%# -*- coding: utf-8-unix -*-
%%==================================================

\chapter{流形上的微积分}

\section{张量的语言}

\subsection{线性代数回顾}
\label{线性代数回顾}

\begin{mythm}{线性扩张}{chapter2线性扩张}
	设$\{\alpha_{ 1 },\alpha_{ 2 },\cdots,\alpha_{ n } \}$为向量空间$V$的任意一组基，$\{\beta_{ 1 },\beta_{ 2 },\cdots,\beta_{ n } \}$为向量空间$W$中的任意$n$个向量. 则存在唯一的线性映射$\mathscr{A}:V\longrightarrow W$，使得$\mathscr{A}(\alpha_{ i })=\beta_{ i},\;i=1,2,\cdots,n$.
\end{mythm}
\begin{proof}
	对任意$\alpha\in V$，则$\alpha$可表示为
	$
	\alpha=a_{1}\alpha_{ 1 }+a_{2}\alpha_{ 2 }+\cdots+a_{n}\alpha_{ n }
	$
	定义$\mathscr{A}:V\longrightarrow W$为$\mathscr{A}(\alpha)=\sum_{i=1}^{n}a_{i}\beta_{ i}$即可。
\end{proof}
这个定理的重要意义在于，你要定义一个线性映射，只需给$V$的一组基指定值，然后通过线性扩张即可给$V$中每一个向量指定它所对应的值。

假设$T:V\rightarrow W$为一个线性映射，$B,B^{\prime}$分别为向量空间$V$和$W$的基。设$B=\{v_{1},v_{2},\cdots,v_{n} \},B^{\prime}=\{w_{1},w_{2},\cdots,w_{m} \}$，则在这两组基下线性映射$T$具有矩阵形式：
\[
T(v_{1},v_{2},\cdots,v_{n})=(w_{1},w_{2},\cdots,w_{m})\left( \begin{array}{cccc}{a_{11}} & {a_{12}} & {\cdots} & {a_{1 n}} \\ {a_{21}} & {a_{22}} & {\cdots} & {a_{2 n}} \\ {\vdots} & {\vdots} & { } & {\vdots} \\ {a_{m 1}} & {a_{m 2}} & {\cdots} & {a_{m n}}\end{array}\right)
\]
我们把上述矩阵记作$[T]_{B}^{B^{\prime}}$. 特别地，当$V=W,B=B^{\prime}$时，简记为$[T]_{B}$. 

以后我们把矩阵$A$中的第$i$行第$j$列的元素$a_{ij}$重新记为$a^{i}_{j}$，即把标记行的指标写在右上角，把标记列的指标写在右下角。比如说
\[
\left( \begin{array}{cccc}{a_{11}} & {a_{12}} & {\cdots} & {a_{1 n}} \\ {a_{21}} & {a_{22}} & {\cdots} & {a_{2 n}} \\ {\vdots} & {\vdots} & { } & {\vdots} \\ {a_{m 1}} & {a_{m 2}} & {\cdots} & {a_{m n}}\end{array}\right)=\left( \begin{array}{cccc}{a_{1}^{1}} & {a^{1}_{2}} & {\cdots} & {a^{1}_{ n}} \\ {a^{2}_{1}} & {a^{2}_{2}} & {\cdots} & {a^{2}_{ n}} \\ {\vdots} & {\vdots} & { } & {\vdots} \\ {a^{m}_{ 1}} & {a^{m}_{ 2}} & {\cdots} & {a^{m}_{ n}}\end{array}\right)
\]

另一个重要的概念是\textbf{对偶空间}，这在第一章\ref{chapter1section1.2.3}节已经详细讲过，此处不再赘述。

\subsection{配合运算}
\label{配合}

\begin{mydef}{配合运算}{chapter2配合运算}
设$V,W$为两个实向量空间. 
\begin{itemize}
\item 若$\langle \quad ,\quad\rangle:V\times W\longrightarrow \mr$是一个双线性映射，则称$\langle \quad ,\quad\rangle$为向量空间$V$与$W$的一个配合(paring).
\item 若配合运算$\langle \quad ,\quad\rangle:V\times W\longrightarrow \mr$满足以下条件：
		\begin{itemize}
		\item[(1)] 由$\langle v,w\rangle=0,\forall w\in W\Rightarrow v=0$
		\item[(2)] 由$\langle v,w\rangle=0,\forall v\in V\Rightarrow w=0$
		\end{itemize}
	   则称$\langle \quad ,\quad\rangle:V\times W\longrightarrow \mr$为向量空间$V$与$W$的一个非退化的配合.
\end{itemize}
\end{mydef}

比如向量空间上的内积就可以看成是该向量空间与自身的一个配合。配合运算最典型的例子是向量空间与其对偶空间的配合。

\begin{example}[对偶配合]
	\label{chapter2对偶空间的配合}
	任取$f\in V^{\ast},v\in V$，则表达式$f(v)$是有意义的，这种写法正是大家习以为常的“函数型”写法。我们也知道对于有限维向量空间$V$，$V$与$V^{\ast}$是\textbf{互为}对偶空间的，但$f(v)$这一记号更多强调的是“$V^{\ast}$是$V$的对偶空间”的含义而体现不出“$V$也是$V^{\ast}$的对偶空间”这一层含义。要想体现出后一种含义，你需要写成$v^{\ast\ast}(f)$，并且由式\requ{equ:1.12}有$f(v)=v^{\ast\ast}(f)$.

	所以如果你坚持使用“$f(v)$”这样的记号，就体现不出“互相”的意味。为了克服这一点，结合$f(v)=v^{\ast\ast}(f)$，我们可以引入$V$与$V^{\ast}$的配合：
	\begin{equation}
	\langle f,v \rangle\triangleq f(v)=v^{\ast\ast}(f)
	\label{对偶空间的配合运算}
	\end{equation}
	容易验证如此得到的确实是一个配合运算，称为向量空间与其对偶空间的对偶配合。
\soln
 
这种写法的好处显而易见：$\langle f,v \rangle$中$f$与$v$两者地位平等。你既可以侧重于把它看成是$f(v)$，也可以侧重于看成是$v^{\ast\ast}(f)$. 

事实上，如果我们把$f$固定不动，则$\langle {\color{gray!80}{f}},\bullet \rangle:V\longrightarrow \mr$是线性函数，从而$\langle {\color{gray!80}{f}},\bullet\rangle\in V^{\ast}$且$\langle {\color{gray!80}{f}},\bullet\rangle=f$；如果把$v$固定不动，则$\langle \bullet,{\color{gray!80}{v}}\rangle: V^{\ast}\longrightarrow \mr$是线性函数，从而$\langle \bullet,{\color{gray!80}{v}}\rangle\in V^{\ast\ast}$，又根据典范同构\rthm{thm:chapter1典范同构定理}，可以把$V^{\ast\ast}$与$V$等同看待，因此$\langle \bullet ,{\color{gray!80}{v}}\rangle\in V^{\ast\ast}=V$且$\langle \bullet ,{\color{gray!80}{v}}\rangle=v^{\ast\ast}=v$. 这样就完美表达了$V$与$V^{\ast}$互为对偶空间的意味。
\end{example}

以后我们在处理向量空间与其对偶空间的元素的计算时，经常会使用对偶配合“$\langle f,v\rangle$”来代替“$f(v)$”. 读者要记住，一旦你看出来$\langle \quad ,\quad\rangle$两个空位中的元素一个是向量空间里的(比如切空间)，另一个是其对偶空间里的(比如余切空间)，那么这个表达式$\langle \quad ,\quad\rangle$就表示对偶配合的意思，即按照式\requ{对偶空间的配合运算}来理解计算。

\begin{mythm}{}{chapter2配合则同构}
	设$V,W$为两个有限维的实向量空间. \\
	若$\langle \quad ,\quad\rangle:V\times W\longrightarrow \mr$是$V$与$W$的一个非退化的配合，则$V\cong W^{\ast}$且$W\cong V^{\ast}$
\end{mythm}
\begin{proof}

显然，当固定$v\in V$不动时，$\langle {\color{gray!80}{v}},\bullet\rangle\in W^{\ast}$；当固定$w\in W$不动时，$\langle \bullet,{\color{gray!80}{w}} \rangle\in V^{\ast}$. 于是我们定义映射$\sigma_{1}:V\longrightarrow W^{\ast}$和$\sigma_{2}:W\longrightarrow V^{\ast}$如下：
\[
	\sigma_{1}(v)=\langle {\color{gray!80}{v}},\bullet\rangle\;,\forall v\in V\quad\quad\sigma_{2}(w)=\langle \bullet,{\color{gray!80}{w}} \rangle\;,\forall w\in W
\]
因为配合运算是双线性的，所以$\sigma_{1},\sigma_{2}$均为线性映射；再由非退化性，则$\sigma_{1},\sigma_{2}$均为单射。最后只需说明$\sigma_{1},\sigma_{2}$均为满射即可。

由单射性质可知，$\operatorname{dim}(V)\leqslant \operatorname{dim}(W^{\ast}),\operatorname{dim}(W)\leqslant \operatorname{dim}(V^{\ast})$. 由于是有限维向量空间，所以$\operatorname{dim}(V)=\operatorname{dim}(V^{\ast}),\operatorname{dim}(W)=\operatorname{dim}(W^{\ast})$. 于是
\[
	\operatorname{dim}(V)\leqslant \operatorname{dim}(W^{\ast})=\operatorname{dim}(W)\leqslant \operatorname{dim}(V^{\ast})=\operatorname{dim}(V)
\]
因此$\operatorname{dim}(V)\leqslant \operatorname{dim}(W)\leqslant \operatorname{dim}(V)$，从而$\operatorname{dim}(V^{\ast})=\operatorname{dim}(V)=\operatorname{dim}(W)=\operatorname{dim}(W^{\ast})$. 
\[
\begin{split}
\sigma_{1}\text{\kaishu 为单射}&\Leftrightarrow \operatorname{dim}(\operatorname{Im}\sigma_{1})=\operatorname{dim}(V)\\
&\Leftrightarrow \operatorname{dim}(\operatorname{Im}\sigma_{1})=\operatorname{dim}(W^{\ast})\\
&\Leftrightarrow \sigma_{1}\text{\kaishu 为满射}
\end{split}
\]
同理$\sigma_{2}$也是满射。也就是说$\sigma_{1}:V\longrightarrow W^{\ast}$和$\sigma_{2}:W\longrightarrow V^{\ast}$为同构映射。
\end{proof}

上述同构映射$\sigma_{1}:V\longrightarrow W^{\ast}$与$\sigma_{2}:W\longrightarrow V^{\ast}$称为由配合运算$\langle \quad ,\quad\rangle$诱导的同构。

\begin{mythm}{}{chapter2同构则配合}
	设$V,W$为两个有限维的实向量空间. 若$\sigma:V\longrightarrow W^{\ast}$为同构，则可以给$V$与$W$赋予适当的非退化配合运算，使得同构映射$\sigma$恰好由此配合运算所诱导.
\end{mythm}
\begin{proof}
	对任意$v\in V,w\in W$，则$\sigma(v)\in W^{\ast}$. 令$\langle v,w\rangle\triangleq \langle \sigma(v),w\rangle$. 其中$\langle \sigma(v),w\rangle$表示$W$与$W^{\ast}$之间的对偶配合

	可以验证这样定义的映射$\langle \quad,\quad\rangle: V\times W\longrightarrow \mr$是向量空间$V$与$W$的一个非退化的配合运算，且它所诱导的$V\longrightarrow W^{\ast}$的同构恰好就是$\sigma$. 
\end{proof}

\rthm{thm:chapter2配合则同构}与\rthm{thm:chapter2同构则配合}一起表明，倘若我们想要描述$V\cong W^{\ast}$这样的事实，可以等价地去给$V$与$W$约定一个非退化的配合运算。




\subsection{多重线性映射}\label{多重线性映射}

\begin{mydef}{$r$重线性映射}{chapter2多重线性映射}
	设$V_{1},V_{2},\cdots,V_{r},W$都是向量空间. 若映射$f:V_{1}\times V_{2}\times\cdots\times V_{r}\longrightarrow W$关于每个变量都是线性的，则称$f$为$r$重线性映射. 并记$L(V_{1},V_{2},\cdots,V_{r};W)$为全体$V_{1}\times V_{2}\times\cdots\times V_{r}\longrightarrow W$的$r$重线性映射所组成的集合. 
\end{mydef}

线性映射可以看作是多重线性映射的特殊情形：1重线性映射。但是注意，大于1重的多重线性映射与线性映射具有本质上的区别！另外，$L(V_{1},V_{2},\cdots,V_{r};W)$在通常的映射的加法与数乘运算之下构成一个向量空间。

为了以后叙述方便，我们引入多重指标的记号。若数组$\gamma=(\gamma_{1},\gamma_{2},\cdots,\gamma_{m})$的每个分量都取正整数，则称$\gamma$是一个多重指标，并令$|\gamma|\triangleq\gamma_{1}+\gamma_{2}+\cdots+\gamma_{m}$. 定义多重指标集如下
\[
\Gamma(n_{1},n_{2},\cdots,n_{m})\triangleq \{\gamma=(\gamma_{1},\gamma_{2},\cdots,\gamma_{m}) | 1\leqslant\gamma_{i}\leqslant n_{i},i=1,\cdots,m\}
\]
令$|\Gamma(n_{1},n_{2},\cdots,n_{m})|\triangleq n_{1}\times n_{2}\times\cdots\times n_{m}$. 显然数$|\Gamma(n_{1},n_{2},\cdots,n_{m})|$反映了指标集中所含元素的个数. 在不引起歧义时，指标集$\Gamma(n_{1},n_{2},\cdots,n_{m})$可简写为$\Gamma$，$|\Gamma(n_{1},n_{2},\cdots,n_{m})|$可简写为$|\Gamma|$.

 注意到$\Gamma$中的元素可以按照字典排序法从小到大一个不漏地排起来成为有序集合，从而可与数集$\{1,2,3,\cdots, n_{1}\times n_{2}\times\cdots\times n_{m}\}$中的数一一对应，正是因为存在这种一一对应的关系，以后我们也直接用多重指标来表示与它对应的那个正整数。
 
 举个例子，指标集$\Gamma(2,4,3)$在进行字典排序后，多重指标$(2,1,1)$对应于正整数$13$，$(2,3,3)$对应于正整数$21$，$(1,3,2)$对应于正整数$8$，$(2,4,3)$对应于正整数$24$等等，所有这些多重指标一个不落得对应完$\{1,2,\cdots,24 \}$中的所有正整数。如果记$\gamma=(1,3,2)$，则说法“矩阵的第$\gamma$行”实际上是说“矩阵的第$8$行”，因为$\gamma$对应于正整数$8$；同理，说法“矩阵的第$(2,3,3)$行第$(2,1,1)$列”其实意指“矩阵的第$21$行第$13$列”。
 
 在PDE中，用多重指标的记号可以十分简洁地表示高阶混合偏导数，比如若有多元函数$u(x_{1},x_{2},\cdots,x_{m})$，以及多重指标$\gamma=(\gamma_{1},\gamma_{2},\cdots,\gamma_{m})$，则$D^{\gamma}f$就表示偏导数$\dfrac{\partial^{|\gamma|}}{\partial x_{1}^{\gamma_{1}}\cdots\partial x_{i}^{\gamma_{ i }}\cdots\partial x_{m}^{\gamma_{m}}}u$. 我们以后会多次使用与此类似的表示方法.
 
设$V_{1},V_{2},\cdots,V_{r},W$都是向量空间，其中$\operatorname{dim}(V_{i})=n_{i},1\leqslant i\leqslant r$. 再设$\{e_{i1},e_{i2},\cdots,e_{in_{i}} \}$是$V_{i},1\leqslant i\leqslant r$的基。若$f:V_{1}\times V_{2}\times\cdots\times V_{r}\longrightarrow W$是任一个$r$重线性映射，则由定义可知
\begin{equation}
	\begin{split}
	f(v_{1},v_{2},\cdots,v_{r})&=f\left(\sum_{k_{1}=1}^{n_{1}}a_{1k_{1}}e_{1k_{1}},\sum_{k_{2}=1}^{n_{2}}a_{2k_{2}}e_{2k_{2}},\cdots,\sum_{k_{r}=1}^{n_{r}}a_{rk_{r}}e_{rk_{r}} \right)\\
	&=\sum_{k_{1}=1}^{n_{1}}\sum_{k_{2}=1}^{n_{2}}\cdots\sum_{k_{r}=1}^{n_{r}}a_{1k_{1}}a_{2k_{2}}\cdots a_{rk_{r}}f\left(e_{1k_{1}},e_{2k_{2}},\cdots,e_{rk_{r}} \right)\\
	&=\sum_{ \gamma\in \Gamma }a_{1\gamma_{1}}a_{2\gamma_{2}}\cdots a_{r\gamma_{r}}f\left(e_{1\gamma_{1}},e_{2\gamma_{2}},\cdots,e_{r\gamma_{r}} \right)\\
	&=\sum_{ \gamma\in \Gamma }a_{\gamma}f(e_{\gamma})
	\end{split}
\end{equation}
其中约定$(k_{1},k_{2},\cdots,k_{r})=\gamma\in \Gamma=\Gamma(n_{1},\cdots,n_{r})\;, a_{\gamma}\triangleq \prod_{i=1}^{r}a_{i\gamma_{i}}\;,e_{\gamma}\triangleq \left(e_{1\gamma_{1}},e_{2\gamma_{2}},\cdots,e_{r\gamma_{r}} \right)$. 利用多重指标的记号，我们就给出了多重线性映射在给定基下计算的一般表达式。如何使用多重指标全凭个人心意，只有一个原则：怎么方便怎么来！比如我们这里约定的$a_{\gamma}$与$e_{\gamma}$，它们具有不同的含义，具体什么意思就看你怎么约定的。

将上式与线性映射的计算表达式$f(v)=\sum_{i=1}^{n}a_{i}f(e_{i})$类比，你会发现在形式上就是把从$1$到$n$的求和改写为对多重指标集$\Gamma$的求和。
\begin{mythm}{多重线性扩张}{chapter2多重线性扩张}
	设$V_{1},V_{2},\cdots,V_{r},W$都是向量空间，$\{e_{i1},e_{i2},\cdots,e_{in_{i}} \}$是$V_{i},1\leqslant i\leqslant r$的基. 则存在唯一的$r$重线性映射$\varphi:V_{1}\times V_{2}\times\cdots\times V_{r}\longrightarrow W$，使得$\varphi(e_{\gamma})=w_{ \gamma}\in W,\;\gamma\in \Gamma$. 其中$e_{\gamma},w_{\gamma}$约定如上.
\end{mythm}
\begin{proof}
	对于$v_{i}\in V_{i}$，则$v_{i}$可表示为
	$
	v_{i}=a_{i1}e_{i1}+a_{i2}e_{i2}+\cdots+a_{in_{i}}e_{in_{i}}\;,i=1,2,\cdots,r
	$
	定义$\varphi:V_{1}\times V_{2}\times\cdots\times V_{r}\longrightarrow W$为$\varphi(v_{1},v_{2},\cdots,v_{r})=\sum_{ \gamma\in \Gamma }a_{\gamma}w_{\gamma}$即可，其中$a_{\gamma}\triangleq \prod_{i=1}^{r}a_{i\gamma_{i}}$.
\end{proof}

这个定理与线性扩张的意义一样重大，你要定义一个$r$重线性映射$\varphi:V_{1}\times V_{2}\times\cdots\times V_{r}\longrightarrow W$，只需给$\operatorname{dim}(V_{1})\times\operatorname{dim}(V_{2})\times\cdots\times\operatorname{dim}(V_{r})=|\Gamma|$个元素$\{e_{\gamma} \}_{\gamma\in \Gamma}$指定值，然后通过多重线性扩张即可给$V_{1}\times V_{2}\times\cdots\times V_{r}$中每一个元素指定它所对应的值。特别地，$V\times V$上的的双线性映射(函数)可以看成是该向量空间$V$上的一种\textbf{二元运算}。因此，以后如果你想定义一个具有双线性性质的二元运算，那么就只需对$n\times n$个元素$\{(e_{i},e_{j}) \}$定义相应的运算值再做2重线性扩张即可，其中$\{e_{1},\cdots,e_{n} \}$为$V$的基。在张量以及后续的很多讨论中，我们都是按照这种方式来定义所需的新的运算。

我们前面说过，线性映射与多重线性映射有本质的区别，多重线性扩张定理就揭示了一个明显的区别。回忆要确定线性映射$\mathscr{A}:V\longrightarrow W$，你需要给定值的点的数目等于$V$的维数$\operatorname{dim}(V)$. 于是当$\varphi:V_{1}\times V_{2}\times\cdots\times V_{r}\longrightarrow W$为\textbf{线性映射}时，要确定它只要给定$\operatorname{dim}(V_{1}\times V_{2}\times\cdots\times V_{r})=\sum_{i=1}^{r}\operatorname{dim}(V_{i})$个点的值即可；但是当$\varphi:V_{1}\times V_{2}\times\cdots\times V_{r}\longrightarrow W$为\textbf{$r$重线性映射}时，要确定它需要给定值的点的数目为$|\Gamma|=\prod_{i=1}^{r}\operatorname{dim}(V_{i})$，这个数一般来说要比向量空间$V_{1}\times V_{2}\times\cdots\times V_{r}$的维数大得多！

我们又知道线性映射$\mathscr{A}:V\longrightarrow W$的值域$\mathscr{A}(V)=\operatorname{Im}(\mathscr{A})$是$W$的一个子空间，但多重线性映射$\varphi:V_{1}\times V_{2}\times\cdots\times V_{r}\longrightarrow W$的值域$\operatorname{Im}(\varphi)$却不一定是$W$的子空间。这是线性映射与多重线性映射的另一个不太明显但更加本质的区别！

虽然多重线性映射$\varphi:V_{1}\times V_{2}\times\cdots\times V_{r}\longrightarrow W$的值域$\operatorname{Im}(\varphi)$不能像线性映射的值域那样完美地自成一个线性空间，不过也没关系，我们可以用它来生成一个子空间嘛。考虑由$\operatorname{Im}(\varphi)$所生成的$W$的子空间$\langle\operatorname{Im}(\varphi) \rangle$. 
\begin{mythm}{}{chapter2像空间的生成子空间}
	$\langle\operatorname{Im}(\varphi) \rangle=\langle\{\varphi(e_{\gamma}) \}_{\gamma\in \Gamma} \rangle $.
\end{mythm}
\begin{proof}
	显然$\langle\{\varphi(e_{\gamma}) \}_{\gamma\in \Gamma} \rangle \subset \langle\operatorname{Im}(\varphi) \rangle$. 另一方面根据生成子空间的定义，$\langle\operatorname{Im}(\varphi) \rangle$中的元素是$\operatorname{Im}(\varphi)$中元素的线性组合，而由式\requ{equ:}又知，$\operatorname{Im}(\varphi)$中的可以表示为$\{\varphi(e_{\gamma}) \}_{\gamma\in \Gamma}$中元素的线性组合，因此$\langle\operatorname{Im}(\varphi) \rangle$中的元素都可以表示成$\{\varphi(e_{\gamma}) \}_{\gamma\in \Gamma}$中元素的线性组合，从而$\langle\operatorname{Im}(\varphi) \rangle\subset\langle\{\varphi(e_{\gamma}) \}_{\gamma\in \Gamma} \rangle $. 
	
	综上所述，$\langle\operatorname{Im}(\varphi) \rangle=\langle\{\varphi(e_{\gamma}) \}_{\gamma\in \Gamma} \rangle $.
\end{proof}

上述定理告诉我们，$\langle\operatorname{Im}(\varphi) \rangle$可能达到的最大的维数为$|\Gamma|$，即$\operatorname{dim}(\langle\operatorname{Im}(\varphi) \rangle)\leqslant \prod_{i=1}^{r}\operatorname{dim}(V_{i})$. 这是多重线性映射的一个十分关键的特征。

\begin{mydef}{张映射}{chapter2张映射}
	设$\varphi:V_{1}\times V_{2}\times\cdots\times V_{r}\longrightarrow W$为$r$重线性映射. 若$\operatorname{dim}(\langle\operatorname{Im}(\varphi) \rangle)= \prod_{i=1}^{r}\operatorname{dim}(V_{i})$，则称$\varphi$为张映射.
\end{mydef}

也就是说，张映射是张开到最大维数的一种多重线性映射。
\begin{mythm}{}{chapter2张映射的等价条件}
	$\varphi:V_{1}\times V_{2}\times\cdots\times V_{r}\longrightarrow W$为张映射$\Leftrightarrow\{\varphi(e_{\gamma}) \}_{\gamma\in \Gamma}$线性无关.
\end{mythm}
\begin{proof}
	由$\langle\operatorname{Im}(\varphi) \rangle=\langle\{\varphi(e_{\gamma}) \}_{\gamma\in \Gamma} \rangle $以及$|\Gamma|=\prod_{i=1}^{r}\operatorname{dim}(V_{i})$容易看出定理成立。
\end{proof}

张映射的重要性体现在，它可以把任意其他的\textbf{多重线性映射“转换”成线性映射}来研究。
\begin{mythm}{因子化性质}{chapter2因子化性质}
	设$\varphi:V_{1}\times V_{2}\times\cdots\times V_{r}\longrightarrow W$为$r$重线性映射. 则$\varphi$为张映射的充要条件是：\\
	对任意的向量空间$Z$以及任意多重线性映射$f:V_{1}\times V_{2}\times\cdots\times V_{r}\longrightarrow Z$，总存在线性映射$g:W\longrightarrow Z$，使得$f=g\circ\varphi$. 也即下图可交换
	\[
	\begin{tikzcd}
	V_{1}\times V_{2}\times\cdots\times V_{r} \arrow[r, "\varphi"] \arrow[d, "f", left ] & W \arrow[dl,dashrightarrow, "g"]\\
	Z  &
	\end{tikzcd} 
	\]
\end{mythm}

这个定理我们不打算证了，因为相比于证明，结论更加重要！值得一提的是，利用因子化性质我们可以给出张映射的等价定义。事实上，由因子化性质所给出的张映射的定义适用范围比我们这给出的张映射的定义来得更广！我们所给的定义只适用于有限维向量空间；对于无穷维的向量空间，就必须得用因子化性质来定义张映射。不过一般我们所处理的空间还是以有限维为主。

\begin{remark}
	上述定理中的线性映射$g:W\longrightarrow Z$不一定是唯一的。事实上可以进一步证明：线性映射$g:W\longrightarrow Z$是唯一的$\Leftrightarrow \langle\operatorname{Im}(\varphi) \rangle=W$
\end{remark}

最后我们介绍多重线性函数的张量积运算。$r$重线性映射$\varphi:V_{1}\times V_{2}\times\cdots\times V_{r}\longrightarrow {\color{red}{\mr}}$又称为\textbf{$r$重线性函数}。两个多重线性函数可以生成一个新的多重线性函数。
\begin{mydef}{多重线性函数的张量积}{chapter多重线性函数的张量积}
	设$f:V_{1}\times V_{2}\times\cdots\times V_{r}\longrightarrow {\color{red}{\mr}}$是一个$r$重线性函数，$g:W_{1}\times W_{2}\times\cdots\times W_{s}\longrightarrow {\color{red}{\mr}}$是一个$s$重线性函数. 对任意$(v_{1},v_{2},\cdots,v_{r})\in V_{1}\times V_{2}\times\cdots\times V_{r}$，任意$(w_{1},w_{2},\cdots,w_{s})\in W_{1}\times W_{2}\times\cdots\times W_{s}$，令
	\[
	f\otimes g(v_{1},v_{2},\cdots,v_{r},w_{1},w_{2},\cdots,w_{s})\triangleq f(v_{1},v_{2},\cdots,v_{r})\cdot g(w_{1},w_{2},\cdots,w_{s})
	\]
	则$f\otimes g:V_{1}\times V_{2}\times\cdots\times V_{r}\times W_{1}\times W_{2}\times\cdots\times W_{s}\longrightarrow F$是一个$r+s$重线性函数. 我们称$f\otimes g$为$f$与$g$的张量积.
\end{mydef}

多重线性函数的张量积运算有一些简单的性质：
\begin{itemize}
	\item[(1)] $\otimes$是双线性的，即$\otimes $关于$f$与$g$都是线性的.
	\item[(2)] $\otimes$满足结合律，即$(f\otimes g)\otimes h=f\otimes (g\otimes h)$.
	\item[(3)] $f_{1}\otimes f_{2}\otimes \cdots\otimes f_{r}$关于每一个$f_{i}$都是线性的，即
	           \[
	            f_{1}\otimes\cdots\otimes \left(\lambda_{1}h_{i}+\lambda_{2}g_{i} \right)\otimes \cdots\otimes f_{r}=\lambda_{1}f_{1}\otimes \cdots\otimes h_{i}\otimes \cdots\otimes f_{r}+\lambda_{2}f_{1}\otimes \cdots\otimes g_{i}\otimes \cdots\otimes f_{r}
	           \]
\end{itemize}
正是因为结合律的成立，所以对于多个多重线性函数的张量积我们可以不用写括号。第(3)个性质是(1)与(2)的直接推论(当然直接按定义验证也很简单)。

\subsection{张量空间}

\begin{mydef}{张量空间}{chapter2张量空间}
	设$W$为一个向量空间. 若存在张映射$\varphi:V_{1}\times V_{2}\times\cdots\times V_{r}\longrightarrow W$使得$W=\langle\operatorname{Im}(\varphi) \rangle$，则称向量空间$W$为$V_{1},V_{2},\cdots,V_{r}$的带有张映射$\varphi$的张量空间. 此时$W$中的向量就称为张量.
\end{mydef}

由于$\varphi:V_{1}\times V_{2}\times\cdots\times V_{r}\longrightarrow W$是张映射且$W=\langle\operatorname{Im}(\varphi) \rangle$，根据张映射的定义立即可知张量空间$W$的维数是$\prod_{i=1}^{r}\operatorname{dim}(V_{i})$. 下面我们就来给张量空间$W$找出一组具体的基。

因为$W=\langle\operatorname{Im}(\varphi) \rangle $，又由\rthm{thm:chapter2像空间的生成子空间}可知$\langle\operatorname{Im}(\varphi)= \langle\{\varphi(e_{\gamma}) \}_{\gamma\in \Gamma} \rangle $，所以$W=\langle\{\varphi(e_{\gamma}) \}_{\gamma\in \Gamma} \rangle $. 又因为$\varphi$是张映射，所以$\{\varphi(e_{\gamma}) \}_{\gamma\in \Gamma}$线性无关. 这表明{\color{red}{$\{\varphi(e_{\gamma}) \}_{\gamma\in \Gamma}$就是张量空间$W$的一组基}}！

因此，任一张量$z\in W$都可以表示为$z=\sum_{ \gamma\in \Gamma }z_{\gamma}\varphi(e_{\gamma})$，我们把系数$\{z_{\gamma} \}_\gamma\in \Gamma$称为张量$z$在基$\{\varphi(e_{\gamma}) \}_{\gamma\in \Gamma}$下的坐标分量。

\begin{mythm}{唯一因子化性质}{chapter2唯一因子化性质}
	$W$为$V_{1},V_{2},\cdots,V_{r}$的带有张映射$\varphi$的张量空间$\Leftrightarrow$
	对任意的向量空间$Z$以及任意多重线性映射$f:V_{1}\times V_{2}\times\cdots\times V_{r}\longrightarrow Z$，总存在{\color{red}{唯一的}}线性映射$g:W\longrightarrow Z$，使得$f=g\circ\varphi$. 也即下图可交换
	\[
	\begin{tikzcd}
	V_{1}\times V_{2}\times\cdots\times V_{r} \arrow[r, "\varphi"] \arrow[d, "f", left ] & W \arrow[dl,dashrightarrow, "{\color{red}g}"]\\
	Z  &
	\end{tikzcd} 
	\]
\end{mythm}
\begin{proof}
	$\varphi$是张映射$\Leftrightarrow\varphi$具有因子化性质；$\langle\operatorname{Im}(\varphi) \rangle=W\Leftrightarrow$线性映射$g:W\longrightarrow Z$是唯一的，因而定理成立。
\end{proof}

从定义可以看出，张量空间$W$是由$V_{1},V_{2},\cdots,V_{r}$以及张映射$\varphi$所共同确定的，因此关于$V_{1},V_{2},\cdots,V_{r}$的张量空间可以有很多。下面的定理表明，在同构的意义下关于$V_{1},V_{2},\cdots,V_{r}$的张量空间是唯一的。
\begin{mythm}{张量空间的唯一性}{chapter2张量空间的唯一性}
	设$W$为$V_{1},V_{2},\cdots,V_{r}$的带有张映射$\varphi$的张量空间，$V$为$V_{1},V_{2},\cdots,V_{r}$的带有张映射$\phi$的张量空间，则$W$与$V$同构.
\end{mythm}
\begin{proof}
	利用唯一因子化性质易知$g_{1}\circ g_{2}=id_{Z},g_{2}\circ g_{1}=id_{W}$. 从而$g_{1}$与$g_{2}$均为双射，又是线性映射，所以$W$与$V$同构.
	\[
	\begin{tikzcd}
	V_{1}\times V_{2}\times\cdots\times V_{r} \arrow[r, "\varphi"] \arrow[d, "\phi", left ] & W \arrow[dl, dashrightarrow,shift left=1ex, "{\color{red}g_{1}}"]\\
	Z \arrow[ur, dashrightarrow, "{\color{red}g_{2}}"] &
	\end{tikzcd} 
	\]
\end{proof}

由此可知$V_{1},V_{2},\cdots,V_{r}$的所有张量空间都是同构的，故我们只需选取其中一个张量空间作为代表来研究就行了。下面我们就来构造一个具体的张量空间$W$，从而给出张量的一个具体模型。
想要构造张量空间$W$，只需构造张映射$\varphi:V_{1}\times V_{2}\times\cdots\times V_{r}\longrightarrow W$即可，因为$W=\langle\operatorname{Im}(\varphi) \rangle$. 我们试图用一种比较自然的方式来构造出$\varphi$. 

首先我们要换个视角来重新看待每个向量$v$.
我们知道每一个向量空间都与其二次对偶空间典范同构，典范同构的意思就是说你可以把向量空间与其二次对偶空间看作是完全一样的，用不着区分。这样每一个向量$v\in V$，又有另一种表达方式$v^{\ast\ast}\in V^{\ast\ast}$，即$v=v^{\ast\ast}$. 而将$v$改写成$v^{\ast\ast}$的好处就在于，$v^{\ast\ast}$是定义在一次对偶空间$V^{\ast}$上的\textbf{线性函数}\footnote{Instead of working with vectors, it turns out to be more fruitful to adopt the dual point of view and work with linear functions on a vector space. After all, there is only so much that one can do with vectors, which are essentially arrows, but functions, far more flexible, can be added, scalar-multiplied, and composed with other maps.}。

对于任意$(v_{1},v_{2},\cdots,v_{r})\in V_{1}\times V_{2}\times\cdots\times V_{r}$，在上述对偶观点下，我们可以把$(v_{1},v_{2},\cdots,v_{r})$重新改写为$(v^{\ast\ast}_{1},v^{\ast\ast}_{2},\cdots,v^{\ast\ast}_{r})$，即
\[
(v_{1},v_{2},\cdots,v_{r})=(v^{\ast\ast}_{1},v^{\ast\ast}_{2},\cdots,v^{\ast\ast}_{r})
\]

现在注意到每一个$v^{\ast\ast}_{i},\;i=1,\cdots,r$都是定义在相应的一次对偶空间$V^{\ast}_{i}$上的\textbf{1重线性函数}. 对于线性函数我们有很多运算可以做，比如加法、数乘、复合，还有上一节刚刚介绍的多重线性函数的张量积运算！我们考虑将这$r$个1重线性函数$v^{\ast\ast}_{1},v^{\ast\ast}_{2},\cdots,v^{\ast\ast}_{r}$\textbf{做张量积运算}，得到一个$r$重线性函数
\[
v^{\ast\ast}_{1}\otimes v^{\ast\ast}_{2}\otimes\cdots\otimes v^{\ast\ast}_{r}:V^{\ast}_{1}\times V^{\ast}_{2}\times\cdots V^{\ast}_{r}\longrightarrow \mr
\]
把这个$r$重线性函数的作用方式具体写开来就是：
\begin{equation}
	v^{\ast\ast}_{1}\otimes \cdots\otimes v^{\ast\ast}_{r}\left(\sigma_{1},\cdots,\sigma_{r} \right)=v^{\ast\ast}_{1}(\sigma_{1})\cdots v^{\ast\ast}_{r}(\sigma_{r}),\;\sigma_{i}\in V^{\ast}_{i},i=1,\cdots,r
\end{equation}

再捋一下上述想法，也就是说
\begin{center}
	\begin{tikzpicture}[scale=1, node distance=1.1cm, auto]
	\node [property, text width=10em] (MId1) {$(v_{1},v_{2},\cdots,v_{r})$};
	\node [notation, below of=MId1, text width=8em] (ONeg1) {$(v^{\ast\ast}_{1},v^{\ast\ast}_{2},\cdots,v^{\ast\ast}_{r})$};
	\node [operation, below of=ONeg1, text width=10em] (DOS1) {$v^{\ast\ast}_{1}\otimes v^{\ast\ast}_{2}\otimes\cdots\otimes v^{\ast\ast}_{r}$};	
	\path [line] (MId1) -- (ONeg1);
	\path [line] (ONeg1) -- (DOS1);
	\end{tikzpicture}
\end{center}
这启发我们考虑这样的映射$\varphi$：对任意$(v_{1},v_{2},\cdots,v_{r})\in V_{1}\times V_{2}\times\cdots\times V_{r}$，定义
\begin{equation}
	\varphi(v_{1},v_{2},\cdots,v_{r})\triangleq v^{\ast\ast}_{1}\otimes v^{\ast\ast}_{2}\otimes\cdots\otimes v^{\ast\ast}_{r}\in L(V^{\ast}_{1}, V^{\ast}_{2},\cdots ,V^{\ast}_{r};\mr)
\end{equation}

注意到二次对偶的性质$(\lambda_{1}v+\lambda_{2}w)^{\ast\ast}=\lambda_{1}v^{\ast\ast}+\lambda_{2}w^{\ast\ast}$，且张量积运算$\otimes$关于每一个位置都是线性的，所以不难发现如上定义的$\varphi$是$V_{1}\times V_{2}\times\cdots\times V_{r}\longrightarrow L(V^{\ast}_{1}, V^{\ast}_{2},\cdots ,V^{\ast}_{r};\mr)$的一个$r$重线线性映射，是不是有点小激动！接下一个自然问题就是，这样定义的$\varphi$是张映射吗？答案是肯定的！
\begin{mythm}{}{chapter2张映射的构造定理}
	式\requ{equ:}定义的$r$重线性映射$\varphi$是张映射.
\end{mythm}
\begin{proof}
	根据\rthm{thm:chapter2张映射的等价条件}，要证明$\varphi$是张映射，只要证$\{\varphi(e_{\gamma}) \}_{\gamma\in \Gamma}$线性无关. 设
\begin{equation}
	\sum_{ \gamma\in \Gamma }k_{\gamma}\varphi(e_{\gamma})=0
\end{equation}
	其中$k_{\gamma}\in \mr$. 将上式两边作用在$e^{\ast}_{\gamma}$上，其中$e^{\ast}_{\gamma}\triangleq (e^{\ast}_{1\gamma_{1}},e^{\ast}_{2\gamma_{2}},\cdots,e^{\ast}_{r\gamma_{r}})\in V^{\ast}_{1}\times V^{\ast}_{2}\times\cdots V^{\ast}_{r}$. 直接计算
	\[
	\begin{split}
	k_{\alpha}\varphi(e_{\alpha})(e^{\ast}_{\gamma})&=k_{\alpha}e^{\ast\ast}_{1\alpha_{1}}\otimes \cdots\otimes e^{\ast\ast}_{r\alpha_{r}}\left(e^{\ast}_{1\gamma_{1}},\cdots,e^{\ast}_{r\gamma_{r}} \right)\\
	&=k_{\alpha}e^{\ast\ast}_{1\alpha_{ 1 }}\left(e^{\ast}_{1\gamma_{1}} \right)\cdots e^{\ast\ast}_{r\alpha_{ r }}\left(e^{\ast}_{r\gamma_{r}} \right)\\
	&=k_{\alpha}e^{\ast}_{1\gamma_{1}}\left(e_{1\alpha_{ 1 }} \right)\cdots e^{\ast}_{1\gamma_{1}}\left(e_{1\alpha_{ 1 }} \right)\\
	&=k_{\alpha}\cdot \delta_{ \gamma_{1} \alpha_{ 1 } }\cdots \delta_{ \gamma_{r} \alpha_{ r } }
	\end{split}
	\]
	
	对于$\alpha\neq \gamma\in \Gamma$，则至少存在一个$\gamma_{ i }\neq\alpha_{ i }$，从而$\delta_{ \gamma_{i} \alpha_{ i } }=0\Rightarrow k_{\alpha}\varphi(e_{\alpha})(e^{\ast}_{\gamma})=0$.
	
	对于$\alpha=\gamma\in\Gamma$，则$\gamma_{1}=\alpha_{ 1 },\cdots,\gamma_{r}=\alpha_{ r }$，从而$\delta_{ \gamma_{1} \alpha_{ 1 } }=1,\cdots, \delta_{ \gamma_{r} \alpha_{ r } }=1\Rightarrow k_{\gamma}\varphi(e_{\gamma})(e^{\ast}_{\gamma})=k_{\gamma}$. 
	
	因此将\requ{equ:}两边作用在$e^{\ast}_{\gamma}$上后得到$k_{\gamma}=0,\forall \gamma\in \Gamma$. 这表明$\{\varphi(e_{\gamma}) \}_{\gamma\in \Gamma}$线性无关，从而$\varphi$是张映射。
\end{proof}

这样我们便构造出了一个张映射$\varphi$，再令$W=\langle\operatorname{Im}(\varphi) \rangle$就得到了一个具体的张量空间$W$.

 我们来看如此构造出来的$W$中的元素(也就是张量)究竟长什么样？按照$\varphi$的构造方式可知，$\operatorname{Im}(\varphi) $中的元素是形如$v^{\ast\ast}_{1}\otimes v^{\ast\ast}_{2}\otimes\cdots\otimes v^{\ast\ast}_{r}$这个样子的、作用在$V^{\ast}_{1}\times V^{\ast}_{2}\times\cdots V^{\ast}_{r}$上的$r$重线性函数，于是$\langle\operatorname{Im}(\varphi) \rangle$中的元素都是它们的线性组合，自然也是作用在$V^{\ast}_{1}\times V^{\ast}_{2}\times\cdots V^{\ast}_{r}$上的$r$重线性函数。也就是说，$W$中的元素都是作用在$V^{\ast}_{1}\times V^{\ast}_{2}\times\cdots V^{\ast}_{r}$上的$r$重线性函数！用数学符号来表示，即$W\subset L(V^{\ast}_{1}, V^{\ast}_{2},\cdots ,V^{\ast}_{r};\mr)$. 聪明的你肯定会想，反过来对不对呢？
 \begin{mythm}{}{chapter2反过来也对定理}
 	设$\varphi$是由式\requ{equ:}定义的$r$重线性映射，则$\langle\operatorname{Im}(\varphi) \rangle=L(V^{\ast}_{1}, V^{\ast}_{2},\cdots ,V^{\ast}_{r};\mr)$.
 \end{mythm}
\begin{proof}
	由之前的讨论已经知道$\langle\operatorname{Im}(\varphi) \rangle\subset L(V^{\ast}_{1}, V^{\ast}_{2},\cdots ,V^{\ast}_{r};\mr)$. 只需再说明
	\[
		L(V^{\ast}_{1}, V^{\ast}_{2},\cdots ,V^{\ast}_{r};\mr)\subset \langle\operatorname{Im}(\varphi) \rangle
	\]
	为此任取$f\in L(V^{\ast}_{1},V^{\ast}_{2},\cdots,V^{\ast}_{r};\mr)$，则
	\begin{equation}
	\begin{split}
	f(\sigma_{1},\sigma_{2},\cdots,\sigma_{r})&=f\left(\sum_{k_{1}=1}^{n_{1}}\sigma_{1}(e_{1k_{1}})e^{\ast}_{1k_{1}},\sum_{k_{2}=1}^{n_{2}}\sigma_{2}(e_{2k_{2}})e^{\ast}_{2k_{2}},\cdots,\sum_{k_{r}=1}^{n_{r}}\sigma_{r}(e_{rk_{r}})e^{\ast}_{rk_{r}} \right)\\
	&=\sum_{k_{1}=1}^{n_{1}}\sum_{k_{2}=1}^{n_{2}}\cdots\sum_{k_{r}=1}^{n_{r}}\sigma_{1}(e_{1k_{1}})\sigma_{2}(e_{2k_{2}})\cdots \sigma_{r}(e_{rk_{r}})f\left(e^{\ast}_{1k_{1}},e^{\ast}_{2k_{2}},\cdots,e^{\ast}_{rk_{r}} \right)\\
	&=\sum_{ \alpha\in \Gamma }\sigma_{1}(e_{1\alpha_{1}})\sigma_{2}(e_{2\alpha_{2}})\cdots \sigma_{r}(e_{r\alpha_{r}})f\left(e^{\ast}_{1\alpha_{1}},e^{\ast}_{2\alpha_{2}},\cdots,e^{\ast}_{r\alpha_{r}} \right)\\
	&=\sum_{ \alpha\in \Gamma }\sigma_{1}(e_{1\alpha_{1}})\sigma_{2}(e_{2\alpha_{2}})\cdots \sigma_{r}(e_{r\alpha_{r}})f(e^{\ast}_{\alpha})\\
	&=\sum_{ \alpha\in \Gamma }f(e^{\ast}_{\alpha})\sigma_{1}(e_{1\alpha_{1}})\sigma_{2}(e_{2\alpha_{2}})\cdots \sigma_{r}(e_{r\alpha_{r}})\\
	&=\sum_{ \alpha\in \Gamma }f(e^{\ast}_{\alpha})e^{\ast\ast}_{1\alpha_{1}}(\sigma_{1})e^{\ast\ast}_{2\alpha_{2}}(\sigma_{2})\cdots e^{\ast\ast}_{r\alpha_{r}}(\sigma_{r})\\
	&=\sum_{ \alpha\in \Gamma }f(e^{\ast}_{\alpha})e^{\ast\ast}_{1\alpha_{1}}\otimes e^{\ast\ast}_{2\alpha_{2}}\cdots \otimes e^{\ast\ast}_{r\alpha_{r}}\left( \sigma_{1},\sigma_{2},\cdots,\sigma_{r}\right)
	\end{split}
	\end{equation}
	其中约定$(k_{1},k_{2},\cdots,k_{r})=\alpha\in \Gamma=\Gamma(n_{1},\cdots,n_{r})\;,e^{\ast}_{\alpha}\triangleq \left(e^{\ast}_{1\alpha_{1}},e^{\ast}_{2\alpha_{2}},\cdots,e^{\ast}_{r\alpha_{r}} \right)$. 
	
	于是从上式我们得出，$f=\sum_{ \alpha\in \Gamma }f(e^{\ast}_{\alpha})e^{\ast\ast}_{1\alpha_{1}}\otimes e^{\ast\ast}_{2\alpha_{2}}\cdots \otimes e^{\ast\ast}_{r\alpha_{r}}$，从而$f\in \langle\operatorname{Im}(\varphi) \rangle\Rightarrow L(V^{\ast}_{1}, V^{\ast}_{2},\cdots ,V^{\ast}_{r};\mr)\subset \langle\operatorname{Im}(\varphi) \rangle$.
	
	综上所述，我们有$\langle\operatorname{Im}(\varphi) \rangle=L(V^{\ast}_{1}, V^{\ast}_{2},\cdots ,V^{\ast}_{r};\mr)$.
\end{proof}

上面这个定理揭开了张量空间最后的面纱。我们构造的张量空间$W=\langle\operatorname{Im}(\varphi) \rangle$不是别的什么东西，它就是$V^{\ast}_{1}\times V^{\ast}_{2}\times\cdots\times V^{\ast}_{r}$上的全体$r$重线性函数组成的向量空间$L(V^{\ast}_{1}, V^{\ast}_{2},\cdots ,V^{\ast}_{r};\mr)$！所谓的张量也不是新东西，就是作用在$V^{\ast}_{1}\times V^{\ast}_{2}\times\cdots\times V^{\ast}_{r}$上的$r$重线性函数！又因为$V_{1},V_{2},\cdots,V_{r}$的所有张量空间彼此同构，所以我们以后就把$L(V^{\ast}_{1}, V^{\ast}_{2},\cdots ,V^{\ast}_{r};\mr)$作为默认的张量空间，事实上这也是很多教材对于张量的定义。

这样我们就给出了一个十分具体的张量及张量空间的模型。读者要明白这种构造的核心想法：把不太好操作的向量，转换为更容易操作的多重线性函数。具体的转换通过典范同构$v=v^{\ast\ast}$完成。总结一下：
\begin{lstlisting}[language={python}, caption={张量空间的构造}]
已知：$r$个向量空间$V_{1},V_{2},\cdots,V_{r}$
目标：构造$V_{1},V_{2},\cdots,V_{r}$的一个张量空间$W$

for $\forall (v_{1},v_{2},\cdots,v_{r})\in V_{1}\times V_{2}\times\cdots\times V_{r}$，定义:
    $\varphi(v_{1},v_{2},\cdots,v_{r})\triangleq v^{\ast\ast}_{1}\otimes v^{\ast\ast}_{2}\otimes\cdots\otimes v^{\ast\ast}_{r}\in L(V^{\ast}_{1}, V^{\ast}_{2},\cdots ,V^{\ast}_{r};\mr)$
则  $\langle\operatorname{Im}(\varphi) \rangle=L(V^{\ast}_{1}, V^{\ast}_{2},\cdots ,V^{\ast}_{r};\mr)$

最后令$W=L(V^{\ast}_{1}, V^{\ast}_{2},\cdots ,V^{\ast}_{r};\mr)$
\end{lstlisting}

接下来我们把上述构造的\textbf{思想}应用到对偶空间$V^{\ast}_{1},V^{\ast}_{2},\cdots,V^{\ast}_{r}$上。

对于任意$(v^{\ast}_{1},v^{\ast}_{2},\cdots,v^{\ast}_{r})\in V^{\ast}_{1}\times V^{\ast}_{2}\times\cdots\times V^{\ast}_{r}$，注意到这时候每一个$v^{\ast}_{i}$本身就是现成的\textbf{1重线性函数}(定义域为向量空间$V_{i}$)，所以我们不必像之前那样把$v_{i}$改写成它的二次对偶再用，而可以直接用$v^{\ast}_{i}$. 即令
\begin{equation}
\bar{\varphi}(v^{\ast}_{1},v^{\ast}_{2},\cdots,v^{\ast}_{r})\triangleq v^{\ast}_{1}\otimes v^{\ast}_{2}\otimes\cdots\otimes v^{\ast}_{r}\in L(V_{1}, V_{2},\cdots ,V_{r};\mr)
\end{equation}

由于张量积运算$\otimes$关于每一个位置都是线性的，所以如上定义的$\bar{\varphi}$是$V^{\ast}_{1}\times \cdots\times V^{\ast}_{r}\longrightarrow L(V_{1}, V_{2},\cdots ,V_{r};\mr)$的一个$r$重线线性映射。更进一步，我们有如类似结论成立。
\begin{mythm}{}{chapter2另一种张映射}
	式\requ{equ:}定义的$r$重线性映射$\bar{\varphi}$是张映射.
\end{mythm}
\begin{proof}
	根据\rthm{thm:chapter2张映射的等价条件}，要证明$\bar{\varphi}$是张映射，只要证$\{\bar{\varphi}(e^{\ast}_{\gamma}) \}_{\gamma\in \Gamma}$线性无关. 其中约定$e^{\ast}_{\gamma}\triangleq (e^{\ast}_{1\gamma_{1}},e^{\ast}_{2\gamma_{2}},\cdots,e^{\ast}_{r\gamma_{r}})$. 设
	\begin{equation}
	\sum_{ \gamma\in \Gamma }k_{\gamma}\bar{\varphi}(e^{\ast}_{\gamma})=0
	\end{equation}
	其中$k_{\gamma}\in \mr$. 将上式两边作用在$e_{\gamma}$上，其中$e_{\gamma}\triangleq (e_{1\gamma_{1}},e_{2\gamma_{2}},\cdots,e_{r\gamma_{r}})\in V_{1}\times V_{2}\times\cdots V_{r}$. 直接计算
	\[
	\begin{split}
	k_{\alpha}\bar{\varphi}(e^{\ast}_{\alpha})(e_{\gamma})&=k_{\alpha}e^{\ast}_{1\alpha_{1}}\otimes \cdots\otimes e^{\ast}_{r\alpha_{r}}\left(e_{1\gamma_{1}},\cdots,e_{r\gamma_{r}} \right)\\
	&=k_{\alpha}e^{\ast}_{1\alpha_{ 1 }}\left(e_{1\gamma_{1}} \right)\cdots e^{\ast}_{r\alpha_{ r }}\left(e_{r\gamma_{r}} \right)\\
	&=k_{\alpha}\cdot \delta_{ \gamma_{1} \alpha_{ 1 } }\cdots \delta_{ \gamma_{r} \alpha_{ r } }
	\end{split}
	\]
	
	对于$\alpha\neq \gamma\in \Gamma$，则至少存在一个$\gamma_{ i }\neq\alpha_{ i }$，从而$\delta_{ \gamma_{i} \alpha_{ i } }=0\Rightarrow k_{\alpha}\bar{\varphi}(e^{\ast}_{\alpha})(e_{\gamma})=0$.
	
	对于$\alpha=\gamma\in\Gamma$，则$\gamma_{1}=\alpha_{ 1 },\cdots,\gamma_{r}=\alpha_{ r }$，从而$\delta_{ \gamma_{1} \alpha_{ 1 } }=1,\cdots, \delta_{ \gamma_{r} \alpha_{ r } }=1\Rightarrow k_{\gamma}\bar{\varphi}(e^{\ast}_{\gamma})(e_{\gamma})=k_{\gamma}$. 
	
	因此将\requ{equ:}两边作用在$e_{\gamma}$上后得到$k_{\gamma}=0,\forall \gamma\in \Gamma$. 这表明$\{\bar{\varphi}(e^{\ast}_{\gamma}) \}_{\gamma\in \Gamma}$线性无关，从而$\bar{\varphi}$是张映射。
\end{proof}

令$\bar{W}=\langle\operatorname{Im}(\bar{\varphi}) \rangle$就得到了对偶空间$V^{\ast}_{1},V^{\ast}_{2},\cdots,V^{\ast}_{r}$的一个具体的张量空间$\bar{W}$.
按照$\bar{\varphi}$的定义方式可知，$\operatorname{Im}(\bar{\varphi}) $中的元素是形如$v^{\ast}_{1}\otimes v^{\ast}_{2}\otimes\cdots\otimes v^{\ast}_{r}$这个样子的、作用在$V_{1}\times V_{2}\times\cdots V_{r}$上的$r$重线性函数，于是$\langle\operatorname{Im}(\bar{\varphi}) \rangle$中的元素都是它们的线性组合，自然也是作用在$V_{1}\times V_{2}\times\cdots V_{r}$上的$r$重线性函数。也就是说，$\bar{W}$中的元素都是一些作用在$V_{1}\times V_{2}\times\cdots V_{r}$上的$r$重线性函数！用数学符号来表示，即$\bar{W}\subset L(V_{1}, V_{2},\cdots ,V_{r};\mr)$. 事实上与前面类似，我们有
\begin{mythm}{}{chapter2反过来也对定理2}
	设$\bar{\varphi}$是由式\requ{equ:}定义的$r$重线性映射，则$\langle\operatorname{Im}(\bar{\varphi}) \rangle=L(V_{1}, V_{2},\cdots ,V_{r};\mr)$.
\end{mythm}
\begin{proof}
	由之前的讨论已经知道$\langle\operatorname{Im}(\bar{\varphi}) \rangle\subset L(V_{1}, V_{2},\cdots ,V_{r};\mr)$. 只需再说明
	\[
		L(V_{1}, V_{2},\cdots ,V_{r};\mr)\subset \langle\operatorname{Im}(\bar{\varphi}) \rangle
	\]
	为此任取$f\in L(V_{1},V_{2},\cdots,V_{r};\mr)$，则
	\begin{equation}
	\begin{split}
	f(v_{1},v_{2},\cdots,v_{r})&=f\left(\sum_{k_{1}=1}^{n_{1}}e^{\ast}_{1k_{1}}(v_{1})e_{1k_{1}},\sum_{k_{2}=1}^{n_{2}}e^{\ast}_{2k_{2}}(v_{2})e_{2k_{2}},\cdots,\sum_{k_{r}=1}^{n_{r}}e^{\ast}_{rk_{r}}(v_{r})e_{rk_{r}}\right)\\
	&=\sum_{k_{1}=1}^{n_{1}}\sum_{k_{2}=1}^{n_{2}}\cdots\sum_{k_{r}=1}^{n_{r}}e^{\ast}_{1k_{1}}(v_{1})e^{\ast}_{2k_{2}}(v_{2})\cdots e^{\ast}_{rk_{r}}(v_{r})f\left(e_{1k_{1}},e_{2k_{2}},\cdots,e_{rk_{r}} \right)\\
	&=\sum_{ \alpha\in \Gamma }e^{\ast}_{1\alpha_{1}}(v_{1})e^{\ast}_{2\alpha_{2}}(v_{2})\cdots e^{\ast}_{r\alpha_{r}}(v_{r})f\left(e_{1\alpha_{1}},e_{2\alpha_{2}},\cdots,e_{r\alpha_{r}} \right)\\
	&=\sum_{ \alpha\in \Gamma }e^{\ast}_{1\alpha_{1}}(v_{1})e^{\ast}_{2\alpha_{2}}(v_{2})\cdots e^{\ast}_{r\alpha_{r}}(v_{r})f(e_{\alpha})\\
	&=\sum_{ \alpha\in \Gamma }f(e_{\alpha})e^{\ast}_{1\alpha_{1}}(v_{1})e^{\ast}_{2\alpha_{2}}(v_{2})\cdots e^{\ast}_{r\alpha_{r}}(v_{r})\\
	&=\sum_{ \alpha\in \Gamma }\left(f(e_{\alpha})e^{\ast}_{1\alpha_{1}}\otimes e^{\ast}_{2\alpha_{2}}\otimes \cdots \otimes e^{\ast}_{r\alpha_{r}}\right)(v_{1},v_{2},\cdots,v_{r})\\
	&=\left(\sum_{ \alpha\in \Gamma }f(e_{\alpha})e^{\ast}_{1\alpha_{1}}\otimes e^{\ast}_{2\alpha_{2}}\otimes \cdots \otimes e^{\ast}_{r\alpha_{r}}\right)(v_{1},v_{2},\cdots,v_{r})
	\end{split}
	\end{equation}
	其中约定$(k_{1},k_{2},\cdots,k_{r})=\alpha\in \Gamma=\Gamma(n_{1},\cdots,n_{r})\;,e_{\alpha}\triangleq \left(e_{1\alpha_{1}},e_{2\alpha_{2}},\cdots,e_{r\alpha_{r}} \right)$. 
	
	于是从上式我们得出，$f=\sum_{ \alpha\in \Gamma }f(e_{\alpha})e^{\ast}_{1\alpha_{1}}\otimes e^{\ast}_{2\alpha_{2}}\otimes \cdots \otimes e^{\ast}_{r\alpha_{r}}$，从而$f\in \langle\operatorname{Im}(\bar{\varphi}) \rangle\Rightarrow L(V_{1}, V_{2},\cdots ,V_{r};\mr)\subset \langle\operatorname{Im}(\bar{\varphi}) \rangle$.
	
	综上所述，我们有$\langle\operatorname{Im}(\bar{\varphi}) \rangle=L(V_{1}, V_{2},\cdots ,V_{r};\mr)$.
\end{proof}

因此$V^{\ast}_{1}\otimes V^{\ast}_{2}\otimes \cdots\otimes V^{\ast}_{r}=\langle\operatorname{Im}(\bar{\varphi}) \rangle=L(V_{1}, V_{2},\cdots ,V_{r};\mr)$

\begin{lstlisting}[language={python}, caption={张量空间的构造}]
已知：$r$个向量空间$V^{\ast}_{1},V^{\ast}_{2},\cdots,V^{\ast}_{r}$
目标：构造$V^{\ast}_{1},V^{\ast}_{2},\cdots,V^{\ast}_{r}$的一个张量空间$\bar{W}$

for $\forall (v^{\ast}_{1},v^{\ast}_{2},\cdots,v^{\ast}_{r})\in V^{\ast}_{1}\times V^{\ast}_{2}\times\cdots\times V^{\ast}_{r}$，定义:
$\bar{\varphi}(v^{\ast}_{1},v^{\ast}_{2},\cdots,v^{\ast}_{r})\triangleq v^{\ast}_{1}\otimes v^{\ast}_{2}\otimes\cdots\otimes v^{\ast}_{r}\in L(V_{1}, V_{2},\cdots ,V_{r};\mr)$
则  $\langle\operatorname{Im}(\bar{\varphi}) \rangle=L(V_{1}, V_{2},\cdots ,V_{r};\mr)$

最后令$\bar{W}=L(V_{1}, V_{2},\cdots ,V_{r};\mr)$
\end{lstlisting}

最后我们引入一些新的记号，从而统一上述两种张量空间的构造。
\begin{mydef}{}{chapter2张量积运算}
	设$V_{1},V_{2},\cdots,V_{r}$为向量空间. 对任意$v_{1}\in V_{1},v_{2}\in V_{2},\cdots,v_{r}\in V_{r}$，定义向量的张量积运算如下：
	\[
	v_{1}\otimes v_{2}\otimes \cdots\otimes v_{r}\triangleq v^{\ast\ast}_{1}\otimes v^{\ast\ast}_{2}\otimes\cdots\otimes v^{\ast\ast}_{r}
	\]
	则$v_{1}\otimes v_{2}\otimes \cdots\otimes v_{r}\in L(V^{\ast}_{1}, V^{\ast}_{2},\cdots ,V^{\ast}_{r};\mr)$，且对$\forall \sigma_{i}\in V^{\ast}_{i},\;i=1,\cdots,r$有如下计算规则：
	\[
	\begin{split}
    &{\color{red}{v_{1}\otimes v_{2}\otimes \cdots\otimes v_{r}(\sigma_{1},\sigma_{2},\cdots,\sigma_{r})}}=v^{\ast\ast}_{1}\otimes v^{\ast\ast}_{2}\otimes\cdots\otimes v^{\ast\ast}_{r}(\sigma_{1},\sigma_{2},\cdots,\sigma_{r})\\
    &=v^{\ast\ast}_{1}(\sigma_{1})\cdot v^{\ast\ast}_{2}(\sigma_{2})\cdots v^{\ast\ast}_{r}(\sigma_{r})
    ={\color{red}{\sigma_{1}(v_{1})\cdot \sigma_{2}(v_{2})\cdots \sigma_{r}(v_{r})}}
	\end{split}
	\]
\end{mydef}
这个定义不管在形式上还是在本质上都是有意义的：因为$v^{\ast\ast}_{i}=v_{i}$. 特别地，对于对偶空间$V^{\ast}_{1},V^{\ast}_{2},\cdots,V^{\ast}_{r}$，其中的元素$(v^{\ast}_{i})^{\ast\ast}=((v^{\ast}_{i})^{\ast})^{\ast}=(v^{\ast\ast})^{\ast}_{i}=v^{\ast}_{i}$，从而${\color{red}{v^{\ast}_{1}\otimes v^{\ast}_{2}\otimes\cdots\otimes v^{\ast}_{r}=v^{\ast}_{1}\otimes v^{\ast}_{2}\otimes\cdots\otimes v^{\ast}_{r}}}$. 用更通俗(但不严谨)的话来说，你在做张量积运算时，如果参与运算的对象$v$来自一般向量空间，即$v\in V$，那你就用该元素的二次对偶$v^{\ast\ast}$来做运算；如果参与运算的对象$v^{\prime}$本身就取自某个对偶向量空间，即$v^{\prime}\in V^{\prime\ast}$，那就用它本身$v^{\prime}$来做运算即可。

举个例子，设$V$是一个向量空间，$V^{\ast}$为其对偶空间，$v\in V,w\in V^{\ast}$，那么符号$v\otimes w$就表示张量积运算$v^{\ast\ast}\otimes w$，这是一个$V^{\ast}\times V\longrightarrow \mr$的2重线性函数，即$v\otimes w\triangleq v^{\ast\ast}\otimes w\in L(V^{\ast},V;\mr)$. 

有了上述记号，$\varphi(v_{1},v_{2},\cdots,v_{r})= v^{\ast\ast}_{1}\otimes v^{\ast\ast}_{2}\otimes\cdots\otimes v^{\ast\ast}_{r}$就可以重新简记为
\[\varphi(v_{1},v_{2},\cdots,v_{r})=v_{1}\otimes v_{2}\otimes \cdots\otimes v_{r}\]
这就与$\bar{\varphi}(v^{\ast}_{1},v^{\ast}_{2},\cdots,v^{\ast}_{r})=v^{\ast}_{1}\otimes v^{\ast}_{2}\otimes \cdots\otimes v^{\ast}_{r}$\textbf{在形式}上统一起来了！

甚至张映射$\varphi$与张映射$\bar{\varphi}$的记号也可以统一，不用区分了，因为在上述记号的约定下它们的作用无非就是把元素用符号“$\otimes$”连接起来，至于元素用“$\otimes$”连接起来之后所表示的真实含义是什么，只需按照\rdef{def:chapter2张量积运算}来理解即可。因此我们以后把$\varphi$与$\bar{\varphi}$都统一用符号“$\otimes$”来表示。即
\[
\otimes :V_{1}\times V_{2}\times\cdots\times V_{r}\longrightarrow W
\]
其中$W=\langle \operatorname{Im}(\otimes)\rangle=L(V^{\ast}_{1},V^{\ast}_{2},\cdots,V^{\ast}_{r};\mr)$. 同理
\[
\otimes :V^{\ast}_{1}\times V^{\ast}_{2}\times\cdots\times V^{\ast}_{r}\longrightarrow \bar{W}
\]
其中$\bar{W}=\langle \operatorname{Im}(\otimes)\rangle=L(V_{1},V_{2},\cdots,V_{r};\mr)$. 

这种形式上的统一使我们以后不必再区分向量空间$V_{1},V_{2},\cdots,V_{r}$中哪些是一般的向量空间，哪些是对偶空间，反正张映射作用后得到的结果都可以写成$v_{1}\otimes v_{2}\otimes \cdots\otimes v_{r}$这个样子，其中$v_{i}\in V_{i},i=1,\cdots,r$. 举个例子，设$V$是一个向量空间，$V^{\ast}$为其对偶空间. 按照之前描述的构造方式我们可以得到张映射$\otimes :V\times V^{\ast}\times V\longrightarrow W$，其中对$\forall v_{1}\in V,v_{2}\in V^{\ast},v_{3}\in V$，有$\otimes (v_{1},v_{2},v_{3})=v_{1}\otimes v_{2}\otimes v_{3}$. 你看得到的结果在形式上非常简洁统一，但它所表达的真正含义是$v_{1}\otimes v_{2}\otimes v_{3}=v^{\ast\ast}_{1}\otimes v_{2}\otimes v^{\ast\ast}_{3}$. 

最后，我们把张量空间$\langle\operatorname{Im}(\otimes ) \rangle$也用时髦的新记号重新记为
\[
\langle\operatorname{Im}(\otimes) \rangle\triangleq V_{1}\otimes V_{2}\otimes\cdots\otimes V_{r}=\otimes_{i=1}^{r}V_{i}
\] 

\begin{remark}
	新记号有好处也有坏处，比如经常会让初学者误以为张量空间$V_{1}\otimes V_{2}\otimes\cdots\otimes V_{r}$中的元素都是形如$v_{1}\otimes v_{2}\otimes \cdots\otimes v_{r}$，实际上它里面的元素是形如$v_{1}\otimes v_{2}\otimes \cdots\otimes v_{r}$这样的元素的线性组合！因为$V_{1}\otimes V_{2}\otimes\cdots\otimes V_{r}=\langle\operatorname{Im}(\otimes) \rangle$.
\end{remark}

下面我们研究张量空间之间的线性映射。设$V_{1},\cdots,V_{m},W_{1},\cdots,W_{m}$都是向量空间，\textbf{已知条件}是给定$m$个线性映射$T_{i}:V_{i}\longrightarrow W_{i}$. \textbf{我们的目标}是构造出线性映射$T:V_{1}\otimes\cdots\otimes V_{m}\longrightarrow W_{1}\otimes\cdots\otimes W_{m}$. \textbf{我们的工具}是唯一因子性质，即以下交换图
\[
\begin{tikzcd}
V_{1}\times V_{2}\times\cdots\times V_{m} \arrow[r, "\otimes"] \arrow[d, "f", left ] & V_{1}\otimes V_{2}\otimes\cdots\otimes V_{m} \arrow[dl,dashrightarrow, "{\color{red}T}"]\\
Z  &
\end{tikzcd} 
\]
把向量空间$Z$取作$W_{1}\otimes\cdots\otimes W_{m}$，再定义$f$如下：对任意$(v_{1},v_{2},\cdots,v_{m})\in V_{1}\times V_{2}\times\cdots\times V_{m}$，令
\[
f(v_{1},v_{2},\cdots,v_{m})\triangleq T_{1}(v_{1})\otimes T_{2}(v_{2})\otimes\cdots\otimes T_{m}(v_{m})\quad
\]
可以验证如上定义的映射$f:V_{1}\times V_{2}\times\cdots\times V_{m}\longrightarrow W_{1}\otimes\cdots\otimes W_{m}$是$m$重线性映射。

根据唯一因子性质可知，存在唯一的线性映射$T:V_{1}\otimes\cdots\otimes V_{m}\longrightarrow W_{1}\otimes\cdots\otimes W_{m}$，满足$f=T\circ\otimes$. 我们把$T$称为由线性映射$T_{1},T_{2},\cdots,T_{m}$所诱导的线性映射。

\begin{mythm}{}{chapter2诱导的多重线性映射}
	设$T_{i}:V_{i}\longrightarrow W_{i}$是向量空间之间的线性映射$(1\leqslant i\leqslant m)$. 若
	$T:V_{1}\otimes\cdots\otimes V_{m}\longrightarrow W_{1}\otimes\cdots\otimes W_{m}$是由$T_{1},T_{2},\cdots,T_{m}$所诱导的线性映射，则对任意$x\in V_{1}\otimes\cdots\otimes V_{m}$，有
	\[
	T(x)\left(\sigma_{1},\sigma_{2},\cdots,\sigma_{m} \right)=x\left(T^{\ast}_{1}(\sigma_{1}),T^{\ast}_{2}(\sigma_{2}),\cdots,T^{\ast}_{m}(\sigma_{m}) \right)\;,\;\forall \sigma_{i}\in W^{\ast}_{i},1\leqslant i\leqslant m
	\]
	其中$T_{i}^{\ast}:W^{\ast}_{i}\longrightarrow V^{\ast}_{i}$是$T_{i}:V_{i}\longrightarrow W_{i}$所诱导的对偶空间之间的线性映射. 
\end{mythm}

\begin{proof}
	设$x=\sum\limits_{ \alpha }x_{\alpha}v_{\alpha}$，其中$x_{\alpha}\in \mr$，$v_{\alpha}=v_{\alpha_{1}}\otimes\cdots\otimes v_{\alpha_{ m }}\;,v_{\alpha_{ i }}\in V_{i},1\leqslant i\leqslant m$. 根据线性映射$T$的定义，则
	\[
	T(x)=\sum\limits_{ \alpha }x_{\alpha}T\left(v_{\alpha}\right)=\sum\limits_{ \alpha }x_{\alpha}f(v_{\alpha_{1}},\cdots, v_{\alpha_{ m }})=\sum\limits_{ \alpha }x_{\alpha}T_{1}(v_{\alpha_{ 1 }})\otimes\cdots\otimes T_{m}(v_{\alpha_{ m }})
	\]
	于是对任意$\left(\sigma_{1},\sigma_{2},\cdots,\sigma_{m} \right)\in W^{\ast}_{1}\otimes\cdots\otimes W^{\ast}_{m}$，便有
	\[
	\begin{split}
	&T(x)\left(\sigma_{1},\sigma_{2},\cdots,\sigma_{m} \right)\\
	&=\sum\limits_{ \alpha }x_{\alpha}\left[T_{1}(v_{\alpha_{ 1 }})\otimes\cdots\otimes T_{m}(v_{\alpha_{ m }})\left(\sigma_{1},\sigma_{2},\cdots,\sigma_{m} \right)\right]\\
	&=\sum\limits_{ \alpha }x_{\alpha}\left[\sigma_{1}\left(T_{1}(v_{\alpha_{ 1 }})\right)\cdots\sigma_{m}\left(T_{m}(v_{\alpha_{ m }})\right)\right]\\
	\end{split}
	\]
	再根据$T_{i}^{\ast}:W^{\ast}_{i}\longrightarrow V^{\ast}_{i}$的定义可知，$\sigma_{i}\left(T_{i}(v_{\alpha_{i }})\right)=T^{\ast}_{i}(\sigma_{i})(v_{\alpha_{ 1 }})$. 因此上式又可以写为
	\[
	\begin{split}
	&T(x)\left(\sigma_{1},\sigma_{2},\cdots,\sigma_{m} \right)\\
	&=\sum\limits_{ \alpha }x_{\alpha}\left[T^{\ast}_{1}(\sigma_{1})(v_{\alpha_{ 1 }})\cdots T^{\ast}_{m}(\sigma_{m})(v_{\alpha_{ m }})\right]\\
	&=\sum\limits_{ \alpha }x_{\alpha}\left[v_{\alpha_{1}}\otimes\cdots\otimes v_{\alpha_{ m }}\left(T^{\ast}_{1}(\sigma_{1}),\cdots,T^{\ast}_{m}(\sigma_{m}) \right) \right]\\
	&=x\left(T^{\ast}_{1}(\sigma_{1}),\cdots,T^{\ast}_{m}(\sigma_{m}) \right)
	\end{split}
	\]
\end{proof}

下面我们给出线性映射$T:V_{1}\otimes\cdots\otimes V_{m}\longrightarrow W_{1}\otimes\cdots\otimes W_{m}$在有序基下的矩阵形式。为此先引入一个新的矩阵运算。
\begin{mydef}{矩阵的Kronecker积}{chapter2矩阵的Kronecker积}
	设$A$为一个$m\times n$的实矩阵，$B$为一个$p\times q$的实矩阵. 令
	\[
	A\otimes B\triangleq \left( \begin{array}{cccc}{a_{11}B} & {a_{12}B} & {\cdots} & {a_{1 n}B} \\ {a_{21}B} & {a_{22}B} & {\cdots} & {a_{2 n}B} \\ {\vdots} & {\vdots} & { } & {\vdots} \\ {a_{m 1}B} & {a_{m 2}B} & {\cdots} & {a_{m n}B}\end{array}\right)
	\]
    得到的结果是一个$m\times p$行、$n\times q$列的实矩阵. 我们称$A\otimes B$为矩阵$A$与$B$的Kronecker积.
\end{mydef}

矩阵的Kronecker积有一些基本的运算性质。
\begin{mythm}{}{chapter2Kronecker积的性质}
	\begin{itemize}
		\item[(1)] Kronecker积不满足交换律，即$A\otimes B\neq B\otimes A$.
		\item[(2)] Kronecker积满足结合律，即$(A\otimes B)\otimes C=A\otimes (B\otimes C)$.
		\item[(3)] $(A\otimes B)^{T}=A^{T}\otimes B^{T}$.
		\item[(4)] $\forall\lambda\in \mr$，$\lambda(A\otimes B)=A\otimes (\lambda B)=(\lambda A)\otimes B$.
		\item[(5)] 若矩阵$A,B$可以做加法，则$(A+B)\otimes C=A\otimes C+B\otimes C,C\otimes (A+B)=C\otimes A+C\otimes B$.
		\item[(6)] 若矩阵$A,C$可以做乘法，$B,D$可以做乘法，则$(A\otimes B) (C\otimes D)=(AC)\otimes(BD)$.
		\item[(7)] 若$A,B$都可逆，则$A\otimes B$也可逆，且$(A\otimes B)^{-1}=A^{-1}\otimes B^{-1}$.
	\end{itemize}
\end{mythm}

设$A_{i}$是$k_{i}\times n_{i}$的实矩阵，则$A=A_{1}\otimes A_{2}\otimes \cdots\otimes A_{m}$是一个$k_{1}\times k_{2}\times\cdots\times k_{m}$行、$n_{1}\times n_{2}\times \cdots\times n_{m} $列的实矩阵。利用多重指标的记号，我们可以方便的计算矩阵Kronecker积中某一行某一列的元素。将$A$的行数$\{1,2,\cdots,k_{1}\times k_{2}\times\cdots\times k_{m} \}$用多重指标集$\Gamma(k_{1},\cdots,k_{m})$来表示，将$A$的列数$\{1,2,\cdots,n_{1}\times n_{2}\times\cdots\times n_{m} \}$用多重指标集$\Gamma(n_{1},\cdots,n_{m})$来表示，其中$\Gamma(k_{1},\cdots,k_{m})$与$\Gamma(n_{1},\cdots,n_{m})$中的元素都按照字典顺序排列。
\begin{mythm}{}{chapter2Kronecker积中元素的计算}
	若$A=A_{1}\otimes A_{2}\otimes \cdots\otimes A_{m}$，则对$\forall\alpha\in \Gamma(k_{1},\cdots,k_{m})$，$\forall\beta\in \Gamma(n_{1},\cdots,n_{m})$，矩阵$A$的第$\alpha$行第$\beta$列的元素$A_{\alpha\beta}$为
	\[
	A_{\alpha\beta}=\prod_{i=1}^{m}a_{\alpha_{ i }\beta_{ i }}
	\]
	其中$a_{\alpha_{ \color{red}{i} }\beta_{\color{red}{ i} }}$表示矩阵$A_{\color{red}{i}}$的第$\alpha_{ i }$行第$\beta_{ i }$列的元素.
\end{mythm}
这个结论就不证明了，我们举个例子验证一下。
\begin{example}
	\label{chapter2确定Kronecker积元素的例子}
	设
	\[
	A=\left( \begin{array}{cc}{a_{11}} & {a_{12}}  \\ {a_{21}B} & {a_{22}} \end{array}\right)\quad B=\left( \begin{array}{ccc}{b_{11}} & {b_{12}} & {b_{13}} \\ {b_{21}B} & {b_{22}} & {b_{23}} \end{array}\right)
	\]
	则
	\[
	A\otimes B=\left( \begin{array}{cccccc}{a_{11}b_{11}} & {a_{11}b_{12}} & {a_{11}b_{13}} & {a_{12}b_{11}} &{a_{12}b_{12}}& {a_{12}b_{13}}\\ {a_{11}b_{21}} & {a_{11}b_{22}} & {a_{11}b_{23}} & {a_{12}b_{21}} &{a_{12}b_{22}}& {a_{12}b_{23}}\\{a_{21}b_{11}} & {a_{21}b_{12}} & {a_{21}b_{13}} & {a_{22}b_{11}} &{a_{22}b_{12}}& {a_{22}b_{13}}\\ {a_{21}b_{21}} & {a_{21}b_{22}} & {a_{21}b_{23}} & {a_{22}b_{21}} &{a_{22}b_{22}}& {a_{22}b_{23}}
	\end{array}\right)
	\]
	将行数与$\Gamma(2,2)$对应，即
	\[
	\begin{array}{cccccccc}\{{1,} & {2}, & {3}, & {4}\} \\{\updownarrow} & {\updownarrow} & {\updownarrow} &{\updownarrow}\\ \{{(1,1)}, & {(1,2)}, & {(2,1)}, & {(2,2)}\}\end{array}
	\]
	将列数与$\Gamma(2,3)$对应，即
	\[
	\begin{array}{cccccccc}\{{1,} & {2}, & {3}, & {4} &{5}&{6}  \} \\{\updownarrow} & {\updownarrow} & {\updownarrow} &{\updownarrow}&{\updownarrow}&{\updownarrow}\\ \{{(1,1)}, & {(1,2)},& {(1,3)}, & {(2,1)}, & {(2,2)}, & {(2,3)}\}\end{array}
	\]
	
	我们来看$A\otimes B$的第$\alpha=(1,2)$行第$\beta=(2,1)$列的元素(即第2行第4列的元素)，是$a_{12}b_{21}=a_{\alpha_{1}\beta_{ 1 }}b_{\alpha_{ 2 }\beta_{ 2 }}$. 
\end{example}

通过这个具体的例子读者应该清楚上述定理所表达的意思。下面我们就用矩阵的Kronecker积来正式给出线性映射$T:V_{1}\otimes\cdots\otimes V_{m}\longrightarrow W_{1}\otimes\cdots\otimes W_{m}$在基下的矩阵表达形式。

设$A_{i}=\{e_{i1},e_{i2},\cdots,e_{in_{i}} \}$为向量空间$V_{i}$的有序基，其中$\operatorname{dim}(V_{i})=n_{i}\;,1\leqslant i\leqslant m$. 再设$B_{i}=\{f_{i1},f_{i2},\cdots,f_{ik_{i}} \}$为向量空间$W_{i}$的有序基，其中$\operatorname{dim}(W_{i})=k_{i}\;,1\leqslant i\leqslant m$. 则
\[
A=\{e^{\otimes}_{\alpha}=e_{1\alpha_{ 1 }}\otimes e_{2\alpha_{ 2 }}\otimes\cdots\otimes e_{m\alpha_{ m }} \}_{\alpha\in \Gamma(n_{1},\cdots,n_{m})}
\]
就是$V_{1}\otimes \cdots\otimes V_{m}$的一组有序基(按字典序排列)；同样的，
\[
B=\{f^{\otimes}_{\beta}=f_{1\beta_{ 1 }}\otimes f_{2\beta_{ 2 }}\otimes\cdots\otimes f_{m\beta_{ m }} \}_{\beta\in \Gamma(k_{1},\cdots,k_{m})}
\]
就是$W_{1}\otimes \cdots\otimes W_{m}$的一组有序基(按字典序排列)。

\begin{mythm}{}{chapter2诱导多重映射的矩阵表示}
	设线性映射$T_{i}:V_{i}\longrightarrow W_{i}$在有序基$A_{i}$和$B_{i}$下的矩阵为$[T_{i}]^{B_{i}}_{A_{i}}$，诱导线性映射$T:V_{1}\otimes\cdots\otimes V_{m}\longrightarrow W_{1}\otimes\cdots\otimes W_{m}$在有序基$A$和$B$下的矩阵为$[T]_{A}^{B}$. 则有
	\[
	[T]_{A}^{B}=[T_{i}]^{B_{i}}_{A_{i}}\otimes [T_{2}]^{B_{2}}_{A_{2}}\otimes \cdots\otimes[T_{m}]^{B_{m}}_{A_{m}}
	\]
\end{mythm}

\begin{mythm}{}{chapter2张量空间的对偶空间}
$(V_{1}\otimes V_{2}\otimes \cdots\otimes V_{r})^{\ast}\cong V_{1}^{\ast}\otimes V_{2}^{\ast}\otimes \cdots\otimes V_{r}^{\ast}$
\end{mythm}
\begin{proof}
首先注意到$V_{1}^{\ast}\otimes V_{2}^{\ast}\otimes \cdots\otimes V_{r}^{\ast}=L(V_{1},V_{2}, \cdots, V_{r};\mr)$，因此任取$g\in V_{1}^{\ast}\otimes V_{2}^{\ast}\otimes \cdots\otimes V_{r}^{\ast}$，则$g\in L(V_{1},V_{2}, \cdots, V_{r};\mr)$. 根据唯一因子性质(\rthm{thm:chapter2唯一因子化性质})可知，存在唯一的线性映射$f:V_{1}\otimes V_{2}\otimes \cdots\otimes V_{r}\longrightarrow \mr$，使得$g=f\circ \otimes$. 也就是有以下交换图成立。
\[
	\begin{tikzcd}
	V_{1}\times V_{2}\times\cdots\times V_{r} \arrow[r, "\otimes"] \arrow[d, "g", left ] & V_{1}\otimes V_{2}\otimes \cdots\otimes V_{r} \arrow[dl, "f=\iota(g)"]\\
	\mr &
	\end{tikzcd} 
	\]
	注意到$f$是从$V_{1}\otimes V_{2}\otimes \cdots\otimes V_{r}\longrightarrow \mr$的线性映射，从而$f\in (V_{1}\otimes V_{2}\otimes \cdots\otimes V_{r})^{\ast}$. 
	
	于是在$V_{1}^{\ast}\otimes V_{2}^{\ast}\otimes \cdots\otimes V_{r}^{\ast}$与$(V_{1}\otimes V_{2}\otimes \cdots\otimes V_{r})^{\ast}$之间我们可以建立如下一一对应$\iota:V_{1}^{\ast}\otimes V_{2}^{\ast}\otimes \cdots\otimes V_{r}^{\ast}\longrightarrow (V_{1}\otimes V_{2}\otimes \cdots\otimes V_{r})^{\ast}$，其定义为
	\begin{equation}
		 \iota(g)=f\;,\; \forall g\in V_{1}^{\ast}\otimes V_{2}^{\ast}\otimes \cdots\otimes V_{r}^{\ast}
		 \label{张量空间及其对偶空间的同构映射}
	\end{equation}
	容易说明$\iota$是一个线性同构映射，因此$(V_{1}\otimes V_{2}\otimes \cdots\otimes V_{r})^{\ast}\cong V_{1}^{\ast}\otimes V_{2}^{\ast}\otimes \cdots\otimes V_{r}^{\ast}$.
\end{proof}

我们在\ref{配合}节讲过，配合运算可以等价地描述形如“$V\cong W^{\ast}$”这样的事实。下面我们就从配合运算的角度来重新看待\rthm{thm:chapter2张量空间的对偶空间}，为此只需给$V_{1}^{\ast}\otimes V_{2}^{\ast}\otimes \cdots\otimes V_{r}^{\ast}$与$V_{1}\otimes V_{2}\otimes \cdots\otimes V_{r}$规定一个非退化的配合运算，而这由\rthm{thm:chapter2同构则配合}结合式\requ{张量空间及其对偶空间的同构映射}很容易做到。

\begin{mydef}{$V_{1}^{\ast}\otimes  \cdots\otimes V_{r}^{\ast}$与$V_{1}\otimes \cdots\otimes V_{r}$的配合}{chapter2张量空间的配合运算}
对于$\forall g\in V_{1}^{\ast}\otimes V_{2}^{\ast}\otimes \cdots\otimes V_{r}^{\ast} $，$\forall x\in V_{1}\otimes V_{2}\otimes \cdots\otimes V_{r}$，我们约定
\[
\langle g,x\rangle\triangleq \langle{\color{red}{\iota(g)}},x\rangle
\]
其中$\iota:V_{1}^{\ast}\otimes \cdots\otimes V_{r}^{\ast}\longrightarrow (V_{1}\otimes \cdots\otimes V_{r})^{\ast}$是由式(\ref{张量空间及其对偶空间的同构映射})所定义的同构映射.
\end{mydef}

注意到$\langle{\color{red}{\iota(g)}},x\rangle$中，$\iota(g)\in (V_{1}\otimes \cdots\otimes V_{r})^{\ast}$，所以$\langle{\color{red}{\iota(g)}},x\rangle$表示对偶配合，即$\langle{\color{red}{\iota(g)}},x\rangle={\color{red}{\iota(g)}}(x)$.
我们来看上述定义的配合运算具体是如何运作的。

若$x=v_{1}\otimes v_{2}\otimes \cdots\otimes v_{r}$，则
\begin{equation}
\begin{split}
	\langle g,x\rangle=\langle{\color{red}{\iota(g)}},x\rangle&={\color{red}{\iota(g)}}(v_{1}\otimes v_{2}\otimes \cdots\otimes v_{r})\\
	&=\left({\color{red}{\iota(g)}}\circ \otimes\right)  (v_{1}, v_{2}, \cdots , v_{r})\\
	&=g(v_{1}, v_{2}, \cdots , v_{r})
\end{split}
\label{配合运算1}
\end{equation}
最后一个等号是因为唯一因子性质。一般地，设$x=\sum\limits_{\gamma}x_{\gamma}\cdot  v_{\gamma_{1}}\otimes v_{\gamma_{2}}\otimes\cdots\otimes v_{\gamma_{r}}$，则
\begin{equation}
\begin{split}
	\langle g,x\rangle=\langle{\color{red}{\iota(g)}},x\rangle&={\color{red}{\iota(g)}}\left(\sum_{\gamma}x_{\gamma} \cdot v_{\gamma_{1}}\otimes v_{\gamma_{2}}\otimes\cdots\otimes v_{\gamma_{r}} \right)\\
&=\sum_{\gamma}x_{\gamma} \cdot \left[{\color{red}{\iota(g)}}\left(v_{\gamma_{1}}\otimes v_{\gamma_{2}}\otimes\cdots\otimes v_{\gamma_{r}}\right)\right]\\
&=\sum_{\gamma}x_{\gamma}\cdot  g\left(v_{\gamma_{1}},v_{\gamma_{2}},\cdots, v_{\gamma_{r}}\right)
\end{split}
\label{配合运算2}
\end{equation}

有些书上在证明\rthm{thm:chapter2张量空间的对偶空间}时，扔下如上定义的配合运算就完事了，初学者往往会对此感到十分困惑。在学过\ref{配合}节后，读者现在应该能理解这背后的原理：因为一旦给出了配合运算，也就等价于证明了同构，再无需多言。还有些书上会把$(V_{1}\otimes V_{2}\otimes \cdots\otimes V_{r})^{\ast}\cong V_{1}^{\ast}\otimes V_{2}^{\ast}\otimes \cdots\otimes V_{r}^{\ast}$直接写成$(V_{1}\otimes V_{2}\otimes \cdots\otimes V_{r})^{\ast}= V_{1}^{\ast}\otimes V_{2}^{\ast}\otimes \cdots\otimes V_{r}^{\ast}$，这可能也会给初学者带来一些困扰，下面我们来解释一下。

设$\{e_{i1},e_{i2},\cdots,e_{in_{i}} \}$为向量空间$V_{i}$的有序基，$\{e^{\ast}_{i1},e^{\ast}_{i2},\cdots,e^{\ast}_{in_{i}} \}$为$V^{\ast}_{i}$中相应的对偶基，其中$n_{i}=\operatorname{dim}(V_{i}),1\leqslant i\leqslant r$. 则
\[
\{e_{1\alpha_{ 1 }}\otimes e_{2\alpha_{ 2 }}\otimes\cdots\otimes e_{r\alpha_{ r }}|\alpha\in \Gamma(n_{1},\cdots,n_{r}) \}\quad\quad \{e^{\ast}_{1\beta_{ 1 }}\otimes e^{\ast}_{2\beta_{ 2 }}\otimes\cdots\otimes e^{\ast}_{r\beta_{ r }} |\beta\in \Gamma(n_{1},\cdots,n_{r})\}
\]
分别构成$V_{1}\otimes \cdots\otimes V_{r}$与$V^{\ast}_{1}\otimes V^{\ast}_{2}\otimes\cdots\otimes V^{\ast}_{r}$的有序基(按字典顺序排列)。于是我们有
\begin{equation}
\begin{split}
	&\langle e^{\ast}_{1\beta_{ 1 }}\otimes e^{\ast}_{2\beta_{ 2 }}\otimes\cdots\otimes e^{\ast}_{r\beta_{ r }},e_{1\alpha_{ 1 }}\otimes e_{2\alpha_{ 2 }}\otimes\cdots\otimes e_{r\alpha_{ r }}\rangle\\
	&=e^{\ast}_{1\beta_{ 1 }}\otimes e^{\ast}_{2\beta_{ 2 }}\otimes\cdots\otimes e^{\ast}_{r\beta_{ r }}\left(e_{1\alpha_{ 1 }}, e_{2\alpha_{ 2 }},\cdots, e_{r\alpha_{ r }} \right)\\
	&=e^{\ast}_{1\beta_{ 1 }}(e_{1\alpha_{ 1 }})\cdot e^{\ast}_{2\beta_{ 2 }}(e_{2\alpha_{ 2 }})\cdots e^{\ast}_{r\beta_{ r }}(e_{r\alpha_{ r}})\\
	&=\delta^{\beta_{1}}_{\alpha_{1}}\cdot \delta^{\beta_{2}}_{\alpha_{2}}\cdots\delta^{\beta_{r}}_{\alpha_{r}}\\
	&=\begin{cases}
		1 ,&\beta=\alpha\\
		0 ,&\beta\neq \alpha
	  \end{cases}
\end{split}
\label{对偶配合性质}
\end{equation}

上式表明某种程度上，我们给$V_{1}^{\ast}\otimes  \cdots\otimes V_{r}^{\ast}$与$V_{1}\otimes \cdots\otimes V_{r}$赋予的这种配合运算可以看成是对偶配合。因此在指定了该配合运算后，你写成$(V_{1}\otimes V_{2}\otimes \cdots\otimes V_{r})^{\ast}= V_{1}^{\ast}\otimes V_{2}^{\ast}\otimes \cdots\otimes V_{r}^{\ast}$也没毛病。

当然啦，到底该看成是同构还是相等，这完全取决于个人喜好，{\textbf{只要你自己拎得清}}。我个人更喜欢同构的说法，因为同构才是本质，相等只是一种视角而已。

\subsection{$(r,s)$型张量}
\label{rs型张量}

上一节考虑的是一般的张量空间，不过很多时候我们用不着那么一般的情形。在几何或物理中所用到的张量及张量空间经常是由\textbf{一个}向量空间及其对偶空间$V$衍生来的，比如在几何中这个向量空间通常就是切空间。下面具体说明。

设$V$是一个$n$维实向量空间，$V^{\ast}$为其对偶空间. 按照上一节构造张量空间的方法，我们可以得到如下张量空间
\[
V^{r}_{s}\triangleq \underbrace{V\otimes \cdots\otimes V}_{r\text{个}}\otimes \underbrace{V^{\ast}\otimes \cdots\otimes V^{\ast}}_{s\text{个}}=\langle\operatorname{Im}(\otimes) \rangle
\]
这个张量空间中的元素都是形如$v_{1}\otimes \cdots \otimes v_{r}\otimes\theta_{1}\otimes\cdots\otimes \theta_{s}$的线性组合，其中$v_{i}\in V_{i},1\leqslant i\leqslant r,\;\theta_{j}\in V^{\ast},1\leqslant j\leqslant s$. 
根据上一节最后约定的记号，我们知道$v_{1}\otimes \cdots \otimes v_{r}\otimes\theta_{1}\otimes\cdots\otimes \theta_{s}$表达的真正含义是
\[
v_{1}\otimes \cdots \otimes v_{r}\otimes\theta_{1}\otimes\cdots\otimes \theta_{s}=v^{\ast\ast}_{1}\otimes \cdots \otimes v^{\ast\ast}_{r}\otimes\theta_{1}\otimes\cdots\otimes \theta_{s}
\]
因此
\[
v_{1}\otimes \cdots \otimes v_{r}\otimes\theta_{1}\otimes\cdots\otimes \theta_{s}\in L(\underbrace{V^\ast,\cdots,V^{\ast}}_{r\text{个}},\underbrace{V,\cdots,V}_{s\text{个}};\mr)
\]
上一节的结论可以全部照搬过来，比如$V^{r}_{s}=L(V^\ast,\cdots,V^{\ast},V,\cdots,V;\mr)$，$V^{r}_{s}$的维数是$n^{r+s}$. 另外$V^{r}_{s}$的基当然也可以由上一节的方法得到，具体如下：

设$\{e_{1},e_{2},\cdots,e_{n} \}$是$V$的一组基，$\{e^{\ast}_{1},e^{\ast}_{2} ,\cdots,e^{\ast}_{n}\}$为$V^{\ast}$中相应的对偶基，令
\[
	e^{\otimes}_{\gamma}\triangleq\otimes(e_{\gamma})= e_{\gamma_{1}}\otimes\cdots\otimes e_{\gamma_{r}}\otimes e^{\ast}_{\gamma_{ r+1 }} \otimes\cdots\otimes e^{\ast}_{\gamma_{ r+s }} 
\]
则$\{e^{\otimes}_{\gamma} \}_{\gamma\in \Gamma^{r}_{s}}$就构成了$V^{r}_{s}$的一组基. 其中$e_{\gamma}\triangleq (e_{\gamma_{1}},\cdots, e_{\gamma_{r}}, e^{\ast}_{\gamma_{ r+1 }} ,\cdots, e^{\ast}_{\gamma_{ r+s }})$，多重指标集$\Gamma^{r}_{s}$定义为
\[
\Gamma^{r}_{s}=\{\gamma=(\gamma_{1},\cdots,\gamma_{ r+s })  | 1\leqslant \gamma_{ i }\leqslant n,i=1,2,\cdots,n\}
\]

但是上面的写法有点臃肿，要注意到$V^{r}_{s}=V\otimes \cdots\otimes V\otimes V^{\ast}\otimes \cdots\otimes V^{\ast}$只涉及到两种空间$V$和$V^{\ast}$，而且前面$r$个都是$V$，后面$s$个都是$V^{\ast}$. {\color{red}{所以在$V^{r}_{s}$的情景下，我们可以把多重指标设计得更精巧些}}。我们约定
\[
\begin{array}{cccccccc}({\gamma_{1},} & {\gamma_{2}}, & {\cdots}, & {\gamma_{r}},& {\gamma_{r+1}},& {\gamma_{r+2}},& {\cdots},& {\gamma_{r+s}}) \\{\updownarrow} & {\updownarrow} & { \cdots} & {\updownarrow}& {\updownarrow}& {\updownarrow}& {\cdots }& {\updownarrow} \\ ({\overline{\gamma_{1}}}, & {\overline{\gamma_{2}}}, & {\cdots}, & {\overline{\gamma_{r}}},& {\underline{\gamma_{1}}},& {\underline{\gamma_{2}}},& {\cdots},& {\underline{\gamma_{s}}})\end{array}
\]
因此以后我们把多重指标写成这个样子：$\gamma=\left(\boxed{\overline{\gamma_{1}},\cdots,\overline{\gamma_{r}}},\boxed{\underline{\gamma_{1}},\cdots,\underline{\gamma_{s}}} \right)=\left(\overline{\gamma},\underline{\gamma} \right)$. 也就是说
\[
\gamma_{ i }= \left\{\begin{array}{l}\overline{\gamma_{i}} \\ \underline{\gamma_{i-r}}\end{array}\right.\;\;,\;\forall i=1,2,\cdots,r+s
\]
这样写纯粹为了契合$V^{r}_{s}=V\otimes \cdots\otimes V\otimes V^{\ast}\otimes \cdots\otimes V^{\ast}$. 从而使得指标中用来索引$V\otimes\cdots\otimes V$的前$r$个分量与用来索引$V^{\ast}\otimes\cdots\otimes V^{\ast}$的后$s$个分量更加简洁！

在新记号下，$V^{r}_{s}$的基$\{e^{\otimes}_{\gamma} \}_{\gamma\in \Gamma^{r}_{s}}$可以表示为
\[
e^{\otimes}_{\gamma}\triangleq\otimes(e_{\gamma})= e_{\overline{\gamma_{1}}}\otimes\cdots\otimes e_{\overline{\gamma_{r}}}\otimes e^{\ast}_{\underline{\gamma_{1}}} \otimes\cdots\otimes e^{\ast}_{\underline{\gamma_{s}}} 
\]
这里也重新记$e_{\gamma}\triangleq (e_{\overline{\gamma_{1}}},\cdots, e_{\overline{\gamma_{r}}}, e^{\ast}_{\underline{\gamma_{1}}} ,\cdots,e^{\ast}_{\underline{\gamma_{s}}} )$. 另外我们以后默认多重指标集里面的元素都已经按照字典顺序排列好了。

\begin{mydef}{}{chapter2rs型张量}
	设$V$是一个$n$维实向量空间，$V^{\ast}$为其对偶空间. 张量空间
	\[
	V^{r}_{s}\triangleq V\otimes \cdots\otimes V\otimes V^{\ast}\otimes \cdots\otimes V^{\ast}
	\]
	中的元素称为$(r,s)$型张量，称$r$为反变阶数，$s$为协变阶数.
\end{mydef}

在基$\{ e^{\otimes}_{\gamma}\}_{\gamma\in \Gamma^{r}_{s}}$下，任一$(r,s)$型张量$x\in V^{r}_{s} $都可以表示为
\[
x=\sum_{\gamma\in \Gamma}x_{\gamma}e^{\otimes}_{\gamma}\;,\;x_{\gamma}\in \mr
\]
我们称$\left(x_{\gamma} \right)_{\gamma\in\Gamma}$为张量$x$在基$\{ e^{\otimes}_{\gamma}\}_{\gamma\in \Gamma_{r+s}(n)}$下的\textbf{坐标分量}。

两个同为$(r,s)$型的张量可以做加法与数乘运算，所得结果仍是$(r,s)$型张量。这里的本质是因为张量空间$V^{r}_{s}$本身就是一个向量空间，所以它里面的元素当然可以做加法与数乘运算。除了加法与数乘这两种简单的运算之外，张量还有另外两种极其重要的运算：张量乘法和缩并。

首先来看张量的乘法运算。
\begin{mydef}{}{chapter2张量积}
	设$x\in V^{r}_{s},y\in V^{p}_{q}$. 张量$x$和$y$的乘积定义为：{\color{red}{$x\otimes y\in V^{r+p}_{s+q}$}}. 对任意的
	\[
	(\theta_{1},\theta_{2},\cdots,\theta_{r+s},v_{1},v_{2},\cdots,v_{s+q})\in\underbrace{ V^{\ast}\times \cdots\times V^{\ast}}_{r+p\text{个}}\times \underbrace{V\times\cdots\times V}_{s+q\text{个}}
	\]
	按如下规则计算
	\[
	\begin{split}
	&x\otimes y(\theta_{1},\cdots,\theta_{r+p},v_{1},\cdots,v_{s+q})\\
	&= x(\theta_{1},\cdots,\theta_{r},v_{1},\cdots,v_{s })\cdot y(\theta_{r+1},\cdots,\theta_{r+p},v_{s +1},\cdots,v_{s+q})
	\end{split}
	\]
\end{mydef}

由于$(r,s)$型张量本质上是一个$V^{\ast}\times\cdots\times V^{\ast}\times V\times\cdots \times V\longrightarrow \mr$的\textbf{$r+s$重线性函数}，所以我们可以对它\textbf{做多重线性函数的张量积运算}。张量乘法就是在把张量视作多重线性函数再做张量积运算的基础上略微\textbf{调整了一下自变量的顺序}，本质上仍是多重线性函数的张量积运算，所以张量乘法也是用的符号“$\otimes$”来表示。读者看了以上定义再回过头对照多重线性函数的张量积运算的定义，容易看出其中的细微区别。      
       
\begin{mythm}{}{chapter2张量积的性质}
	张量乘法具有下列性质
	\begin{itemize}
		\item 对任意$x\in V^{r}_{s},y\in V^{p}_{q},\lambda\in \mr$，$(\lambda x)\otimes y=\lambda(x\otimes y)=x\otimes (\lambda y)$.
		\item 结合律：对任意$x\in V^{r}_{s },y\in V^{p}_{q},z\in V^{k}_{l}$，$(x\otimes y)\otimes z=x\otimes (y\otimes z)$.
		\item 右分配律：对任意$x\in V^{r }_{s },y\in V^{r }_{s },z\in V^{p}_{q}$，$(x+y)\otimes z=x\otimes z+y\otimes z$.
		\item 左分配律：对任意$x\in V^{r }_{s },y\in V^{p}_{q},z\in V^{p}_{q}$，$x\otimes (y+z)=x\otimes y+x\otimes z$.
	\end{itemize}
\end{mythm}
证明很简单，按照定义验证即可。但是注意，张量乘法不满足交换律，即$x\otimes y\neq y\otimes x$.

\begin{mylem}{}{chapter2基的张量积定理}
	设$\{ e^{\otimes}_{\alpha}\}_{\alpha\in \Gamma^{r}_{s}}$为$V^{r}_{s}$的基，$\{ e^{\otimes}_{\beta}\}_{\beta\in \Gamma^{p}_{q}}$为$V^{p}_{q}$的基，$\{ e^{\otimes}_{\gamma}\}_{\gamma\in \Gamma^{r+p}_{s+q}}$为$V^{r+p}_{s+q}$的基.\\ 则对任意多重指标$\alpha=(\overline{\alpha},\underline{\alpha})\in \Gamma^{r}_{s}\;,\beta=(\overline{\beta},\underline{\beta})\in \Gamma^{p}_{q}\;,\gamma\in \Gamma^{r+p}_{s+q}$，有
\[
e^{\otimes}_{\gamma}=e^{\otimes}_{\alpha}\otimes e^{\otimes}_{\beta}\;\Leftrightarrow\; \gamma=(\overline{\alpha},\overline{\beta},\underline{\alpha},\underline{\beta})
\]
\end{mylem}
\begin{proof}
	直接计算
	\[
	\begin{split}
	&e^{\otimes}_{\alpha}\otimes e^{\otimes}_{\beta}(\theta_{1},\cdots,\theta_{r+p},v_{1},\cdots,v_{s+q})\\
    =&{\color{red}{e^{\otimes}_{\alpha}(\theta_{1},\cdots,\theta_{r},v_{1},\cdots,v_{s})}}\cdot {\color{blue}{e^{\otimes}_{\beta}(\theta_{r+1},\cdots,\theta_{r+p},v_{s+1},\cdots,v_{s+q})}}\\
    =&{\color{red}{\theta_{1}(e_{\overline{\alpha_{ 1 }}})\cdots\theta_{r}(e_{\overline{\alpha_{ r }}})\cdot e^{\ast}_{\underline{\alpha_{ 1}}}(v_{1})\cdots e^{\ast}_{\underline{\alpha_{s}}}(v_{s})}}\cdot{\color{blue}{\theta_{r+1}(e_{\overline{\beta_{ 1 }}})\cdots\theta_{r+p}(e_{\overline{\beta_{ p }}})\cdot e^{\ast}_{\underline{\beta_{1}}}(v_{s+1})\cdots e^{\ast}_{\underline{\beta_{ q}}}(v_{s+q})}}\\
    =&{\color{red}{\theta_{1}(e_{\overline{\alpha_{ 1 }}})\cdots\theta_{r}(e_{\overline{\alpha_{ r }}})}}\cdot
    {\color{blue}{\theta_{r+1}(e_{\overline{\beta_{ 1 }}})\cdots\theta_{r+p}(e_{\overline{\beta_{ p }}})}}\cdot
    {\color{red}{e^{\ast}_{\underline{\alpha_{ 1}}}(v_{1})\cdots e^{\ast}_{\underline{\alpha_{s}}}(v_{s})}}
    \cdot{\color{blue}{ e^{\ast}_{\underline{\beta_{1}}}(v_{s+1})\cdots e^{\ast}_{\underline{\beta_{ q}}}(v_{s+q})}}\\
    =&e_{\overline{\alpha_{ 1 }}}\otimes\cdots\otimes e_{\overline{\alpha_{ r }}}\otimes e_{\overline{\beta_{ 1 }}}\otimes\cdots\otimes e_{\overline{\beta_{ p }}}\otimes e^{\ast}_{\underline{\alpha_{ 1}}}\otimes \cdots \otimes e^{\ast}_{\underline{\alpha_{ s}}}\otimes e^{\ast}_{\underline{\beta_{ 1}}}\otimes\cdots\otimes e^{\ast}_{\underline{\beta_{q}}}(\theta_{1},\cdots,\theta_{r+p},v_{1},\cdots,v_{s+q})\\
	\end{split}
	\]
	从而，$e^{\otimes}_{\alpha}\otimes e^{\otimes}_{\beta}=e_{\overline{\alpha_{ 1 }}}\otimes\cdots\otimes e_{\overline{\alpha_{ r }}}\otimes e_{\overline{\beta_{ 1 }}}\otimes\cdots\otimes e_{\overline{\beta_{ p }}}\otimes e^{\ast}_{\underline{\alpha_{ 1}}}\otimes \cdots \otimes e^{\ast}_{\underline{\alpha_{ s}}}\otimes e^{\ast}_{\underline{\beta_{ 1}}}\otimes\cdots\otimes e^{\ast}_{\underline{\beta_{q}}}$. 于是
	\[
	e^{\otimes}_{\alpha}\otimes e^{\otimes}_{\beta}=e^{\otimes}_{\gamma}
	\Leftrightarrow\gamma=(\overline{\alpha_{ 1 }},\cdots,\overline{\alpha_{ r }},\overline{\beta_{ 1 }},\cdots,\overline{\beta_{ p }},\underline{\alpha_{1}},\cdots,\underline{\alpha_{s}},\underline{\beta_{1}},\cdots,\underline{\beta_{q}})
	\Leftrightarrow\gamma=(\overline{\alpha},\overline{\beta},\underline{\alpha},\underline{\beta})
	\]
	上述“$\Leftrightarrow$”之所以能成立，是因为由指标$\alpha,\beta$得到指标$\gamma$的整个过程都是\textbf{一一对应}的。
	而且你只要把$\gamma$按照$r,p,s,q$的长度切分成四块，比如说$\gamma=(\text{\ding{202}},\text{\ding{203}} ,\text{\ding{204}},\text{\ding{205}})$，再分别令$\alpha=(\text{\ding{202}},\text{\ding{204}}),\beta=(\text{\ding{203}},\text{\ding{205}})$，即可由 $\gamma\in \Gamma^{r+p}_{s+q}$获得$\alpha\in \Gamma^{r}_{s}$与$\beta \in \Gamma^{p}_{q}$.
\end{proof}

\begin{mythm}{}{chapter2张量积的坐标分量}
	设$\{ e^{\otimes}_{\alpha}\}_{\alpha\in \Gamma^{r}_{s}}$为$V^{r}_{s}$的基，$\{ e^{\otimes}_{\beta}\}_{\beta\in \Gamma^{p}_{q}}$为$V^{p}_{q}$的基，$\{ e^{\otimes}_{\gamma}\}_{\gamma\in \Gamma^{r+p}_{s+q}}$为$V^{r+p}_{s+q}$的基.\\
	 则$x\otimes y$在基下的第$\gamma$个坐标分量为
	\[
	(x\otimes y)_{\gamma}=x_{\alpha}\cdot y_{\beta}
	\]
	其中指标满足$\gamma=(\overline{\alpha},\overline{\beta},\underline{\alpha},\underline{\beta})\in \Gamma^{r+p}_{s+q}\;,\alpha=(\overline{\alpha},\underline{\alpha})\in \Gamma^{r}_{s}\;,\beta=(\overline{\beta},\underline{\beta})\in \Gamma^{p}_{q}$.
\end{mythm}

\begin{proof}
	分别设$x=\sum_{ \alpha }x_{\alpha}e^{\otimes}_{\alpha}\;,y=\sum_{\beta}y_{\beta}e^{\otimes}_{\beta}$. 则根据张量乘法的运算性质可知
	\[
	\begin{split}
	x\otimes y&=\left(\sum_{ \alpha }x_{\alpha}e^{\otimes}_{\alpha} \right)\otimes \left( \sum_{\beta}y_{\beta}e^{\otimes}_{\beta}\right)\\
	&=\sum_{ \alpha }\sum_{\beta}x_{\alpha}y_{\beta}e^{\otimes}_{\alpha}\otimes e^{\otimes}_{\beta}
	\end{split}
	\]
	由于$\gamma=(\overline{\alpha},\overline{\beta},\underline{\alpha},\underline{\beta})$，则由\rlem{lem:chapter2基的张量积定理}可知$e^{\otimes}_{\alpha}\otimes e^{\otimes}_{\beta}=e^{\otimes}_{\gamma}$. 于是上式变为
	\[
	x\otimes y
	={\color{blue}{\sum_{ \alpha }\sum_{\beta}x_{\alpha}y_{\beta}e^{\otimes}_{\gamma}}}
	\]
	注意到$\alpha,\beta$与$\gamma$是一一对应的，所以当$\alpha$跑遍多重指标集$\Gamma^{r}_{s}$，且$\beta$跑遍多重指标集$\Gamma^{p}_{q}$时，意味着由$\alpha,\beta$唯一确定的$\gamma=(\overline{\alpha},\overline{\beta},\underline{\alpha},\underline{\beta})$也同时跑遍了多重指标集$ \Gamma^{r+p}_{s+q}$. 所以$\sum_{ \alpha }\sum_{\beta}$所表示的求和等同于$\sum_{ \gamma }$表示的求和。即
	\[
	x\otimes y
	={\color{blue}{\sum_{ \gamma }x_{\alpha}y_{\beta}e^{\otimes}_{\gamma}}}
	\]
	而$e^{\otimes}_{\gamma}$前面的系数就是$x\otimes y$在基下的第$\gamma$个坐标分量$(x\otimes y)_{\gamma}$，因此$(x\otimes y)_{\gamma}=x_{\alpha}\cdot y_{\beta}$.
\end{proof}

接下来我们介绍张量的缩并运算。

\begin{mydef}{}{chapter2张量的缩并}
	设$\{e_{1},e_{2},\cdots,e_{n} \}$为向量空间$V$的基，$\{e^{\ast}_{1},e^{\ast}_{2},\cdots,e^{\ast}_{n} \}$为$V^{\ast}$中相应的对偶基. \\
	若$x\in V^{r}_{s}$，取定$1\leqslant k\leqslant r,1\leqslant l\leqslant s$，则称由下式确定的$(r-1,s-1)$型张量$C^{k}_{l}(x)$为$x$的缩并：
	\[
		(\theta_{1},\cdots,\theta_{r-1},v_{1},\cdots,v_{s-1})\in\underbrace{ V^{\ast}\times \cdots\times V^{\ast}}_{r-1\text{个}}\times \underbrace{V\times\cdots\times V}_{s-1\text{个}}
	\]
	按如下规则计算
	\[
	\begin{split}
	&C^{k}_{l}(x)(\theta_{1},\cdots,\theta_{r-1},v_{1},\cdots,v_{s-1})\\
	=& \sum_{i=1}^{n}x(\theta_{1},\cdots,\theta_{\color{red}{k-1}},{\color{blue}{e^{\ast}_{i}}},\theta_{\color{red}{k}},\cdots,\theta_{r-1},v_{1},\cdots,v_{\color{red}{l-1}},{\color{blue}{e_{i}}},v_{\color{red}{l}},\cdots,v_{s-1})
	\end{split}
	\]
	我们称映射$C^{k}_{l}:V^{r}_{s}\longrightarrow V^{r-1}_{s-1}$为张量的缩并运算.
\end{mydef}

容易验证$C^{k}_{l}:V^{r}_{s}\longrightarrow V^{r-1}_{s-1}$为\textbf{线性映射}，所以缩并运算是线性运算。特别地，当张量$x\in V^{r}_{s}$本身是一个\textbf{单项式}时，即形如$x=w_{1}\otimes \cdots \otimes w_{r}\otimes\sigma_{1}\otimes\cdots\otimes \sigma_{s}$，其中$w_{i}\in V,\sigma_{j}\in V^{\ast} \;, i=1,\cdots,r\;, j=1,\cdots,s$. 我们有如下结论。
\begin{mythm}{}{chapter2单项式的缩并}
	对任一单项式$x=w_{1}\otimes \cdots \otimes w_{r}\otimes\sigma_{1}\otimes\cdots\otimes \sigma_{s}\in V^{r}_{s}$，对其做缩并运算所得到的$(r-1,s-1)$型张量$C^{k}_{l}(x)$为
	\[
	C^{k}_{l}(x)={\color{blue}{\sigma_{l}(w_{k})}} w_{1}\otimes \cdots\otimes \hat{w_{\color{blue}{k}}}\otimes \cdots\otimes w_{r}\otimes \sigma_{1}\otimes\cdots\otimes \hat{\sigma_{\color{blue}{l}}}\otimes\cdots\otimes  \sigma_{s}
	\]
	其中其中$\hat{w_{k}}$与$\hat{\sigma_{l}}$表示删去$w_{k}$与$\sigma_{l} $这两项.
\end{mythm}
\begin{proof}
	直接计算
	\[
	\begin{split}
	&C^{k}_{l}(x)(\theta_{1},\cdots,\theta_{r-1},v_{1},\cdots,v_{s-1})\\
	=&\sum_{i=1}^{n}x(\theta_{1},\cdots,\theta_{\color{red}{k-1}},{\color{blue}{e^{\ast}_{i}}},\theta_{\color{red}{k}},\cdots,\theta_{r-1},v_{1},\cdots,v_{\color{red}{l-1}},{\color{blue}{e_{i}}},v_{\color{red}{l}},\cdots,v_{s-1})\\
	=&\sum_{i=1}^{n}\left(\theta_{1}(w_{1})\cdots \theta_{\color{red}{k-1}}(w_{\color{red}{k-1}})\cdot{\color{blue}{e^{\ast}_{i}(w_{k})}}\cdot\theta_{\color{red}{k}}(w_{\color{red}{k+1}})\cdots \cdots \sigma_{\color{red}{l-1}}(v_{\color{red}{l-1}})\cdot{\color{blue}{\sigma_{l}(e_{i})}}\cdot\sigma_{\color{red}{l+1}}(v_{\color{red}l})\cdots \sigma_{s}(v_{s-1})\right)\\
	=&\left(\theta_{1}(w_{1})\cdots \theta_{\color{red}{k-1}}(w_{\color{red}{k-1}})\cdot\theta_{\color{red}{k}}(w_{\color{red}{k+1}})\cdots \cdots \sigma_{\color{red}{l-1}}(v_{\color{red}{l-1}})\cdot\sigma_{\color{red}{l+1}}(v_{\color{red}l})\cdots\sigma_{s}(v_{s-1})\right)\cdot  \left(\sum_{i=1}^{n}{\color{blue}{e^{\ast}_{i}(w_{k})}}{\color{blue}{\sigma_{l}(e_{i})}}\right)\\
	=&w_{1}\otimes \cdots\otimes \hat{w_{\color{blue}{k}}}\otimes \cdots\otimes w_{r}\otimes \sigma_{1}\otimes\cdots\otimes \hat{\sigma_{\color{blue}{l}}}\otimes\cdots\otimes  \sigma_{s}(\theta_{1},\cdots,\theta_{r-1},v_{1},\cdots,v_{s-1})\cdot {\color{blue}{\sigma_{l}(w_{k})}}\\
	=&\left({\color{blue}{\sigma_{l}(w_{k})}}\cdot w_{1}\otimes \cdots\otimes \hat{w_{\color{blue}{k}}}\otimes \cdots\otimes w_{r}\otimes \sigma_{1}\otimes\cdots\otimes \hat{\sigma_{\color{blue}{l}}}\otimes\cdots\otimes  \sigma_{s}\right)(\theta_{1},\cdots,\theta_{r-1},v_{1},\cdots,v_{s-1})\\
	\end{split}
	\]
	
	其中$\hat{w_{k}}$与$\hat{\sigma_{l}}$表示删去这两项。第四个等号成立则是因为$w_{k}=\sum_{i=1}^{n}e^{\ast}_{i}(w_{k})e_{i}$且$\sigma_{l}$是线性函数. 从而
	$
	C^{k}_{l}(x)={\color{blue}{\sigma_{l}(w_{k})}} \cdot w_{1}\otimes \cdots\otimes \hat{w_{\color{blue}{k}}}\otimes \cdots\otimes w_{r}\otimes \sigma_{1}\otimes\cdots\otimes \hat{\sigma_{\color{blue}{l}}}\otimes\cdots\otimes  \sigma_{s}
	$.
\end{proof}

\begin{mythm}{}{chapter2缩并的坐标分量}
	设$\{ e^{\otimes}_{\alpha}\}_{\alpha\in \Gamma^{r}_{s}}$为$V^{r}_{s}$的基，$\{ e^{\otimes}_{\gamma}\}_{\gamma\in \Gamma^{r-1}_{s-1}}$为$V^{r-1}_{s-1}$的基. 取定$1\leqslant k\leqslant r,1\leqslant l\leqslant s$. \\
	则对任意$x\in V^{r}_{s}$，$C^{k}_{l}(x)$在基下的第$\gamma$个坐标分量为
	\[
	 \left(C^{k}_{l}(x)\right)_{\gamma}=\sum_{ i = 1 }^{n}x_{i(\gamma)}
	\]
	其中$i:\Gamma^{r-1}_{s-1}\longrightarrow \Gamma^{r}_{s}$定义为：对$\forall\; \gamma=(\overline{\gamma_{1}},\cdots,\overline{\gamma_{r-1}},\underline{\gamma_{1}},\cdots,\underline{\gamma_{s-1}})\in \Gamma^{r-1}_{s-1}$
	\[
	i(\gamma)\triangleq \left(\overline{\gamma_{1}},\cdots,\overline{\gamma_{\color{red}{k-1}}},{\color{blue}{i}},\overline{\gamma_{\color{red}{k}}},\cdots,\overline{\gamma_{r-1}},\underline{\gamma_{1}},\cdots,\underline{\gamma_{\color{red}{l-1}}},{\color{blue}{i}},\underline{\gamma_{\color{red}{l}}},\cdots,\underline{\gamma_{s-1}}\right)\in \Gamma^{r}_{s}
	\]
\end{mythm}
\begin{proof}
	设$x=\sum_{ \alpha\in \Gamma^{r}_{s} }x_{\alpha}e^{\otimes}_{\alpha}$. 则根据缩并运算的线性性质可知
	\begin{equation}
		C^{k}_{l}(x)=\sum_{ \alpha\in \Gamma^{r}_{s} }x_{\alpha}C^{k}_{l}\left(e^{\otimes}_{\alpha}\right)
	\end{equation}
	
	而根据\rthm{thm:chapter2单项式的缩并}可知
	\[
	C^{k}_{l}\left(e^{\otimes}_{\alpha}\right)={\color{blue}{e^{\ast}_{\underline{\alpha_{ l }}}} (e_{\overline{\alpha_{ k }}})}e_{\overline{\alpha_{ 1 }}}\otimes \cdots\otimes \hat{e_{\color{blue}{\overline{\alpha_{ k }}}}}\otimes \cdots\otimes e_{\overline{\alpha_{ r }}}\otimes e^{\ast}_{\underline{\alpha_{ 1 }}}\otimes\cdots\otimes \hat{e^{\ast}_{\color{blue}{\underline{\alpha_{ l }}}}}\otimes\cdots\otimes  e^{\ast}_{\underline{\alpha_{ s }}}
	\]
	由对偶基的性质可知，${\color{blue}{e^{\ast}_{\underline{\alpha_{ l }}}} (e_{\overline{\alpha_{ k }}})}$当$\underline{\alpha_{ l }}\neq \overline{\alpha_{ k }}$时为$0$；当$\underline{\alpha_{ l }}= \overline{\alpha_{ k }}$时为$1$. 所以式\requ{equ:}的求和范围缩小到$\tilde{  \Gamma  }^{r}_{s}\triangleq \{\alpha\in \Gamma^{r}_{s} | \underline{\alpha_{ l }}= \overline{\alpha_{ k }} \}$. 即
	\[
	\begin{split}
	C^{k}_{l}(x)&=\sum_{ \alpha\in \Gamma^{r}_{s} }x_{\alpha}C^{k}_{l}\left(e^{\otimes}_{\alpha}\right)\\
	&=\sum_{ \alpha\in \tilde{  \Gamma  }^{r}_{s} }x_{\alpha}e_{\overline{\alpha_{ 1 }}}\otimes \cdots\otimes \hat{e_{\color{blue}{\overline{\alpha_{ k }}}}}\otimes \cdots\otimes e_{\overline{\alpha_{ r }}}\otimes e^{\ast}_{\underline{\alpha_{ 1 }}}\otimes\cdots\otimes \hat{e^{\ast}_{\color{blue}{\underline{\alpha_{ l }}}}}\otimes\cdots\otimes  e^{\ast}_{\underline{\alpha_{ s }}}\\
	\end{split}
	\]
	注意到$\tilde{  \Gamma  }^{r}_{s}\triangleq \{\alpha\in \Gamma^{r}_{s} | \underline{\alpha_{ l }}= \overline{\alpha_{ k }} \}=\tilde{  \Gamma  }^{r}_{s}(1)\cup\tilde{  \Gamma  }^{r}_{s}(2)\cup\cdots\cup \tilde{  \Gamma  }^{r}_{s}(n)$. 其中$\tilde{  \Gamma  }^{r}_{s}(i)\triangleq \{\alpha\in \Gamma^{r}_{s} | \overline{\alpha_{ k }}=\underline{\alpha_{ l }}=i\}$. 因此上式又可以写成
	\[
	\begin{split}
	C^{k}_{l}(x)&=\sum_{ \alpha\in \Gamma^{r}_{s} }x_{\alpha}C^{k}_{l}\left(e^{\otimes}_{\alpha}\right)\\
	&=\sum_{ \alpha\in \tilde{  \Gamma  }^{r}_{s} }x_{\alpha}e_{\overline{\alpha_{ 1 }}}\otimes \cdots\otimes \hat{e_{\color{blue}{\overline{\alpha_{ k }}}}}\otimes \cdots\otimes e_{\overline{\alpha_{ r }}}\otimes e^{\ast}_{\underline{\alpha_{ 1 }}}\otimes\cdots\otimes \hat{e^{\ast}_{\color{blue}{\underline{\alpha_{ l }}}}}\otimes\cdots\otimes  e^{\ast}_{\underline{\alpha_{ s }}}\\
	&=\sum_{ i = 1 }^{n}\sum_{ \alpha\in \tilde{  \Gamma  }^{r}_{s}(i) }x_{\alpha}e_{\overline{\alpha_{ 1 }}}\otimes \cdots\otimes \hat{e_{\color{blue}{\overline{\alpha_{ k }}}}}\otimes \cdots\otimes e_{\overline{\alpha_{ r }}}\otimes e^{\ast}_{\underline{\alpha_{ 1 }}}\otimes\cdots\otimes \hat{e^{\ast}_{\color{blue}{\underline{\alpha_{ l }}}}}\otimes\cdots\otimes  e^{\ast}_{\underline{\alpha_{ s }}}\\
	\end{split}
	\]
	注意到$\tilde{  \Gamma  }^{r}_{s}(i)=i(  \Gamma  ^{r-1}_{s-1})$. 这里$i:\Gamma^{r-1}_{s-1}\longrightarrow \Gamma^{r}_{s}$就是条件中所给的映射。所以对于任意$\alpha\in\tilde{  \Gamma  }^{r}_{s}(i) $，都存在唯一的$\gamma\in\Gamma  ^{r-1}_{s-1} $，使得$\alpha=i(\gamma)$. 因此我们得到
	\[
	\begin{split}
	C^{k}_{l}(x)&=\sum_{ \alpha\in \Gamma^{r}_{s} }x_{\alpha}C^{k}_{l}\left(e^{\otimes}_{\alpha}\right)\\
	&=\sum_{ \alpha\in \tilde{  \Gamma  }^{r}_{s} }x_{\alpha}e_{\overline{\alpha_{ 1 }}}\otimes \cdots\otimes \hat{e_{\color{blue}{\overline{\alpha_{ k }}}}}\otimes \cdots\otimes e_{\overline{\alpha_{ r }}}\otimes e^{\ast}_{\underline{\alpha_{ 1 }}}\otimes\cdots\otimes \hat{e^{\ast}_{\color{blue}{\underline{\alpha_{ l }}}}}\otimes\cdots\otimes  e^{\ast}_{\underline{\alpha_{ s }}}\\
	&=\sum_{ i = 1 }^{n}\sum_{ \alpha\in \tilde{  \Gamma  }^{r}_{s}(i) }x_{\alpha}e_{\overline{\alpha_{ 1 }}}\otimes \cdots\otimes \hat{e_{\color{blue}{\overline{\alpha_{ k }}}}}\otimes \cdots\otimes e_{\overline{\alpha_{ r }}}\otimes e^{\ast}_{\underline{\alpha_{ 1 }}}\otimes\cdots\otimes \hat{e^{\ast}_{\color{blue}{\underline{\alpha_{ l }}}}}\otimes\cdots\otimes  e^{\ast}_{\underline{\alpha_{ s }}}\\
	&=\sum_{ i = 1 }^{n}\sum_{ \gamma\in \Gamma^{r-1}_{s-1} }x_{i(\gamma)}e_{\overline{\gamma_{ 1 }}}\otimes \cdots\otimes e_{\color{red}{\overline{\gamma_{ k-1 }}}}\otimes e_{\color{red}{\overline{\gamma_{ k }}}}\otimes \cdots\otimes e_{\overline{\gamma_{ r-1 }}}\otimes e^{\ast}_{\underline{\gamma_{ 1 }}}\otimes\cdots\otimes e^{\ast}_{\color{red}{\underline{\gamma_{ l-1 }}}}\otimes e^{\ast}_{\color{red}{\underline{\gamma_{ l }}}}\otimes\cdots\otimes  e^{\ast}_{\underline{\gamma_{ s-1 }}}\\
	&=\sum_{ \gamma\in \Gamma^{r-1}_{s-1} }\left(\sum_{ i = 1 }^{n}x_{i(\gamma)}\right)e^{\otimes}_{\gamma}\\
	\end{split}
	\]
	而$e^{\otimes}_{\gamma}$前面的系数就是$C^{k}_{l}(x)$在基下的第$\gamma$个坐标分量$\left(C^{k}_{l}(x)\right)_{\gamma}$，因此$\left(C^{k}_{l}(x)\right)_{\gamma}=\sum_{ i = 1 }^{n}x_{i(\gamma)}$.
\end{proof}

在线性代数中我们研究过当向量空间的基变化时，向量的坐标如何变化的问题。现在我们对于$(r,s)$型张量空间$V^{r}_{s}$也提出这个问题：当向量空间$V$中的基发生变化时，$(r,s)$型张量的坐标分量会随之如何变化？

此前接触过张量的读者一定知道张量分量的变换公式写出来十分复杂，为什么看起来那么复杂呢？因为大多数教材上的推导过程都比较干净直接，省略了很多背后的原理。实际上张量分量的变换公式的本质极其简单，早在线性代数课程中我们就已经学过：
\begin{itemize}
	\item[$\bullet$] 若向量空间$V$中有基变换$(\overline{e}_{1},\overline{e}_{2},\cdots,\overline{e}_{n})=(e_{1},e_{2},\cdots,e_{n})A$，则向量$v$在相应的基下的坐标满足$\overline{x}=A^{-1}x$.
\end{itemize}
这也是我们下面建立张量分量变化关系的第一个基本事实。另一个基本事实也很简单，简述如下：
\begin{itemize}
	\item[$\bullet$] 从基$(e_{1},e_{2},\cdots,e_{n})$到基$(\overline{e}_{1},\overline{e}_{2},\cdots,\overline{e}_{n})$的过渡矩阵$A$，可以视为恒等变换$id:\left(V,\{\overline{e}_{1},\cdots,\overline{e}_{n} \} \right)\longrightarrow \left(V,\{e_{1},e_{2},\cdots,e_{n} \} \right)$在基下的矩阵.
\end{itemize}
只要学过线性代数，就肯定看得懂这两个基本事实。下面我们基于以上两个十分简单的事实来推导$(r,s)$型张量的分量在基变换下的坐标变换公式。

设$(e_{1},e_{2},\cdots,e_{n})$为向量空间$V$的一组基，$(\overline{e}_{1},\overline{e}_{2},\cdots,\overline{e}_{n})$为向量空间$V$的另一组基。假设$(\overline{e}_{1},\overline{e}_{2},\cdots,\overline{e}_{n})=(e_{1},e_{2},\cdots,e_{n})A$，根据基本事实，过渡矩阵$A$可以视为恒等变换
\[id:\left(V,\{\overline{e}_{1},\cdots,\overline{e}_{n} \} \right)\longrightarrow \left(V,\{e_{1},e_{2},\cdots,e_{n} \} \right)\]
在基下的矩阵.

\textbf{\underline{Claim1}}：{\kaishu{$id:\left(V,\{\overline{e}_{1},\cdots,\overline{e}_{n} \} \right)\longrightarrow \left(V,\{e_{1},e_{2},\cdots,e_{n} \} \right)$所诱导的对偶空间之间的线性映射
$id^{\ast}:\left(V^{\ast},\{e^{\ast}_{1},e^{\ast}_{2},\cdots,e^{\ast}_{n} \} \right)\longrightarrow \left(V^{\ast}, \{\overline{e}^{\ast}_{1},\cdots,\overline{e}^{\ast}_{n} \}\right)$也是恒等映射. }}

事实上，根据$id^{\ast}$的定义可知，$id^{\ast}(\theta)(v)=\theta(id(v))=\theta(v)\Rightarrow id^{\ast}(\theta)=\theta\;,\forall \theta\in V^{\ast} $. 这表明$id^{\ast}$为恒等映射。

\textbf{\underline{Claim2}}：{\kaishu{$Id=(id^{\ast})^{-1}:\left(V^{\ast}, \{\overline{e}^{\ast}_{1},\cdots,\overline{e}^{\ast}_{n} \}\right)\longrightarrow \left(V^{\ast},\{e^{\ast}_{1},e^{\ast}_{2},\cdots,e^{\ast}_{n} \} \right)$为恒等映射且$Id$在两组基下的矩阵为$(A^{T})^{-1}$. }}

由于$id^{\ast}$为恒等映射，那么显然$id^{\ast}$可逆并且逆映射仍为恒等映射。又因为$id^{\ast}$在两组基下的矩阵为$A^{T}$，即
\[
\begin{split}
&id^{\ast}\left(e^{\ast}_{1},e^{\ast}_{2},\cdots,e^{\ast}_{n}  \right)=\left(\overline{e}^{\ast}_{1},\cdots,\overline{e}^{\ast}_{n} \right) A^{T}\\
\Leftrightarrow& \left(e^{\ast}_{1},e^{\ast}_{2},\cdots,e^{\ast}_{n}  \right)=\left(\overline{e}^{\ast}_{1},\cdots,\overline{e}^{\ast}_{n} \right) A^{T}\\
\Leftrightarrow& (\overline{e}^{\ast}_{1},\overline{e}^{\ast}_{2},\cdots,\overline{e}^{\ast}_{n})=(e^{\ast}_{1},e^{\ast}_{2},\cdots,e^{\ast}_{n})(A^{T})^{-1}\\
\Leftrightarrow& Id\left(\overline{e}^{\ast}_{1},\cdots,\overline{e}^{\ast}_{n} \right)=\left(e^{\ast}_{1},e^{\ast}_{2},\cdots,e^{\ast}_{n}  \right) (A^{T})^{-1}\\
\end{split}
\]

\textbf{\underline{Claim3}}：{\kaishu{由$r$个线性映射$id:\left(V,\{\overline{e}_{1},\cdots,\overline{e}_{n} \} \right)\longrightarrow \left(V,\{e_{1},e_{2},\cdots,e_{n} \} \right)$与$s$个线性映射$Id:\left(V^{\ast}, \{\overline{e}^{\ast}_{1},\cdots,\overline{e}^{\ast}_{n} \}\right)\longrightarrow \left(V^{\ast},\{e^{\ast}_{1},e^{\ast}_{2},\cdots,e^{\ast}_{n} \} \right)$所诱导的张量空间之间的线性映射$T:\left(V^{r}_{s},\{ \overline{e}^{\otimes}_{\alpha}\}_{\alpha\in \Gamma^{r}_{s}} \right)\longrightarrow \left(V^{r}_{s},\{e^{\otimes}_{\beta} \}_{\beta\in \Gamma^{r}_{s}} \right)$也是恒等映射.}}

事实上，对于单项式$v_{1}\otimes\cdots\otimes v_{r}\otimes \sigma_{1}\otimes\cdots\otimes \sigma_{s}\in V^{r}_{s}$，我们有
\[
\begin{split}
&T(v_{1}\otimes\cdots\otimes v_{r}\otimes \sigma_{1}\otimes\cdots\otimes \sigma_{s})\\
&=id(v_{1})\otimes\cdots\otimes id(v_{r})\otimes Id(\sigma_{1})\otimes\cdots\otimes Id(\sigma_{s})\\
&=v_{1}\otimes\cdots\otimes v_{r}\otimes \sigma_{1}\otimes\cdots\otimes \sigma_{s}
\end{split}
\]
因此，对于任意$x\in V^{r}_{s}$，设$x=\sum_{ \gamma\in \Gamma^{r}_{s} }x_{\alpha}v_{\overline{\alpha_{ 1 }}}\otimes\cdots\otimes v_{\overline{\alpha_{ r }}}\otimes \sigma_{\underline{\alpha_{ 1 }}}\otimes\cdots\otimes \sigma_{\underline{\alpha_{ s }}}$，其中$x_{\alpha}\in \mr$. 则
\[
\begin{split}
T(x)&=\sum_{ \gamma\in \Gamma^{r}_{s} }x_{\alpha}T\left(v_{\overline{\alpha_{ 1 }}}\otimes\cdots\otimes v_{\overline{\alpha_{ r }}}\otimes \sigma_{\underline{\alpha_{ 1 }}}\otimes\cdots\otimes \sigma_{\underline{\alpha_{ s }}}\right)\\
&=\sum_{ \gamma\in \Gamma^{r}_{s} }x_{\alpha}v_{\overline{\alpha_{ 1 }}}\otimes\cdots\otimes v_{\overline{\alpha_{ r }}}\otimes \sigma_{\underline{\alpha_{ 1 }}}\otimes\cdots\otimes \sigma_{\underline{\alpha_{ s }}}\\
&=x
\end{split}
\]
因此$T:\left(V^{r}_{s},\{ \overline{e}^{\otimes}_{\alpha}\}_{\alpha\in \Gamma^{r}_{s}} \right)\longrightarrow \left(V^{r}_{s},\{e^{\otimes}_{\beta} \}_{\beta\in \Gamma^{r}_{s}} \right)$是恒等映射.

根据基本事实可知，恒等映射$T$在基$\{ \overline{e}^{\otimes}_{\alpha}\}_{\alpha\in \Gamma^{r}_{s}}$与基$\{e^{\otimes}_{\beta} \}_{\beta\in \Gamma^{r}_{s}} $下的矩阵$C$就是从基$\{e^{\otimes}_{\beta} \}_{\beta\in \Gamma^{r}_{s}} $到基$\{ \overline{e}^{\otimes}_{\alpha}\}_{\alpha\in \Gamma^{r}_{s}}$的过渡矩阵。即
\[
\{ \overline{e}^{\otimes}_{\alpha}\}_{\alpha\in \Gamma^{r}_{s}}=T\{ \overline{e}^{\otimes}_{\alpha}\}_{\alpha\in \Gamma^{r}_{s}}=\{e^{\otimes}_{\beta} \}_{\beta\in \Gamma^{r}_{s}}\cdot C
\]
其中基里面的元素都按照字典顺序排列。显然$C$是一个$n^{r+s}$行、$n^{r+s}$列的方阵。由\rthm{thm:chapter2诱导多重映射的矩阵表示}可知
\[
C=\underbrace{A\otimes \cdots \otimes A}_{r\text{个}}\otimes \underbrace{(A^{T})^{-1}\otimes \cdots\otimes (A^{T})^{-1}}_{s\text{个}}
\]

综上所述，我们有以下定理。
\begin{mythm}{}{chapter2张量空间之间的基过渡矩阵}
	若在向量空间$V$中，从基$(e_{1},e_{2},\cdots,e_{n})$到基$(\overline{e}_{1},\overline{e}_{2},\cdots,\overline{e}_{n})$的过渡矩阵为$A$，则在张量空间$V^{r}_{s}$中，从基$\{e^{\otimes}_{\beta} \}_{\beta\in \Gamma^{r}_{s}} $到基$\{ \overline{e}^{\otimes}_{\alpha}\}_{\alpha\in \Gamma^{r}_{s}}$的过渡矩阵$C$可由矩阵$A$的Kronecker积表示：
	\[
	C=\underbrace{A\otimes \cdots \otimes A}_{r\text{个}}\otimes \underbrace{(A^{T})^{-1}\otimes \cdots\otimes (A^{T})^{-1}}_{s\text{个}}
	\]
\end{mythm}

既然过渡矩阵都找出来了，我们实际上也就知道了$(r,s)$型张量在基下的坐标分量的变换规律。对任意$x\in V^{r}_{s}$，设$x$在基$\{ \overline{e}^{\otimes}_{\alpha}\}_{\alpha\in \Gamma^{r}_{s}}$下的第$\alpha$个坐标分量为$\overline{x}_{\alpha}$，在基$\{e^{\otimes}_{\beta} \}_{\beta\in \Gamma^{r}_{s}} $下的第$\beta$个坐标分量为$x_{\beta}$，则根据基本事实可知
\[
\left( \begin{array}{ccccc}{\vdots} \\ {\overline{x}_{\alpha}}  \\ {\vdots} \\ {\vdots}\\ {\vdots}\end{array}\right)=C^{-1}\left( \begin{array}{ccccc}{\vdots} \\ {\vdots}  \\ {\vdots} \\ {x_{\beta}}\\ {\vdots}\end{array}\right)
\]
而根据Kronecker积的运算性质又知道
$
C^{-1}=\underbrace{A^{-1}\otimes \cdots \otimes A^{-1}}_{r\text{个}}\otimes \underbrace{A^{T}\otimes \cdots\otimes A^{T}}_{s\text{个}}
$.
设
\[
A=\left( \begin{array}{cccc}{a_{1}^{1}} & {a^{1}_{2}} & {\cdots} & {a^{1}_{ n}} \\ {a^{2}_{1}} & {a^{2}_{2}} & {\cdots} & {a^{2}_{ n}} \\ {\vdots} & {\vdots} & { } & {\vdots} \\ {a^{n}_{ 1}} & {a^{n}_{ 2}} & {\cdots} & {a^{n}_{ n}}\end{array}\right)\quad\quad A^{-1}=\left( \begin{array}{cccc}{b_{1}^{1}} & {b^{1}_{2}} & {\cdots} & {b^{1}_{ n}} \\ {b^{2}_{1}} & {b^{2}_{2}} & {\cdots} & {b^{2}_{ n}} \\ {\vdots} & {\vdots} & { } & {\vdots} \\ {b^{n}_{ 1}} & {b^{n}_{ 2}} & {\cdots} & {b^{n}_{ n}}\end{array}\right)
\]
这里我们用上指标标识元素所在的行，用下指标标识元素所在的列。根据\rthm{thm:chapter2Kronecker积中元素的计算}可知，矩阵$C^{-1}$中第$\alpha$行第$\beta$列的元素为$b^{\overline{\alpha_{1}}}_{\overline{\beta_{ 1 }}}\cdots b^{\overline{\alpha_{r}}}_{\overline{\beta_{ r }}}\cdot a^{\underline{\beta_{ 1 }}}_{\underline{\alpha_{ 1 }}}\cdots a^{\underline{\beta_{ s }}}_{\underline{\alpha_{ s }}}$. 注意这里是转置矩阵$A^{T}$.
\[
\left( \begin{array}{ccccc}{\vdots} \\ {\overline{x}_{\alpha}}  \\ {\vdots} \\ {\vdots}\\ {\vdots}\end{array}\right)=\left( \begin{array}{ccccc}{ } & { } & { } & { }&{} \\ { \cdots} & {\cdots} & {\cdots} & { \boxed{b^{\overline{\alpha_{1}}}_{\overline{\beta_{ 1 }}}\cdots b^{\overline{\alpha_{r}}}_{\overline{\beta_{ r }}}\cdot a^{\underline{\beta_{ 1 }}}_{\underline{\alpha_{ 1 }}}\cdots a^{\underline{\beta_{ s }}}_{\underline{\alpha_{ s }}}}} &{\cdots}\\ { } & { } & { } & { } &{}\\ { } & { } & { } & { }&{}\\ { } & { } & { } & { }&{}\end{array}\right)\left( \begin{array}{ccccc}{\vdots} \\ {\vdots}  \\ {\vdots} \\ {\boxed{x_{\beta}}}\\ {\vdots}\end{array}\right)
\]
因此
\[
\overline{x}_{\alpha}=\sum_{ \beta\in \Gamma^{r}_{s} }\boxed{b^{\overline{\alpha_{1}}}_{\overline{\beta_{ 1 }}}\cdots b^{\overline{\alpha_{r}}}_{\overline{\beta_{ r }}}\cdot a^{\underline{\beta_{ 1 }}}_{\underline{\alpha_{ 1 }}}\cdots a^{\underline{\beta_{ s }}}_{\underline{\alpha_{ s }}}}\cdot \boxed{x_{\beta}}
\]

把\rthm{thm:chapter2张量空间的对偶空间}照搬过来，我们得到如下定理。

\begin{mythm}{}{}
$(\underbrace{V\otimes V\otimes \cdots\otimes V}_{k\text{个}})^{\ast}\cong \underbrace{V^{\ast}\otimes V^{\ast}\otimes \cdots\otimes V^{\ast}}_{k\text{个}}  $
\end{mythm}

再把\rdef{def:chapter2张量空间的配合运算}给出的配合运算搬到这里，我们就获得了$V^{0}_{k}$与$V^{k}_{0}$的之间配合。

\begin{mydef}{$V^{0}_{k}$与$V^{k}_{0}$的配合}{chapter2张量空间的配合运算2}
	对于$\forall g\in \underbrace{V^{\ast}\otimes V^{\ast}\otimes \cdots\otimes V^{\ast}}_{k\text{个}} $，$\forall x\in \underbrace{V\otimes V\otimes \cdots\otimes V}_{k\text{个}}$，我们约定
	\[
		\langle g,x\rangle\triangleq \langle{\color{red}{\iota(g)}},x\rangle
	\]
	其中$\iota:V^{\ast}\otimes \cdots\otimes V^{\ast}\longrightarrow (V\otimes \cdots\otimes V)^{\ast}$是由式(\ref{张量空间及其对偶空间的同构映射})所定义的同构映射.
\end{mydef}

相比于上一节定义的一般情形的配合运算，这里的配合运算在计算上要简单许多，因为我们需要处理的张量空间只由一个向量空间$V$生成。

\subsection{对称性与反对称性}

下面我们以
\[
	V^{r}_{0}=\underbrace{V\otimes V\otimes \cdots\otimes V}_{r\text{个}}
\]
也就是反变张量，作为讨论张量的对称性与反对称性的主要对象。

\begin{mydef}{置换群在$V^{r}_{0}$上的作用}{chapter2置换群的作用}
记$G(r)$为$\{1,2,\cdots,r\}$的置换群. 对$ \forall x\in V^{r}_{0}\;,\forall\sigma\in G(r)$，定义$\sigma(x)\in V^{r}_{0}$如下：
\[
\sigma(x)(\theta_{1},\theta_{2},\cdots,\theta_{r})\triangleq x \left( \theta_{\sigma(1)},\theta_{\sigma(2)},\cdots,\theta_{\sigma(r)}\right)
\]
其中$\theta_{i}\in V^{\ast},i=1,2,\cdots,r$. 可以验证这是一个群作用.
\end{mydef}

直接由上述定义出发，我们可以得到一个很有用的小结论。

\begin{myprop}{}{chapter2置换在单项式上的作用}
任取$\sigma\in G(r)$. 若$x=v_{1}\otimes v_{2}\otimes \cdots\otimes v_{r}$，则$\sigma(x)=v_{\sigma^{-1}(1)}\otimes v_{\sigma^{-1}(2)}\otimes\cdots\otimes v_{\sigma^{-1}(r)}$.
\end{myprop}

\begin{proof}

\end{proof}




\begin{mydef}{}{chapter2对称与反对称张量}
设$x\in V^{r}_{0}$.
\begin{itemize}
\item 若$\forall \sigma\in G(r)$，有$\sigma(x)=x$，则称$x$是对称的.
\item 若$\forall \sigma\in G(r)$，有$\sigma(x)=\operatorname{sgn}(\sigma)\cdot x$，则称$x$是反对称的.
\end{itemize}
\end{mydef}

\begin{myprop}{}{chapter2对称与反对称的等价条件}
	\begin{itemize}
	\item[(1)] 张量$x\in V^{r}_{0}$是对称的$\Leftrightarrow x(\cdots,\theta_{i},\cdots,\theta_{j}\cdots)=x(\cdots,		\theta_{j},\cdots,\theta_{i}\cdots),\forall i,j$.
	\item[(2)] 张量$x\in V^{r}_{0}$是反对称的$\Leftrightarrow x(\cdots,\theta_{i},\cdots,\theta_{j}\cdots)={\color{red}{-}}x(\cdots,\theta_{j},\cdots,\theta_{i}\cdots),\forall i,j$.
	\end{itemize}
\end{myprop}
\begin{proof}
	$(1)\Rightarrow\;$因为$x$为对称张量，所以$\forall \sigma\in G(r),\sigma(x)=x$. 特别地取$\sigma$为对换$(ij)$，即
\[
	\sigma=\left(\begin{array}{lllll}
		\cdots & i &\cdots & j & \cdots  \\
	    \cdots & j &\cdots & i & \cdots 
		\end{array}\right)
\]
则
\[
\begin{split}
	x(\cdots,\theta_{i},\cdots,\theta_{j}\cdots)=\sigma(x)(\cdots,\theta_{i},\cdots,\theta_{j}\cdots)&=x(\cdots,\theta_{\sigma(i)},\cdots,\theta_{\sigma(j)}\cdots)\\
	&=x(\cdots,\theta_{j},\cdots,\theta_{i}\cdots)
\end{split}
\]

$\Leftarrow\;$ 任取$\sigma\in G(r)$，则根据代数知识可知$\sigma$可以分解成若干个对换的乘积。而已知条件相当于是说$x$在对换作用下保持不变，从而对$x$连续施加若干次对换作用仍然保持不变，即$\sigma(x)=x$. 这表明$x$是对称的。

$(2)\Rightarrow\;$因为$x$为反对称张量，所以$\forall \sigma\in G(r),\sigma(x)=\operatorname{sgn}(\sigma) \cdot x$. 特别地取$\sigma$为对换$(ij)$，即
\[
	\sigma=\left(\begin{array}{lllll}
		\cdots & i &\cdots & j & \cdots  \\
	    \cdots & j &\cdots & i & \cdots 
		\end{array}\right)
\]
而对换总是奇置换，即$\operatorname{sgn}(\sigma)=-1$. 从而
\[
\begin{split}
x(\cdots,\theta_{i},\cdots,\theta_{j}\cdots)&=\operatorname{sgn}(\sigma)\cdot \sigma(x)(\cdots,\theta_{i},\cdots,\theta_{j}\cdots)\\
&=-x(\cdots,\theta_{\sigma(i)},\cdots,\theta_{\sigma(j)}\cdots)\\
&=-x(\cdots,\theta_{j},\cdots,\theta_{i}\cdots)
\end{split}
\]

$\Leftarrow\;$ 任取$\sigma\in G(r)$，则根据代数知识可知$\sigma$可以分解成若干个对换的乘积，不妨设它是$I_{\sigma}$个对换的乘积。已知条件是说每施加一次对换，值就改变一次符号，连续施加这$I_{\sigma}$个对换就有
\[
\begin{split}
x(\theta_{1},\theta_{2},\cdots,\theta_{r})&=(-1)^{I_{\sigma}}x(\theta_{\sigma(1)},\theta_{\sigma(2)},\cdots,\theta_{\sigma(r)})\\
&=\operatorname{sgn}(\sigma)\cdot \sigma(x)(\theta_{1},\theta_{2},\cdots,\theta_{r})
\end{split}
\]
因此$\sigma(x)=\operatorname{sgn}(\sigma)\cdot x$. 这表明$x$是反对称的。
\end{proof}

\begin{mydef}{}{chapter2系数的对称与反对称}
	取定$V$的一组基$\{e_{1} ,e_{2},\cdots,e_{n}\}$. 设$x=\tensor{x}{^{i_{1}i_{2}\cdots i_{r}}}e_{i_{1}}\otimes e_{i_{2}}\otimes \cdots \otimes e_{i_{r}}\in V^{r}_{0}$.
	\begin{itemize}
	\item 若$\forall \sigma\in G(r)$，有$\tensor{x}{^{i_{1}i_{2}\cdots i_{r}}}=\tensor{x}{^{i_{\sigma(1)}i_{\sigma(2)}\cdots i_{\sigma(r)}}}$，则称系数$\tensor{x}{^{i_{1}i_{2}\cdots i_{r}}}$关于各指标是对称的.
	\item 若$\forall \sigma\in G(r)$，有$\tensor{x}{^{i_{1}i_{2}\cdots i_{r}}}=\operatorname{sgn}(\sigma)\cdot \tensor{x}{^{i_{\sigma(1)}i_{\sigma(2)}\cdots i_{\sigma(r)}}}$，则称系数$\tensor{x}{^{i_{1}i_{2}\cdots i_{r}}}$关于各指标是反对称的.
	\end{itemize}
\end{mydef}

\begin{myprop}{}{chapter2系数对称与反对称的等价条件}
	取定$V$的一组基$\{e_{1} ,e_{2},\cdots,e_{n}\}$. 设$x=\tensor{x}{^{i_{1}i_{2}\cdots i_{r}}}e_{i_{1}}\otimes e_{i_{2}}\otimes \cdots \otimes e_{i_{r}}\in V^{r}_{0}$.
	\begin{itemize}
	\item[(1)] 系数$\tensor{x}{^{i_{1}i_{2}\cdots i_{r}}}$关于各指标是对称的$\Leftrightarrow \tensor{x}{^{\cdots i_{k}\cdots i_{l}\cdots }}=\tensor{x}{^{\cdots i_{l}\cdots i_{k}\cdots }},\forall k,l$.
	\item[(2)] 系数$\tensor{x}{^{i_{1}i_{2}\cdots i_{r}}}$关于各指标是反对称的$\Leftrightarrow \tensor{x}{^{\cdots i_{k}\cdots i_{l}\cdots }}={\color{red}{-}}\tensor{x}{^{\cdots i_{l}\cdots i_{k}\cdots }},\forall k,l$.
	\end{itemize}
\end{myprop}

\begin{proof}
	$(1)\Rightarrow$因为$\tensor{x}{^{i_{1}i_{2}\cdots i_{r}}}$关于各指标是对称的，所以$\forall \sigma\in G(r)$，$\tensor{x}{^{i_{1}i_{2}\cdots i_{r}}}=\tensor{x}{^{i_{\sigma(1)}i_{\sigma(2)}\cdots i_{\sigma(r)}}}$. 特别地取$\sigma$为对换$(kl)$，即
\[
	\sigma=\left(\begin{array}{lllll}
		\cdots & k &\cdots & l & \cdots  \\
	    \cdots & l &\cdots & k & \cdots 
		\end{array}\right)
\]
则$\tensor{x}{^{\cdots i_{k}\cdots i_{l}\cdots }}=\tensor{x}{^{\cdots i_{\sigma(k)}\cdots i_{\sigma(l)}\cdots }}=\tensor{x}{^{\cdots i_{l}\cdots i_{k}\cdots }}$.

$\Leftarrow$任取$\sigma\in G(r)$，首先根据代数知识可知$\sigma$可以分解成若干个对换的乘积，将$\tensor{x}{^{i_{1}i_{2}\cdots i_{r}}}$的指标按上述分解得到的对换逐次进行变换，最终得到$\tensor{x}{^{i_{\sigma(1)}i_{\sigma(2)}\cdots i_{\sigma(r)}}}$. 再根据已知条件即有$\tensor{x}{^{i_{1}i_{2}\cdots i_{r}}}=\tensor{x}{^{i_{\sigma(1)}i_{\sigma(2)}\cdots i_{\sigma(r)}}}$. 这表明$\tensor{x}{^{i_{1}i_{2}\cdots i_{r}}}$是对称的。

$(2)\Rightarrow$因为$x$为反对称张量，所以$\forall \sigma\in G(r),\tensor{x}{^{i_{1}i_{2}\cdots i_{r}}}=\operatorname{sgn}(\sigma)\cdot \tensor{x}{^{i_{\sigma(1)}i_{\sigma(2)}\cdots i_{\sigma(r)}}}$. 特别地取$\sigma$为对换$(kl)$，即
\[
	\sigma=\left(\begin{array}{lllll}
		\cdots & k &\cdots & l & \cdots  \\
	    \cdots & l &\cdots & k & \cdots 
		\end{array}\right)
\]
而对换总是奇置换，即$\operatorname{sgn}(\sigma)=-1$. 从而$\tensor{x}{^{\cdots i_{k}\cdots i_{l}\cdots }}=\operatorname{sgn}(\sigma)\cdot \tensor{x}{^{\cdots i_{\sigma(k)}\cdots i_{\sigma(l)}\cdots }}={\color{red}{-}}\tensor{x}{^{\cdots i_{l}\cdots i_{k}\cdots }}$.

$\Leftarrow$任取$\sigma\in G(r)$，首先根据代数知识可知$\sigma$可以分解成若干个对换的乘积，不妨设它是$I_{\sigma}$个对换的乘积。将$\tensor{x}{^{i_{1}i_{2}\cdots i_{r}}}$的指标按上述分解得到的对换逐次进行变换，最终得到$\tensor{x}{^{i_{\sigma(1)}i_{\sigma(2)}\cdots i_{\sigma(r)}}}$. 再根据已知条件，每交换两个指标就改变一次符号，所以$\tensor{x}{^{i_{1}i_{2}\cdots i_{r}}}=(-1)^{I_{\sigma}}\cdot \tensor{x}{^{i_{\sigma(1)}i_{\sigma(2)}\cdots i_{\sigma(r)}}}=\operatorname{sgn}(\sigma)\cdot \tensor{x}{^{i_{\sigma(1)}i_{\sigma(2)}\cdots i_{\sigma(r)}}}$. 这表明$\tensor{x}{^{i_{1}i_{2}\cdots i_{r}}}$是对称的。
\end{proof}

\begin{mythm}{}{chapter2张量对称等价于系数对称}
取定$V$的一组基$\{e_{1} ,e_{2},\cdots,e_{n}\}$. 设$x=\tensor{x}{^{i_{1}i_{2}\cdots i_{r}}}e_{i_{1}}\otimes e_{i_{2}}\otimes \cdots \otimes e_{i_{r}}\in V^{r}_{0}$.
\begin{itemize}
\item[(1)] 张量$x$是对称的$\Leftrightarrow$系数$\tensor{x}{^{i_{1}i_{2}\cdots i_{r}}}$关于各指标是对称的.
\item[(2)] 张量$x$是反对称的$\Leftrightarrow$系数$\tensor{x}{^{i_{1}i_{2}\cdots i_{r}}}$关于各指标是反对称的.
\end{itemize}
\end{mythm}
\begin{proof}
$(1)\Rightarrow$因为$x$对称，所以$\forall \sigma\in G(r),\sigma(x)=x$. 于是
\[
\begin{split}
\tensor{x}{^{i_{1}i_{2}\cdots i_{r}}}&=x(e^{\ast}_{i_{1}},e^{\ast}_{i_{2}},\cdots,e^{\ast}_{i_{r}})\\
&= \sigma(x)(e^{\ast}_{i_{1}},e^{\ast}_{i_{2}},\cdots,e^{\ast}_{i_{r}})\\
&= x(e^{\ast}_{i_{\sigma(1)}},e^{\ast}_{i_{\sigma(2)}},\cdots,e^{\ast}_{i_{\sigma(r)}})\\
&= \tensor{x}{^{i_{\sigma(1)}i_{\sigma(2)}\cdots i_{\sigma(r)}}}
\end{split}
\]

$\Leftarrow$注意到$x=\tensor{x}{^{\cdots i_{k}\cdots i_{l}\cdots }}\cdots \otimes e_{i_{k}}\otimes \cdots\otimes e_{i_{l}}\otimes\cdots=\tensor{x}{^{\cdots i_{l}\cdots i_{k}\cdots }}\cdots \otimes e_{i_{l}}\otimes \cdots\otimes e_{i_{k}}\otimes\cdots$.
所以我们直接计算
\[
\begin{split}
	x(\cdots,\theta_{l},\cdots,\theta_{k},\cdots)&=\tensor{x}{^{\cdots i_{l}\cdots i_{k}\cdots }}\cdots\theta_{l}(e_{i_{l}})\cdots\theta_{k}(e_{i_{k}})\cdots\\
&=\tensor{x}{^{\cdots i_{k}\cdots i_{l}\cdots }}\cdots  \theta_{k}(e_{i_{k}}) \cdots \theta_{l}(e_{i_{l}})\cdots\\
&=x(\cdots,\theta_{k},\cdots,\theta_{l},\cdots)
\end{split}
\]
所以$x$是对称的。

$(2)\Rightarrow$因为$x$反对称，所以$\forall \sigma\in G(r),\operatorname{sgn}(\sigma) \cdot \sigma(x)=x$. 于是
\[
\begin{split}
\tensor{x}{^{i_{1}i_{2}\cdots i_{r}}}&=x(e^{\ast}_{i_{1}},e^{\ast}_{i_{2}},\cdots,e^{\ast}_{i_{r}})\\
&=\operatorname{sgn}(\sigma) \cdot \sigma(x)(e^{\ast}_{i_{1}},e^{\ast}_{i_{2}},\cdots,e^{\ast}_{i_{r}})\\
&= \operatorname{sgn}(\sigma) \cdot x(e^{\ast}_{i_{\sigma(1)}},e^{\ast}_{i_{\sigma(2)}},\cdots,e^{\ast}_{i_{\sigma(r)}})\\
&=\operatorname{sgn}(\sigma) \cdot \tensor{x}{^{i_{\sigma(1)}i_{\sigma(2)}\cdots i_{\sigma(r)}}}
\end{split}
\]

$\Leftarrow$注意到$x=\tensor{x}{^{\cdots i_{k}\cdots i_{l}\cdots }}\cdots \otimes e_{i_{k}}\otimes \cdots\otimes e_{i_{l}}\otimes\cdots=\tensor{x}{^{\cdots i_{l}\cdots i_{k}\cdots }}\cdots \otimes e_{i_{l}}\otimes \cdots\otimes e_{i_{k}}\otimes\cdots$.
所以我们直接计算
\[
\begin{split}
	x(\cdots,\theta_{l},\cdots,\theta_{k},\cdots)&=\tensor{x}{^{\cdots i_{l}\cdots i_{k}\cdots }}\cdots\theta_{l}(e_{i_{l}})\cdots\theta_{k}(e_{i_{k}})\cdots\\
&={\color{red}{-}}\tensor{x}{^{\cdots i_{k}\cdots i_{l}\cdots }}\cdots  \theta_{k}(e_{i_{k}}) \cdots \theta_{l}(e_{i_{l}})\cdots\\
&={\color{red}{-}}x(\cdots,\theta_{k},\cdots,\theta_{l},\cdots)
\end{split}
\]
所以$x$是反对称的。
\end{proof}

\begin{mydef}{}{chapter2对称空间与反对称空间}
\begin{itemize}
\item $\operatorname{P}^{r}(V)\triangleq\{x |x\in V^{r}_{0}\text{且} x\text{为对称张量} \}$.
\item $\Lambda^{r}(V)\triangleq \{x |x\in V^{r}_{0}\text{且} x\text{为反对称张量} \}$. 另外约定$\Lambda^{1}(V)=V,\Lambda^{0}(V)=\mr$.
\end{itemize}
\end{mydef}

容易验证，$\operatorname{P}^{r}(V)$与$\Lambda^{r}(V)$作为向量空间都是$V^{r}_{0}$的子空间。也就是说，$V^{r}_{0}$中的对称(或反对称)张量的线性组合仍是对称(或反对称)张量。

\begin{mydef}{}{chapter2对称化与反对称化算子}
\begin{itemize}
\item 对$\forall x\in V^{r}_{0}$，令
         \[
			S_{r}(x)\triangleq \frac{1}{r!}\sum_{\sigma\in G(r)}\sigma(x)
		 \]
		 我们称$S_{r}$为对称化算子.
\item 对$\forall x\in V^{r}_{0}$，令
         \[
			A_{r}(x)\triangleq \frac{1}{r!}\sum_{\sigma\in G(r)}\operatorname{sgn}(\sigma)\cdot \sigma(x)
		 \]
		 我们称$A_{r}$为反对称化算子.
\end{itemize}
\end{mydef}

容易验证，$S_{r}$与$A_{r}$都是$V^{r}_{0}\longrightarrow V^{r}_{0}$的线性变换。

\begin{mythm}{}{chapter2对称化与反对称化性质}
	\begin{itemize}
	\item[(1)] $S_{r}\left( V^{r}_{0} \right)=\operatorname{P}^{r}(V)$.
	\item[(2)] $A_{r}\left( V^{r}_{0} \right)=\Lambda^{r}(V)$.
	\end{itemize}
\end{mythm}
\begin{proof}
首先由定义
\[
\begin{split}
	&\forall x\in \operatorname{P}^{r}(V),i.e.,\sigma(x)=x,\forall \sigma\in G(r)\Rightarrow {\color{red}{x=S_{r}(x)}}\Rightarrow \operatorname{P}^{r}(V)\subset S_{r}(V^{r}_{0})\\
	&\forall x\in \Lambda^{r}(V),i.e., \operatorname{sgn}(\sigma) \cdot\sigma(x)= x,\forall \sigma\in G(r) \Rightarrow {\color{red}{x=A_{r}(x)}}\Rightarrow \Lambda^{r}(V)\subset A_{r}(V^{r}_{0})
\end{split}
\]

另一方面，对$\forall S_{r}(x)\in S_{r}(V^{r}_{0})$. 任取$\tau\in G(r)$，则
\[
	\tau(S_{r}(x))=\tau \left( \frac{1}{r!}\sum_{\sigma\in G(r)}\sigma(x) \right)=\frac{1}{r!}\sum_{\sigma\in G(r)}\tau(\sigma(x))=S_{r}(x)
\]
最后一个等号成立是因为，当$\sigma$遍历$G(r)$时，相应的$\tau\sigma$也遍历$G(r)$. 因此$S_{r}(x)\in \operatorname{P}^{r}(V)$，从而$ S_{r}(V^{r}_{0})\subset \operatorname{P}^{r}(V)$.

同理，对$\forall A_{r}(x)\in A_{r}(V^{r}_{0})$. 任取$\tau\in G(r)$，则
\[
\begin{split}
\tau(A_{r}(x))&=\tau \left(\frac{1}{r!}\sum_{\sigma\in G(r)}\operatorname{sgn}(\sigma)\cdot \sigma(x)  \right)=\frac{1}{r!}\sum_{\sigma\in G(r)}{\color{red}{1}}\cdot \operatorname{sgn}(\sigma)\cdot \tau(\sigma(x))\\
&=\frac{1}{r!}\sum_{\sigma\in G(r)}{\color{red}{\operatorname{sgn}(\tau)\cdot \operatorname{sgn}(\tau)}}\cdot \operatorname{sgn}(\sigma)\cdot \tau(\sigma(x))\\
&={\color{red}{\operatorname{sgn}(\tau)}}\cdot \left( \frac{1}{r!}\sum_{\sigma\in G(r)}{\color{red}{\operatorname{sgn}(\tau)}}\cdot \operatorname{sgn}(\sigma)\cdot \tau(\sigma(x)) \right)\\
&={\color{red}{\operatorname{sgn}(\tau)}}\cdot \left( \frac{1}{r!}\sum_{\sigma\in G(r)}\operatorname{sgn}({\color{red}{\tau}}\sigma) \cdot \tau(\sigma(x)) \right)\\
&=\operatorname{sgn} (\tau)\cdot A_{r}(x)
\end{split}
\]
最后一个等号成立是因为，当$\sigma$遍历$G(r)$时，相应的$\tau\sigma$也遍历$G(r)$. 因此$A_{r}(x)\in \Lambda^{r}(V)$，从而$ A_{r}(V^{r}_{0})\subset \Lambda^{r}(V)$.
\end{proof}

以上所有这些讨论对于$V^{0}_{s}$也同样适用，因为从代数结构来看，$V$和$V^{\ast}$没有本质上的差别，只是具体解释不同而已。当以$V^{0}_{s}$为对象讨论对称性与反对称性时，我们约定如下记号。
\begin{mydef}{}{chapter2斜边张量的对称与反对称空间}
	\begin{itemize}
	\item $\operatorname{P}^{s}(V^{\ast})\triangleq\{x |x\in V^{0}_{s}\text{且} x\text{为对称张量} \}$.
	\item $\Lambda^{s}(V^{\ast})\triangleq \{x |x\in V^{0}_{s}\text{且} x\text{为反对称张量} \}$. 另外约定$\Lambda^{1}(V^{\ast})=V^{\ast},\Lambda^{0}(V^{\ast})=\mr$.
	\end{itemize}
\end{mydef}

在几何里面，反对称张量具有特殊而重要的地位，所以我们有必要给这一类张量重新取名。我们把$\Lambda^{r}(V)$称为\textbf{$r$ 次外向量空间}，其中的元素(即$V^{r}_{0}$中的反对称张量)称为\textbf{$r$ 次外向量}；把$\Lambda^{s}(V^{\ast})$称为\textbf{$s$次外形式空间}，其中的元素(即$V^{0}_{s}$中的反对称张量)称为\textbf{$s$ 次外形式}。

\subsection{对称积与外积}

上一节我们讨论了张量的对称性与反对称性。一个自然的问题是:
\begin{itemize}
\item 对于$x\in \operatorname{P}^{k}(V),y\in \operatorname{P}^{l}(V)$，它们的张量积$x\otimes 		  y$是否仍对称？
\item 对于$x\in \Lambda^{k}(V),y\in \Lambda^{l}(V)$，它们的张量积$x\otimes y$是否仍反对称？
\end{itemize}

一般而言，上面两个问题的答案都是否定的。但我们可以利用对称化算子与反对称化算子，重新定义出满足上述性质的张量运算。
\begin{mydef}{对称积}{chapter2对称积}
设$x\in \operatorname{P}^{k}(V)\;,y\in \operatorname{P}^{l}(V)$. 令
\[
xy\triangleq S_{k+l}\left( x\otimes y \right)=\frac{1}{(k+l)!}\sum_{\sigma\in G(k+l)}\sigma(x\otimes y)
\]
则$xy\in \operatorname{P}^{k+l}(V)$，称为$x$与$y$的对称积(symmetric product).
\end{mydef}

\begin{mydef}{外积}{chapter2外积}
	设$x\in \Lambda^{k}(V)\;,y\in \Lambda^{l}(V)$. 令
	\[
	x\wedge y\triangleq \frac{(k+l)!}{k!l!}A_{k+l}\left( x\otimes y \right)=\frac{1}{k!l!}\sum_{\sigma\in G(k+l)}\operatorname{sgn}(\sigma)\cdot \sigma(x\otimes y)
	\]
	则$x\wedge\in \Lambda^{k+l}(V)$，称为$x$与$y$的外积(exterior product)或楔积(wedge product).
\end{mydef}

明显可以看出，在形式上外积的定义比对称积稍显复杂——前面多了一个系数，这个系数你可以简单地理解为是为了方便计算才添上去的，更深层次的原因是为了排除和式中的重复项\footnote{参看Loring Tu P26 3.7节。}。如果说对称积是为了保证两个对称张量“乘起来”之后仍是对称张量而设计的，那么外积就是为了保证两个反对称张量“乘起来”之后仍是反对称张量而设计的，所以某种程度上我们也可以把外积叫作“反对称积”。

很多书上都不提对称积，因为这个概念有点鸡肋，但是我觉得还是有必要写一下。一方面与外积相对比，读者可能更好理解外积这一概念；另一方面，在黎曼几何中经常会把度量表达成$g_{ij}\dd u^{i}\dd u^{j}$的形式，这里的$\dd u^{i}\dd u^{j}$在对称积意义下才是严格的表达式。

\begin{myprop}{对称积的运算规律}{chapter2对称积的运算规律}
	设$x,w\in \operatorname{P}^{k}(V)\;,y\in \operatorname{P}^{l}(V)\;,z\in \operatorname{P}^{h}(V)\;,k_{1}\in\mr\;,k_{2}\in \mr$.
	\begin{itemize}
	\item[(1)] 分配律：\[
						\left \{ \begin{array} { l }
							{ (k_{1}x+k_{2}w)z=k_{1}xz+k_{2}wz } \\
							{  z(k_{1}x+k_{2}w)=k_{1}zx+k_{2}zw } 
	   					\end{array} \right.
					   \]
	\item[(2)] 交换律：$xy=yx$.
	\item[(3)] 结合律：$(xy)z=x(yz)\xlongequal{\text{记作}}xyz$.
	\end{itemize}
\end{myprop}	

对称积的这些性质我们就不证明了，读者可以自行验证，我们最关心的还是外积。

\begin{myprop}{外积的运算规律}{chapter2外积的运算规律}
	设$x,w\in \Lambda^{k}(V)\;,y\in \Lambda^{l}(V)\;,z\in \Lambda^{h}(V)\;,k_{1}\in\mr\;,k_{2}\in \mr$.
	\begin{itemize}
	\item[(1)] 分配律：\[
						\left \{ \begin{array} { l }
							{ (k_{1}x+k_{2}w)\wedge z=k_{1}x\wedge z+k_{2}w\wedge z } \\
							{  z\wedge (k_{1}x+k_{2}w)=k_{1}z\wedge x+k_{2}z\wedge w } 
	   					\end{array} \right.
					   \]
	\item[(2)] 反交换律：$x\wedge y=(-1)^{kl} y\wedge x$.
	\item[(3)] 结合律：$(x\wedge y)\wedge z=\frac{(k+l+h)!}{k!l!h!}A_{k+l+h}\left( x\otimes y \otimes z\right)=x\wedge (y\wedge z)\xlongequal{\text{记作}}x\wedge y\wedge z $.
	\end{itemize}
\end{myprop}
\begin{proof}
(1)因为反对称化算子是线性的，且张量积运算$x\otimes y$具有分配律，所以外积运算自然也有分配律。

(2)根据定义$x\wedge y=\frac{1}{k!l!}\sum_{\sigma\in G(k+l)}\operatorname{sgn}(\sigma)\cdot \sigma(x\otimes y)$. 所以要先看$x\otimes y$与$y\otimes x$之间的关系。显然$x\otimes y\neq y\otimes x$. 这是因为它们定义域中的自变量顺序不同，具体来说，即
\[
\begin{split}
x\otimes y(\theta_{1},\cdots,\theta_{k},\theta_{k+1},\cdots,\theta_{k+l})&=x(\theta_{1},\cdots,\theta_{k})\cdot y(\theta_{k+1},\cdots,\theta_{k+l})\\
&=y(\theta_{k+1},\cdots,\theta_{k+l})\cdot x(\theta_{1},\cdots,\theta_{k})\\
&=y\otimes x(\theta_{k+1},\cdots,\theta_{k+l},\theta_{1},\cdots,\theta_{k})
\end{split}
\]
我们需要用一个置换把自变量的顺序重新调一下。根据上式，取$\tau\in G(k+l)$如下：
\[
	\left(\begin{array}{llllll}
		\theta_{1} & \cdots  & \theta_{l} & \theta_{l+1} & \cdots & \theta_{k+l} \\
		\theta_{k+1} & \cdots & \theta_{k+l} & \theta_{1} & \cdots  & \theta_{k} 
		\end{array}\right)\Leftrightarrow \tau=\left(\begin{array}{llllll}
			1 & \cdots  & l & l+1 & \cdots & k+l \\
			k+1 & \cdots & k+l & 1 & \cdots  & k 
			\end{array}\right)
\]
于是
\[
\begin{split}
y\otimes x(\theta_{k+1},\cdots,\theta_{k+l},\theta_{1},\cdots,\theta_{k})
&=y\otimes x(\theta_{\tau(1)},\cdots,\theta_{\tau(l)},\theta_{\tau(l+1)},\cdots,\theta_{\tau(k+l)})\\
&=\tau(y\otimes x)(\theta_{1},\cdots,\theta_{k},\theta_{k+1},\cdots,\theta_{k+l})
\end{split}
\]
从而$x\otimes y=\tau(y\otimes x)$. 因此
\[
\begin{split}
	x\wedge y&=\frac{1}{k!l!}\sum_{\sigma\in G(k+l)}\operatorname{sgn}(\sigma)\cdot \sigma(x\otimes y)\\
	&=\frac{1}{k!l!}\sum_{\sigma\in G(k+l)}\operatorname{sgn}(\sigma)\cdot \sigma({\color{red}{\tau}}(y\otimes x))\\
	&={\color{red}{\operatorname{sgn}(\tau)}}\cdot \frac{1}{k!l!}\sum_{\sigma\in G(k+l)}\operatorname{sgn}(\sigma)\cdot {\color{red}{\operatorname{sgn}(\tau)}}\cdot \sigma({\color{red}{\tau}}(y\otimes x))\\
	&={\color{red}{\operatorname{sgn}(\tau)}}\cdot \frac{1}{k!l!}\sum_{\sigma\in G(k+l)}\operatorname{sgn}(\sigma{\color{red}{\tau}})\cdot (\sigma{\color{red}{\tau}})(y\otimes x)\\
	&={\color{red}{\operatorname{sgn}(\tau)}}\cdot y\wedge x
\end{split}
\]
最后数一下置换$\tau$中逆序对的个数，得到逆序数为$kl$. 因此$\operatorname{sgn}(\tau)=(-1)^{kl}$. 也就得到了$x\wedge y=(-1)^{kl} y\wedge x$.

(3)证明有点繁琐，不想证了。
\end{proof}
以上性质是我们在做外积运算时最主要的工具。特别地对于1次外向量，我们有如下推论。
\begin{myprop}{}{chapter21次外向量的外积}
	若$x,y\in \Lambda^{1}(V)=V$，则$x\wedge y=-y\wedge x$. 特别地，$x\wedge x=0$.
\end{myprop}
\begin{proof}
根据反对称性立得$x\wedge y=-y\wedge x$. 注意到$x\wedge x=-x\wedge x$，从而$x\wedge x=0$.
\end{proof}

\begin{myprop}{反对称性}{chapter2外积的反对称性}
设$e_{i_{1}},e_{i_{2}},\cdots,e_{i_{k}}\in \Lambda^{1}(V)=V$，则对$\forall \sigma\in G(k)$有
\[
	e_{i_{1}}\wedge e_{i_{2}}\wedge \cdots\wedge e_{i_{k}}=\operatorname{sgn}(\sigma)\cdot e_{i_{\sigma(1)}}\wedge e_{i_{\sigma(2)}}\wedge \cdots\wedge e_{i_{\sigma(k)}}
\]
\end{myprop}
\begin{proof}
	根据代数知识可知$\sigma$可以分解成若干个对换的乘积，不妨设它是$I_{\sigma}$个对换的乘积。将$e_{i_{1}}\wedge e_{i_{2}}\wedge \cdots\wedge e_{i_{k}}$的指标按上述分解得到的对换逐次进行交换，最终得到$e_{i_{\sigma(1)}}\wedge e_{i_{\sigma(2)}}\wedge \cdots\wedge e_{i_{\sigma(k)}}$. 再根据\rprop{prop:chapter21次外向量的外积}，每交换两个指标就改变一次符号，所以
	\[
	\begin{split}
		e_{i_{1}}\wedge e_{i_{2}}\wedge \cdots\wedge e_{i_{k}}&=(-1)^{I_{\sigma}}\cdot e_{i_{\sigma(1)}}\wedge e_{i_{\sigma(2)}}\wedge \cdots\wedge e_{i_{\sigma(k)}}\\
		&=\operatorname{sgn}(\sigma)\cdot e_{i_{\sigma(1)}}\wedge e_{i_{\sigma(2)}}\wedge \cdots\wedge e_{i_{\sigma(k)}}
	\end{split}
	\]
\end{proof}

利用结合律，我们还可以给出多个外向量做外积的计算表达式。
\begin{myprop}{}{chapter2多个外向量的外积}
设$x_{1}\in \Lambda^{r_{1}}(V),x_{2}\in \Lambda^{r_{2}}(V),\cdots,x_{k}\in \Lambda^{r_{k}}(V)$. 则有：
\[
	x_{1}\wedge x_{2}\wedge\cdots\wedge x_{k}=\frac{(r_{1}+r_{2}+\cdots+r_{k})!}{r_{1}!r_{2}!\cdots r_{k}!}A_{r_{1}+r_{2}+\cdots+r_{k}}\left( x_{1}\otimes x_{2}\otimes\cdots\otimes x_{k}\right) 
\]
\end{myprop}

这个性质在反对称张量的外积与张量积之间架起了桥梁。特别地，若$\{e_{1} ,e_{2},\cdots,e_{n}\}$是$V$的一组基，根据\rprop{prop:chapter2多个外向量的外积}我们有
\begin{equation}
	e_{i_{1}}\wedge e_{i_{2}}\wedge \cdots\wedge e_{i_{r}}=r!A_{r}\left( e_{i_{1}}\otimes e_{i_{2}}\otimes\cdots\otimes e_{i_{r}}\right) 
	\label{1次外向量的外积与张量积}
\end{equation}
由此出发我们可以弄清楚向量空间$\lambda^{r}(V)$的结构，并给出它的一组基。
\begin{mythm}{}{chapter2反对称空间的结构}
	设$\operatorname{dim}(V)=n$. 取定$V$的一组基$\{e_{1} ,e_{2},\cdots,e_{n}\}$. 考虑$r$次外向量空间$\Lambda^{r}(V)$.
\begin{itemize}
\item[(1)] 若$r>n$，则$\Lambda^{r}(V)=\{0\}$.
\item[(2)] 若$r\leqslant n$，则$\{e_{i_{1}}\wedge e_{i_{2}}\wedge \cdots\wedge e_{i_{r}} | 1\leqslant i_{1}<i_{2}<\cdots<i_{r}\leqslant n \}$是$\Lambda^{r}(V)$的一组基，并且对$\forall x\in \Lambda^{r}(V)$，有
\[
x=\sum_{1\leqslant i_{1}<i_{2}<\cdots<i_{r}\leqslant n}\tensor{x}{^{i_{1}i_{2}\cdots i_{r}}}e_{i_{1}}\wedge e_{i_{2}}\wedge \cdots\wedge e_{i_{r}}
\]
其中$\tensor{x}{^{i_{1}i_{2}\cdots i_{r}}}$关于各指标是反对称的.
\end{itemize}
\end{mythm}
\begin{proof}
任取$x\in \Lambda^{r}(V)\subset V^{r}_{0}$，则可设$x=\tensor{x}{^{i_{1}i_{2}\cdots i_{r}}}e_{i_{1}}\otimes e_{i_{2}}\otimes \cdots \otimes e_{i_{r}}$. 因为$x$是反对称的，根据\rthm{thm:chapter2张量对称等价于系数对称}与\rthm{thm:chapter2对称化与反对称化性质}可知，系数$x^{i_{1}i_{2}\cdots i_{r}}$关于各指标是反对称的，并且$A_{r}(x)=x$. 
\begin{equation}
\begin{split}
	x=A_{r}(x)&=A_{r}(\tensor{x}{^{i_{1}i_{2}\cdots i_{r}}}e_{i_{1}}\otimes e_{i_{2}}\otimes \cdots \otimes e_{i_{r}})\\
	&=\tensor{x}{^{i_{1}i_{2}\cdots i_{r}}}A_{r}(e_{i_{1}}\otimes e_{i_{2}}\otimes \cdots \otimes e_{i_{r}})\\
	&\xlongequal{\requ{1次外向量的外积与张量积}}\frac{1}{r!}\tensor{x}{^{i_{1}i_{2}\cdots i_{r}}}e_{i_{1}}\wedge e_{i_{2}}\wedge \cdots\wedge e_{i_{r}}
\end{split}
\label{外向量的表达式}
\end{equation}
这里求和指标$i_{1},i_{2},\cdots,i_{r}$各自分别跑遍$\{1,2,\cdots,n \}$. 

(1)当$r>n$时，由鸽巢原理，$e_{i_{1}}\wedge e_{i_{2}}\wedge \cdots\wedge e_{i_{r}}$中必有重复项，则$x=0\Rightarrow \Lambda^{r}(V)=\{0\}$. 

(2)当$r\leqslant n$时，注意到\requ{外向量的表达式}所表达的和式中，并非每一项都对求和有贡献——只有那些指标$i_{1},i_{2},\cdots,i_{r}$两两不同的才对求和有贡献，否则该项等于0. 我们把这些等于0的项扔掉得到
\[
\begin{split}
x&=\frac{1}{r!}\tensor{x}{^{i_{1}i_{2}\cdots i_{r}}}e_{i_{1}}\wedge e_{i_{2}}\wedge \cdots\wedge e_{i_{r}}\\
&=\sum_{\substack{ 1\leqslant i_{1},i_{2},\cdots,i_{r}\leqslant n\\ \text{\tiny 且两两不同}}}\left[\frac{1}{r!}\tensor{x}{^{i_{1}i_{2}\cdots i_{r}}}e_{i_{1}}\wedge e_{i_{2}}\wedge \cdots\wedge e_{i_{r}}\right]\\
&=\sum_{1\leqslant j_{1}<j_{2}<\cdots<j_{n}\leqslant n}\left[ \sum_{\sigma\in G(r)}\frac{1}{r!}\tensor{x}{^{j_{\sigma(1)}j_{\sigma(2)}\cdots j_{\sigma(r)}}}e_{j_{\sigma(1)}}\wedge e_{j_{\sigma(2)}}\wedge \cdots\wedge e_{j_{\sigma(r)}}\right]
\end{split}
\]
因为$x^{i_{1}i_{2}\cdots i_{r}}$关于各指标是反对称的，根据\rdef{def:chapter2系数的对称与反对称}可知$\tensor{x}{^{j_{\sigma(1)}j_{\sigma(2)}\cdots j_{\sigma(r)}}}=\operatorname{sgn}(\sigma)\cdot \tensor{x}{^{j_{1}j_{2}\cdots j_{r}}}$；另外根据\rprop{prop:chapter2外积的反对称性}我们又有$e_{j_{\sigma(1)}}\wedge e_{j_{\sigma(2)}}\wedge \cdots\wedge e_{j_{\sigma(r)}}=\operatorname{sgn}(\sigma)\cdot e_{j_{1}}\wedge e_{j_{2}}\wedge \cdots\wedge e_{j_{r}}$. 从而
\[
\begin{split}
x&=\sum_{1\leqslant j_{1}<j_{2}<\cdots<j_{n}\leqslant n}\left[ \sum_{\sigma\in G(r)}\frac{1}{r!}\tensor{x}{^{j_{\sigma(1)}j_{\sigma(2)}\cdots j_{\sigma(r)}}}e_{j_{\sigma(1)}}\wedge e_{j_{\sigma(2)}}\wedge \cdots\wedge e_{j_{\sigma(r)}}\right]\\
&=\sum_{1\leqslant j_{1}<j_{2}<\cdots<j_{n}\leqslant n}\left[ \sum_{\sigma\in G(r)}\frac{1}{r!}\cdot {\color{red}{\left(\operatorname{sgn}(\sigma)\right)^{2}}}\cdot \tensor{x}{^{j_{1}j_{2}\cdots j_{r}}}e_{j_{1}}\wedge e_{j_{2}}\wedge \cdots\wedge e_{j_{r}}\right]\\
&=\sum_{1\leqslant j_{1}<j_{2}<\cdots<j_{n}\leqslant n}\left[r!\cdot \left( \frac{1}{r!}\cdot {\color{red}{1}}\cdot \tensor{x}{^{j_{1}j_{2}\cdots j_{r}}}e_{j_{1}}\wedge e_{j_{2}}\wedge \cdots\wedge e_{j_{r}}\right)\right]\\
&=\sum_{1\leqslant j_{1}<j_{2}<\cdots<j_{n}\leqslant n}\tensor{x}{^{j_{1}j_{2}\cdots j_{r}}}e_{j_{1}}\wedge e_{j_{2}}\wedge \cdots\wedge e_{j_{r}}
\end{split}
\]
这表明$\forall x\in \lambda^{r}(V)$，可由$\{e_{i_{1}}\wedge e_{i_{2}}\wedge \cdots\wedge e_{i_{r}} | 1\leqslant i_{1}<i_{2}<\cdots<i_{r}\leqslant n \}$线性表出。

下面再证线性无关性。首先导出$e_{i_{1}}\wedge e_{i_{2}}\wedge \cdots\wedge e_{i_{r}} $的求值公式。根据式\requ{1次外向量的外积与张量积}，我们有
\[
	e_{i_{1}}\wedge e_{i_{2}}\wedge \cdots\wedge e_{i_{r}}=r!A_{r}\left( e_{i_{1}}\otimes e_{i_{2}}\otimes\cdots\otimes e_{i_{r}}\right)=\sum_{\sigma\in G(r)}\operatorname{sgn}(\sigma)\cdot \sigma\left( e_{i_{1}}\otimes e_{i_{2}}\otimes\cdots\otimes e_{i_{r}}\right)
\]
于是对任意$\theta_{1},\theta_{2},\cdots,\theta_{r}\in V^{\ast}$，则
\begin{equation}
\begin{split}
	&e_{i_{1}}\wedge e_{i_{2}}\wedge \cdots\wedge e_{i_{r}}(\theta_{1},\theta_{2},\cdots,\theta_{r})\\
	&=\sum_{\sigma\in G(r)}\operatorname{sgn}(\sigma)\cdot \sigma\left( e_{i_{1}}\otimes e_{i_{2}}\otimes\cdots\otimes e_{i_{r}}\right) (\theta_{1},\theta_{2},\cdots,\theta_{r})\\
	&=\sum_{\sigma\in G(r)}\operatorname{sgn}(\sigma)\cdot \left( e_{i_{1}}\otimes e_{i_{2}}\otimes\cdots\otimes e_{i_{r}}\right) (\theta_{\sigma(1)},\theta_{\sigma(2)},\cdots,\theta_{\sigma(r)})\\
	&=\sum_{\sigma\in G(r)}\operatorname{sgn}(\sigma)\cdot \theta_{\sigma(1)}(e_{i_{1}})\cdot \theta_{\sigma(2)}(e_{i_{2}})\cdots\theta_{\sigma(r)}(e_{i_{r}})\\
	&=\left|\begin{array}{cccc}\theta_{1}(e_{i_{1}}) & \theta_{2}(e_{i_{1}}) & \cdots & \theta_{r}(e_{i_{1}}) \\ \theta_{1}(e_{i_{2}}) & \theta_{2}(e_{i_{2}}) & \cdots & \theta_{r}(e_{i_{2}}) \\ \vdots & \vdots & \ddots& \vdots \\ \theta_{1}(e_{i_{r}}) & \theta_{2}(e_{i_{r}}) & \cdots & \theta_{r}(e_{i_{r}})\end{array}\right|
\end{split}
\label{外积求值公式}
\end{equation}
上式称为\textbf{求值公式}。特别地，$e_{1}\wedge e_{2}\wedge \cdots\wedge e_{n}(e^{\ast}_{1},e^{\ast}_{2},\cdots,e^{\ast}_{n})=1$. 这表明${\color{red}{e_{1}\wedge e_{2}\wedge \cdots\wedge e_{n}\neq 0}}$.

假设$\{e_{i_{1}}\wedge e_{i_{2}}\wedge \cdots\wedge e_{i_{r}} | 1\leqslant i_{1}<i_{2}<\cdots<i_{r}\leqslant n \}$线性相关，则存在不全为0的一组数$\{ \tensor{a}{^{i_{1}i_{2}\cdots i_{r}}}\}$，使得
\begin{equation*}
	\sum_{1\leqslant i_{1}<i_{2}<\cdots<i_{r}\leqslant n}\tensor{a}{^{i_{1}i_{2}\cdots i_{r}}}e_{i_{1}}\wedge e_{i_{2}}\wedge \cdots\wedge e_{i_{r}}=0 \eqno{(\star)}
	\label{线性相关}
\end{equation*}

不妨设$\tensor{a}{^{j_{1}j_{2}\cdots j_{r}}}\neq 0$. 对指标$\{j_{1},j_{2},\cdots,j_{r}\}$，我们将其补全，即存在$\{j_{r+1},j_{r+2},\cdots,j_{n}\}$，使得$\{j_{1},j_{2},\cdots,j_{r},j_{r+1},j_{r+2},\cdots,j_{n}\}$是$\{1,2,\cdots,n\}$的一个排列。

注意到$\{j_{r+1},j_{r+2},\cdots,j_{n}\}$只跟$\{j_{1},j_{2},\cdots,j_{r}\}$不重复，与其他的$\{i_{1},i_{2},\cdots,i_{r}\}$至少有一个指标会重复，因此考虑将$e_{j_{r+1}}\wedge e_{j_{r+2}}\wedge\cdots\wedge e_{j_{n}}$与式$(\star)$做外积运算，这样就只剩下
\[
	\tensor{a}{^{j_{1}j_{2}\cdots j_{r}}}e_{j_{1}}\wedge e_{j_{2}}\wedge \cdots\wedge e_{j_{r}}\wedge e_{j_{r+1}}\wedge e_{j_{r+2}}\wedge\cdots\wedge e_{j_{n}}=0
\]
将外积式中的指标重新排列为自然顺序，则上式左边最多会产生一个符号的差别，即
\[
\pm \tensor{a}{^{j_{1}j_{2}\cdots j_{r}}}e_{1}\wedge e_{2}\wedge \cdots\wedge e_{n}=0
\]
而我们已经知道${\color{red}{e_{1}\wedge e_{2}\wedge \cdots\wedge e_{n}\neq 0}}$，因此只能$\tensor{a}{^{j_{1}j_{2}\cdots j_{r}}}=0$，而这与$\tensor{a}{^{j_{1}j_{2}\cdots j_{r}}}\neq 0$相矛盾！这表明$\{e_{i_{1}}\wedge e_{i_{2}}\wedge \cdots\wedge e_{i_{r}} | 1\leqslant i_{1}<i_{2}<\cdots<i_{r}\leqslant n \}$线性无关。
\end{proof}

值得注意的是，证明中出现的式\requ{外向量的表达式}比定理给出的表达式更常用，读者要记住。

易知满足$1\leqslant i_{1}<i_{2}<\cdots <i_{r}\leqslant n$的有序对$(i_{1},i_{2},\cdots,i_{r})$一共有$\binom{n}{r}$个，从而上述定理的一个简单推论是：$\operatorname{dim}(\Lambda^{r}(V))=\binom{n}{r}$. 特别地，$\operatorname{dim}(\Lambda^{n}(V))=\binom{n}{n}=1$，这表明$\Lambda^{n}(V)$中的元素都呈现为线性形式，即形如$k\cdot e_{1}\wedge e_{2}\wedge\cdots\wedge e_{n}$.

为了以后能更方便地描述$\Lambda^{r}(V)$，我们给它量身定制一套带有顺序的多重指标语言。若数组$I=(i_{1},i_{2},\cdots,i_{r})$的每个分量都取正整数，并且满足$1\leqslant i_{1}<i_{2}<\cdots <i_{r}\leqslant n$，则称$I$为有序多重指标，并定义有序多重指标集如下
\[
Q(r,n)\triangleq \{I=(i_{1},i_{2},\cdots,i_{r}) | 1\leqslant i_{1}<i_{2}<\cdots <i_{r}\leqslant n  \}
\]
显然$|Q(r,n)|=\binom{n}{r}$. 与一般的多重指标集一样，$Q(r,n)$的元素可以按照字典排序法从小到大一个不漏地排起来成为有序集合，从而实现与数集$\{1,2,\cdots,\binom{n}{r} \}$中的数字一一对应。正是因为存在这种一一对应的关系，以后我们也直接用有序多重指标来表示与它对应的那个正整数。这一点和\ref{多重线性映射}节介绍的一般多重指标完全一样。

利用有序多重指标，我们可以很简洁地描述\rthm{thm:chapter2反对称空间的结构}所给出的$\Lambda^{r}(V)$的基。\textbf{约定：}
\[
e^{\wedge}_{I}\triangleq e_{i_{1}}\wedge e_{i_{2}}\wedge \cdots\wedge e_{i_{r}}\;,\text{\kaishu 其中}I=(i_{1},i_{2},\cdots,i_{r})\in Q(r,n)
\]
于是$\{e_{i_{1}}\wedge e_{i_{2}}\wedge \cdots\wedge e_{i_{r}} | 1\leqslant i_{1}<i_{2}<\cdots<i_{r}\leqslant n \}$就可以简写为$\{e^{\wedge}_{I} |I\in Q(r,n) \}$.

\rthm{thm:chapter2反对称空间的结构}给出了$r$次外向量空间的基底的构造，接下来自然就要考虑不同基底之间是如何变换的问题，也就是给出不同基底之间的过渡矩阵。

首先把线性代数中行列式的子式这一概念搬过来。给定$n$阶方阵$A=(a_{ij})_{n\times n}$，从它的行列式$\operatorname{det}(A)$中取出第$i_{1},i_{2},\cdots,i_{r}$行与第$j_{1},j_{2},\cdots,j_{r}$列的交叉位置上的元素，其中$1\leqslant i_{1}<i_{2}<\cdots <i_{r}\leqslant n,1\leqslant j_{1}<j_{2}<\cdots <j_{r}\leqslant n$，并按这些元素在原来行列式$\operatorname{det}(A)$中的次序排成一个$r$阶行列式，它称为行列式$\operatorname{det}(A)$的\textbf{$r$阶子式}，记为
\[
A\left(\begin{array}{cccc}i_{1} & i_{2} & \cdots & i_{r} \\j_{1} & j_{2} & \cdots & j_{r}\end{array}\right)=\left|\begin{array}{cccc}
	a_{i_{1} j_{1}} & a_{i_{1} j_{2}} & \cdots & a_{i_{1} j_{r}} \\
	a_{i_{2} j_{1}} & a_{i_{2} j_{2}} & \cdots & a_{i_{2} j_{r}} \\
	\vdots & \vdots & \ddots & \vdots \\
	a_{i_{r} j_{1}} & a_{i_{r} j_{2}} & \cdots & a_{i_{r}} j_{r}
	\end{array}\right|
\]

关于行列式的子式，有一条十分著名的性质，也是我们下一个定理需要用到的。
\begin{mylem}{Binet-Cauchy公式}{chapter2Binet-Cauchy公式}
设$A$是一个$m\times n$矩阵，$B$是一个$n\times l$矩阵，并且记$C=AB$. 若$r\leqslant n$，则矩阵$C$的$r$阶子式
\[
C\left(\begin{array}{cccc}i_{1} & i_{2} & \cdots & i_{r} \\j_{1} & j_{2} & \cdots & j_{r}\end{array}\right)=\sum\limits_{1\leqslant k_{1}<\cdots<k_{r}\leqslant n}A\left(\begin{array}{cccc}i_{1} & i_{2} & \cdots & i_{r} \\k_{1} & k_{2} & \cdots & k_{r}\end{array}\right)\cdot B\left(\begin{array}{cccc}k_{1} & k_{2} & \cdots & k_{r} \\j_{1} & j_{2} & \cdots & j_{r}\end{array}\right)
\]
\end{mylem}
\begin{proof}
参阅李炯生《线性代数》。我看的是李清老师重排的版本，在这一版本中此结论的证明位于第94页定理3.2.3.
\end{proof}

\begin{mythm}{}{chapter2反对称空间的基过渡矩阵}
	已知在向量空间$V$中，从基$(e_{1},e_{2},\cdots,e_{n})$到基$(\overline{e}_{1},\overline{e}_{2},\cdots,\overline{e}_{n})$的过渡矩阵为$A$. 在$r$次外向量空间$\Lambda^{r}(V)$中，若从基$\{e^{\wedge}_{\beta} \}_{\beta\in Q(r,n) } $到基$\{ \overline{e}^{\wedge}_{\alpha}\}_{\alpha\in Q(r,n)}$的过渡矩阵为$C$. 则
	\begin{itemize}
	\item $C$可由$\operatorname{det}(A)$的$r$阶子式确定：$\forall I=(i_{1},\cdots,i_{r}),J=(j_{1},\cdots,j_{r})\in Q(r,n)$
	\[
	C_{IJ}=A\left(\begin{array}{cccc}i_{1} & i_{2} & \cdots & i_{r} \\j_{1} & j_{2} & \cdots & j_{r}\end{array}\right)
	\]
	其中$C_{IJ}$表示矩阵$C$的第$I$行第$J$列处的元素.
	\item $C^{-1}$可由$\operatorname{det}(A^{-1})$的$r$阶子式确定：$\forall \alpha=(\alpha_{1},\cdots,\alpha_{r}),\beta=(\beta_{1},\cdots,\beta_{r})\in Q(r,n)$
	\[
	C^{-1}_{\alpha\beta}=A^{-1}\left(\begin{array}{cccc}\alpha_{1} & \alpha_{2} & \cdots & \alpha_{r} \\\beta_{1} & \beta_{2} & \cdots & \beta_{r}\end{array}\right)
	\]
	其中$C^{-1}_{\alpha\beta}$表示矩阵$C^{-1}$的第$\alpha$行第$\beta$列处的元素.
	\end{itemize}
\end{mythm}

\begin{proof}
因为从基$(e_{1},e_{2},\cdots,e_{n})$到基$(\overline{e}_{1},\overline{e}_{2},\cdots,\overline{e}_{n})$的过渡矩阵为$A$，不妨设$A=(a_{ij})_{n\times n}$. 则$\overline{e}_{j}=\sum\limits_{i=1}^{n}a_{ij}e_{i}$. 对任意$1\leqslant j_{1}<j_{2}<\cdots <j_{r}\leqslant n$，我们来计算$\overline{e}_{j_{1}}\wedge \overline{e}_{j_{2}}\wedge\cdots\wedge \overline{e}_{j_{r}}$.
\[
\begin{split}
	\overline{e}_{j_{1}}\wedge \overline{e}_{j_{2}}\wedge\cdots\wedge \overline{e}_{j_{r}}&=\left( \sum_{k_{1}=1}^{n}a_{k_{1}j_{1}}e_{k_{1}} \right)\wedge \left( \sum_{k_{2}=1}^{n}a_{k_{2}j_{2}}e_{k_{2}} \right)\wedge\cdots\wedge \left( \sum_{k_{r}=1}^{n}a_{k_{r}j_{r}}e_{k_{r}} \right)\\
	&=\sum_{1\leqslant k_{1},\cdots,k_{r}\leqslant n}a_{k_{1}j_{1}}a_{k_{2}j_{2}}\cdots a_{k_{r}j_{r}}\cdot e_{k_{1}}\wedge e_{k_{2}}\wedge \cdots \wedge e_{k_{r}}\\
	&=\sum_{1\leqslant i_{1}<\cdots< i_{r}\leqslant n}\left[\sum_{\sigma\in G(r)}a_{i_{\sigma(1)}j_{1}}a_{i_{\sigma(2)}j_{2}}\cdots a_{i_{\sigma(r)}j_{r}}\cdot e_{i_{\sigma(1)}}\wedge e_{i_{\sigma(2)}}\wedge \cdots e_{i_{\sigma(r)}}\right]\\
	&=\sum_{1\leqslant i_{1}<\cdots< i_{r}\leqslant n}\left[\sum_{\sigma\in G(r)}\operatorname{sgn}(\sigma)\cdot  a_{i_{\sigma(1)}j_{1}}a_{i_{\sigma(2)}j_{2}}\cdots a_{i_{\sigma(r)}j_{r}}\cdot  e_{i_{1}}\wedge e_{i_{2}}\wedge \cdots e_{i_{r}}\right]\\
	&=\sum_{1\leqslant i_{1}<\cdots< i_{r}\leqslant n}e_{i_{1}}\wedge e_{i_{2}}\wedge \cdots e_{i_{r}}\left[\sum_{\sigma\in G(r)}\operatorname{sgn}(\sigma)\cdot  a_{i_{\sigma(1)}j_{1}}a_{i_{\sigma(2)}j_{2}}\cdots a_{i_{\sigma(r)}j_{r}}  \right]
\end{split}
\]
注意到
\[
\begin{split}
	\sum_{\sigma\in G(r)}\operatorname{sgn}(\sigma)\cdot  a_{i_{\sigma(1)}j_{1}}a_{i_{\sigma(2)}j_{2}}\cdots a_{i_{\sigma(r)}j_{r}}&=\left|\begin{array}{cccc}
		a_{i_{1} j_{1}} & a_{i_{2} j_{1}} & \cdots & a_{i_{r} j_{1}} \\
		a_{i_{1} j_{2}} & a_{i_{2} j_{2}} & \cdots & a_{i_{r} j_{2}} \\
		\vdots & \vdots & \ddots & \vdots \\
		a_{i_{1} j_{r}} & a_{i_{2} j_{r}} & \cdots & a_{i_{r}} j_{r}
		\end{array}\right|=\left|\begin{array}{cccc}
		a_{i_{1} j_{1}} & a_{i_{1} j_{2}} & \cdots & a_{i_{1} j_{r}} \\
		a_{i_{2} j_{1}} & a_{i_{2} j_{2}} & \cdots & a_{i_{2} j_{r}} \\
		\vdots & \vdots & \ddots & \vdots \\
		a_{i_{r} j_{1}} & a_{i_{r} j_{2}} & \cdots & a_{i_{r}} j_{r}
		\end{array}\right|\\
		&=A\left(\begin{array}{cccc}i_{1} & i_{2} & \cdots & i_{r} \\j_{1} & j_{2} & \cdots & j_{r}\end{array}\right)
\end{split}
\]
因此
\[
\begin{split}
	\overline{e}_{j_{1}}\wedge \overline{e}_{j_{2}}\wedge\cdots\wedge \overline{e}_{j_{r}}&=\sum_{1\leqslant i_{1}<\cdots< i_{r}\leqslant n}A\left(\begin{array}{cccc}i_{1} & i_{2} & \cdots & i_{r} \\j_{1} & j_{2} & \cdots & j_{r}\end{array}\right)\cdot e_{i_{1}}\wedge e_{i_{2}}\wedge \cdots e_{i_{r}}
\end{split}
\]
也就是说
\begin{equation}
\overline{e}^{\wedge}_{J}=\sum_{I\in Q(r,n)}A\left(\begin{array}{cccc}i_{1} & i_{2} & \cdots & i_{r} \\j_{1} & j_{2} & \cdots & j_{r}\end{array}\right)\cdot e^{\wedge}_{I}
\end{equation}
这就表明过渡矩阵$C$的第$I$行第$J$列处的元素$C_{IJ}$是
\[
	C_{IJ}=A\left(\begin{array}{cccc}i_{1} & i_{2} & \cdots & i_{r} \\j_{1} & j_{2} & \cdots & j_{r}\end{array}\right)
	\]

再利用Binet-Cauchy公式(\rlem{lem:chapter2Binet-Cauchy公式})，计算可得
\begin{equation}
\begin{split}
&\sum\limits_{K\in Q(r,n)}A\left(\begin{array}{cccc}i_{1} & i_{2} & \cdots & i_{r} \\k_{1} & k_{2} & \cdots & k_{r}\end{array}\right)\cdot A^{-1}\left(\begin{array}{cccc}k_{1} & k_{2} & \cdots & k_{r} \\j_{1} & j_{2} & \cdots & j_{r}\end{array}\right)\\
&=\left( AA^{-1} \right)\left(\begin{array}{cccc}i_{1} & i_{2} & \cdots & i_{r} \\j_{1} & j_{2} & \cdots & j_{r}\end{array}\right)=E\left(\begin{array}{cccc}i_{1} & i_{2} & \cdots & i_{r} \\j_{1} & j_{2} & \cdots & j_{r}\end{array}\right)=
	\begin{cases}
		0,& I\neq J\\
		1,& I=J
	\end{cases}
\end{split}
\label{Binet-Cauchy公式计算}
\end{equation}
若记
\[
B_{KJ}=A^{-1}\left(\begin{array}{cccc}k_{1} & k_{2} & \cdots & k_{r} \\j_{1} & j_{2} & \cdots & j_{r}\end{array}\right)
\]
则式\requ{Binet-Cauchy公式计算}可以写成$\sum\limits_{K\in Q(r,n)}C_{IK}B_{KJ}=\delta_{IJ}$. 这表明以$B_{KJ}$作为第$K$行第$J$列处的元素的矩阵恰好就是矩阵$C$的逆矩阵，因此对$\forall \alpha=(\alpha_{1},\cdots,\alpha_{r}),\beta=(\beta_{1},\cdots,\beta_{r})\in Q(r,n)$，有
\[
C^{-1}_{\alpha\beta}=B_{\alpha\beta}=A^{-1}\left(\begin{array}{cccc}\alpha_{1} & \alpha_{2} & \cdots & \alpha_{r} \\\beta_{1} & \beta_{2} & \cdots & \beta_{r}\end{array}\right)
\]
\end{proof}

\begin{example}[]
设$\operatorname{dim}(V)=3$，考虑$\Lambda^{2}(V)$. 假设$\{e_{1},e_{2},e_{3} \}$与$\{\overline{e}_{1},\overline{e}_{2},\overline{e}_{3}\}$都是$V$的基，且
\[
\left( \overline{e}_{1},\overline{e}_{2},\overline{e}_{3} \right)=\left( e_{1},e_{2},e_{3} \right)\left[\begin{array}{ccc}
	a_{11} & a_{12} & a_{13} \\
	a_{21} & a_{22} & a_{23} \\
	a_{31} & a_{32} & a_{33} 
	\end{array}\right]
\]
则$\{e_{1}\wedge e_{2},e_{1}\wedge e_{3},e_{2}\wedge e_{3}\}$和$\{\overline{e}_{1}\wedge \overline{e}_{2},\overline{e}_{1}\wedge \overline{e}_{3},\overline{e}_{2}\wedge \overline{e}_{3}\}$都是$\Lambda^{2}(V)$的基。下面我们来写$\{e_{1}\wedge e_{2},e_{1}\wedge e_{3},e_{2}\wedge e_{3}\}$到$\{\overline{e}_{1}\wedge \overline{e}_{2},\overline{e}_{1}\wedge \overline{e}_{3},\overline{e}_{2}\wedge \overline{e}_{3}\}$的过渡矩阵$C$.
\soln

首先将有序多重指标集$Q(2,3)$中的元素都按照字典顺序排列，得到
\[
	Q(2,3)=\{(1,2),(1,3),(2,3) \}
\]
因此我们说“$C$的第$(1,2)$行”实际上是指$C$的第$1$行，因为指标$(1,2)$在字典顺序下排在第一位；同理，所谓“$C$的第$(1,3)$行第$(2,3)$列处的元素”是指$C$的第$2$行第$3$列处的元素，因为指标$(1,3)$在字典顺序下排在第二位，而指标$(2,3)$在字典顺序下排在第三位。

根据\rthm{thm:chapter2反对称空间的基过渡矩阵}，结合指标的字典顺序，我们来逐个计算矩阵$C$中的元素。比如
\[
C_{(1,2)(1,2)}=A\left(\begin{array}{cc} 1 & 2  \\ 1 & 2 \end{array}\right)=\left|\begin{array}{ccc}
	a_{11} & a_{12}   \\
	a_{21} & a_{22}  
	\end{array}\right|\quad\quad C_{(1,2)(1,3)}=A\left(\begin{array}{cc} 1 & 2  \\ 1 & 3 \end{array}\right)=\left|\begin{array}{ccc}
		a_{11} & a_{13}   \\
		a_{21} & a_{23}  
		\end{array}\right|
\]
这里$C_{(1,2)(1,2)}$表示第$1$行第$1$列处的元素，$C_{(1,2)(1,3)}$表示第$1$行第$2$列处的元素。其余元素可同理计算，最终得到过渡矩阵$C$如下，它是一个$\binom{3}{2}\times \binom{3}{2}=3\times 3$的矩阵：
\[
	C=\left[\begin{array}{ccc}
		{\left|\begin{array}{ccc}
			a_{11} & a_{12}   \\
			a_{21} & a_{22}  
		\end{array}\right|} & {\left|\begin{array}{ccc}
				a_{11} & a_{13}   \\
				a_{21} & a_{23}  
		\end{array}\right|} & {\left|\begin{array}{ccc}
				a_{12} & a_{13}   \\
				a_{22} & a_{23}  
				\end{array}\right|} \\
		{\left|\begin{array}{ccc}
				a_{11} & a_{12}   \\
				a_{31} & a_{32}  
		\end{array}\right|} & {\left|\begin{array}{ccc}
				a_{11} & a_{13}   \\
				a_{31} & a_{33}  
		\end{array}\right|} & {\left|\begin{array}{ccc}
				a_{12} & a_{13}   \\
				a_{32} & a_{33}  
		\end{array}\right|} \\
		{\left|\begin{array}{ccc}
				a_{21} & a_{22}   \\
				a_{31} & a_{32}  
		\end{array}\right|} & {\left|\begin{array}{ccc}
				a_{21} & a_{23}   \\
				a_{31} & a_{33}  
		\end{array}\right|} & {\left|\begin{array}{ccc}
				a_{22} & a_{23}   \\
				a_{32} & a_{33}  
		\end{array}\right|} 
		\end{array}\right]
\]
\end{example}


在\rdef{def:chapter2张量空间的配合运算2}中，我们给$V^{r}_{0}$与$V_{r}^{0}$赋予了配合运算。作为$V^{r}_{0}$与$V_{r}^{0}$的线性子空间，$\Lambda^{r}(V)$与$\Lambda^{r}(V^{\ast})$当然也继承了其上的配合运算。注意到配合运算是双线性的，再结合式\requ{外向量的表达式}，所以我们只需挑一个简单的情形看看计算结果是怎样的即可。

取$\theta_{1}\wedge \theta_{2}\wedge\cdots\wedge \theta_{r}\in \Lambda^{r}(V^{\ast})$，取$v_{1}\wedge v_{2}\wedge \cdots\wedge v_{r}\in \Lambda^{r}(V)$. 首先有
\[
\begin{split}
&\theta_{1}\wedge \theta_{2}\wedge\cdots\wedge \theta_{r}=\sum_{\sigma\in G(r)}\operatorname{sgn}(\sigma)\cdot \sigma(\theta_{1}\otimes  \theta_{2}\otimes\cdots\otimes \theta_{r})\\
&v_{1}\wedge v_{2}\wedge \cdots\wedge v_{r}=\sum_{\tau\in G(r)}\operatorname{sgn}(\tau)\cdot \tau(v_{1}\otimes  v_{2}\otimes\cdots\otimes v_{r})
\end{split}
\]
按照$V^{0}_{r}$与$V_{0}^{r}$上约定的配合运算，以及配合的双线性性质，我们有
\[
\begin{split}
	&\bigl\langle\theta_{1}\wedge \theta_{2}\wedge\cdots\wedge \theta_{r} ,v_{1}\wedge v_{2}\wedge \cdots\wedge v_{r} \bigr\rangle\\
	&=\biggl\langle \sum_{\sigma\in G(r)}\operatorname{sgn}(\sigma)\cdot \sigma(\theta_{1}\otimes  \theta_{2}\otimes\cdots\otimes \theta_{r}),\sum_{\tau\in G(r)}\operatorname{sgn}(\tau)\cdot \tau(v_{1}\otimes  v_{2}\otimes\cdots\otimes v_{r}) \biggr\rangle \\
	&=\sum_{\sigma\in G(r)}\operatorname{sgn}(\sigma)\left[\sum_{\tau\in G(r)}\operatorname{sgn}(\tau)\cdot \bigl\langle \sigma(\theta_{1}\otimes  \theta_{2}\otimes\cdots\otimes \theta_{r}), \tau(v_{1}\otimes  v_{2}\otimes\cdots\otimes v_{r}) \bigr\rangle \right]
\end{split}
\]
再利用\rprop{prop:chapter2置换在单项式上的作用}与式\requ{配合运算1}可得
\[
\begin{split}
	&\bigl\langle\theta_{1}\wedge \theta_{2}\wedge\cdots\wedge \theta_{r} ,v_{1}\wedge v_{2}\wedge \cdots\wedge v_{r} \bigr\rangle\\
	&=\sum_{\sigma\in G(r)}\operatorname{sgn}(\sigma)\left[\sum_{\tau\in G(r)}\operatorname{sgn}(\tau)\cdot \bigl\langle \theta_{\sigma^{-1}(1)}\otimes  \theta_{\sigma^{-1}(2)}\otimes\cdots\otimes \theta_{\sigma^{-1}(r)}, v_{\tau^{-1}(1)}\otimes  v_{\tau^{-1}(2)}\otimes\cdots\otimes v_{\tau^{-1}(r)} \bigr\rangle \right]\\
	&\xlongequal{\requ{配合运算1}}\sum_{\sigma\in G(r)}\operatorname{sgn}(\sigma)\left[\sum_{\tau\in G(r)}\operatorname{sgn}(\tau)\cdot  \theta_{\sigma^{-1}(1)}\otimes  \theta_{\sigma^{-1}(2)}\otimes\cdots\otimes \theta_{\sigma^{-1}(r)} \left(v_{\tau^{-1}(1)},  v_{\tau^{-1}(2)},\cdots, v_{\tau^{-1}(r)} \right) \right]\\
	&=\sum_{\sigma\in G(r)}\operatorname{sgn}(\sigma)\left[\sum_{\tau\in G(r)}\operatorname{sgn}(\tau)\cdot  \theta_{\sigma^{-1}(1)}\left(v_{\tau^{-1}(1)}\right)\cdot   \theta_{\sigma^{-1}(2)}\left(v_{\tau^{-1}(2)}\right)\cdots \theta_{\sigma^{-1}(r)}\left(v_{\tau^{-1}(r)}\right)\right]
\end{split}
\]
注意到当$\sigma,\tau$遍历$G(r)$时，$\sigma^{-1},\tau^{-1}$也同时遍历$G(r)$，而且置换与它的逆置换总是同号，即$\operatorname{sgn}(\sigma)=\operatorname{sgn}(\sigma^{-1}),\operatorname{sgn}(\tau)=\operatorname{sgn}(\tau^{-1})$. 若令$\sigma^{\prime}=\sigma^{-1},\tau^{\prime}=\tau^{-1}$，则上式又可改写为
\[
\begin{split}
	&\bigl\langle\theta_{1}\wedge \theta_{2}\wedge\cdots\wedge \theta_{r} ,v_{1}\wedge v_{2}\wedge \cdots\wedge v_{r} \bigr\rangle\\
	&=\sum_{\sigma^{-1}\in G(r)}\operatorname{sgn}(\sigma^{-1})\left[\sum_{\tau^{-1}\in G(r)}\operatorname{sgn}(\tau^{-1})\cdot  \theta_{\sigma^{-1}(1)}\left(v_{\tau^{-1}(1)}\right)\cdot   \theta_{\sigma^{-1}(2)}\left(v_{\tau^{-1}(2)}\right)\cdots \theta_{\sigma^{-1}(r)}\left(v_{\tau^{-1}(r)}\right)\right]\\
	&=\sum_{\sigma^{\prime}\in G(r)}\operatorname{sgn}(\sigma^{\prime})\left[\sum_{\tau^{\prime}\in G(r)}\operatorname{sgn}(\tau^{\prime})\cdot  \theta_{\sigma^{\prime}(1)}\left(v_{\tau^{\prime}(1)}\right)\cdot   \theta_{\sigma^{\prime}(2)}\left(v_{\tau^{\prime}(2)}\right)\cdots \theta_{\sigma^{\prime}(r)}\left(v_{\tau^{\prime}(r)}\right)\right]
\end{split}
\]
注意到上式的方括号里面恰好可以看成是一个行列式的展开式，即
\[
\bigl\langle\theta_{1}\wedge \theta_{2}\wedge\cdots\wedge \theta_{r} ,v_{1}\wedge v_{2}\wedge \cdots\wedge v_{r} \bigr\rangle=\sum_{\sigma^{\prime}\in G(r)}\operatorname{sgn}(\sigma^{\prime}) \left|\begin{array}{cccc}\theta_{\sigma^{\prime}(1)}(v_{1}) & \theta_{\sigma^{\prime}(1)}(v_{2}) & \cdots & \theta_{\sigma^{\prime}(1)}(v_{r}) \\ \theta_{\sigma^{\prime}(2)}(v_{1}) & \theta_{\sigma^{\prime}(2)}(v_{2}) & \cdots & \theta_{\sigma^{\prime}(2)}(v_{r}) \\ \vdots & \vdots & \ddots& \vdots \\ \theta_{\sigma^{\prime}(r)}(v_{1}) & \theta_{\sigma^{\prime}(r)}(v_{2}) & \cdots & \theta_{\sigma^{\prime}(r)}(v_{r})\end{array}\right|
\]
依次交换行列式的两行，将行列式的行指标调整为自然顺序。注意这样不断交换两行后，行列式的值会改变一个符号$\operatorname{sgn}(\sigma^{\prime})$. 从而
\[
\begin{split}
	\bigl\langle\theta_{1}\wedge \theta_{2}\wedge\cdots\wedge \theta_{r} ,v_{1}\wedge v_{2}\wedge \cdots\wedge v_{r} \bigr\rangle&=\sum_{\sigma^{\prime}\in G(r)}\underbrace{\operatorname{sgn}(\sigma^{\prime}) {\color{red}{\operatorname{sgn}(\sigma^{\prime})}} }_{=1} \left|\begin{array}{cccc}\theta_{1}(v_{1}) & \theta_{1}(v_{2}) & \cdots & \theta_{1}(v_{r}) \\ \theta_{2}(v_{1}) & \theta_{2}(v_{2}) & \cdots & \theta_{2}(v_{r}) \\ \vdots & \vdots & \ddots& \vdots \\ \theta_{r}(v_{1}) & \theta_{r}(v_{2}) & \cdots & \theta_{r}(v_{r})\end{array}\right|\\
	&=\sum_{\sigma^{\prime}\in G(r)}\left|\begin{array}{cccc}\theta_{1}(v_{1}) & \theta_{1}(v_{2}) & \cdots & \theta_{1}(v_{r}) \\ \theta_{2}(v_{1}) & \theta_{2}(v_{2}) & \cdots & \theta_{2}(v_{r}) \\ \vdots & \vdots & \ddots& \vdots \\ \theta_{r}(v_{1}) & \theta_{r}(v_{2}) & \cdots & \theta_{r}(v_{r})\end{array}\right|\\
	&=r! \left|\begin{array}{cccc}\theta_{1}(v_{1}) & \theta_{1}(v_{2}) & \cdots & \theta_{1}(v_{r}) \\ \theta_{2}(v_{1}) & \theta_{2}(v_{2}) & \cdots & \theta_{2}(v_{r}) \\ \vdots & \vdots & \ddots& \vdots \\ \theta_{r}(v_{1}) & \theta_{r}(v_{2}) & \cdots & \theta_{r}(v_{r})\end{array}\right|
\end{split}
\]
特别地，若取定$V$的一组基$\{e_{1} ,e_{2},\cdots,e_{n}\}$，则$\bigl\langle e^{\ast}_{i_{1}}\wedge e^{\ast}_{i_{2}}\wedge \cdots\wedge e^{\ast}_{i_{r}} ,e_{i_{1}}\wedge e_{i_{2}}\wedge \cdots\wedge e_{i_{r}} \bigr\rangle=r!$.

因此，按照$V^{0}_{r}$与$V^{r}_{0}$上的配合运算来计算$\Lambda^{r}(V^{\ast})$与$\Lambda^{r}(V)$中元素的配合，会多出一个常数$r!$，但我们希望$\Lambda^{r}(V^{\ast})$与$\Lambda^{r}(V)$之间的配合也能够具有对偶配合的性质。所以不宜直接继承$V^{0}_{r}$与$V^{r}_{0}$上的配合运算，还需要对系数进行规范化。
\begin{mydef}{$\Lambda^{r}(V^{\ast})$与$\Lambda^{r}(V)$的配合}{chapter2反对称空间之间的配合}
对于$\forall g\in \Lambda^{r}(V^{\ast})$，$\forall x\in \Lambda^{r}(V)$，约定配合运算$\bigl(\quad,\quad \bigr):\Lambda^{r}(V^{\ast})\times \Lambda^{r}(V)\longrightarrow \mr$如下：
\[
\bigl(g,x \bigr) \triangleq \frac{1}{r!} \langle g,x \rangle
\]
其中$\langle\quad,\quad\rangle$表示$V^{0}_{r}$与$V^{r}_{0}$之间的配合运算(即\rdef{def:chapter2张量空间的配合运算2}). 
\end{mydef}

注意，$\Lambda^{r}(V^{\ast})$与$\Lambda^{r}(V)$上的配合并非直接继承自$V^{0}_{r}$与$V^{r}_{0}$上的配合运算，而是相差了一个系数$r!$. 以后在做$\Lambda^{r}(V^{\ast})$与$\Lambda^{r}(V)$中的元素之间的配合运算时，都按照上述定义进行。既然有了配合运算，那么根据\rthm{thm:chapter2配合则同构}，立刻就有如下结论成立。
\begin{mythm}{}{chapter2反对称空间的对偶同构}
$\left( \Lambda^{r}(V) \right)^{\ast}\cong \Lambda^{r}(V^{\ast})$
\end{mythm}

对于形如$\theta_{1}\wedge \theta_{2}\wedge\cdots\wedge \theta_{r}\in \Lambda^{r}(V^{\ast})$，$v_{1}\wedge v_{2}\wedge \cdots\wedge v_{r}\in \Lambda^{r}(V)$的元素，在新的配合下就可按如下方式计算
\begin{equation}
\begin{split}
	\bigl(\theta_{1}\wedge \theta_{2}\wedge\cdots\wedge \theta_{r} ,v_{1}\wedge v_{2}\wedge \cdots\wedge v_{r} \bigr)&=\frac{1}{r!}\bigl\langle\theta_{1}\wedge \theta_{2}\wedge\cdots\wedge \theta_{r} ,v_{1}\wedge v_{2}\wedge \cdots\wedge v_{r} \bigr\rangle\\
&=\left|\begin{array}{cccc}\theta_{1}(v_{1}) & \theta_{1}(v_{2}) & \cdots & \theta_{1}(v_{r}) \\ \theta_{2}(v_{1}) & \theta_{2}(v_{2}) & \cdots & \theta_{2}(v_{r}) \\ \vdots & \vdots & \ddots& \vdots \\ \theta_{r}(v_{1}) & \theta_{r}(v_{2}) & \cdots & \theta_{r}(v_{r})\end{array}\right|
\end{split}
\end{equation}
特别地，若$\{e_{1} ,e_{2},\cdots,e_{n}\}$是$V$的一组基，则
\begin{equation}
\bigl(e^{\ast}_{j_{1}}\wedge e^{\ast}_{j_{2}}\wedge \cdots\wedge e^{\ast}_{j_{r}} ,e_{i_{1}}\wedge e_{i_{2}}\wedge \cdots\wedge e_{i_{r}}\bigr)=\left|\begin{array}{cccc}\delta^{j_{1}}_{i_{1}} & \delta^{j_{1}}_{i_{2}} & \cdots & \delta^{j_{1}}_{i_{r}} \\ \delta^{j_{2}}_{i_{1}} & \delta^{j_{2}}_{i_{2}} & \cdots & \delta^{j_{2}}_{i_{r}} \\ \vdots & \vdots & \ddots& \vdots \\ \delta^{j_{r}}_{i_{1}} & \delta^{j_{r}}_{i_{2}} & \cdots & \delta^{j_{r}}_{i_{r}}\end{array}\right|
	\label{反对称空间的对偶基}
\end{equation}

\begin{mydef}{广义Kronecker符号}{chapter2广义Kronecker符号}
\[
\delta^{j_{1}j_{2}\cdots j_{r}}_{i_{1}i_{2}\cdots i_{r}}\triangleq \left|\begin{array}{cccc}\delta^{j_{1}}_{i_{1}} & \delta^{j_{1}}_{i_{2}} & \cdots & \delta^{j_{1}}_{i_{r}} \\ \delta^{j_{2}}_{i_{1}} & \delta^{j_{2}}_{i_{2}} & \cdots & \delta^{j_{2}}_{i_{r}} \\ \vdots & \vdots & \ddots& \vdots \\ \delta^{j_{r}}_{i_{1}} & \delta^{j_{r}}_{i_{2}} & \cdots & \delta^{j_{r}}_{i_{r}}\end{array}\right|
\]
我们称$\delta^{j_{1}j_{2}\cdots j_{r}}_{i_{1}i_{2}\cdots i_{r}}$为广义Kronecker符号.
\end{mydef}

$\delta^{j_{1}j_{2}\cdots j_{r}}_{i_{1}i_{2}\cdots i_{r}}$是Kronecker符号$\delta^{j}_{i}$的推广，它的定义方式看起来比较复杂，不过其实还是很容易把握的。
\begin{myprop}{}{chapter2广义Kronecker符号的性质}
\begin{itemize}
\item[(1)] 若存在$j_{k}=j_{l}$，则$\delta^{j_{1}j_{2}\cdots j_{r}}_{i_{1}i_{2}\cdots i_{r}}=0$；若存在$i_{k}=i_{l}$，则$\delta^{j_{1}j_{2}\cdots j_{r}}_{i_{1}i_{2}\cdots i_{r}}=0$.
\item[(2)] 若存在$j_{k}\notin \{i_{1},\cdots,i_{r} \}$，则$\delta^{j_{1}j_{2}\cdots j_{r}}_{i_{1}i_{2}\cdots i_{r}}=0$；若存在$i_{k}\notin \{j_{1},\cdots,j_{r} \}$，则$\delta^{j_{1}j_{2}\cdots j_{r}}_{i_{1}i_{2}\cdots i_{r}}=0$.
\item[(3)] 若
			\[
				\left(\begin{array}{lllll}
					j_{1} & j_{2} &j_{3} & \cdots & j_{r}  \\
					i_{1} & i_{2} & i_{3} & \cdots & i_{r} 
					\end{array}\right)
			\]
			形成一个置换$\sigma$，则$\delta^{j_{1}j_{2}\cdots j_{r}}_{i_{1}i_{2}\cdots i_{r}}=\operatorname{sgn}(\sigma)$.
\end{itemize}
\end{myprop}

\begin{proof}
证明其实很简单，抓住行列式的特点即可。

(1)若$j_{k}=j_{l}$，则行列式的第$k$行与第$l$行是一样的，从而行列式值为0；同理，若$i_{k}=i_{l}$，则行列式的第$k$列于第$l$列是一样的，从而行列式值为0.

(2)若$j_{k}\notin \{i_{1},i_{2},\cdots,i_{r} \}$，则行列式的第$k$行全是0，从而行列式值为0；同理，若$i_{k}\notin \{j_{1},j_{2},\cdots,j_{r} \}$，则行列式的第$k$列全是0，从而行列式值为0.

(3)当指标集形成置换时，注意到
		\[
			\left(\begin{array}{llll}
					j_{1} & j_{2}  & \cdots & j_{r}  \\
					i_{1} & i_{2}  & \cdots & i_{r} 
			\end{array}\right)=
			\left(\begin{array}{llll}
				j_{1} & j_{2}  & \cdots & j_{r}  \\
				1 & 2  & \cdots & r 
			\end{array}\right)
			\left(\begin{array}{llll}
				1 & 2  & \cdots & r  \\
				i_{1} & i_{2}  & \cdots & i_{r} 
		\end{array}\right)=\sigma_{1}\sigma_{2}
		\]
从$r$阶单位方阵$E$出发，对行施加$\sigma_{1}^{-1}$对应的置换，再对列施加$\sigma_{2}$对应的置换，最终得到$\delta^{j_{1}j_{2}\cdots j_{r}}_{i_{1}i_{2}\cdots i_{r}}$. 于是
\[
	\delta^{j_{1}j_{2}\cdots j_{r}}_{i_{1}i_{2}\cdots i_{r}}=\operatorname{sgn}(\sigma_{1}^{-1})\cdot \operatorname{sgn}(\sigma_{2})\cdot |E|
	=\operatorname{sgn}(\sigma_{1})\cdot \operatorname{sgn}(\sigma_{2})=\operatorname{sgn}(\sigma_{1}\sigma_{2})=\operatorname{sgn}(\sigma)
\]
其中第二个等号的成立是因为置换与其逆置换总是同号的。
\end{proof}

利用广义Kronecker符号，我们来看$\Lambda^{r}(V^{\ast})$与$\Lambda^{r}(V)$的基底之间的配合。
取定$V$的一组基$\{e_{1} ,e_{2},\cdots,e_{n}\}$，则由\rthm{thm:chapter2反对称空间的结构}可知，$\{e_{i_{1}}\wedge e_{i_{2}}\wedge \cdots\wedge e_{i_{r}} | 1\leqslant i_{1}<i_{2}<\cdots<i_{r}\leqslant n \}$是$\Lambda^{r}(V)$的一组基，$\{e^{\ast}_{j_{1}}\wedge e^{\ast}_{j_{2}}\wedge \cdots\wedge e^{\ast}_{j_{r}} | 1\leqslant j_{1}<j_{2}<\cdots<j_{r}\leqslant n \}$是$\Lambda^{r}(V^{\ast})$的一组基。再根据式\requ{反对称空间的对偶基}可知
\[
	\bigl(e^{\ast}_{j_{1}}\wedge e^{\ast}_{j_{2}}\wedge \cdots\wedge e^{\ast}_{j_{r}} ,e_{i_{1}}\wedge e_{i_{2}}\wedge \cdots\wedge e_{i_{r}}\bigr)=\delta^{j_{1}j_{2}\cdots j_{r}}_{i_{1}i_{2}\cdots i_{r}}
\]
注意此时$1\leqslant j_{1}<j_{2}<\cdots<j_{r}\leqslant n,1\leqslant i_{1}<i_{2}<\cdots<i_{r}\leqslant n $. 因为它们是作为$\Lambda^{r}(V^{\ast})$与$\Lambda^{r}(V)$的基底的下标出现的，所以有顺序。
显然当$j_{1}=i_{1},j_{2}=i_{2},\cdots,j_{r}=i_{r}$时，$\delta^{j_{1}j_{2}\cdots j_{r}}_{i_{1}i_{2}\cdots i_{r}}=1$.
否则的话，假设$j_{k}\neq i_{k}$是从左往右第一个出现的不相等序对，即$j_{1}=i_{1},\cdots,j_{k-1}=i_{k-1}$，但$j_{k}\neq i_{k}$. 
\begin{itemize}
\item 若$j_{k}<i_{k}$，则$j_{k}\notin \{i_{1},i_{2},\cdots,i_{r} \}$，从而$\delta^{j_{1}j_{2}\cdots j_{r}}_{i_{1}i_{2}\cdots i_{r}}=0$；
\item 若$j_{k}>i_{k}$，则$i_{k}\notin \{j_{1},j_{2},\cdots,j_{r} \}$，从而$\delta^{j_{1}j_{2}\cdots j_{r}}_{i_{1}i_{2}\cdots i_{r}}=0$.
\end{itemize}
这表明$\bigl(e^{\ast}_{j_{1}}\wedge e^{\ast}_{j_{2}}\wedge \cdots\wedge e^{\ast}_{j_{r}} ,e_{i_{1}}\wedge e_{i_{2}}\wedge \cdots\wedge e_{i_{r}}\bigr)=1\Leftrightarrow j_{1}=i_{1},j_{2}=i_{2},\cdots,j_{r}=i_{r}$. 

也就是说，我们给$\Lambda^{r}(V^{\ast})$与$\Lambda^{r}(V)$赋予的配合运算$\bigl(\quad,\quad \bigr)$可以看成是对偶配合。因此，在指定了该配合运算后，你也可以直接把\rthm{thm:chapter2反对称空间的对偶同构}写成$\left( \Lambda^{r}(V) \right)^{\ast}=\Lambda^{r}(V^{\ast})$. 当然啦还是那句话，无论看成是$\left( \Lambda^{r}(V) \right)^{\ast}\cong\Lambda^{r}(V^{\ast})$还是$\left( \Lambda^{r}(V) \right)^{\ast}=\Lambda^{r}(V^{\ast})$，只要你自己拎得清就行！

\begin{remark}
	以后在不引起歧义时，我们也用$\langle\quad,\quad\rangle$来表示$\bigl(\quad,\quad \bigr)$. 实际上你只要看参与配合的元素来自哪儿，就能明白此时的配合运算$\langle\quad,\quad\rangle$采用的是\rdef{def:chapter2张量空间的配合运算2}还是\rdef{def:chapter2反对称空间之间的配合}。
\end{remark}

\subsection{张量代数与外代数}

在开始介绍这一节内容之前，先跟大家声明：本节行文中有些说法并不妥当，更谈不上严格，不必与我较真，毕竟这只是一本笔记。如果你想要更深入严谨地了解张量代数等方面的内容，还请查阅专业教材，比如李文威老师的书。不过我并不是很推荐这么做，因为这是一个大坑。更何况，咱们学的是几何呀！

首先读者要明确一点，这里所谓的“代数(algebra)”，不是我们通常说的“分析、代数、几何”这一语境中的意思，而是一个跟“群、环、域”等等类似的有着确切定义的概念。不过我不打算介绍它的严格定义，在大多数时候你都可以把代数简单地理解为带有乘法运算的向量空间，当然这个乘法运算要满足一些规定。也就是说，一个代数上有三种运算：加法、乘法与数乘。就像群和环一样，这三种运算也要满足一些基本公理。

其次我们粗略说一下张量代数和外代数的名称由来。它们得名于其上所采用的乘法：所谓张量代数，简单地说就是\textbf{以张量积运算作为向量空间上的乘法}；所谓外代数，就是\textbf{以外积运算作为乘法}。记性好的读者可能记得我们上一节还介绍过一种称为对称积的运算，那如果把乘法取作为对称积，形成的代数(如果存在的话)是不是就叫“对称代数”呢？答案是Yes! 不过我们不准备讲对称代数，甚至于张量代数也只是一笔带过，因为在几何中外代数才是更加常用的。

我们最最熟悉、最最常见的一个代数是多项式代数，我想多项式的加法、数乘以及乘法运算大家应该都很熟练。只要你会多项式的这些运算，那么你自然而然也就会计算跟张量代数和外代数有关的运算了。唯一要注意到的就只有张量积和外积所独有的一些运算性质，比如计算多项式相乘时你可以随意交换变元$x,y$，但是当以外积为乘法时，就不能随意交换了，因为外积没有交换律，取而代之的是反交换律。总之，涉及具体计算时，你就以多项式为参考模型，辅以张量积和外积特有的运算性质即可。

以下内容其实只是介绍一些术语和记号，我真正想说的就是上面那三段话。令
\[
T(V)\triangleq \bigoplus_{r,s\geqslant 0}^{\infty} V^{r}_{s}=\mr\oplus V\oplus V^{\ast}\oplus V^{1}_{1}\oplus V^{2}_{0}\oplus V^{0}_{2}\oplus\cdots
\]

解释一下，上面所作的张量空间的直和只是形式和(严格说来叫外直和)。具体一点，也就是说我们约定其中的元素长这个样子：$\forall w\in T(V),$
\[
	w=w^{0}_{0}+w^{1}_{0}+w^{0}_{1}+w^{1}_{1}+\cdots
\]
其中诸$ w^{r}_{s}\in V^{r}_{s}$，并且右边的形式和中只有有限多项不为0，其余皆为0. 

初学者可能会困惑这里“$+$”的意义，其实不必深究\footnote{深究下去你会发现，$w$也写成$w=(w^{0}_{0},w^{1}_{0},w^{0}_{1},w^{1}_{1},\cdots)$的形式。事实上这正是一般的外直和定义中所采用的写法。}，你把它当成连接符就行了，因为我们做的只是形式和。就像多元多项式$x^{2}y+y^{3}+xy$，你有想过其中$+$的意义吗？(笑)

$T(V)$按照自然的方式形成一个无穷维向量空间，这里所谓的“自然的方式”，你可以想象成多元多项式的加法与数乘运算：数乘就是和式的每一项都乘以该数；加法就是“合并同类项”，能加的就合并在一起加起来(即来自同一个$V^{r}_{s}$的元素，就按照$V^{r}_{s}$中的加法相加)，不能加的就放着构成形式和。也就是说“$+$”对于能相加的项就表示加法的含义，对于无法相加的项(即来自不同$V^{r}_{s}$的元素)就仅仅起到连接符的作用。

\begin{example}[$ T(V) $上的加法与数乘]
设$k\in \mr$，再取$ w,u\in T(V) $为
\[
	w=w_{0}^{4}+ w_{1}^{3}+ w^{12}_{7}\quad
	\quad\quad
	u=u_{2}^{2}+ u_{2}^{3}+ u^{12}_{7}+ u^{5}_{25}
\]
\soln

 \[
 \begin{split}
 w+u&=w_{0}^{4}+ w_{1}^{3}+ (w^{12}_{7}+ u^{12}_{7})+ u_{2}^{2}+  u_{2}^{3}+ u^{5}_{25}\\
kw&=kw_{0}^{4}+ kw_{1}^{3}+ kw^{12}_{7}
 \end{split}
 \]
\end{example}

除了有加法与乘法，$T(V)  $还可以引入乘法运算。我们把$T(V)$中\textbf{两个元素相乘规定为张量的张量积运算通过分配律扩充得到的}，仍用记号$\otimes $来表示$T(V)$上的这种乘法。这种乘法就跟多项式乘多项式差不多，但是要注意张量积没有交换律，所以在做这个乘法时不可随意交换位置。看看下面的例子就明白了。

\begin{example}[$ T(V) $上的乘法]
	取$ w,u\in T(V) $为
\[
	w=w_{0}^{4}+ w_{1}^{3}+ w^{12}_{7}\quad
	\quad\quad
	u=u_{2}^{2}+ u_{2}^{3}+ u^{12}_{7}+ u^{5}_{25}
\]
\soln

\[
\begin{split}
	w\otimes u&=\left( w_{0}^{4}+ w_{1}^{3}+ w^{12}_{7} \right)\otimes \left( u_{2}^{2}+ u_{2}^{3}+ u^{12}_{7}+ u^{5}_{25} \right)\\
	&=w_{0}^{4}\otimes u_{2}^{2}+ w_{0}^{4}\otimes u_{2}^{3}+ w_{0}^{4}\otimes u^{12}_{7}+ w_{0}^{4}\otimes u^{5}_{25}\\
	&\quad + w_{1}^{3}\otimes u_{2}^{2}+ w_{1}^{3}\otimes u_{2}^{3}+ w_{1}^{3}\otimes u^{12}_{7}+ w_{1}^{3}\otimes u^{5}_{25}\\
	&\quad + w^{12}_{7}\otimes u_{2}^{2}+ w^{12}_{7}\otimes u_{2}^{3}+ w^{12}_{7}\otimes u^{12}_{7}+ w^{12}_{7}\otimes u^{5}_{25}
\end{split}
\]
\end{example}

在此乘法下，$T(V)$成为$\mr$上的一个代数，称为$V$上的\textbf{张量代数}。特别地，如果令
\[
T^{A}(V)\triangleq \bigoplus_{r= 0}^{\infty} V^{r}_{0} \quad\quad\quad T^{B}(V)\triangleq \bigoplus_{s=0}^{\infty} V^{0}_{s}
\]
则作为向量空间来看，$T^{A}(V)$与$T^{B}(V)$都是$T(V)$的子空间。如果给它们也配上张量积运算作为乘法，那么$T^{A}(V)$与$T^{B}(V)$各自也都形成代数，而且可以看作是$T(V)$的两个子代数。很多书上也把$T^{A}(V)$或者$T^{B}(V)$叫作$V$上的张量代数。至此我们说完了张量代数，下面再来看外代数。

前面我们已经提到过，所谓外代数就是把乘法取作为外积运算。而我们只对$\Lambda^{r}(V)$这一类空间定义过外积运算，所以令
\[
\Lambda(V)\triangleq \bigoplus_{r=0}^{\infty} \Lambda^{r}(V)
\]
如果$\operatorname{dim}(V)=n$，则根据\rthm{thm:chapter2反对称空间的结构}可知，当$r>n$时，$\Lambda^{r}(V)=\{0\}$. 所以实际上$\Lambda(V)$并不像张量代数那样是无穷多个空间的直和，它本质上只是有限个$\Lambda^{r}(V)$的直和，即
\[
	\Lambda(V)= \bigoplus_{r=0}^{n} \Lambda^{r}(V)\;,\;\text{\kaishu 其中}n=\operatorname{dim}(V)
\]
$\Lambda(V)$中的元素长这个样子：$\forall \xi\in \Lambda(V),\xi=\xi^{0}+\xi^{1}+\cdots +\xi^{n}$，其中$\xi^{i}\in \Lambda^{i}(V),i=1,\cdots,n$.

与$T(V)$一样，$\Lambda(V)$上的加法与数乘运算也是十分自然的。加法就是“合并同类项”，能加的就合并在一起加起来(即来自同一个$\Lambda^{r}(V)$的元素，就按照$\Lambda^{r}(V)$中的加法相加)，不能加的就放着构成形式和。也就是说“$+$”对于能相加的项就表示加法的含义，对于无法相加的项(即来自不同$\Lambda^{r}(V)$的元素)就仅仅起到连接符的作用。因此，$\Lambda(V)$构成一个$\binom{n}{0}+\binom{n}{1}+\cdots+\binom{n}{n}=2^{n}$维的向量空间，$\Lambda(V)$作为向量空间的组基为：
\[
\begin{split}
	\Lambda(V)\text{\kaishu 的基}&=\{\Lambda^{0}(V)\text{\kaishu 的基}\}\cup \{\Lambda^{1}(V)\text{\kaishu 的基}\}\cup \{\Lambda^{2}(V)\text{\kaishu 的基}\}\cup\cdots \cup \{\Lambda^{n}(V)\text{\kaishu 的基}\}\\
	&=\{1\}\cup\{ e_{1},\cdots,e_{n}\}\cup\{e_{i_{1}}\wedge e_{i_{2}} \}_{1\leqslant i_{1}<i_{2}\leqslant n}\cup\cdots \cup\{ e_{1}\wedge\cdots\wedge e_{n}\}
\end{split}
\]

然后我们再\textbf{把外积运算}添加进向量空间$\Lambda(V)$中\textbf{作为其上的乘法}。在外积乘法下，$\Lambda(V)$形成一个代数，称为$V$上的\textbf{外代数}或\textbf{Grassman代数}。
\begin{example}[$\Lambda(V)$上的加法、数乘和乘法]
设$k\in\mr$，再取$\xi,\eta\in \Lambda(V)$为
\[
\xi=\xi^{1}+\xi^{7}-\xi^{10}\quad\quad\quad \eta=\eta^{1}-\eta^{2}+\eta^{6}
\]
\soln

 \[
 \begin{split}
	\xi+\eta&=(\xi^{1}+\eta^{1})-\eta^{2}+\eta^{6}+\xi^{7}-\xi^{10}\\
	k\xi&=k\xi^{1}+k\xi^{7}-k\xi^{10}\\
	\xi\wedge\eta&=(\xi^{1}+\xi^{7}-\xi^{10})\wedge (\eta^{1}-\eta^{2}+\eta^{6})\\
	&=\xi^{1}\wedge \eta^{1}-\xi^{1}\wedge\eta^{2}+\xi^{1}\wedge\eta^{6}\\
	&\quad +\xi^{7}\wedge \eta^{1}-\xi^{7}\wedge\eta^{2}+\xi^{7}\wedge\eta^{6}\\
	&\quad -\xi^{10}\wedge \eta^{1}+\xi^{10}\wedge\eta^{2}-\xi^{10}\wedge\eta^{6}
 \end{split}
 \]
\end{example}

以上是以$r$次外向量空间$\Lambda^{r}(V)$构建的外代数，我们也可以用$s$次外形式空间$\Lambda^{s}(V^{\ast})$来构建相应的外代数$\Lambda(V^{\ast})$。令
\[
	\Lambda(V^{\ast})\triangleq \bigoplus_{s=0}^{\infty} \Lambda^{s}(V^{\ast})= \bigoplus_{s=0}^{n} \Lambda^{s}(V^{\ast})\;,\;\text{\kaishu 其中}n=\operatorname{dim}(V^{\ast})
\]
然后$\Lambda(V)$上的所有讨论照搬到$\Lambda(V^{\ast})$上即可。

\subsection{拉回映射}

这一节我们在$V^{0}_{r}$以及外形式空间$\Lambda^{r}(V^{\ast})$上讨论，因为主要是为后面的微分形式提供服务。之前所有关于外向量的性质均可照搬到外形式上。

给定线性映射$f:V\longrightarrow W$，我们在前面\ref{chapter1section1.2.3}节的式\requ{线性映射的诱导映射}定义了诱导映射
\[
	f^{\ast}:W^{\ast}\longrightarrow V^{\ast}
\]
注意到对偶空间$W^{\ast}$可以看作是$(0,1)$型张量空间$W^{0}_{1}$，$V^{\ast}$可以看作是$(0,1)$型张量空间$V^{0}_{1}$. 所以诱导映射$f^{\ast}$可以看作是$W^{0}_{1}\longrightarrow V^{0}_{1}$的线性映射。把空间进一步推广至$(0,r)$型张量空间，同时移植式\requ{线性映射的诱导映射}的想法，我们便得到以下定义。

\begin{mydef}{拉回映射}{chapter2拉回映射}
设$f:V\longrightarrow W$为向量空间之间的线性映射，定义$f^{\ast}:W^{0}_{r}\longrightarrow V^{0}_{r}$如下：对$\forall g\in W^{0}_{r}$
\[
f^{\ast}(g)(v_{1},v_{2},\cdots,v_{r})\triangleq g \left( f(v_{1}),f(v_{2}),\cdots,f(v_{r}) \right)\;,\forall v_{1},v_{2},\cdots,v_{r}\in V
\]
则$f^{\ast}$本身也是一个线性映射. 我们称$f^{\ast}$为$f$诱导的拉回映射(pull-back).
\end{mydef}

当$r=1$时，上述定义即化为\ref{chapter1section1.2.3}节中由式\requ{线性映射的诱导映射}所定义的诱导映射。

\begin{mythm}{}{chapter2保持张量积不变}
	设$f:V\longrightarrow W$为线性映射，$f^{\ast}:W^{0}_{r}\longrightarrow V^{0}_{r}$为$f$诱导的拉回映射. 则：\\
对$\forall \phi\in W^{0}_{k}$，$\forall \psi\in W^{0}_{l}$，均有$f^{\ast}(\phi\otimes\psi)=f^{\ast}(\phi)\otimes f^{\ast}(\psi)$. 
\end{mythm}

要说明的是，定理中拉回映射$f^{\ast}:W^{0}_{r}\longrightarrow V^{0}_{r}$的$r$不是固定的一个数，它是可自动调节的，随着作用对象所处空间的改变而改变。比如在表达式$f^{\ast}(\phi\otimes \psi)$中，$r=k+l$；而在表达$f^{\ast}(\phi)$时，$r$取作$k$；在表达$f^{\ast}(\psi)$时，$r$取作$l$.

\begin{proof}
因为$\phi\in W^{0}_{k},\psi\in W^{0}_{l}$，所以$\phi\otimes \psi\in W^{0}_{k+l},f^{\ast}(\phi\otimes \psi)\in V^{0}_{k+l}$. 任取$v_{1},\cdots,v_{k},v_{k+1},\cdots,v_{k+l}\in V$，则
\[
\begin{split}
	&f^{\ast}(\phi\otimes \psi)\left( v_{1},\cdots,v_{k},v_{k+1},\cdots,v_{k+l} \right)\\
	&=\phi\otimes \psi\left( f(v_{1}),\cdots,f(v_{k}),f(v_{k+1}),\cdots,f(v_{k+l}) \right)\\
	&=\phi \left(  f(v_{1}),\cdots,f(v_{k})\right)\cdot \psi \left( f(v_{k+1}),\cdots,f(v_{k+l}) \right)\\
	&=f^{\ast}(\phi) \left( v_{1},\cdots,v_{k} \right)\cdot f^{\ast}(\psi) \left( v_{k+1},\cdots,v_{k+l} \right)\\
	&=f^{\ast}(\phi) \otimes f^{\ast}(\psi) \left( v_{1},\cdots,v_{k},v_{k+1},\cdots,v_{k+l} \right)
\end{split}
\]

因此，$f^{\ast}(\phi\otimes\psi)=f^{\ast}(\phi)\otimes f^{\ast}(\psi)$. 
\end{proof}

上述定理的意义在于，它给出了张量代数$T^{B}(W)$到$T^{B}(V)$的\textbf{代数同态}。让我们来解释一下什么叫代数同态？学过抽象代数的读者一定知道“群同态”的概念，更进一步在环论中还有“环同态”，我们这里所谓的代数同态正是一个与之类似的概念。无论是群同态还是环同态，其核心都是保持结构不变的映射，比如群同态保持的是群之间的乘法关系不变，环同态保持是环之间的加法与乘法关系不变。代数同态亦是如此，它保持两个代数之间的加法、数乘与乘法关系都不变。

由于拉回映射$f^{\ast}:W^{0}_{r}\longrightarrow V^{0}_{r}$是线性映射，所以我们完全可以把$f^{\ast}$线性扩充为$T^{B}(W)\longrightarrow T^{B}(V)$的映射(本质是因为张量代数是分次代数)。这样$f^{\ast}:T^{B}(W)\longrightarrow T^{B}(V)$就自动保持加法与数乘关系不变，即对$\forall x,y\in T^{B}(W),\forall k\in \mr$，均有
\[
\begin{split}
	f^{\ast}(x+y)&=f^{\ast}(x)+f^{\ast}(y)\\
	f^{\ast}(kx)&=kf^{\ast}(x)
\end{split}
\]
现在结合\rthm{thm:chapter2保持张量积不变}，$f^{\ast}:T^{B}(W)\longrightarrow T^{B}(V)$还保持张量积运算，而我们知道张量代数上规定的乘法正好是张量积，从而$f^{\ast}$保持张量代数上的乘法关系不变(本质仍是因为张量代数是分次代数)，即对$\forall x,y\in T^{B}(W)$，均有
\[
f^{\ast}(x\otimes y)=f^{\ast}(x)\otimes f^{\ast}(y)
\]
按照我们上面介绍的代数同态的概念，这表明$f^{\ast}:T^{B}(W)\longrightarrow T^{B}(V)$是张量代数之间的一个代数同态。

下面我们把目光转向更为重要的外形式空间与外代数。作为$W^{0}_{r}$与$V^{0}_{r}$的线性子空间，我们可以把拉回映射$f^{\ast}:W^{0}_{r}\longrightarrow V^{0}_{r}$限制在$\Lambda^{r}(W^{\ast})$与$\Lambda^{r}(V^{\ast})$上，从而得到$f^{\ast}:\Lambda^{r}(W^{\ast})\longrightarrow \Lambda^{r}(V^{\ast})$. 具体的计算方式仍遵循\rdef{def:chapter2拉回映射}，只不过是把作用对象的范围缩小到$(0,r)$型{\color{purple}{反对称}}张量上而已。为此，我们先准备一个引理。

\begin{mylem}{}{chapter2拉回反对称交换引理}
	设$f:V\longrightarrow W$为线性映射，$f^{\ast}:W^{0}_{r}\longrightarrow V^{0}_{r}$为$f$诱导的拉回映射. 则有
\begin{itemize}
\item[(1)] $S_{r}\circ f^{\ast}=f^{\ast}\circ S_{r}$，等号左边的$S_{r}$作用在$V^{0}_{r}$上，等号右边的$S_{r}$作用在$W^{0}_{r}$上.
\item[(2)] $A_{r}\circ f^{\ast}=f^{\ast}\circ A_{r}$，等号左边的$A_{r}$作用在$V^{0}_{r}$上，等号右边的$A_{r}$作用在$W^{0}_{r}$上.
\end{itemize}
\end{mylem}

\begin{proof}
我们只证明第(2)个结论，(1)同理可得。任取$g\in W^{0}_{r}$，则
\[
\begin{split}
	A_{r}\left( f^{\ast}(g) \right)&=\frac{1}{r!}\sum_{\sigma\in G(r)}\operatorname{sgn}(\sigma)\cdot \sigma \left( f^{\ast}(g) \right) \in V^{0}_{r}\\
	f^{\ast} \left( A_{r}(g) \right)&=\frac{1}{r!}\sum_{\sigma\in G(r)}\operatorname{sgn}(\sigma)\cdot f^{\ast}\left(\sigma(g)\right) \in V^{0}_{r}
\end{split}
\]
于是对任意$v_{1},v_{2},\cdots,v_{r}\in V$，直接计算
\[
\begin{split}
	A_{r}\left( f^{\ast}(g) \right) \left( v_{1},v_{2},\cdots,v_{r} \right)&=\frac{1}{r!}\sum_{\sigma\in G(r)}\operatorname{sgn}(\sigma)\cdot f^{\ast}(g)\left( v_{\sigma(1)},v_{\sigma(2)},\cdots,v_{\sigma(r)} \right)\\
	&=\frac{1}{r!}\sum_{\sigma\in G(r)}\operatorname{sgn}(\sigma)\cdot g\left( f(v_{\sigma(1)}),f(v_{\sigma(2)}),\cdots,f(v_{\sigma(r)}) \right)\\
	&=\frac{1}{r!}\sum_{\sigma\in G(r)}\operatorname{sgn}(\sigma)\cdot \sigma(g)\left( f(v_{1}),f(v_{2}),\cdots,f(v_{r}) \right)\\
	&=\frac{1}{r!}\sum_{\sigma\in G(r)}\operatorname{sgn}(\sigma)\cdot f^{\ast}(\sigma(g))\left( v_{1},v_{2},\cdots,v_{r} \right)\\
	&=f^{\ast} \left( A_{r}(g) \right)\left( v_{1},v_{2},\cdots,v_{r} \right)
\end{split}
\]
因此，$A_{r}\left( f^{\ast}(g) \right)=f^{\ast} \left( A_{r}(g) \right)$. 由$g$的任意性便知，$A_{r}\circ f^{\ast}=f^{\ast}\circ A_{r}$.
\end{proof}

用人话来叙述\rlem{lem:chapter2拉回反对称交换引理}，它实际上是说拉回映射保持对称性与反对称性不变：若$\phi$是对称的，则$f^{\ast}(\phi)$也是对称的；若$\phi$是反对称的，则$f^{\ast}(\phi)$也是反对称的。

这种保持张量的对称性与反对称性的特点使得拉回映射是一个极其常用的工具，比如在黎曼几何中，我们可以借助于拉回映射把一个流形上的度量拉回到另一个流形上，得到所谓的诱导度量。

\begin{mythm}{}{chapter2保持外积不变}
	设$f:V\longrightarrow W$为线性映射，$f^{\ast}:\Lambda^{r}(W^{\ast})\longrightarrow \Lambda^{r}(V^{\ast})$为$f$诱导的拉回映射. 则：\\
对$\forall \phi\in \Lambda^{k}(W^{\ast})$，$\forall \psi\in \Lambda^{l}(W^{\ast})$，均有$f^{\ast}(\phi\wedge\psi)=f^{\ast}(\phi)\wedge f^{\ast}(\psi)$. 
\end{mythm}

先说明一下，拉回映射$f^{\ast}:\Lambda^{r}(W^{\ast})\longrightarrow \Lambda^{r}(V^{\ast})$的$r$不是固定的一个数，它是可自动调节的，随着作用对象所处空间的改变而改变。比如在表达式$f^{\ast}(\phi\wedge\psi)$中，$r=k+l$；而在表达$f^{\ast}(\phi)$时，$r$取作$k$；在表达$f^{\ast}(\psi)$时，$r$取作$l$.

\begin{proof}
	\[
	\begin{split}
	f^{\ast}\left( \phi\wedge\psi \right)&=f^{\ast}\left( \frac{(k+l)!}{k!l!} A_{k+l}(\phi\otimes \psi)\right)\\
	&=\frac{(k+l)!}{k!l!} f^{\ast}\left(A_{k+l}(\phi\otimes \psi)\right)\\
	&\xlongequal{\rlem{lem:chapter2拉回反对称交换引理}}\frac{(k+l)!}{k!l!} A_{k+l}\left(f^{\ast}(\phi\otimes \psi)\right)\\
	&\xlongequal{\rthm{thm:chapter2保持张量积不变}}\frac{(k+l)!}{k!l!} A_{k+l}\left(f^{\ast}(\phi)\otimes f^{\ast}(\psi)\right)\\
	&=f^{\ast}(\phi)\wedge f^{\ast}(\psi)
	\end{split}
	\]
\end{proof}

这个定理与\rthm{thm:chapter2保持张量积不变}的意义一样，它给出了外代数$\Lambda(W^{\ast})$与$\Lambda(V^{\ast})$之间的代数同态。这是因为外代数上采用的乘法正好是外积运算，所以在将$f^{\ast}:\Lambda^{r}(W^{\ast})\longrightarrow \Lambda^{r}(V^{\ast})$线性扩充为$\Lambda(W^{\ast})\longrightarrow \Lambda(V^{\ast})$的映射后，$f^{\ast}$保持外代数上的加法、数乘与乘法运算，从而
$f^{\ast}:\Lambda(W^{\ast})\longrightarrow \Lambda(V^{\ast})$便构成一个代数同态。

因此，拉回映射可同时看作是张量代数与外代数上的代数同态。


\section{张量丛}\label{张量丛}

给定一个光滑流形$M^{m}$，对于每个点$p\in M$，在这一点处有切空间$T_{p}M$和余切空间$T^{\ast}_{p}M$. 因此在流形$M$的每一点$p$都可以构造一个$(r,s)$型张量空间：
\[
T^{r}_{s}(p)\triangleq \underbrace{T_{p}M\otimes \cdots\otimes T_{p}M}_{r\text{个}}\otimes \underbrace{T^{\ast}_{p}M\otimes\cdots\otimes T^{\ast}_{p}M}_{s\text{个}}
\]

根据上一节张量空间的有关知识易知，这是一个$m^{r+s}$维向量空间。下面我们要把各点处的$T^{r}_{s}(p)$粘合成一个整体，形成一个新的光滑流形，即所谓的$(r,s)$型张量丛。为了搞清楚什么是张量丛，三个最基本的问题是：
\begin{itemize}
	\item[(1)] 张量丛长什么样子？
	\item[(2)] 张量丛上的拓扑结构是什么？
	\item[(3)] 张量丛上的微分结构是什么？
\end{itemize}
下面我们来详细回答这三个问题。

\textbf{\underline{外貌}}：张量丛在外貌上就是流形$M$上全体$(r,s)$型张量组成的集合，即
		\[
		T^{r}_{s}\triangleq \bigcup\limits_{p\in M}T^{r}_{s}(p)
		\]

实际上更严格来说，这个“并”应该是指“不交并”。设$\{A_{i} \}_{i\in I}$为一族集合，所谓不交并是指
\[
\bigsqcup_{i \in I} A_{i}\triangleq\bigcup_{i \in I}\left(\{i\} \times A_{i}\right)
\]
原本的诸集合$A_{i}$可能会相交，但是在做不交并时，每个$A_{i}$都被$\{i\} \times A_{i}$所替代，而后者总是互不相交的。

对于张量丛来说，普通的并集$\bigcup\limits_{p\in M}T^{r}_{s}(p)$实际上与不交并$\bigsqcup\limits_{p\in M}T^{r}_{s}(p)$是一样的。因为对于流形上两个不同的点$p,q$，它们的切空间$T_{p}M,T_{q}M$已经就是不相交的了，从而张量空间$T^{r}_{s}(p),T^{r}_{s}(q)$也都天生是不相交的。所以在定义$T^{r}_{s}$时，无论是用$\cup$还是用$\sqcup$，都是可以的。不同的书上可能会有不同的说法，读者要明白两种说法其实是一样的。

\textbf{\underline{拓扑结构}}：{\color{red}{设$\cx$为流形$M$的任一坐标卡，局部坐标系为$(x_{1},\cdots,x_{m})$}}. 则对$\forall p\in \x(U)$，切空间$T_{p}M$有自然基底
\[
\left\{ \left( \frac { \partial } { \partial x _ { 1 } } \right) _ { p },\left( \frac { \partial } { \partial x _ { 2 } } \right) _ { p } , \dots , \left( \frac { \partial } { \partial x _ { m } } \right) _ { p } \right\}
\]
在余切空间$T^{\ast}_{p}M$中有相应的对偶基
\[
\left\{(\dd x_{1})_{p},(\dd x_{2})_{p},\cdots,(\dd x_{m})_{p} \right\}
\]
注意，如果按照我们之前的写法，上述对偶基应当严格地写成$\left\{(\dd c_{1})_{p},(\dd c_{2})_{p},\cdots,(\dd c_{n})_{p} \right\}$的形式. 其中$c_{i}=p_{i}\circ\x^{-1}$表示流形上的第$i$个坐标函数(参看\rexa{exa14})，写成上面这个样子只是为了更漂亮一些。读者要牢记，在我们这本笔记的体系下，$x_{i}=p_{i}$是欧氏空间的坐标函数，并不是流形上的坐标函数；流形上的坐标函数还需转换一下，即$c_{i}=p_{i}\circ\x^{-1}$. 

于是我们可以写出张量空间$T^{r}_{s}(p)$的基$\{T^{\otimes}_{\gamma} \}_{\gamma\in \Gamma^{r}_{s}}$，符号的意义如下：
\[
T^{\otimes}_{\gamma}\triangleq \left( \frac { \partial } { \partial x _ { \overline{\gamma_{1}} } } \right) _ { p } \otimes \cdots \otimes \left( \frac { \partial } { \partial x _ {  \overline{\gamma_{r}} } }\right) _ { p }\otimes (\dd x_{\underline{\gamma_{1}}})_{p}\otimes\cdots\otimes(\dd x_{\underline{\gamma_{s}}})_{p} 
\]
这里多重指标集合约定为$\Gamma^{r}_{s}=\{\gamma=( \overline{\gamma_{1}},\cdots, \overline{\gamma_{r}},\underline{\gamma_{1}},\cdots,\underline{\gamma_{s}}) | 1\leqslant\text{每个分量}\leqslant m\}$. 

每一个$(r,s)$型张量$X\in T^{r}_{s}(p)$，在上述基$\{T^{\otimes}_{\gamma} \}_{\gamma\in \Gamma^{r}_{s}}$下都有唯一的坐标. 利用这一点，我们{\color{red}{定义映射$\theta^{U}_{p}:\mr^{m^{r+s}}\longrightarrow T^{r}_{s}(p)$}}如下
\begin{equation}
	\theta^{U}_{p}(\underbrace{a_{\alpha},\cdots,a_{\gamma},\cdots,a_{\beta}}_{m^{r+s}\text{个}})\triangleq a_{\alpha} T^{\otimes}_{\alpha}+\cdots+a_{\gamma} T^{\otimes}_{\gamma}+\cdots+a_{\beta} T^{\otimes}_{\beta}\in T^{r}_{s}(p)
\end{equation}
其中$(a_{\alpha},\cdots,a_{\gamma},\cdots,a_{\beta})\in \mr^{m^{r+s}}$. 索引的指标$\alpha,\cdots\,\gamma,\cdots,\beta$取自$\Gamma^{r}_{s}$. 还有一点要注意，我们这里的符号标识有点多，读者不要搞混了：映射$\theta^{U}_{p}$的{\color{blue}{下标$p$用来标识值域，上标$U$用来标识值域中所用的基底}}。

也许$(\theta^{U}_{p})^{-1}$更容易想象：$(\theta^{U}_{p})^{-1}$就是获取$T^{r}_{s}(p)$中的每一个张量在(由$U$中的局部坐标系所给出的)基底下的坐标。显然这种坐标映射$(\theta^{U}_{p})^{-1}$是向量空间$T^{r}_{s}(p)$与$ \mr^{m^{r+s}}$之间的一个同构，取这个坐标同构的逆映射就是我们这里的$\theta^{U}_{p}$. 

令${\color{red}{\tilde{  U  }=U\times \mr^{m^{r+s}}}}$，显然$\tilde{  U  }$是$\mr^{m}\times\mr^{m^{r+s}}=\mr^{m+m^{r+s}}$中的开集。再定义映射${\color{red}{\tilde{  \x  }:\tilde{  U  }\longrightarrow T^{r}_{s}}}$如下
\begin{equation}
	\tilde{  \x  }\left(a,b \right)=\theta^{U}_{\x(a)}(b)\in T^{r}_{s}(\x(a))
\end{equation}
这里$a=(a_{1},\cdots,a_{m})\in U\subset\mr^{m} \;, b=(a_{\alpha},\cdots,a_{\gamma},\cdots,a_{\beta})\in \mr^{m^{r+s}}\in  \mr^{m^{r+s}}$，所以$\left(a,b \right)$实际上一共有$m+m^{r+s}$个分量。前$m$个分量给出的是张量的位置信息，即得到的是$T^{r}_{s}(\x(a))$这个张量空间里的张量；后$m^{r+s}$个分量给出的是张量的坐标信息，即所得的张量在(由$U$中的局部坐标系所给出的)基底下的坐标就是$(a_{\alpha},\cdots,a_{\gamma},\cdots,a_{\beta})$.

{\color{red}{每给定流形$M$上的一个坐标卡$\cx$，就可以按照上述方式确定相应的$(\tilde{  U  },\tilde{  \x })$}}.

容易验证上述定义的$\tilde{  \x  }$是单射，且$\tilde{  \x  }(\tilde{  U  })=\theta^{U}_{\x(U)}(\mr^{m^{r+s}})=\bigcup\limits_{p\in \x(U)}T^{r}_{s}(p)$. 也就是说$\tilde{  \x  }:\tilde{  U  }\longrightarrow \tilde{  \x  }(\tilde{  U  })$为双射！既然是双射，就意味着我们可以用$\tilde{  \x  }$把$\tilde{  U  }=U\times \mr^{m^{r+s}}$上的拓扑搬运到$\tilde{  U  }$上，规定：
\[
A\subset \tilde{  \x  }(\tilde{  U  })\;\text{为}\;\tilde{  \x  }(\tilde{  U  })\;\text{中的开集}\Leftrightarrow \tilde{  \x  }^{-1}(A)\subset \tilde{  U  }\;\text{为}\;\tilde{  U  }\;\text{中的开集}
\]
其中$\tilde{  U  }=U\times \mr^{m^{r+s}}$上的拓扑是它作为$\mr^{m+m^{r+s}}$的子空间的拓扑。在这样的定义下，映射$\tilde{  \x  }:\tilde{  U  }\longrightarrow \tilde{  \x  }(\tilde{  U  })$就成为了同胚映射。

但是这样只是在$T^{r}_{s}$的局部范围上$\bigcup\limits_{p\in \x(U)}T^{r}_{s}(p)$定义了拓扑，我们需要的是整体上的拓扑。不要慌，问题不大，我们可以用(拓扑)基的方式赋予$T^{r}_{s}$一个整体的拓扑，并且能够保证兼容上述局部定义的拓扑。设$\Sigma$为流形$M$的一个光滑图册，令
\[
\mathscr{B}=\bigcup_{\cx\in \Sigma}\{A\;|\;A\;\text{为}\;\tilde{  \x  }(\tilde{  U  })\;\text{的开集}  \}
\]

可以验证$T^{r}_{s}$的这族子集$\mathscr{B}$满足成为(拓扑)基的两个条件：
\begin{itemize}
	\item $T^{r}_{s}=\bigcup\limits_{A\in \mathscr{B}}A$.
	\item 若$A_{1}\in \mathscr{B},A_{2}\in \mathscr{B}$，则$A_{1}\cap A_{2}$可以表示成$\mathscr{B}$中成员之并.
\end{itemize}
因此，利用子集族$\mathscr{B}$就可以在集合$T^{r}_{s}$上生成一个拓扑$\tau$：
\[
\tau=\{W\subset T^{r}_{s}\;|\;\text{W is a union of members of}\;\mathscr{B} \}
\]
这样我们成功地给$T^{r}_{s}$赋予了一个拓扑结构，并且当流形$M$是$T_{2},A_{2}$的拓扑空间时，这样得到的拓扑空间$(T^{r}_{s},\tau)$也是$T_{2},A_{2}$的。

\begin{mythm}{}{chapter2Trs是拓扑流形}
	$(T^{r}_{s},\tau)$是一个$m+m^{r+s}$维的拓扑流形. 其中$m=\operatorname{dim}(M)$.
\end{mythm}
\begin{proof}
	由上面讨论可知，$(T^{r}_{s},\tau)$为$T_{2},A_{2}$的拓扑空间。
      对任意$X\in T^{r}_{s}$，则存在某个点$p\in M$，使得$X\in T^{r}_{s}(p)$. 而$M$是一个光滑流形，所以存在点$p$处的参数化$\cx$. 再按照上面的式子\requ{equ:}得到$(\tilde{  U  },\tilde{  \x })$. 其中$\tilde{  U  }$是$\mr^{m+m^{r+s}}$中的开集，显然$X\in \tilde{  \x  }(\tilde{  U  })$. 再由拓扑$\tau$的定义方式可知，$\tilde{  \x  }:\tilde{  U  }\longrightarrow \tilde{  \x  }(\tilde{  U  })$为同胚映射。
      这就表明$(T^{r}_{s},\tau)$是一个$m+m^{r+s}$维的拓扑流形。
\end{proof}

\textbf{\underline{微分结构}}：设$\Sigma$为$M$的光滑图册，对任意$\cx\in\Sigma$，按照上面的做法可以得到与之相应的$(\tilde{  U  },\tilde{  \x  })$. 我们来证明$\tilde{  \Sigma  }=\{(\tilde{  U  },\tilde{  \x  }) | \cx\in \Sigma\}$是$T^{r}_{s}$的光滑图册。

首先，由于$\tilde{  \x  }(\tilde{  U  })=\bigcup\limits_{p\in \x(U)}T^{r}_{s}(p)$，所以
\[
\bigcup\limits_{(\tilde{  U  },\tilde{  \x  })\in \tilde{  \Sigma  }}\tilde{  \x  }(\tilde{  U  })=\bigcup\limits_{\cx\in\Sigma}\bigcup\limits_{p\in \x(U)}T^{r}_{s}(p)=\bigcup\limits_{p\in M}T^{r}_{s}(p)=T^{r}_{s}
\]
这表明$\tilde{  \Sigma  }$确实能覆盖$T^{r}_{s}$，从而是$T^{r}_{s}$的一个图册。

其次，选取图册$\tilde{  \Sigma  }$的任意两个成员$(\tilde{  U  },\tilde{  \x  })$和$(\tilde{  V  },\tilde{  \y  }) $，它们分别对应于$M$的两个坐标卡$\cx$和$\cy$. 我们需要说明若$\tilde{  \x  }(\tilde{  U  })\cap\tilde{  \y  }(\tilde{  V  })\neq\varnothing$时，则转换映射$\tilde{  \y  }^{-1}\circ \tilde{  \x}$和$\tilde{  \x  }^{-1}\circ \tilde{  \y  }$都是$C^{\infty}$映射。注意$\tilde{  \x  }(\tilde{  U  })\cap\tilde{  \y  }(\tilde{  V  })\neq\varnothing\Leftrightarrow\x(U)\cap\y(V)\neq\varnothing$.

下面以$\tilde{  \y  }^{-1}\circ \tilde{  \x }$为例来说明它是$C^{\infty}$的。
为此跟踪数据在映射$\tilde{  \y  }^{-1}\circ \tilde{  \x }$之下的变化情况。输入的数据$(\x^{-1}(p),a_{\alpha},\cdots,a_{\gamma},\cdots,a_{\beta})\in U\times\mr^{m^{r+s}}$，首先被$\tilde{  \x  }$合成为一个位于张量空间$T^{r}_{s}(p)$中的张量$X_{p}$，这里$p\in\x(U)\cap\y(V)$. 然后$\tilde{  \y  }^{-1}$解析出该张量$X_{p}$所蕴含的位置信息以及{\color{red}{在由$V$中的局部坐标系所给出的基底下的坐标信息}}并输出。用流程图来看更直观。
\begin{center}
	\begin{tikzpicture}[scale=1, node distance=1.5cm, auto]
	\node [property, text width=25em] (MId1) {\textbf{Input}：$(\x^{-1}(p),a_{\alpha},\cdots,a_{\gamma},\cdots,a_{\beta})\in U\times\mr^{m^{r+s}}$};
	\node [notation, below of=MId1, text width=30em] (ONeg1) {$a_{\alpha} T^{\otimes}_{\alpha}+\cdots+a_{\gamma} T^{\otimes}_{\gamma}+\cdots+a_{\beta} T^{\otimes}_{\beta}=X_{p}=b_{\alpha} \overline{T}^{\otimes}_{\alpha}+\cdots+b_{\gamma} \overline{T}^{\otimes}_{\gamma}+\cdots+b_{\beta} \overline{T}^{\otimes}_{\beta}$};
	\node [operation, below of=ONeg1, text width=25em] (DOS1) {\textbf{Output}：$(\y^{-1}(p),b_{\alpha} ,\cdots,b_{\gamma} ,\cdots,b_{\beta} )\in V\times \mr^{m^{r+s}}$};	
	\path [line] (MId1) -- (ONeg1);
	\path [line] (ONeg1) -- (DOS1);
	\end{tikzpicture}
\end{center}
因此，不难看出$\tilde{  \y  }^{-1}\circ \tilde{  \x }$可以写成如下形式
\[
	\tilde{  \y  }^{-1}\circ \tilde{  \x }(a_{1},\cdots,a_{m},a_{\alpha},\cdots,\cdots,a_{\beta})
	=\left(\y^{-1}\circ\x(a_{1},\cdots,a_{m}) ,(\theta^{V}_{p})^{-1}\circ\theta^{U}_{p}(a_{\alpha},\cdots,\cdots,a_{\beta}) \right)
\]
其中$p=\x(a_{1},\cdots,a_{m})\in \x(U)\cap\y(V)$. 

注意到$\y^{-1}\circ\x$显然已经是$C^{\infty}$的了，所以为了说明$\tilde{  \y  }^{-1}\circ \tilde{  \x }$是$C^{\infty}$的，只要说明$\left(\theta^{V}_{(\cdot)}\right)^{-1}\circ\theta^{U}_{(\cdot)}$在$\x(U)\cap\y(V)$上是$C^{\infty}$映射即可。

\begin{mylem}{}{chapter2张量丛过渡函数的矩阵形式}
	设$\cx,\cy$为光滑流形$M$的任意两个参数化且$\x(U)\cap\y(V)\neq\varnothing$. $\forall p\in\x(U)\cap\y(V),\forall(a_{\alpha},\cdots,a_{\gamma},\cdots,a_{\beta})\in \mr^{m^{r+s}}$，若记$(\theta^{V}_{p})^{-1}\circ\theta^{U}_{p}(a_{\alpha},\cdots,a_{\gamma},\cdots,a_{\beta})=(b_{\alpha},\cdots,b_{\gamma},\cdots,b_{\beta})$，则有
	\[
	\left( \begin{array}{ccccc}{\vdots} \\ {b_{\gamma}}  \\ {\vdots} \\ {\vdots}\\ {\vdots}\end{array}\right)=\underbrace{J(\y^{-1}\circ\x)\otimes \cdots\otimes J(\y^{-1}\circ\x)}_{r\text{个}}\otimes \underbrace{\left(J(\x^{-1}\circ\y) \right)^{T}\otimes\cdots \otimes \left(J(\x^{-1}\circ\y)\right)^{T}}_{s\text{个}} \left( \begin{array}{ccccc}{\vdots} \\ {\vdots}  \\ {\vdots} \\ {a_{\lambda}}\\ {\vdots}\end{array}\right)
	\]
	这里$J(\y^{-1}\circ\x)$表示映射$\y^{-1}\circ\x$点$\x^{-1}(p)$处的Jacobi矩阵，$J(\x^{-1}\circ\y) $表示映射$x^{-1}\circ\y$在点$\y^{-1}(p)$为Jacobi矩阵. 
\end{mylem}
\begin{proof}
	$(a_{\alpha},\cdots,a_{\gamma},\cdots,a_{\beta})$可看作是$T^{r}_{s}(p)$中某个张量$X_{p}$在由$U$中的局部坐标系所给出的$T^{r}_{s}(p)$的基底$\{T^{\otimes}_{\gamma}\}_{\gamma\in \Gamma^{r}_{s}}$下的坐标；$(b_{\alpha},\cdots,b_{\gamma},\cdots,b_{\beta})$是同一个张量$X_{p}$在由$V$中的局部坐标系所给出的$T^{r}_{s}(p)$的基底$\{\overline{T}^{\otimes}_{\gamma}\}_{\gamma\in \Gamma^{r}_{s}}$下的坐标. 这里基$\{T^{\otimes}_{\gamma}\}_{\gamma\in \Gamma^{r}_{s}}$与基$\{\overline{T}^{\otimes}_{\gamma}\}_{\gamma\in \Gamma^{r}_{s}}$默认都是按照字典顺序排列的有序基，因此坐标$(a_{\alpha},\cdots,a_{\gamma},\cdots,a_{\beta})$与$(b_{\alpha},\cdots,b_{\gamma},\cdots,b_{\beta})$的分量也都是按照字典顺序从左到右排列的。
	
	所以{\color{red}{$(\theta^{V}_{p})^{-1}\circ\theta^{U}_{p}$的功能就是对张量空间$T^{r}_{s}(p)$中的元素在不同基底下的坐标进行转换}}。而张量的坐标转换我们熟啊，前面刚刚研究过！我们把那里的结论直接拿来套用就ok了。
	
	设从有序基$\{T^{\otimes}_{\gamma}\}_{\gamma\in \Gamma^{r}_{s}}$到有序基$\{\overline{T}^{\otimes}_{\gamma}\}_{\gamma\in \Gamma^{r}_{s}}$的过渡矩阵为$C$，即
	\[
	\{\overline{T}^{\otimes}_{\gamma}\}_{\gamma\in \Gamma^{r}_{s}}=\{T^{\otimes}_{\gamma}\}_{\gamma\in \Gamma^{r}_{s}}\cdot C
	\]
	则
	\[
	\left( \begin{array}{ccccc}{\vdots} \\ {b_{\gamma}}  \\ {\vdots} \\ {\vdots}\\ {\vdots}\end{array}\right)=C^{-1}\left( \begin{array}{ccccc}{\vdots} \\ {\vdots}  \\ {\vdots} \\ {a_{\lambda}}\\ {\vdots}\end{array}\right)
	\]
	
	注意到$T^{r}_{s}(p)$中的基之所以会从$\{T^{\otimes}_{\gamma}\}_{\gamma\in \Gamma^{r}_{s}}$变成$\{\overline{T}^{\otimes}_{\gamma}\}_{\gamma\in \Gamma^{r}_{s}}$，是因为底层的向量空间$T_{p}M$中的基从$\left\{ \left( \frac { \partial } { \partial x _ { 1 } } \right) _ { p },\left( \frac { \partial } { \partial x _ { 2 } } \right) _ { p } , \dots , \left( \frac { \partial } { \partial x _ { m } } \right) _ { p } \right\}$变成了$\left\{ \left( \frac { \partial } { \partial y _ { 1 } } \right) _ { p },\left( \frac { \partial } { \partial y _ { 2 } } \right) _ { p } , \dots , \left( \frac { \partial } { \partial y _ { m } } \right) _ { p } \right\}$. 因此，若设
	\[
	\left\{ \left( \frac { \partial } { \partial y _ { 1 } } \right) _ { p },\left( \frac { \partial } { \partial y _ { 2 } } \right) _ { p } , \dots , \left( \frac { \partial } { \partial y _ { m } } \right) _ { p } \right\}=\left\{ \left( \frac { \partial } { \partial x _ { 1 } } \right) _ { p },\left( \frac { \partial } { \partial x _ { 2 } } \right) _ { p } , \dots , \left( \frac { \partial } { \partial x _ { m } } \right) _ { p } \right\}\cdot A
	\]
	则根据上一节的\rthm{thm:chapter2张量空间之间的基过渡矩阵}可知，$C=A\otimes\cdots\otimes A\otimes (A^{T})^{-1}\otimes\cdots\otimes (A^{T})^{-1}$. 从而
	\[
	C^{-1}=A^{-1}\otimes \cdots\otimes A^{-1}\otimes A^{T}\otimes\cdots \otimes A^{T}
	\]
	我们又知道
	\[
	{\color{red}{A=J(\x^{-1}\circ\y)}}=\left( \begin{array}{ccc}{\frac { \partial x_{1} } { \partial y _ { 1 } } }  & {\cdots} & {\frac { \partial x_{1} } { \partial y _ { m } } } \\ {\frac { \partial x_{2} } { \partial y _ { 1 } } }  & {\cdots} & {\frac { \partial x_{2} } { \partial y _ { m } } } \\ {\vdots}  & { \cdots} & {\vdots} \\ {\frac { \partial x_{m} } { \partial y _ { 1 } } }  & {\cdots} & {\frac { \partial x_{m} } { \partial y _ { m } } }\end{array}\right)_{\y^{-1}(p)}{\color{red}{A^{-1}=J(\y^{-1}\circ\x)}}=\left( \begin{array}{ccc}{\frac { \partial y_{1} } { \partial x _ { 1 } } }  & {\cdots} & {\frac { \partial y_{1} } { \partial x _ { m } } } \\ {\frac { \partial y_{2} } { \partial x _ { 1 } } }  & {\cdots} & {\frac { \partial y_{2} } { \partial x _ { m } } } \\ {\vdots}  & {\cdots } & {\vdots} \\ {\frac { \partial y_{m} } { \partial x _ { 1 } } }  & {\cdots} & {\frac { \partial y_{m} } { \partial x _ { m } } }\end{array}\right)_{\x^{-1}(p)}
	\]
	
	所以
	$
	C^{-1}=J(\y^{-1}\circ\x)\otimes \cdots\otimes J(\y^{-1}\circ\x)\otimes \left(J(\x^{-1}\circ\y) \right)^{T}\otimes\cdots \otimes \left(J(\x^{-1}\circ\y)\right)^{T}
	$.
\end{proof}

假设$\tilde{  U  }$中的局部坐标系为$(x_{1},\cdots,x_{m},a_{\alpha},\cdots,a_{\gamma},\cdots,a_{\beta})$，$\tilde{  V  }$中的局部坐标系为$(y_{1},\cdots,y_{m},b_{\alpha},\cdots,b_{\gamma},\cdots,b_{\beta})$. 结合上述引理，我们可以把局部坐标系之间的变换关系具体写出来：
首先，显然有
\begin{equation}
	(y_{1},\cdots,y_{m})=\y^{-1}\circ\x(x_{1},\cdots,x_{m})
\end{equation}
其次，由矩阵张量积的性质可知，$C^{-1}$中第$\gamma$行第$\lambda$列的元素为
\[\left(\frac { \partial y_{\overline{\gamma_{1}}} } { \partial x _ { \overline{\lambda_{1}} } }\right)\cdot \left(\frac { \partial y_{\overline{\gamma_{2}}} } { \partial x _ { \overline{\lambda_{2}} } }\right)\cdots\left(\frac { \partial y_{\overline{\gamma_{r}}} } { \partial x _ { \overline{\lambda_{r}} } }\right)\left(\frac { \partial x_{\underline{\lambda_{1}}} } { \partial y _ { \underline{\gamma_{1}} } }\right)\cdot\left(\frac { \partial x_{\underline{\lambda_{2}}} } { \partial y _ { \underline{\gamma_{2}} } }\right)\cdots\left(\frac { \partial x_{\underline{\lambda_{s}}} } { \partial y _ { \underline{\gamma_{s}} } }\right)\]
也就有
\[
\left( \begin{array}{ccccc}{\vdots} \\ {b_{\gamma}}  \\ {\vdots} \\ {\vdots}\\ {\vdots}\end{array}\right)=\left( \begin{array}{ccccc}{ } & { } & { } & { }&{} \\ { \cdots} & {\cdots} & {\cdots} & { \boxed{\left(\frac { \partial y_{\overline{\gamma_{1}}} } { \partial x _ { \overline{\lambda_{1}} } }\right)\cdots\left(\frac { \partial y_{\overline{\gamma_{r}}} } { \partial x _ { \overline{\lambda_{r}} } }\right)\left(\frac { \partial x_{\underline{\lambda_{1}}} } { \partial y _ { \underline{\gamma_{1}} } }\right)\cdots\left(\frac { \partial x_{\underline{\lambda_{s}}} } { \partial y _ { \underline{\gamma_{s}} } }\right)}} &{\cdots}\\ { } & { } & { } & { } &{}\\ { } & { } & { } & { }&{}\\ { } & { } & { } & { }&{}\end{array}\right)\left( \begin{array}{ccccc}{\vdots} \\ {\vdots}  \\ {\vdots} \\ {a_{\lambda}}\\ {\vdots}\end{array}\right)
\]
所以
\begin{equation}
	b_{\gamma}=\sum\limits_{\lambda\in \Gamma^{r}_{s}}\left(\frac { \partial y_{\overline{\gamma_{1}}} } { \partial x _ { \overline{\lambda_{1}} } }\right)\cdots\left(\frac { \partial y_{\overline{\gamma_{r}}} } { \partial x _ { \overline{\lambda_{r}} } }\right)\left(\frac { \partial x_{\underline{\lambda_{1}}} } { \partial y _ { \underline{\gamma_{1}} } }\right)\cdots\left(\frac { \partial x_{\underline{\lambda_{s}}} } { \partial y _ { \underline{\gamma_{s}} } }\right) a_{\lambda}
\end{equation}

这样我们就把$(\tilde{  U  },\tilde{  \x  })$与$(\tilde{  V  },\tilde{  \y  })$在交集部分$\tilde{  \x  }(\tilde{  U  })\cap \tilde{  \y  }(\tilde{  V  })$的局部坐标变换完全写了出来，即式\requ{equ:}与式\requ{equ:}. 由此可知，$\tilde{  \y  }^{-1}\circ\tilde{  \x  }$是$C^{\infty}$映射。

我们将以上讨论总结为如下定理。

\begin{mythm}{}{chapter2Trs的微分结构}
	设$\Sigma$为$M$的光滑图册. 对任意$\cx\in\Sigma$，定义与之相应的$(\tilde{  U  },\tilde{  \x })$如下：
	\[
	\tilde{  U  }\triangleq U\times \mr^{m^{r+s}}
	\]
	$\forall (a_{1},\cdots,a_{m},a_{\alpha},\cdots,a_{\gamma},\cdots,a_{\beta})\in \tilde{  U  }$，定义$\tilde{ \x  }:\tilde{  U  }\longrightarrow T^{r}_{s}$为
	\[
	\tilde{  \x  }(a_{1},\cdots,a_{m},a_{\alpha},\cdots,a_{\gamma},\cdots,a_{\beta})\triangleq \theta^{U}_{p}(a_{\alpha},\cdots,a_{\gamma},\cdots,a_{\beta})
	\]
	其中$p=\x(a_{1},\cdots,a_{m})$. 将所有这样的$(\tilde{  U  },\tilde{  \x })$收集起来得到$\tilde{  \Sigma  }=\{(\tilde{  U  },\tilde{  \x  }) | \cx\in \Sigma\}$，则$\tilde{  \Sigma  }$是$T^{r}_{s}$的光滑图册，再把它极大化后便确定了$T^{r}_{s}$的微分结构.
\end{mythm}


\subsection{切丛}\label{切丛}

简单来说，切丛就是$(1,0)$型张量丛。不过相比于一般$(r,s)$型张量丛的复杂性，初学者可能较难把握，切丛显得更为直观简单。如果你被上一节那些乱糟糟的符号搞得一头雾水的话，不妨先从这一小节开始读。看起来很复杂的张量丛与较为简单的切丛在精神内核上是一脉相承的。本节并不准备展开像上一节那样细致的讨论，而主要是罗列一些实用的概念和结论以备查找参考。

给定一个$m$维流形$M$. 对于每个点$p\in M$，在这一点处有切空间$T_{p}M$. 将所有这样的切空间做并集得到
\[
TM\triangleq\bigcup_{p\in M}T_{p}M
\]
这就是切丛的外在长相，它是流形$M$上全体切向量组成的集合。切丛是最容易直观想象的一种丛，\rfig{圆的切丛}展示了单位圆的切丛。

\begin{figure}[!htp]
\centering
\newcommand\pgfmathsinandcos[3]{%
\pgfmathsetmacro#1{sin(#3)}%
\pgfmathsetmacro#2{cos(#3)}%
}
\newcommand\LongitudePlane[3][current plane]{%
\pgfmathsinandcos\sinEl\cosEl{#2} % elevation
\pgfmathsinandcos\sint\cost{#3} % azimuth
\tikzset{#1/.estyle={cm={\cost,\sint*\sinEl,0,\cosEl,(0,0)}}}
}
\newcommand\LatitudePlane[3][current plane]{%
\pgfmathsinandcos\sinEl\cosEl{#2} % elevation
\pgfmathsinandcos\sint\cost{#3} % latitude
\pgfmathsetmacro\yshift{\cosEl*\sint}
\tikzset{#1/.style={cm={\cost,0,0,\cost*\sinEl,(0,\yshift)}}} %
}
\def\angPhiOne{-130} % longitude of point P
\def\angPhiTwo{-50} % longitude of point Q
\newcommand\DrawLongitudeCircle[2][1]{
    \LongitudePlane{\angEl}{#2}
    \tikzset{current plane/.estyle={cm={\cost,\sint*\sinEl,0,\cosEl,(0,0)},scale=#1}}
    \pgfmathsetmacro\angVis{atan(sin(#2)*cos(\angEl)/sin(\angEl))} %
    \draw[current plane,thin,black] (\angVis:1) arc (\angVis:\angVis+180:1);
    \draw[current plane,thin,dashed] (\angVis-180:1) arc (\angVis-180:\angVis:1);
}
\newcommand\DrawLongitudeCirclered[2][1]{
    \LongitudePlane{\angEl}{#2}
    \tikzset{current plane/.estyle={cm={\cost,\sint*\sinEl,0,\cosEl,(0,0)},scale=#1}}
    \pgfmathsetmacro\angVis{atan(sin(#2)*cos(\angEl)/sin(\angEl))} %
    \draw[current plane,red,thick] (90:1) arc (90:180:1);   
}
\newcommand\DrawLatitudeCircle[2][1]{
    \LatitudePlane{\angEl}{#2}
    \tikzset{current plane/.prefix style={scale=#1}}
    \pgfmathsetmacro\sinVis{sin(#2)/cos(#2)*sin(\angEl)/cos(\angEl)}
    \pgfmathsetmacro\angVis{asin(min(1,max(\sinVis,-1)))}
    \draw[current plane,ultra thick,blue] (\angVis:1) arc (\angVis:-\angVis-180:1);
    \draw[current plane,ultra thick,blue] (180-\angVis:1) arc (180-\angVis:\angVis:1);
}
\newcommand\DrawLatitudeCirclered[2][1]{
    \LatitudePlane{\angEl}{#2}
    \tikzset{current plane/.prefix style={scale=#1}}
    \pgfmathsetmacro\sinVis{sin(#2)/cos(#2)*sin(\angEl)/cos(\angEl)}
    \pgfmathsetmacro\angVis{asin(min(1,max(\sinVis,-1)))}
    \draw[current plane,red,thick] (\angPhiTwo:1) node[below right] {$$} arc (\angPhiTwo:\angPhiOne:1) node[below left] {$$}; %Point Q suivi du point P
      \foreach \r in {-250,-240,-230,-220,-210,-200,-190,-180,-170,-160,-150,-140,-130,-120,-110,...,-50,-40,-30,-20,-10,0,10,20,30,40,50,60,70,80,90}{
    \draw[current plane,red] (\r:1) -- ([turn]90:2);
    \draw[current plane,red] (\r:1) -- ([turn]-90:2);
    }
}
\tikzset{%
>=latex,
inner sep=0pt,%
outer sep=2pt,%
mark coordinate/.style={inner sep=0pt,outer sep=0pt,minimum size=3pt,
fill=black,circle}%
}
\begin{tikzpicture}[scale=1,every node/.style={minimum size=1cm}]
\def\R{1.25} % sphere radius
\def\angEl{100} % 倾斜角度
\def\angAz{-100} % azimuth angle
\def\angBeta{30} % latitude of point P and Q
\DrawLatitudeCircle[\R]{0} % equator
\DrawLatitudeCirclered[\R]{0}
\end{tikzpicture}
\caption{单位圆的切丛}
\label{圆的切丛}
\end{figure}
接下来给$TM$赋予拓扑结构和微分结构，使之也成为一个微分流形。根据上一节对张量丛上的拓扑结构和微分结构的讨论，切丛上也可以按完全相同的方式赋予拓扑和微分结构，这里不赘述相关的细节，只给出最重要的局部坐标系的构造，这是最实用的。

任取光滑流形$M$的一个局部坐标系$\cx$. 对于点$p\in \x(U)$，我们知道
\[
	\left\{ \left( \frac { \partial } { \partial x _ { 1 } } \right) _ { p },\left( \frac { \partial } { \partial x _ { 2 } } \right) _ { p } , \dots , \left( \frac { \partial } { \partial x _ { m } } \right) _ { p } \right\}
\]
是$p$处切空间$T_{p}M$的基。在此基底下，任一切向量$X\in T_{p}M$可以表为$X=\sum\limits_{i=1}^{m}X_{i}\left( \frac { \partial } { \partial x _ { i } } \right) _ { p }$. 也就是说$X$在此自然基底下的坐标是$(X_{1},X_{2},\cdots, X_{m})$. 

令$\theta^{U}_{p}:\mr^{m}\longrightarrow T_{p}M$定义为：对$\forall \left( b_{1},b_{2},\cdots, b_{m} \right)\in \mr^{m} $
\begin{equation}
\begin{split}
	\theta^{U}_{p} \left( b_{1},b_{2},\cdots, b_{m} \right)\triangleq b_{1}\left( \frac { \partial } { \partial x _ { 1 } } \right) _ { p }+b_{2}\left( \frac { \partial } { \partial x _ { 2 } } \right) _ { p }+\cdots +b_{m}\left( \frac { \partial } { \partial x _ { m } } \right) _ { p }
\end{split}
\end{equation}

读者还记不记得我们在线性代数中如何证明“数域$\mathbb{F}$上的任意$n$维向量空间都与$\mathbb{F}^{n}$同构”这一结论？那里我们就是通过给向量空间取定基底，然后把向量映射到它在基下的坐标来证明的。这里$\theta^{U}_{p}$实际上也就是向量空间$\mr^{m}$与向量空间$T_{p}M$之间的坐标同构映射。

再令$\tilde{U}=U\times \mr^{m}$，则$\tilde{U}$显然是$\mr^{m}\times \mr^{m}=\mr^{2m}$的开集。令
$\tilde{\x}:\tilde{U}\longrightarrow TM$定义为：对$\forall (a_{1},\cdots,a_{m},b_{1},\cdots,b_{m})\in \mr^{2m}$
\begin{equation}
\begin{split}
\tilde{\x}({\color{red}{a_{1},\cdots,a_{m}}},b_{1},\cdots,b_{m})\triangleq \theta^{U}_{\color{red}{p}}(b_{1},\cdots,b_{m})\;,\;\text{\kaishu 这里} {\color{red}{p=\x(a_{1},\cdots,a_{m})}}
\end{split}
\end{equation}

{\color{red}{每给定流形$M$上的一个坐标卡$\cx$，就可以按照上述方式确定相应的$(\tilde{  U  },\tilde{  \x })$}}. 这些$(\tilde{  U  },\tilde{  \x })$就是切丛$TM$的\textbf{局部坐标卡}，将这些局部坐标卡收集起来再极大化就是$TM$的\textbf{微分结构}，于是$TM$成为一个$2m$维光滑流形。

切丛的不同局部坐标卡之间的转换映射也是一个经常要用到的性质。设$(\tilde{  U  },\tilde{  \x })$与$(\tilde{  V  },\tilde{  \y })$是$TM$的任意两个局部坐标卡，它们分别对应于$M$的两个坐标卡$\cx$与$\cy$. 如若$\tilde{  \x }(\tilde{  U })\cap \tilde{  \y }(\tilde{  V })\neq\varnothing$，那么也有$\x(U)\cap \y(V)\neq\varnothing$. 此时我们有
\[
	\tilde{  \y }^{-1}\circ\tilde{  \x }=\left( \y^{-1}\circ\x,(\theta^{V}_{p} )^{-1}\circ \theta^{U}_{p}\right)
\]
其中$p=\x(U)\cap \y(V)$. 

所以为了描述转换映射$\tilde{  \y }^{-1}\circ\tilde{  \x }$，我们得先知道$(\theta^{V}_{p} )^{-1}\circ \theta^{U}_{p}$该如何刻画。

\begin{mylem}{}{chapter2切丛过渡函数的矩阵形式}
	$\forall p\in\x(U)\cap\y(V),\forall(b_{1},b_{2},\cdots,b_{m})\in \mr^{m}$，若记$(\theta^{V}_{p})^{-1}\circ\theta^{U}_{p}(b_{1},b_{2},\cdots,b_{m})=(b^{\prime}_{1},b^{\prime}_{2},\cdots,b^{\prime}_{m})$，则有
	\begin{equation}
	\left( \begin{array}{cccc}{b^{\prime}_{1}} \\ {b^{\prime}_{2}}  \\ {\vdots} \\ {b^{\prime}_{m}}\end{array}\right)=\left[\begin{array}{ccccc}
		\frac { \partial y _ { 1 } } { \partial x _ { 1 } }& \frac { \partial y _ { 1 } } { \partial x _ { 2 } }& \cdots &\frac { \partial y _ { 1 } } { \partial x _ { m } }\\
		\frac { \partial y _ { 2 } } { \partial x _ { 1 } }& \frac { \partial y _ { 2 } } { \partial x _ { 2 } }& \cdots &\frac { \partial y _ { 2 } } { \partial x _ { m } }\\
		\vdots & \vdots &\ddots & \vdots\\
		\frac { \partial y _ { m } } { \partial x _ { 1 } }& \frac { \partial y _ { m } } { \partial x _ { 2 } }& \cdots &\frac { \partial y _ { m} } { \partial x _ { m } }
		\end{array}\right] \left( \begin{array}{cccc}{b_{1}} \\ {b_{2}}  \\ {\vdots} \\ {b_{m}}\end{array}\right)
			\label{转化式}
	\end{equation}
\end{mylem}
\begin{proof}
在\rlem{lem:chapter2张量丛过渡函数的矩阵形式}中，取$r=1,s=0$即得。
\end{proof}

因此，对$\forall (a_{1},\cdots,a_{m},b_{1},\cdots,b_{m})\in \mr^{2m}$，我们有
\[
\begin{split}
	\tilde{\y }^{-1}\circ\tilde{\x }( a_{1},\cdots,a_{m},b_{1},\cdots,b_{m})&= \bigl(\y^{-1}\circ\x (a_{1},\cdots,a_{m}),(\theta^{V}_{p} )^{-1}\circ \theta^{U}_{p}(b_{1},\cdots,b_{m}) \bigr)\\
	&=\bigl(\y^{-1}\circ\x (a_{1},\cdots,a_{m}),b^{\prime}_{1},b^{\prime}_{2},\cdots,b^{\prime}_{m} \bigr)
\end{split}
\]
其中$(b^{\prime}_{1},b^{\prime}_{2},\cdots,b^{\prime}_{m})$与$(b_{1},b_{2}\cdots,b_{m})$之间的转换关系由\rlem{lem:chapter2切丛过渡函数的矩阵形式}中的式\requ{转化式}所刻画。

\subsection{余切丛}

所谓余切丛就是$(0,1)$型张量丛。余切丛没有切丛那么直观的图像，不过比起一般$(r,s)$型张量丛又要简单得多。本节与\ref{切丛}节一样，仍是以罗列一些实用的概念和结论为主，方便以后来查找和参考。

给定一个$m$维流形$M$. 对于每个点$p\in M$，在这一点处有余切空间$T^{\ast}_{p}M$. 将所有这样的余切空间做并集得到
\[
T^{\ast}M\triangleq\bigcup_{p\in M}T^{\ast}_{p}M
\]
这就是余切丛的外在长相，它是流形$M$上全体余切向量组成的集合。

接下来给$T^{\ast}M$赋予拓扑结构和微分结构，使之也成为一个微分流形。
根据上一节对张量丛上的拓扑结构和微分结构的讨论，余切丛上也可以按完全相同的方式赋予拓扑和微分结构，这里不赘述相关的细节，只给出最重要的局部坐标系的构造，这是最实用的。

任取光滑流形$M$的一个局部坐标系$\cx$. 对于点$p\in \x(U)$，我们知道
\[
\left\{(\dd x_{1})_{p},(\dd x_{2})_{p},\cdots,(\dd x_{m})_{p} \right\}
\]
是$p$处余切空间$T^{\ast}_{p}M$的基。在此基底下，任一余切向量$\omega\in T^{\ast}_{p}M$可以表为$\omega=\sum\limits_{i=1}^{m}\omega_{i}\cdot (\dd x_{i})_{p}$. 也就是说$\omega$在此基底下的坐标是$(\omega_{1},\omega_{2},\cdots, \omega_{m})$. 

令$\eta^{U}_{p}:\mr^{m}\longrightarrow T^{\ast}_{p}M$定义为：对$\forall \left( b_{1},b_{2},\cdots, b_{m} \right)\in \mr^{m} $
\begin{equation}
\begin{split}
	\eta^{U}_{p} \left( b_{1},b_{2},\cdots, b_{m} \right)\triangleq b_{1}(\dd x_{1})_{p}+b_{2}(\dd x_{2})_{p}+\cdots +b_{m}(\dd x_{m})_{p}
\end{split}
\end{equation}

再令$\bar{U}=U\times \mr^{m}$，则$\bar{U}$显然是$\mr^{m}\times \mr^{m}=\mr^{2m}$的开集。令
$\bar{\x}:\bar{U}\longrightarrow T^{\ast}M$定义为：对$\forall (a_{1},\cdots,a_{m},b_{1},\cdots,b_{m})\in \mr^{2m}$
\begin{equation}
\begin{split}
	\bar{\x}({\color{red}{a_{1},\cdots,a_{m}}},b_{1},\cdots,b_{m})\triangleq \eta^{U}_{\color{red}{p}}(b_{1},\cdots,b_{m})\;,\;\text{\kaishu 这里} {\color{red}{p=\x(a_{1},\cdots,a_{m})}}
\end{split}
\end{equation}

{\color{red}{每给定流形$M$上的一个坐标卡$\cx$，就可以按照上述方式确定相应的$(\bar{  U  },\bar{  \x })$}}. 这些$(\bar{  U  },\bar{  \x })$就是余切丛$T^{\ast}M$的\textbf{局部坐标卡}，将这些局部坐标卡收集起来再极大化就是$T^{\ast}M$的\textbf{微分结构}，于是$T^{\ast}M$成为一个$2m$维光滑流形。

余切丛的不同局部坐标卡之间的转换映射也是一个经常要用到的性质。设$(\bar{  U  },\bar{  \x })$与$(\bar{  V  },\bar{  \y })$是$T^{\ast}M$的任意两个局部坐标卡，它们分别对应于$M$的两个坐标卡$\cx$与$\cy$. 如若$\bar{  \x }(\bar{  U })\cap \bar{  \y }(\bar{  V })\neq\varnothing$，那么也有$\x(U)\cap \y(V)\neq\varnothing$. 此时我们有
\[
	\bar{  \y }^{-1}\circ\bar{  \x }=\left( \y^{-1}\circ\x,(\eta^{V}_{p} )^{-1}\circ \eta^{U}_{p}\right)
\]
其中$p=\x(U)\cap \y(V)$. 

所以为了描述转换映射$\bar{  \y }^{-1}\circ\bar{  \x }$，我们得先知道$(\eta^{V}_{p} )^{-1}\circ \eta^{U}_{p}$该如何刻画。

\begin{mylem}{}{chapter2余切丛过渡函数的矩阵形式}
	$\forall p\in\x(U)\cap\y(V),\forall(b_{1},b_{2},\cdots,b_{m})\in \mr^{m}$，若记$(\eta^{V}_{p})^{-1}\circ\eta^{U}_{p}(b_{1},b_{2},\cdots,b_{m})=(b^{\prime}_{1},b^{\prime}_{2},\cdots,b^{\prime}_{m})$，则有
	\begin{equation}
	\left( \begin{array}{cccc}{b^{\prime}_{1}} \\ {b^{\prime}_{2}}  \\ {\vdots} \\ {b^{\prime}_{m}}\end{array}\right)=\left[\begin{array}{ccccc}
		\frac { \partial x _ { 1 } } { \partial y _ { 1 } }& \frac { \partial x _ { 2 } } { \partial y _ { 1 } }& \cdots &\frac { \partial x _ { m } } { \partial y _ { 1 } }\\
		\frac { \partial x _ { 1 } } { \partial y _ {2 } }& \frac { \partial x _ { 2 } } { \partial y _ { 2 } }& \cdots &\frac { \partial x _ { m } } { \partial y _ { 2 } }\\
		\vdots & \vdots &\ddots & \vdots\\
		\frac { \partial x _ { 1 } } { \partial y _ { m } }& \frac { \partial x _ { 2 } } { \partial y _ { m } }& \cdots &\frac { \partial x _ { m} } { \partial y _ { m } }
		\end{array}\right] \left( \begin{array}{cccc}{b_{1}} \\ {b_{2}}  \\ {\vdots} \\ {b_{m}}\end{array}\right)
			\label{转化式2}
	\end{equation}
\end{mylem}
\begin{proof}
在\rlem{lem:chapter2张量丛过渡函数的矩阵形式}中，取$r=0,s=1$即得。
\end{proof}

因此，对$\forall (a_{1},\cdots,a_{m},b_{1},\cdots,b_{m})\in \mr^{2m}$，我们有
\[
\begin{split}
	\bar{\y }^{-1}\circ\bar{\x }( a_{1},\cdots,a_{m},b_{1},\cdots,b_{m})&= \bigl(\y^{-1}\circ\x (a_{1},\cdots,a_{m}),(\eta^{V}_{p} )^{-1}\circ \eta^{U}_{p}(b_{1},\cdots,b_{m}) \bigr)\\
	&=\bigl(\y^{-1}\circ\x (a_{1},\cdots,a_{m}),b^{\prime}_{1},b^{\prime}_{2},\cdots,b^{\prime}_{m} \bigr)
\end{split}
\]
其中$(b^{\prime}_{1},b^{\prime}_{2},\cdots,b^{\prime}_{m})$与$(b_{1},b_{2}\cdots,b_{m})$之间的转换关系由\rlem{lem:chapter2余切丛过渡函数的矩阵形式}中的式\requ{转化式2}所刻画。

\section{$s$次外形式丛}\label{s次外形式丛}

我们看到张量丛$T^{r}_{s}$实际上就是给流形的每一点$p$赋予$(r,s)$型张量空间$T^{r}_{s}(p)$，特别地有切丛和余切丛。本节我们来构造一类新的丛，构造方法与构造张量丛时完全一样！这类丛\textbf{非常非常非常}重要。

给定一个光滑流形$M^{m}$，对于每个点$p\in M$，在这一点处有张量空间$T^{0}_{s}(p)$和$T^{r}_{0}(p)$. 考虑由它们中的反对称张量构成的子空间，也即$s$次外形式空间$\Lambda^{s}(T^{\ast}_{p}M)$或者$r$次外向量空间$\Lambda^{r}(T_{p}M)$. 特别地，根据\rdef{def:chapter2斜边张量的对称与反对称空间}我们约定，$\Lambda^{1}(T^{\ast}_{p}M)=T^{\ast}_{p}M,\Lambda^{0}(T^{\ast}_{p}M)=\mr$. 

下面我们要把各点处的$\Lambda^{s}(T^{\ast}_{p}M)$粘合成一个整体，形成一个新的光滑流形，即所谓的$s$次外形式丛。为了搞清楚什么是$s$次外形式丛，三个最基本的问题是：
\begin{itemize}
	\item[(1)] $s$次外形式丛长什么样子？
	\item[(2)] $s$次外形式丛上的拓扑结构是什么？
	\item[(3)] $s$次外形式丛上的微分结构是什么？
\end{itemize}

在正式回答这三个问题之前，我们强调一下：从原理上来说，你也可以把各点处的$r$次外向量空间$\Lambda^{r}(T_{p}M)$粘和为一个整体，构造所谓的\textbf{$r$次外向量丛}。不过比起外向量，更常用的还是外形式，所以在几何里面我们基本只考虑外形式，本节也只对$s$次外形式丛做出构造。

\textbf{\underline{外貌}}：$s$次外形式丛在外貌上就是流形$M$上全体$s$阶反对称协变张量组成的集合：
		\[
			\Lambda^{s}(T^{\ast}M)\triangleq \bigcup\limits_{p\in M}\Lambda^{s}(T^{\ast}_{p}M)
		\]

\textbf{\underline{拓扑结构}}：{\color{red}{设$\cx$为$M$的任一坐标卡，局部坐标系为$(x_{1},\cdots,x_{m})$}}. 则对$\forall p\in \x(U)$，余切空间$T^{\ast}_{p}M$有基底$\left\{(\dd x_{1})_{p},(\dd x_{2})_{p},\cdots,(\dd x_{m})_{p} \right\}$. 于是我们可以写出$\Lambda^{s}(T^{\ast}_{p}M)$的基$\{\left( \dd \x \right)^{\wedge}_{I} |I\in Q(s,m) \}$，符号的意义如下：
		\[
			\left( \dd \x \right)^{\wedge}_{I}\triangleq (\dd x_{i_{1}})_{p}\wedge (\dd x_{i_{2}})_{p}\wedge\cdots\wedge(\dd x_{i_{s}})_{p},\; \text{\kaishu 其中}I=(i_{1},\cdots, i_{s})
		\]
		这里$Q(s,m)=\{I=(i_{1},i_{2},\cdots, i_{s})|1\leqslant i_{1}<\leqslant i_{2}<\cdots<i_{s}\leqslant m \}$.
		
		$\Lambda^{s}(T^{\ast}_{p}M)$中的每一个反对称张量$X$在基$\{\left( \dd \x \right)^{\wedge}_{I} |I\in Q(s,m) \}$下有唯一的坐标，可简写为$X=\sum\limits_{I\in Q(s,m)}X_{I}\left( \dd \x \right)^{\wedge}_{I}$. 于是我们{\color{red}{定义映射$\theta^{U}_{p}:\mr^{\binom{m}{s}}\longrightarrow \Lambda^{s}(T^{\ast}_{p}M)$}}如下：
		\begin{equation}
			\theta^{U}_{p}(\underbrace{a_{\alpha},\cdots,a_{\gamma},\cdots,a_{\beta}}_{\binom{m}{s}\text{个}})\triangleq a_{\alpha} \left( \dd \x \right)^{\wedge}_{\alpha}+\cdots+a_{\gamma} \left( \dd \x \right)^{\wedge}_{\gamma}+\cdots+a_{\beta} \left( \dd \x \right)^{\wedge}_{\beta}
		\end{equation}
		其中$(a_{\alpha},\cdots,a_{\gamma},\cdots,a_{\beta})\in \mr^{\binom{m}{s}}$. 索引的指标$\alpha,\cdots\,\gamma,\cdots,\beta$取自$Q(s,m)$. 映射$\theta^{U}_{p}$的{\color{blue}{下标$p$用来标识值域，上标$U$用来标识值域中所用的基底}}。
		
		容易看出，$(\theta^{U}_{p})^{-1}$就是获取$\Lambda^{s}(T^{\ast}_{p}M)$中的每一个张量在(由$U$中的局部坐标系所给出的)基底下的坐标。显然这种坐标映射$(\theta^{U}_{p})^{-1}$是向量空间$\Lambda^{s}(T^{\ast}_{p}M)$与$ \mr^{\binom{m}{s}}$之间的一个同构，取这个坐标同构的逆映射就是我们这里的$\theta^{U}_{p}$. 
		
		令${\color{red}{\tilde{  U  }=U\times \mr^{\binom{m}{s}}}}$，显然$\tilde{  U  }$是$\mr^{m}\times\mr^{\binom{m}{s}}=\mr^{m+\binom{m}{s}}$中的开集。再定义映射${\color{red}{\tilde{  \x  }:\tilde{  U  }\longrightarrow \Lambda^{s}(T^{\ast}M)}}$如下：$\forall(a_{1},\cdots,a_{m})\in U,\forall (a_{\alpha},\cdots,\alpha_{\gamma},\cdots,\alpha_{\beta})\in \mr^{\binom{m}{s}}$
		\begin{equation}
			\tilde{  \x  }\left(a_{1},\cdots,a_{m},a_{\alpha},\cdots,a_{\gamma},\cdots,a_{\beta} \right)=\theta^{U}_{p}(a_{\alpha},\cdots,a_{\gamma},\cdots,a_{\beta})
		\end{equation}
		这里$p=\x(a_{1},\cdots,a_{m})$，指标$\alpha,\cdots\,\gamma,\cdots,\beta\in Q(s,m)$.

		该映射定义域的前$m$个分量给出的是张量的位置信息，即得到的是$\Lambda^{s}(T^{\ast}_{p}M)$这个张量空间里的张量；后$\binom{m}{s}$个分量给出的是张量的坐标信息，即所得的张量在(由$U$中的局部坐标系所给出的)基底下的坐标就是$(a_{\alpha},\cdots,a_{\gamma},\cdots,a_{\beta})$.
		{\color{red}{每给定流形$M$上的一个坐标卡$\cx$，就可以按照上述方式确定相应的$(\tilde{  U  },\tilde{  \x })$}}.
		
		容易验证上述定义的$\tilde{  \x  }$是单射，且$\tilde{  \x  }(\tilde{  U  })=\theta^{U}_{\x(U)}(\mr^{\binom{m}{s}})=\bigcup\limits_{p\in \x(U)}\Lambda^{s}(T^{\ast}_{p}M)$. 也就是说$\tilde{  \x  }:\tilde{  U  }\longrightarrow \tilde{  \x  }(\tilde{  U  })$为双射！双射意味着我们可以用$\tilde{  \x  }$把$\tilde{  U  }=U\times \mr^{m^{\binom{m}{s}}}$上的拓扑搬运到$\tilde{  U  }$上，规定：
		\[
		A\subset \tilde{  \x  }(\tilde{  U  })\;\text{\kaishu 为}\;\tilde{  \x  }(\tilde{  U  })\;\text{\kaishu 中的开集}\Leftrightarrow \tilde{  \x  }^{-1}(A)\subset \tilde{  U  }\;\text{\kaishu 为}\;\tilde{  U  }\;\text{\kaishu 中的开集}
		\]
		其中$\tilde{  U  }=U\times \mr^{\binom{m}{s}}$上的拓扑是它作为$\mr^{m+\binom{m}{s}}$的子空间的拓扑。在这样的定义下，映射$\tilde{  \x  }:\tilde{  U  }\longrightarrow \tilde{  \x  }(\tilde{  U  })$就成为了同胚映射。
		
		但这样只是在局部范围$\bigcup\limits_{p\in \x(U)}\Lambda^{s}(T^{\ast}_{p}M)$上定义了拓扑，我们需要的是整体上的拓扑。不要慌，问题不大，我们可以用(拓扑)基的方式赋予$\Lambda^{s}(T^{\ast}M)$一个整体的拓扑，并且能够保证兼容上述局部定义的拓扑。设$\Sigma$为流形$M$的一个光滑图册，令
		\[
		\mathscr{B}=\bigcup_{\cx\in \Sigma}\{A\;|\;A\;\text{为}\;\tilde{  \x  }(\tilde{  U  })\;\text{\kaishu 的开集}  \}
		\]
		可以验证$\Lambda^{s}(T^{\ast}M)$的这族子集$\mathscr{B}$满足成为(拓扑)基的两个条件：
		\begin{itemize}
			\item $\Lambda^{s}(T^{\ast}M)=\bigcup\limits_{A\in \mathscr{B}}A$.
			\item 若$A_{1}\in \mathscr{B},A_{2}\in \mathscr{B}$，则$A_{1}\cap A_{2}$可以表示成$\mathscr{B}$中成员之并.
		\end{itemize}
		因此利用子集族$\mathscr{B}$就可以在集合$\Lambda^{s}(T^{\ast}M)$上生成一个拓扑$\tau$：
		\[
		\tau=\{W\subset \Lambda^{s}(T^{\ast}M)\;|\;\text{W is a union of members of}\;\mathscr{B} \}
		\]
		这样我们成功地给$\Lambda^{s}(T^{\ast}M)$赋予了一个拓扑结构，并且当流形$M$是$T_{2},A_{2}$的拓扑空间时，如上得到的拓扑空间$(\Lambda^{s}(T^{\ast}M),\tau)$也是$T_{2},A_{2}$的。这表明$(\Lambda^{s}(T^{\ast}M),\tau)$是一个$m+\binom{m}{s}$维的拓扑流形。

\textbf{\underline{微分结构}}：设$\Sigma$为$M$的光滑图册，对任意$\cx\in\Sigma$，按照上面的做法可以得到与之相应的$(\tilde{  U  },\tilde{  \x  })$. 我们来证明$\tilde{  \Sigma  }=\{(\tilde{  U  },\tilde{  \x  }) | \cx\in \Sigma\}$是$\Lambda^{s}(T^{\ast}M)$的光滑图册。

首先，由于$\tilde{  \x  }(\tilde{  U  })=\bigcup\limits_{p\in \x(U)}\Lambda^{s}(T^{\ast}_{p}M)$，所以
	\[
		\bigcup\limits_{(\tilde{  U  },\tilde{  \x  })\in \tilde{  \Sigma  }}\tilde{  \x  }(\tilde{  U  })=\bigcup\limits_{\cx\in\Sigma}\bigcup\limits_{p\in \x(U)}\Lambda^{s}(T^{\ast}_{p}M)=\bigcup\limits_{p\in M}\Lambda^{s}(T^{\ast}_{p}M)=\Lambda^{s}(T^{\ast}M)
		\]
这表明$\tilde{  \Sigma  }$确实能覆盖$\Lambda^{s}(T^{\ast}M)$，从而是$\Lambda^{s}(T^{\ast}M)$的一个图册。
		
其次，选取图册$\tilde{  \Sigma  }$的任意两个成员$(\tilde{  U  },\tilde{  \x  })$和$(\tilde{  V  },\tilde{  \y  }) $，它们分别对应于$M$的两个坐标卡$\cx$和$\cy$. 如若$\tilde{  \x  }(\tilde{  U  })\cap\tilde{  \y  }(\tilde{  V  })\neq\varnothing$，那么也有$\x(U)\cap\y(V)\neq\varnothing$. 下面以$\tilde{  \y  }^{-1}\circ \tilde{  \x }$为例来说明转换映射是$C^{\infty}$的。不难看出$\tilde{  \y  }^{-1}\circ \tilde{  \x }$可以写成如下形式
	\[
		\tilde{  \y  }^{-1}\circ \tilde{  \x }
		=\left(\y^{-1}\circ\x,(\theta^{V}_{p})^{-1}\circ\theta^{U}_{p}\right)
	\]
其中$p\in \x(U)\cap\y(V)$. 
		
注意到$\y^{-1}\circ\x$已经是$C^{\infty}$的了，所以为了说明$\tilde{  \y  }^{-1}\circ \tilde{  \x }$是$C^{\infty}$的，只要说明$\left(\theta^{V}_{p}\right)^{-1}\circ\theta^{U}_{p}$在$\x(U)\cap\y(V)$上是$C^{\infty}$映射即可。
		
\begin{mylem}{}{chapter2s次外形式丛过渡函数的矩阵形式}
		$\forall p\in\x(U)\cap\y(V),\forall(\cdots,a_{I},\cdots)\in \mr^{\binom{m}{s}}$，若记$(\theta^{V}_{p})^{-1}\circ\theta^{U}_{p}(\cdots,a_{I},\cdots)=(\cdots,b_{J},\cdots)$，则有
		\[
		b_{J}=\sum_{I\in Q(s,m)}A \left(\begin{array}{cccc} j_{1} & j_{2} & \cdots & j_{r} \\i_{1} & i_{2} & \cdots & i_{r}\end{array}\right)\cdot a_{I}
		\]
		这里$A=\left( J(\x^{-1}\circ \y ) \right)^{T} $表示映射$\x^{-1}\circ\y$在点$\y^{-1}(p)$的Jacobi矩阵的转置矩阵. 
\end{mylem}
\begin{proof}
	$(\cdots,a_{I},\cdots)$可看作是$\Lambda^{s}(T^{\ast}_{p}M)$中某个张量$X_{p}$在由$U$中的局部坐标系所给出的基底$\{\left( \dd \x \right)^{\wedge}_{\gamma} |\gamma\in Q(s,m)\}$下的坐标；$(\cdots,b_{J},\cdots)$是同一个张量$X_{p}$在由$V$中的局部坐标系所给出的基底$\{\left( \dd \y \right)^{\wedge}_{\gamma} |\gamma\in Q(s,m)\}$下的坐标. 这里基$\{\left( \dd \x \right)^{\wedge}_{\gamma} |\gamma\in Q(s,m)\}$与基$\{\left( \dd \y \right)^{\wedge}_{\gamma} |\gamma\in Q(s,m)\}$默认都是按照字典顺序排列的有序基，因此坐标$(\cdots,a_{I},\cdots)$与$(\cdots,b_{J},\cdots)$的分量也都是按照字典顺序从左到右排列的。
			
	由此可知{\color{red}{$(\theta^{V}_{p})^{-1}\circ\theta^{U}_{p}$的功能就是对张量空间$\Lambda^{s}(T^{\ast}_{p}M)$中的元素在不同基底下的坐标进行转换}}。只要找出从基$\{\left( \dd \x \right)^{\wedge}_{\gamma} |\gamma\in Q(s,m)\}$到基$\{\left( \dd \y \right)^{\wedge}_{\gamma} |\gamma\in Q(s,m)\}$的过渡矩阵$C$，就有
	\[
		\left( \begin{array}{ccccc}{\vdots} \\ {b_{J}}  \\ {\vdots} \\ {\vdots}\\ {\vdots}\end{array}\right)=C^{-1}\left( \begin{array}{ccccc}{\vdots} \\ {\vdots}  \\ {\vdots} \\ {a_{I}}\\ {\vdots}\end{array}\right)
	\]

	设$\left( (\dd y_{1})_{p},(\dd y_{2})_{p},\cdots,(\dd y_{m})_{p} \right)=\left( (\dd x_{1})_{p},(\dd x_{2})_{p},\cdots,(\dd x_{m})_{p} \right)B$，则我们知道
	\[
		B=\left( J(\y^{-1}\circ \x ) \right)^{T}=\left[ \begin{array}{cccc}{\frac { \partial y_{1} } { \partial x _ { 1 } } }  & {\frac { \partial y_{2} } { \partial x _ { 1 } } }  & \cdots & {\frac { \partial y_{m} } { \partial x _ { 1 } } }  \\ {\frac { \partial y_{1} } { \partial x _ { 2 } } }  & {\frac { \partial y_{2} } { \partial x _ { 2 } } }  & \cdots & {\frac { \partial y_{m} } { \partial x _ {2 } } }  \\ {\vdots}  & {\vdots } & {\ddots} & {\vdots}\\ {\frac { \partial y_{1} } { \partial x _ { m } } }  & {\frac { \partial y_{2} } { \partial x _ { m } } }  & \cdots & {\frac { \partial y_{m} } { \partial x _ { m } } }  \end{array}\right]
	\]
根据\rthm{thm:chapter2反对称空间的基过渡矩阵}可知，矩阵$C^{-1}$第$\alpha$行第$\beta$列处的元素为
\[
	C^{-1}_{\alpha\beta}=B^{-1}\left(\begin{array}{cccc}\alpha_{1} & \alpha_{2} & \cdots & \alpha_{s} \\\beta_{1} & \beta_{2} & \cdots & \beta_{s}\end{array}\right)=\left( J(\x^{-1}\circ \y ) \right)^{T}\left(\begin{array}{cccc}\alpha_{1} & \alpha_{2} & \cdots & \alpha_{s} \\\beta_{1} & \beta_{2} & \cdots & \beta_{s}\end{array}\right)
\]
于是
\begin{equation}
\begin{split}
b_{J}=\sum_{I\in Q(s,m)}C^{-1}_{JI}\cdot a_{I}&=\sum_{I\in Q(s,m)}\left( J(\x^{-1}\circ \y ) \right)^{T}\left(\begin{array}{cccc} j_{1} & j_{2} & \cdots & j_{r} \\i_{1} & i_{2} & \cdots & i_{r}\end{array}\right)\cdot a_{I}\\
	&=\sum_{I\in Q(s,m)}
		\left| \begin{array}{ccc}{\frac { \partial x_{i_{1}} } { \partial y _ { j_{1} } } }  & {\cdots} & {\frac { \partial x_{i_{s}} } { \partial y _ { j_{1} } } }  \\ {\vdots}  & { \ddots} & {\vdots} \\ {\frac { \partial x_{i_{1}} } { \partial y _ { j_{s} } } }  & {\cdots} & {\frac { \partial x_{i_{s}} } { \partial y _ { j_{s} } } }\end{array}\right|\cdot a_{I}
\end{split}
\label{外形式丛的坐标转换}
\end{equation}
令$A=\left( J(\x^{-1}\circ \y ) \right)^{T} $即得。
\end{proof}

从式\requ{外形式丛的坐标转换}显然可以看出，$(\theta^{V})^{-1}\circ \theta^{U}$在$\x(U)\cap \y(V)$中是$C^{\infty}$映射。因此，$\tilde{\y}^{-1}\circ \tilde{\x}$也是$C^{\infty}$的。

\begin{remark}
	上述引理的证明也可通过外积的求值公式\requ{外积求值公式}直接计算得到，具体做法如下：设$X_{p}=\sum\limits_{I\in Q(s,m)}a_{I}\left( \dd\x \right)^{\wedge}_{I}=\sum\limits_{J\in Q(s,m)}b_{J}\left( \dd\y \right)^{\wedge}_{J}$. 首先由求值公式\requ{外积求值公式}以及广义Kronecker符号的性质，我们有
	\[
	a_{I}=a_{i_{1}\cdots i_{s}}=X \left[ \left( \frac { \partial  } { \partial x_{i_{1}}} \right)_{p},\cdots,\left( \frac { \partial  } { \partial x_{i_{s}}} \right)_{p} \right]\quad\quad b_{J}=b_{j_{1}\cdots j_{s}}=X \left[ \left( \frac { \partial  } { \partial y_{j_{1}}} \right)_{p},\cdots,\left( \frac { \partial  } { \partial y_{j_{s}}} \right)_{p} \right]
	\]
	于是
	\[
	\begin{split}
	b_{J}&=X \left[ \left( \frac { \partial  } { \partial y_{j_{1}}} \right)_{p},\cdots,\left( \frac { \partial  } { \partial y_{j_{s}}} \right)_{p} \right]\\
	&=X \left[  \sum_{k_{1}=1}^{m}\frac { \partial x_{k_{1}} } { \partial y_{j_{1}}}\left( \frac { \partial  } { \partial x_{k_{1}}} \right)_{p},\cdots, \sum_{k_{s}=1}^{m}\frac { \partial x_{k_{s}} } { \partial y_{j_{s}}}\left( \frac { \partial  } { \partial x_{k_{s}}} \right)_{p} \right]\\
	&=\sum_{1\leqslant k_{1},\cdots,k_{s}\leqslant m}\frac { \partial x_{k_{1}} } { \partial y_{j_{1}}}\cdots \frac { \partial x_{k_{s}} } { \partial y_{j_{s}}} X \left[ \left( \frac { \partial  } { \partial x_{k_{1}}} \right)_{p},\cdots,\left( \frac { \partial  } { \partial x_{k_{s}}} \right)_{p} \right]\\
	&=\sum_{1\leqslant i_{1}<\cdots<i_{s}\leqslant m}\left[\sum_{\sigma\in G(s)} \frac { \partial x_{i_{\sigma(1)}} } { \partial y_{j_{1}}}\cdots \frac { \partial x_{i_{\sigma(s)}} } { \partial y_{j_{s}}} X \left[ \left( \frac { \partial  } { \partial x_{i_{\sigma(1)}}} \right)_{p},\cdots,\left( \frac { \partial  } { \partial x_{i_{\sigma(s)}}} \right)_{p} \right]\right]\\
	&=\sum_{1\leqslant i_{1}<\cdots<i_{s}\leqslant m}\left[\sum_{\sigma\in G(s)}\operatorname{sgn}(\sigma) \frac { \partial x_{i_{\sigma(1)}} } { \partial y_{j_{1}}}\cdots \frac { \partial x_{i_{\sigma(s)}} } { \partial y_{j_{s}}} X \left[ \left( \frac { \partial  } { \partial x_{i_{1}}} \right)_{p},\cdots,\left( \frac { \partial  } { \partial x_{i_{s}}} \right)_{p} \right]\right]\\
	&=\sum_{I\in Q(s,m)}
		\left| \begin{array}{ccc}{\frac { \partial x_{i_{1}} } { \partial y _ { j_{1} } } }  & {\cdots} & {\frac { \partial x_{i_{s}} } { \partial y _ { j_{1} } } }  \\ {\vdots}  & { \ddots} & {\vdots} \\ {\frac { \partial x_{i_{1}} } { \partial y _ { j_{s} } } }  & {\cdots} & {\frac { \partial x_{i_{s}} } { \partial y _ { j_{s} } } }\end{array}\right|\cdot a_{I}
	\end{split}
	\]
\end{remark}

综上所述，$\tilde{  \Sigma  }=\{(\tilde{  U  },\tilde{  \x  }) | \cx\in \Sigma\}$是$\Lambda^{s}(T^{\ast}M)$的光滑图册，将这个图册极大化后便得到$\Lambda^{s}(T^{\ast}M)$的光滑结构。也就是说，对于流形$M$上的每一个局部坐标卡$\cx$，构建一个与之对应的$(\tilde{U},\tilde{\x})$，把所有这些$(\tilde{U},\tilde{\x})$收集起来然后极大化，就是$\Lambda^{s}(T^{\ast}M)$上的光滑结构。

以上讨论表明，$\Lambda^{s}(T^{\ast}M)$是一个$m+\binom{m}{s}$维的光滑流形，其中$m=\operatorname{dim}(M)$. 特别地，当$s=0$时，$\Lambda^{0}(T^{\ast}M)=M\times \mr$，即$0$次外形式丛就是乘积流形$M\times \mr$；当$s=1$时，$\Lambda^{1}(T^{\ast}M)=T^{\ast}M$，即$1$次外形式丛就是余切丛。

\section{向量丛}

在前面两节中，我们分别介绍了张量丛和$s$次外形式丛。注意到无论是张量丛，还是$s$次外形式丛，本质都是在光滑流形$M$的\textbf{每一点粘上一个向量空间}：$(r,s)$型张量丛粘上的是$(r,s)$型张量空间，$s$次外形式丛粘上的是$s$次外形式空间。注意到这一点之后，我们可以把张量丛的构造进一步抽象化：直接给流形的每一点配上一个向量空间。

\begin{mydef}{向量丛}{chapter2向量丛}
	设$E,M$为两个光滑流形，$\pi:E\longrightarrow M$为光滑的满射. 若$(E,M,\pi)$满足以下条件：
	\begin{itemize}
		\item[(1)] 对$\forall p\in M$，集合$E_{p}\triangleq \pi^{-1}(p)$具有$k$维实向量空间的结构.
		\item[(2)] 对$\forall p\in M$，存在包含点$p$的开集$W\subset M$和映射$\varPhi:W\times \mr^{k}\longrightarrow \pi^{-1}(W)$，使得
		\begin{itemize}
			\item[$\bullet$] $\varPhi:W\times \mr^{k}\longrightarrow \pi^{-1}(W)$是光滑的微分同胚映射；
			\item[$\bullet$] $\forall q\in W,\forall a\in \mr^{k}$，均有$\varPhi(q,a)\in E_{q}$；
			\item[$\bullet$]  任意冻结一点$q\in W$，映射$\varPhi(q,\;{\color{red}{\cdot}}\;): \mr^{k}\longrightarrow E_{q}$是向量空间之间的同构.
		\end{itemize}
	\end{itemize}
则称$(E,M,\pi)$为光滑流形$M$上的一个秩为$k$的光滑实向量丛. 
\end{mydef}

我们称$E$为丛的\textbf{全空间}，$M$为丛的\textbf{底空间}，$\pi$为\textbf{丛投影}，$E_{p}=\pi^{-1}(p)$称为丛在点$p$处的\textbf{纤维}，$\mr^{k}$称为\textbf{纤维型}，微分同胚$\varPhi:W\times \mr^{k}\longrightarrow \pi^{-1}(W)$称为\textbf{向量丛的局部平凡化}。秩为$1$的光滑实向量丛通常称为\textbf{线丛}，这个名称很直观。

\begin{example}
	\label{chapter2张量丛是向量丛}
    光滑流形$M^{m}$上的$(r,s)$型张量丛$T^{r}_{s}$是一个秩为$m^{r+s}$的光滑实向量丛.
    \soln
    
    丛投影$\pi:T^{r}_{s}\longrightarrow M$为：对$\forall X\in T^{r}_{s}$，则存在某个唯一的$p\in M$使得$X\in T^{r}_{s}(p)$，令$\pi(X)=p$. 下面给出张量丛的局部平凡化.
    
    对任意一点$p\in M$，有$M$在点$p$处的参数化$\cx$. 令$W=\x(U)\subset M$. 再定义$\varPhi:W\times \mr^{m^{r+s}}\longrightarrow \pi^{-1}(W)$为：$\forall (a_{\eta},\cdots,a_{\gamma},\cdots,a_{\xi})\in \mr^{m^{r+s}}$
    \[
    \varPhi(p,a_{\eta},\cdots,a_{\gamma},\cdots,a_{\xi})\triangleq \theta^{U}_{p}(a_{\eta},\cdots,a_{\gamma},\cdots,a_{\xi})
    \]
    其中指标$\eta,\gamma,\xi\in \Gamma^{r}_{s}$，并按照字典顺序排列. 则可以验证$\{\varPhi:W\times \mr^{m^{r+s}}\longrightarrow \pi^{-1}(W) \}$就是张量丛的局部参数化.
\end{example}

就像切丛、余切丛和张量丛一样，我们也可以通过给流形$M$的每点指定一个向量空间的方式来构造向量丛：先把各点处指定的向量空间作不交并，然后给得到的不交并集合安装拓扑结构和微分结构。
\begin{mythm}{}{chapter2向量丛的等价定义}
	设$M$是一个光滑流形，每一点$p\in M$处都给定了一个$k$维实向量空间$E_{p}$. 令$E=\bigcup\limits_{p\in M}E_{p}$，再定义映射$\pi:E\longrightarrow M$为：对$\forall X\in E$，则存在某个唯一的$p\in M$使得$X\in E_{p}$，令$\pi(X)=p$. 如果存在$M$的开覆盖$\{U_{\alpha} \}_{\alpha\in \Lambda}$，并且对每一个$U_{\alpha}$，存在相应的{\color{red}{双射}}$\varPhi_{\alpha}:U_{\alpha}\times \mr^{k}\longrightarrow \pi^{-1}(U_{\alpha})$，满足以下条件：
	\begin{itemize}
		\item[(1)] $\forall q\in U_{\alpha},\forall a\in \mr^{k}$，均有$\varPhi_{\alpha}(q,a)\in E_{q}$；
		\item[(2)]  任意冻结一点$q\in U_{\alpha}$，映射$\varPhi_{\alpha}(q,\;{\color{red}{\cdot}}\;): \mr^{k}\longrightarrow E_{q}$是向量空间之间的同构；
		\item[(3)] 当$U_{\alpha}\cap U_{\beta}\neq\varnothing$时，存在光滑映射$t_{\beta\alpha}:U_{\beta}\cap U_{\alpha}\longrightarrow GL_{k}(\mr)$，使得对$\;\forall p\in U_{\beta}\cap U_{\alpha}\;,\forall(a_{1},\cdots,a_{k})\in \mr^{k}$，$\varPhi_{\beta}^{-1}\circ\varPhi_{\alpha}:(U_{\beta}\cap U_{\alpha})\times\mr^{k}\longrightarrow (U_{\beta}\cap U_{\alpha})\times\mr^{k}$具有以下形式
		\[
		\varPhi_{\beta}^{-1}\circ\varPhi_{\alpha}(p,a_{1},\cdots,a_{k})=(p,b_{1},\cdots,b_{k})
		\]
		其中
		\[
		\left( \begin{array}{ccc}{b_{1}} \\ {\vdots}  \\ {b_{k}} \end{array}\right)=t_{\beta\alpha}(p) \left( \begin{array}{ccc}{a_{1}} \\ {\vdots}  \\ {a_{k}} \end{array}\right)
		\]
	\end{itemize}
	则$E$上可赋予拓扑结构和微分结构，使得$(E,M,\pi)$成为$M$上的一个秩为$k$的光滑向量丛，并且以$\{\varPhi_{\alpha}:U_{\alpha}\times \mr^{k}\longrightarrow \pi^{-1}(U_{\alpha}) \}_{\alpha\in \Lambda}$作为它的局部平凡化.
\end{mythm}

\begin{proof}
	\textbf{\underline{拓扑结构}}：对任意$p\in M$，则存在某个包含点$p$的$U_{\alpha}$. 选取流形$M$在点$p$处的参数化$\cx$，使得$p\in \x(U)\subset U_{\alpha}$. 
	
	令${\color{red}{\tilde{  U  }=U\times \mr^{k}}}$，显然$\tilde{  U  }$是$\mr^{m}\times\mr^{k}=\mr^{m+k}$中的开集。再定义映射${\color{red}{\tilde{  \x  }:\tilde{  U  }\longrightarrow  E}}$如下
	\begin{equation}
	\tilde{  \x  }\left(x_{1},\cdots,x_{m},a \right)=\varPhi_{\alpha}\left(q,a \right)\in E_{q}
	\end{equation}
	这里$q=\x(x_{1},\cdots,x_{m})\in \x(U)\subset U_{\alpha},a\in \mr^{k}$. 
	
	我们断言$\tilde{  \x  }:\tilde{  U  }\longrightarrow \tilde{  \x  }(\tilde{  U  })=\pi^{-1}\left(\x(U) \right)$为双射。
	\begin{itemize}
		\item 若$\tilde{  \x  }\left(x_{1},\cdots,x_{m},a \right)=\tilde{  \x  }\left(x'_{1},\cdots,x'_{m},b \right)$. 设$q=\x(x_{1},\cdots,x_{m}),q^{\prime}=\x(x'_{1},\cdots,x'_{m})$，则$\varPhi_{\alpha}\left(q,a \right)=\varPhi_{\alpha}\left(q^{\prime},b \right)$. 注意到$q\in \x(U)\subset U_{\alpha}, q^{\prime}\in \x(U)\subset U_{\alpha}$，且$\varPhi_{\alpha}:U_{\alpha}\times \mr^{k}\longrightarrow \pi^{-1}(U_{\alpha})$是双射，所以$\left(q,a \right)=\left(q^{\prime},b \right)$. 于是有$q=q^{\prime},a=b$. 又因为$\x:U\longrightarrow \x(U)$是双射，所以$\x(x_{1},\cdots,x_{m})=q=q^{\prime}=\x(x'_{1},\cdots,x'_{m})\Rightarrow (x_{1},\cdots,x_{m})=(x'_{1},\cdots,x'_{m})$. 这表明$\tilde{  \x }$是单射。
		\item 任取${\color{red}{X}}\in \pi^{-1}\left(\x(U) \right)$，则存在某个唯一(因为是不交并)的$q\in \x(U)$，使得$X\in E_{q}$. 冻结这个点$q\in\x(U)\subset U_{\alpha}$，因为映射$\varPhi_{\alpha}(q,\underline{\hspace{0.3cm}}\;): \mr^{k}\longrightarrow E_{q}$是向量空间之间的同构，所以在此同构映射下，必存在$a\in \mr^{k}$，使得$\varPhi_{\alpha}(q,{\color{red}{a}})={\color{red}{X}}$. 假设$\x^{-1}(q)=(x_{1},\cdots,x_{m})$，则$\tilde{  \x  }(x_{1},\cdots,x_{m},a)=\varPhi_{\alpha}(q,{\color{red}{a}})={\color{red}{X}}$. 这表明$\tilde{  \x  }$是满射。
	\end{itemize}
    既然是双射，就意味着我们可以用$\tilde{  \x  }$把$\tilde{  U  }=U\times \mr^{k}$上的拓扑搬运到$\tilde{  U  }$上。规定：
	\[
	A\subset \tilde{  \x  }(\tilde{  U  })\;\text{为}\;\tilde{  \x  }(\tilde{  U  })\;\text{中的开集}\Leftrightarrow \tilde{  \x  }^{-1}(A)\subset \tilde{  U  }\;\text{为}\;\tilde{  U  }\;\text{中的开集}
	\]
	其中$\tilde{  U  }=U\times \mr^{m^{r+s}}$上的拓扑是它作为$\mr^{m+k}$的子空间的拓扑。在这样的定义下，映射$\tilde{  \x  }:\tilde{  U  }\longrightarrow \tilde{  \x  }(\tilde{  U  })$就成为了同胚映射。注意到{\color{red}{每取定流形$M$上的一点$p$，就可以按照上述方式确定相应的$(\tilde{  U  },\tilde{  \x })$}}，于是我们再令
	\[
	\mathscr{B}=\bigcup_{p\in M}\{A\;|\;A\;\text{为}\;\tilde{  \x  }(\tilde{  U  })\;\text{的开集}  \}
	\]
	
	可以验证$E$的这族子集$\mathscr{B}$满足成为(拓扑)基的两个条件：
	\begin{itemize}
		\item $E=\bigcup\limits_{A\in \mathscr{B}}A$.
		\item 若$A_{1}\in \mathscr{B},A_{2}\in \mathscr{B}$，则$A_{1}\cap A_{2}$可以表示成$\mathscr{B}$中成员之并.
	\end{itemize}
	因此，利用子集族$\mathscr{B}$就可以在集合$E$上生成一个拓扑$\tau$：
	\[
	\tau=\{W\subset E\;|\;\text{W is a union of members of}\;\mathscr{B} \}
	\]
	这样我们给集合$T^{r}_{s}$赋予了一个拓扑结构，并且当流形$M$是$T_{2},A_{2}$的拓扑空间时，这样得到的拓扑空间$(E,\tau)$也是$T_{2},A_{2}$的。
	
	\textbf{\underline{微分结构}}：{\color{red}{每取定流形$M$上的一点$p$，就可以按照上述方式确定一个与之相应的$(\tilde{  U  },\tilde{  \x })$}}. 我们来证明$\tilde{  \Sigma  }=\{(\tilde{  U  },\tilde{  \x  }) | p\in M\}$是$E$的光滑图册。
	
	首先，由于$\tilde{  \x  }(\tilde{  U  })=\pi^{-1}(\x(U))$，并且随着点$p$而取的参数化$\cx$在$p$遍历$M$的同时，覆盖了$M$. 所以
	\[
	\bigcup\limits_{p\in M}\tilde{  \x  }(\tilde{  U  })
	=\bigcup\limits_{p\in M}\pi^{-1}(\x(U))=\pi^{-1}\left(\bigcup\limits_{p\in M}\x(U) \right)= \pi^{-1}(M)=E
	\]
	这表明$\tilde{  \Sigma  }$能覆盖$E$，从而是$E$的一个图册。
	
	其次，选取图册$\tilde{  \Sigma  }$的任意两个成员$(\tilde{  U  },\tilde{  \x  })$和$(\tilde{  V  },\tilde{  \y  }) $，它们分别对应于$M$的两个坐标卡$\cx$和$\cy$，并且假设$\x(U)\subset U_{\alpha},\y(V)\subset U_{\beta}$. 我们需要说明若$\tilde{  \x  }(\tilde{  U  })\cap\tilde{  \y  }(\tilde{  V  })\neq\varnothing$时，则转换映射$\tilde{  \y  }^{-1}\circ \tilde{  \x}$和$\tilde{  \x  }^{-1}\circ \tilde{  \y  }$都是$C^{\infty}$映射。
	
	下面以$\tilde{  \y  }^{-1}\circ \tilde{  \x }$为例来说明它是$C^{\infty}$的。
	为此跟踪数据在映射$\tilde{  \y  }^{-1}\circ \tilde{  \x }$之下的变化情况。
	\begin{center}
		\begin{tikzpicture}[scale=1, node distance=1.5cm, auto]
		\node [property, text width=15em] (MId1) {\textbf{Input}：$(\x^{-1}(p),a_{1},\cdots,a_{k})\in U\times\mr^{k}$};
		\node [notation, below of=MId1, text width=20em] (ONeg1) {$\varPhi_{\alpha}(p,a_{1},\cdots,a_{k})=X_{p}=\varPhi_{\beta}(p,b_{1},\cdots,b_{k})\in E_{p} $};
		\node [operation, below of=ONeg1, text width=15em] (DOS1) {\textbf{Output}：$(\y^{-1}(p),b_{1},\cdots,b_{k})\in V\times \mr^{k}$};	
		\path [line] (MId1) -- (ONeg1);
		\path [line] (ONeg1) -- (DOS1);
		\end{tikzpicture}
	\end{center}
	假设$\tilde{  U  }$的局部坐标系为$(x_{1},\cdots,x_{m},a_{1},\cdots,a_{k})$，$\tilde{  V  }$的局部坐标系为$(y_{1},\cdots,y_{m},b_{1},\cdots,b_{k})$. 结合上述流程图所展示的数据变化情况以及定理本身提供的条件，我们可以把局部坐标系之间的变换关系具体写出来：首先，显然有
	\begin{equation}
		(y_{1},\cdots,y_{m})=\y^{-1}\circ\x(x_{1},\cdots,x_{m})
	\end{equation}
	其次，若设$p=\x(x_{1},\cdots,x_{m})$. 则有$\varPhi_{\alpha}(p,a_{1},\cdots,a_{k})=\varPhi_{\beta}(p,b_{1},\cdots,b_{k})$. 也即
	\[
	\varPhi_{\beta}^{-1}\circ\varPhi_{\alpha}(p,a_{1},\cdots,a_{k})=(p,b_{1},\cdots,b_{k})
	\]
	注意到$\tilde{  \x  }(\tilde{  U  })\cap\tilde{  \y  }(\tilde{  V  })\neq\varnothing\Leftrightarrow\x(U)\cap\y(V)\neq\varnothing$，而$\x(U)\subset U_{\alpha},\y(V)\subset U_{\beta}$，所以当$\tilde{  \x  }(\tilde{  U  })\cap\tilde{  \y  }(\tilde{  V  })\neq\varnothing$时，也有$U_{\alpha}\cap U_{\beta}\neq \varnothing$. 因此由所给条件可知，存在{\color{red}{光滑映射$t_{\beta\alpha}:U_{\beta}\cap U_{\alpha}\longrightarrow GL_{k}(\mr)$}}使得
    \begin{equation}
        \left( \begin{array}{ccc}{b_{1}} \\ {\vdots}  \\ {b_{k}} \end{array}\right)=t_{\beta\alpha}\left(\x(x_{1},\cdots,x_{m}) \right) \left( \begin{array}{ccc}{a_{1}} \\ {\vdots}  \\ {a_{k}} \end{array}\right)
    \end{equation}
    
  这样我们就把$(\tilde{  U  },\tilde{  \x  })$与$(\tilde{  V  },\tilde{  \y  })$在交集部分$\tilde{  \x  }(\tilde{  U  })\cap \tilde{  \y  }(\tilde{  V  })$的局部坐标变换完全写了出来，即式\requ{equ:}与式\requ{equ:}. 由此可知，$\tilde{  \y  }^{-1}\circ\tilde{  \x  }$是$C^{\infty}$映射。同理可证$\tilde{  \x  }^{-1}\circ \tilde{  \y }$也是$C^{\infty}$映射。
  最后，只要将$\tilde{  \Sigma }$极大化便可给出$E$上的微分结构。
	
	至此，我们给$E$赋予了拓扑结构和微分结构，使其成为一个$m+k$维光滑流形。容易验证$(E,M,\pi)$满足向量丛的\rdef{def:chapter2向量丛}，因此$(E,M,\pi)$成为$M$上的一个秩为$k$的光滑向量丛，并且以$\{\varPhi_{\alpha}:U_{\alpha}\times \mr^{k}\longrightarrow \pi^{-1}(U_{\alpha}) \}_{\alpha\in \Lambda}$作为它的局部平凡化。	
\end{proof}

事实上我们可以进一步证明，上述定理与向量丛的定义是\textbf{等价的}。定理本身已经给出了等价性的其中一个方向的证明，为了证明另一个方向，我们先给出以下这个重要的引理。

\begin{mylem}{}{chapter2向量丛的过渡函数的矩阵表示}
	设$(E,M,\pi)$为光滑流形$M$上的一个秩为$k$的光滑实向量丛. $\varPhi:U\times\mr^{k}\longrightarrow \pi^{-1}(U)$和$\varPsi:V\times \mr^{k}\longrightarrow \pi^{-1}(V)$是两个局部平凡化. 若$U\cap V\neq\varnothing$，则存在{\color{red}{光滑映射}}$t_{VU}:V\cap U\longrightarrow GL_{k}(\mr)$，使得对$\;\forall p\in V\cap U\;,\forall(a_{1},\cdots,a_{k})\in \mr^{k}$，$\varPsi^{-1}\circ\varPhi:(V\cap U)\times\mr^{k}\longrightarrow (V\cap U)\times\mr^{k}$具有以下形式
	\[
	\varPsi^{-1}\circ\varPhi(p,a_{1},\cdots,a_{k})=(p,b_{1},\cdots,b_{k})
	\]
	其中
	\[
	\left( \begin{array}{ccc}{b_{1}} \\ {\vdots}  \\ {b_{k}} \end{array}\right)=t_{VU}(p) \left( \begin{array}{ccc}{a_{1}} \\ {\vdots}  \\ {a_{k}} \end{array}\right)
	\]
\end{mylem}
\begin{proof}
	任意取定一点$p\in U\cap V$. 令$\varPhi_{p}\triangleq\varPhi(p,\;{\color{red}{\cdot}}\;),\varPsi_{p}\triangleq\varPsi(p,\;{\color{red}{\cdot}}\;)$. 由于$\varPhi(p,\;{\color{red}{\cdot}}\;): \mr^{k}\longrightarrow E_{p}$是向量空间之间的同构，由于$\varPsi(p,\;{\color{red}{\cdot}}\;): \mr^{k}\longrightarrow E_{p}$也是向量空间之间的同构，所以
	\[
	\varPsi^{-1}_{p}\circ\varPhi_{p}:\mr^{k}\longrightarrow \mr^{k}
	\]
	也是向量空间之间的同构。正因为$\varPsi^{-1}_{p}\circ\varPhi_{p}$是$\mr^{k}\longrightarrow \mr^{k}$的同构映射，所以必存在$k\times k$的可逆矩阵，不妨记为$t_{VU}(p)$，使得
	对$\forall(x_{1},\cdots,x_{k})\in \mr^{k}$，若$\varPsi^{-1}_{p}\circ\varPhi_{p}(x_{1},\cdots,x_{k})=(y_{1},\cdots,y_{k})$，则
	\[
	\left( \begin{array}{ccc}{y_{1}} \\ {\vdots}  \\ {y_{k}} \end{array}\right)=t_{VU}(p) \left( \begin{array}{ccc}{x_{1}} \\ {\vdots}  \\ {x_{k}}\end{array}\right)
	\]
    另一方面，注意到$\forall(a_{1},\cdots,a_{k})\in \mr^{k}$，我们有
	\[
	\varPsi^{-1}\circ\varPhi(p,a_{1},\cdots,a_{k})=(p,	\varPsi^{-1}_{p}\circ\varPhi_{p}(a_{1},\cdots,a_{k}))
	\]
	
	所以对任意$p\in U\cap V$，存在相应的可逆矩阵$t_{VU}(p)\in GL_{k}(\mr)$，使得对$\forall(a_{1},\cdots,a_{k})\in \mr^{k}$，若
	\[
	\varPsi^{-1}\circ\varPhi(p,a_{1},\cdots,a_{k})=(p,b_{1},\cdots,b_{k})
	\]
	则
	\[
	\left( \begin{array}{ccc}{b_{1}} \\ {\vdots}  \\ {b_{k}} \end{array}\right)=\varPsi^{-1}_{p}\circ\varPhi_{p}(a_{1},\cdots,a_{k})=t_{VU}(p) \left( \begin{array}{ccc}{a_{1}} \\ {\vdots}  \\ {a_{k}} \end{array}\right)
	\]
	于是映射$t_{VU}:U\cap V\longrightarrow GL_{k}(\mr)$的存在性得证！下面只要再说明它的光滑性。
	
	注意到映射$t_{VU}$吃进去一个$U\cap V$中的点$p$，吐出来一个$k\times k$的实矩阵，所以我们可以把它想象成一个$k\times k$的矩阵函数：
	\[
	t_{VU}=\left( \begin{array}{cccc}{f_{11}} & {f_{12}} & {\cdots} & {f_{1 k}} \\ {f_{21}} & {f_{22}} & {\cdots} & {f_{2 k}} \\ {\vdots} & {\vdots} & { } & {\vdots} \\ {f_{k 1}} & {f_{k 2}} & {\cdots} & {f_{kk}}\end{array}\right)
	\]
	其中诸$f_{ij}:U\cap V\longrightarrow \mr$. 于是，为了说明$t_{VU}$是光滑映射，只要说明每一个分量函数$f_{ij}$是光滑的即可。
	
	对任意$1\leqslant j\leqslant k$，记$e_{j}\in \mr^{k}$是第$j$个分量为1，其余分量都为0的单位向量。则通过计算可知
	\[
	\varPsi^{-1}\circ\varPhi(p,e_{j})=(p,f_{1j}(p),f_{2j}(p),\cdots,f_{kj}(p))\;,\;\forall p\in U\cap V
	\]
	于是可令
	\[
	h_{j}\triangleq\varPsi^{-1}\circ\varPhi(\;\cdot\;,e_{j})=(\;\cdot\;,f_{1j},f_{2j},\cdots,f_{kj})
	\]
	由于$\varPhi:U\times\mr^{k}\longrightarrow \pi^{-1}(U)$和$\varPsi:V\times \mr^{k}\longrightarrow \pi^{-1}(V)$是局部平凡化，所以它们都是光滑的微分同胚。因此$h_{j}$也是光滑映射。
	
	对$1\leqslant i\leq k$，令$P_{i}:M\times \mr^{k}\longrightarrow\mr$为投影映射，即$
	P_{i}(p,x_{1},\cdots,x_{k})=x_{i}$. 显然$P_{i}$是光滑函数。将投影映射作用到$h_{j}$，我们便得到
	\[
	P_{i}\circ h_{j}= f_{ij}
	\]
	由于$P_{i},h_{j}$都是光滑的，所以$f_{ij}$也是光滑的。由$i,j$的任意性可知，每一个分量函数$f_{ij}$都是光滑的，因此$t_{VU}:U\cap V\longrightarrow GL_{k}(\mr)$是光滑映射。
\end{proof}

\begin{mythm}{}{chapter2向量丛的性质}
	若$(E,M,\pi)$是$M$上的一个秩为$k$的光滑向量丛，则存在$M$的开覆盖$\{U_{\alpha} \}_{\alpha\in \Lambda}$，并且对每一个$U_{\alpha}$，存在相应的{\color{red}{双射}}$\varPhi_{\alpha}:U_{\alpha}\times \mr^{k}\longrightarrow \pi^{-1}(U_{\alpha})$，满足以下条件：
	\begin{itemize}
		\item[(1)] $\forall q\in U_{\alpha},\forall a\in \mr^{k}$，均有$\varPhi_{\alpha}(q,a)\in E_{q}$；
		\item[(2)]  任意冻结一点$q\in U_{\alpha}$，映射$\varPhi_{\alpha}(q,\;{\color{red}{\cdot}}\;): \mr^{k}\longrightarrow E_{q}$是向量空间之间的同构；
		\item[(3)] 当$U_{\alpha}\cap U_{\beta}\neq\varnothing$时，存在光滑映射$t_{\beta\alpha}:U_{\beta}\cap U_{\alpha}\longrightarrow GL_{k}(\mr)$，使得对$\;\forall p\in U_{\beta}\cap U_{\alpha}\;,\forall(a_{1},\cdots,a_{k})\in \mr^{k}$，$\varPhi_{\beta}^{-1}\circ\varPhi_{\alpha}:(U_{\beta}\cap U_{\alpha})\times\mr^{k}\longrightarrow (U_{\beta}\cap U_{\alpha})\times\mr^{k}$具有以下形式
		\[
		\varPhi_{\beta}^{-1}\circ\varPhi_{\alpha}(p,a_{1},\cdots,a_{m})=(p,b_{1},\cdots,b_{m})
		\]
		其中
		\[
		\left( \begin{array}{ccc}{b_{1}} \\ {\vdots}  \\ {b_{k}} \end{array}\right)=t_{\beta\alpha}(p) \left( \begin{array}{ccc}{a_{1}} \\ {\vdots}  \\ {a_{k}} \end{array}\right)
		\]
	\end{itemize}
\end{mythm}
\begin{proof}
	根据向量丛的\rdef{def:chapter2向量丛}和\rlem{lem:chapter2向量丛的过渡函数的矩阵表示}立即可知定理成立！
\end{proof}

因此，我们也可以把\rthm{thm:chapter2向量丛的等价定义}看作是向量丛的等价定义。这种定义就与前面张量丛的构造产生了呼应，在具体构作向量丛的时候也往往更加方便。

上述定理和引理中出现的光滑映射$t_{\alpha\beta}:U_{\alpha}\cap U_{\beta}\longrightarrow GL_{k}(\mr)$与$t_{UV}:U\cap V\longrightarrow GL_{k}(\mr)$，我们称之为向量丛的\textbf{过渡函数}。注意，虽然$t_{UV}$叫“函数”，但它的值域并非$\mr$，叫“过渡函数”只是习惯而已。全体过渡函数的集合$\{t_{UV}:U\cap V\longrightarrow GL_{k}(\mr) \}$称为向量丛的\textbf{过渡函数族}。
\begin{example}
	\label{张量丛的过渡函数族}
	张量丛$T^{r}_{s}$的过渡函数族.
	\soln
	
	若$\x_{\alpha}(U_{\alpha})=W_{\alpha}\cap W_{\beta}=\x_{\beta}(U_{\beta})\neq\varnothing$，假设
	\[
	\varPhi_{\beta}^{-1}\circ\varPhi_{\alpha}(p,a_{\eta},\cdots,a_{\gamma},\cdots,a_{xi})=(p,b_{\eta},\cdots,b_{\gamma},\cdots,b_{xi})
	\]
	也就是
	\[
	\begin{split}
	&\varPhi_{\alpha}(p,a_{\eta},\cdots,a_{\gamma},\cdots,a_{xi})=\varPhi_{\beta}(p,b_{\eta},\cdots,b_{\gamma},\cdots,b_{xi})\\
	&\Leftrightarrow \theta^{\alpha}_{p}(a_{\eta},\cdots,a_{\gamma},\cdots,a_{xi})=\theta^{\beta}_{p}(b_{\eta},\cdots,b_{\gamma},\cdots,b_{xi})
	\end{split}
	\]
	则根据前面的\rlem{lem:chapter2张量丛过渡函数的矩阵形式}可知
	\[
	\left( \begin{array}{ccccc}{\vdots} \\ {b_{\gamma}}  \\ {\vdots} \\ {\vdots}\\ {\vdots}\end{array}\right)=J(\x^{-1}_{\beta}\circ\x_{\alpha})\otimes\cdots\otimes J(\x^{-1}_{\beta}\circ\x_{\alpha})\otimes [J(\x^{-1}_{\alpha}\circ\x_{\beta})]^{T}\otimes \cdots\otimes [J(\x^{-1}_{\alpha}\circ\x_{\beta})]^{T} \left( \begin{array}{ccccc}{\vdots} \\ {\vdots}  \\ {\vdots} \\ {a_{\lambda}}\\ {\vdots}\end{array}\right)
	\]
	这里$J(\y^{-1}\circ\x)$表示映射$\y^{-1}\circ\x$点$\x^{-1}(p)$处的Jacobi矩阵，$J(\x^{-1}\circ\y) $表示映射$x^{-1}\circ\y$在点$\y^{-1}(p)$为Jacobi矩阵. 所以令
	\[
	t_{\beta\alpha}\triangleq \underbrace{J(\x^{-1}_{\beta}\circ\x_{\alpha})\otimes\cdots\otimes J(\x^{-1}_{\beta}\circ\x_{\alpha})}_{r\text{个}}\otimes \underbrace{[J(\x^{-1}_{\alpha}\circ\x_{\beta})]^{T}\otimes \cdots\otimes [J(\x^{-1}_{\alpha}\circ\x_{\beta})]^{T}}_{s\text{个}}
	\]
	则$\{t_{\beta\alpha} \}$就构成了张量丛的过渡函数族.
\end{example}

过渡函数有如下基本性质。
\begin{mythm}{}{chapter2过渡函数的性质}
	\begin{itemize}
		\item[(1)] $t_{UU}=I$，其中$I$为$k$阶单位矩阵. 
		\item[(2)] $t_{UV}\cdot t_{VW}=t_{UW}$. 
		\item[(3)] $t_{UV}=\left(t_{VU} \right)^{-1}$.
	\end{itemize}
\end{mythm}
\begin{proof}
	(1)设$\varPhi:U\times\mr^{k}\longrightarrow \pi^{-1}(U)$为局部平凡化，显然对任意$p\in U$，任意$(a_{1},\cdots,a_{k})\in \mr^{k}$，均有
	\[
	\varPhi^{-1}\circ\varPhi(p,a_{1},\cdots,a_{k})=(p,a_{1},\cdots,a_{k})
	\]
	即$t_{UU}(p)=I,\forall p\in U\Rightarrow t_{UU}=I$. 
	
	(2)设$\varPhi:U\times\mr^{k}\longrightarrow \pi^{-1}(U),\varPsi:V\times\mr^{k}\longrightarrow \pi^{-1}(V),\varUpsilon:W\times\mr^{k}\longrightarrow \pi^{-1}(W)$为局部平凡化，对任意$p\in U\cap V\cap W$，任意$(x_{1},\cdots,x_{k})\in \mr^{k}$，设
	\[
	\varPsi^{-1}\circ\varUpsilon(p,x_{1},\cdots.x_{k})=(p,a_{1},\cdots,a_{k})\quad\varPhi^{-1}\circ\varPsi(p,a_{1},\cdots,a_{k})=(p,y_{1},\cdots,y_{k})
	\]
	则$\varPhi^{-1}\circ\varUpsilon(p,x_{1},\cdots,x_{k})=(p,y_{1},\cdots,y_{k})$. 注意到
	\[
	t_{UW}(p) \left( \begin{array}{ccc}{x_{1}} \\ {\vdots}  \\ {x_{k}} \end{array}\right)=\left( \begin{array}{ccc}{y_{1}} \\ {\vdots}  \\ {b_{k}} \end{array}\right)=t_{UV}(p) \left( \begin{array}{ccc}{a_{1}} \\ {\vdots}  \\ {a_{k}} \end{array}\right)=t_{UV}(p)\cdot t_{VW}(p) \left( \begin{array}{ccc}{x_{1}} \\ {\vdots}  \\ {x_{k}} \end{array}\right)
	\]
	所以$t_{UW}(p)=t_{UV}(p)\cdot t_{VW}(p),\forall p\in U\cap V\cap W\Rightarrow t_{UV}\cdot t_{VW}=t_{UW}$.
	
	(3)根据性质(1)和性质(2)$\Rightarrow t_{UV}\cdot t_{VU}=t_{UU}=I$，所以$t_{UV}=\left(t_{VU} \right)^{-1}$. 
\end{proof}

实际上，满足上述定理的过渡函数族可唯一确定一个向量丛，因此，过渡函数才是向量丛的本质所在。
\begin{mythm}{光滑向量丛的结构定理}{chapter2向量丛的结构定理}
	设$M^{m}$是一个光滑流形，$\{U_{\alpha} \}_{\alpha\in \Lambda}$是$M$的一个开覆盖. 若对每一对指标$\alpha,\beta$，当$U_{\alpha}\cap U_{\beta}\neq\varnothing$时，都给定了一个光滑映射$t_{\alpha\beta}:U_{\alpha}\cap U_{\beta}\longrightarrow GL_{k}(\mr)$，并且这些光滑映射满足：
		\begin{itemize}
		\item[(1)] $t_{\alpha\alpha}=I$. 
		\item[(2)] $t_{\alpha\beta}\cdot t_{\beta\gamma}=t_{\alpha\gamma}$. 
	\end{itemize}
则$M$上存在秩为$k$的光滑向量丛$(E,M,\pi)$，它以$\{t_{\alpha\beta}:U_{\alpha}\cap U_{\beta}\longrightarrow GL_{k}(\mr) \}$为过渡函数族.
\end{mythm}
\begin{proof}
	我们只把关键的构造写出来，细节就不验证了。首先令
	\[
	S\triangleq \bigsqcup_{\alpha\in \Lambda} U_{\alpha}\times \mr^{k}=\bigcup_{\alpha\in \Lambda}\{ \alpha\}\times U_{\alpha}\times \mr^{k}
	\]
	其次在集合$S$上定义等价关系如下
	\[
	(\alpha,p,x_{1},\cdots,x_{k})\sim (\beta,q,y_{1},\cdots,y_{k})\Leftrightarrow p=q\text{且}
	\left( \begin{array}{ccc}{y_{1}} \\ {\vdots}  \\ {b_{k}} \end{array}\right)=t_{\beta\alpha}(p) \left( \begin{array}{ccc}{x_{1}} \\ {\vdots}  \\ {x_{k}} \end{array}\right)
	\]
	最后令$E=S/\sim$，丛投影$\pi$定义为
	\[\pi:E\longrightarrow M\Leftrightarrow\pi([\alpha,p,x_{1},\cdots,x_{k}])=p
	\] 
	局部平凡化$\varPhi_{\alpha}:U_{\alpha}\times \mr^{k}\longrightarrow \pi^{-1}(U_{\alpha})$定义为
	\[
	\varPhi_{\alpha}(p,a_{1},\cdots,a_{k})=[\alpha,p,a_{1},\cdots,a_{k}]
	\]
	这样构造出来的$(E,M,\pi)$即为$M$上的一个秩为$k$的、以$\{t_{\alpha\beta}:U_{\alpha}\cap U_{\beta}\longrightarrow GL_{k}(\mr) \}$为过渡函数族的光滑向量丛。
\end{proof}







\subsection{对偶}

假设已经给定一个秩为$k$的向量丛$(E,M,\pi)$. 我们来构造一个新的向量丛$(E^{\ast},M,\pi^{\ast})$.

对任意$p\in M$，令$E^{\ast}_{p}\triangleq \left(E_{p} \right)^{\ast}$. 即给流形$M$的每一点赋予向量空间$E_{p}$的对偶空间。再令$E^{\ast}\triangleq \bigcup\limits_{p\in M}E^{\ast}_{p}$. $\pi^{\ast}:E^{\ast}\longrightarrow M$定义为$\pi^{\ast}(X_{p})=p,\forall X_{p}\in E^{\ast}_{p}$. 这些都是很自然的构造。

由于$(E,M,\pi)$为向量丛，根据\rthm{thm:chapter2向量丛的性质}可知，存在$M$的开覆盖$\{U_{\alpha} \}_{\alpha\in \Lambda}$以及双射族$\{\varPhi_{\alpha}:U_{\alpha}\times \mr^{k}\longrightarrow \pi^{-1}(U_{\alpha}) \}$，满足以下性质：
\begin{itemize}
	\item[(1)] $\varPhi_{\alpha}(q,a)\in E_{q},\forall q\in U_{\alpha},\forall a\in \mr^{k}$.
	\item[(2)] $\varPhi_{\alpha}(q,\;{\color{red}{\cdot}}\;):\mr^{k}\longrightarrow E_{q}$为线性同构.
	\item[(3)] 当$U_{\alpha}\cap U_{\beta}\neq\varnothing$时，存在光滑映射$t_{\beta\alpha}:U_{\alpha}\cap U_{\beta}\longrightarrow GL_{k}(\mr)$，使得对$\forall p\in U_{\alpha}\cap U_{\beta},\forall (a_{1},\cdots,a_{k})\in \mr^{k}$，若$\varPhi_{\beta}^{-1}\circ\varPhi_{\alpha}(p,a_{1},\cdots,a_{k})=(p,b_{1},\cdots,b_{k})$，则
	\[
	\left( \begin{array}{ccc}{b_{1}} \\ {\vdots}  \\ {b_{k}} \end{array}\right)=t_{\beta\alpha}(p) \left( \begin{array}{ccc}{a_{1}} \\ {\vdots}  \\ {a_{k}} \end{array}\right)
	\]
\end{itemize}
	
对于这个开覆盖$\{U_{\alpha} \}_{\alpha\in \Lambda}$，我们定义一族新的映射$\{\varPhi^{\ast}_{\alpha} :U_{\alpha}\times \mr^{k}\longrightarrow \pi^{\ast -1}(U_{\alpha})\}$如下：$\forall p\in U_{\alpha},\forall (x_{1},\cdots,x_{k})\in\mr^{k},\varPhi^{\ast}_{\alpha}(p,x_{1},\cdots,x_{k})\in E^{\ast}_{p}$. 其中
\[
\varPhi^{\ast}_{\alpha}(p,x_{1},\cdots,x_{k})\left(X_{p} \right)=\varPhi^{\ast}_{\alpha}(p,x_{1},\cdots,x_{k})\left(\varPhi_{\alpha}(p,a_{1},\cdots,a_{k}) \right)=\sum_{ i = 1 }^{k}x_{i}a_{i}
\]
这里$X_{p}=\varPhi_{\alpha}(p,a_{1},\cdots,a_{k})\in E_{p}$. 显然$\varPhi^{\ast}_{\alpha} :U_{\alpha}\times \mr^{k}\longrightarrow \pi^{\ast -1}(U_{\alpha})$是双射。

当$U_{\alpha}\cap U_{\beta}\neq\varnothing$时，设$\varPhi^{\ast -1}_{\beta}\circ\varPhi^{\ast}_{\alpha}(p,x_{1},\cdots,x_{k})=(p,y_{1},\cdots,y_{k})$. 则
\[
\begin{split}
&\varPhi^{\ast -1}_{\beta}\circ\varPhi^{\ast}_{\alpha}(p,x_{1},\cdots,x_{k})=(p,y_{1},\cdots,y_{k})\\
&\Leftrightarrow \varPhi^{\ast}_{\alpha}(p,x_{1},\cdots,x_{k})=\varPhi^{\ast}_{\beta}(p,y_{1},\cdots,y_{k})\\
&\Leftrightarrow \varPhi^{\ast}_{\alpha}(p,x_{1},\cdots,x_{k})(X_{p})=\varPhi^{\ast}_{\beta}(p,y_{1},\cdots,y_{k})(X_{p}),\forall X_{p}\in E_{p}\\
\end{split}
\]
设$X_{p}=\varPhi_{\alpha}(p,a_{1},\cdots,a_{k})=\varPhi_{\beta}(p,b_{1},\cdots,b_{k})$. 
\[
\begin{split}
&\varPhi^{\ast}_{\alpha}(p,x_{1},\cdots,x_{k})(X_{p})=\varPhi^{\ast}_{\beta}(p,y_{1},\cdots,y_{k})(X_{p}),\forall X_{p}\in E_{p}\\
&\Leftrightarrow(x_{1},\cdots,x_{k})\cdot \left( \begin{array}{ccc}{a_{1}} \\ {\vdots}  \\ {a_{k}} \end{array}\right)=(y_{1},\cdots,y_{k})\cdot \left( \begin{array}{ccc}{b_{1}} \\ {\vdots}  \\ {b_{k}} \end{array}\right)\\
&\Leftrightarrow (x_{1},\cdots,x_{k})\cdot \left( \begin{array}{ccc}{a_{1}} \\ {\vdots}  \\ {a_{k}} \end{array}\right)=(y_{1},\cdots,y_{k})\cdot t_{\beta\alpha}(p) \left( \begin{array}{ccc}{a_{1}} \\ {\vdots}  \\ {a_{k}} \end{array}\right)
\end{split}
\]
从而
\[
(x_{1},\cdots,x_{k})=(y_{1},\cdots,y_{k})\cdot t_{\beta\alpha}(p)
\]
改写一下，也就是
\[
\left( \begin{array}{ccc}{y_{1}} \\ {\vdots}  \\ {y_{k}} \end{array}\right)=\left(t^{-1}_{\beta\alpha}(p)\right)^{T} \left( \begin{array}{ccc}{x_{1}} \\ {\vdots}  \\ {x_{k}} \end{array}\right)
\]
 
 因此，当$U_{\alpha}\cap U_{\beta}\neq\varnothing$时，我们定义$t^{\ast}_{\beta\alpha}:U_{\alpha}\cap U_{\beta}\longrightarrow GL_{k}(\mr)$为
 \begin{equation}
 	t^{\ast}_{\beta\alpha}\triangleq \left(t^{-1}_{\beta\alpha}\right)^{T}
 \end{equation}
 
 由于$t_{\beta\alpha}:U_{\alpha}\cap U_{\beta}\longrightarrow GL_{k}(\mr)$是光滑映射，所以如此定义的$t^{\ast}_{\beta\alpha}:U_{\alpha}\cap U_{\beta}\longrightarrow GL_{k}(\mr)$也是光滑映射。
 
 可以验证以上给出的$\{\varPhi^{\ast}_{\alpha} :U_{\alpha}\times \mr^{k}\longrightarrow \pi^{\ast -1}(U_{\alpha})\}$以及$\{ t^{\ast}_{\beta\alpha} \}$满足\rthm{thm:chapter2向量丛的等价定义}的所有要求，从而使得$(E^{\ast},M,\pi^{\ast})$成为$M$的一个新的秩为$k$的向量丛，式\requ{equ:}给出了它的过渡函数族。我们称向量丛$(E^{\ast},M,\pi^{\ast})$为向量丛$(E,M,\pi)$的\textbf{对偶丛}。

\begin{remark}
	对偶丛的纤维型也可以看成是${\color{red}{(\mr^{k})^{\ast}}}$. 即我们也可以定义映射$\varPhi^{\ast}_{\alpha}$为
	$\varPhi^{\ast}_{\alpha} :U_{\alpha}\times {\color{red}{(\mr^{k})^{\ast}}}\longrightarrow \pi^{\ast -1}(U_{\alpha})$，对任意$p\in U_{\alpha}$，任意$f\in {\color{red}{(\mr^{k})^{\ast}}}$，$\varPhi^{\ast}_{\alpha}(p,f)\in (E_{p})^{\ast}$. 其中
	\[
	\varPhi^{\ast}_{\alpha}(p,f)(X_{p})\xlongequal{X_{p}=\varPhi_{\alpha}(p,a_{1},\cdots,a_{k})\in E_{p}}f(a_{1},\cdots,a_{k})
	\]
	实际上，根据泛函分析中的$Riesz$表示定理，任意$f\in {\color{red}{(\mr^{k})^{\ast}}}$，存在唯一的$(x_{1},\cdots,x_{k})\in\mr^{k}$，使得
	\[
	f(a_{1},\cdots,a_{k})=\sum_{ i = 1 }^{k}x_{i}a_{i}
	\]
	因此这两种定义映射$\varPhi^{\ast}_{\alpha}$的方式完全等价！用我们介绍的那种定义更直观一些。
\end{remark}

\subsection{直和}

假设已经给定两个秩分别为$k_{1},k_{2}$的向量丛$(\overline{E},M,\overline{\pi}),(\tilde{E},M,\tilde{\pi})$. 我们来构造一个新的向量丛$(E,M,\pi)$.

对任意$p\in M$，令$E_{p}\triangleq \overline{E}_{p}\oplus \tilde{E}_{p}$. 即给流形$M$的每一点赋予向量空间$\overline{E}_{p}$与$\tilde{E}_{p}$的\textbf{外直和}。再令$E\triangleq \bigcup\limits_{p\in M}E_{p}$. $\pi:E\longrightarrow M$定义为$\pi(X_{p})=p,\forall X_{p}\in E_{p}$. 这些都是很自然的构造。

由于$(\overline{E},M,\overline{\pi})$为向量丛，根据\rthm{thm:chapter2向量丛的性质}可知，存在$M$的开覆盖$\{\overline{U}_{\alpha} \}_{\alpha\in \Lambda}$以及双射族$\{\overline{\varPhi}_{\alpha}:\overline{U}_{\alpha}\times \mr^{k}\longrightarrow \overline{\pi}^{-1}(\overline{U}_{\alpha}) \}$，满足以下性质：
\begin{itemize}
	\item[(1)] $\overline{\varPhi}_{\alpha}(q,a)\in \overline{E}_{q},\forall q\in \overline{U}_{\alpha},\forall a\in \mr^{k_{1}}$.
	\item[(2)] $\overline{\varPhi}_{\alpha}(q,\;{\color{red}{\cdot}}\;):\mr^{k_{1}}\longrightarrow \overline{E}_{q}$为线性同构.
	\item[(3)] 当$\overline{U}_{\alpha}\cap \overline{U}_{\beta}\neq\varnothing$时，存在光滑映射$\overline{t}_{\beta\alpha}:\overline{U}_{\alpha}\cap \overline{U}_{\beta}\longrightarrow GL_{k_{1}}(\mr)$，使得对$\forall p\in \overline{U}_{\alpha}\cap \overline{U}_{\beta},\forall (a_{1},\cdots,a_{k_{1}})\in \mr^{k_{1}}$，若$\overline{\varPhi}_{\beta}^{-1}\circ\overline{\varPhi}_{\alpha}(a_{1},\cdots,a_{k_{1}})=(p,b_{1},\cdots,b_{k_{1}})$，则
	\[
	\left( \begin{array}{ccc}{b_{1}} \\ {\vdots}  \\ {b_{k_{1}}} \end{array}\right)=\overline{t}_{\beta\alpha}(p) \left( \begin{array}{ccc}{a_{1}} \\ {\vdots}  \\ {a_{k_{1}}} \end{array}\right)
	\]
\end{itemize}
同样的，由于$(\tilde{E},M,\tilde{\pi})$为向量丛，根据\rthm{thm:chapter2向量丛的性质}可知，存在$M$的开覆盖$\{\tilde{U}_{\alpha} \}_{\alpha\in \Lambda'}$以及双射族$\{\tilde{\varPhi}_{\alpha}:\tilde{U}_{\alpha}\times \mr^{k}\longrightarrow \tilde{\pi}^{-1}(\tilde{U}_{\alpha}) \}$，满足以下性质：
\begin{itemize}
	\item[(1)] $\tilde{\varPhi}_{\alpha}(q,a)\in \tilde{E}_{q},\forall q\in \tilde{U}_{\alpha},\forall a\in \mr^{k_{2}}$.
	\item[(2)] $\tilde{\varPhi}_{\alpha}(q,\;{\color{red}{\cdot}}\;):\mr^{k_{2}}\longrightarrow \tilde{E}_{q}$为线性同构.
	\item[(3)] 当$\tilde{U}_{\alpha}\cap \tilde{U}_{\beta}\neq\varnothing$时，存在光滑映射$\tilde{t}_{\beta\alpha}:\tilde{U}_{\alpha}\cap \tilde{U}_{\beta}\longrightarrow GL_{k_{2}}(\mr)$，使得对$\forall p\in \tilde{U}_{\alpha}\cap \tilde{U}_{\beta},\forall (a_{1},\cdots,a_{k_{2}})\in \mr^{k_{2}}$，若$\tilde{\varPhi}_{\beta}^{-1}\circ\tilde{\varPhi}_{\alpha}(a_{1},\cdots,a_{k_{2}})=(p,b_{1},\cdots,b_{k_{2}})$，则
	\[
	\left( \begin{array}{ccc}{b_{1}} \\ {\vdots}  \\ {b_{k_{2}}} \end{array}\right)=\tilde{t}_{\beta\alpha}(p) \left( \begin{array}{ccc}{a_{1}} \\ {\vdots}  \\ {a_{k_{2}}} \end{array}\right)
	\]
\end{itemize}

对$\forall p\in M$，存在包含点$p$的开集$\overline{U}_{\alpha}$和$\tilde{U}_{\alpha}$. 令$U_{\alpha}=\overline{U}_{\alpha}\cap \tilde{U}_{\alpha}$，并将映射$\overline{\varPhi}_{\alpha}$和$\tilde{\varPhi}_{\alpha}$限制在$U_{\alpha}$上来考虑。定义$\varPhi_{\alpha}:U_{\alpha}\times \mr^{k_{1}+k_{2}} \longrightarrow \pi^{-1}(U_{\alpha})$为
\[
\varPhi_{\alpha}(p,a_{1},\cdots,a_{k_{1}},a_{k_{1}+1},\cdots,a_{k_{1}+k_{2}})\triangleq  \overline{\varPhi}_{\alpha}(p,a_{1},\cdots,a_{k_{1}})+\tilde{\varPhi}_{\alpha}(p,a_{k_{1}+1},\cdots,a_{k_{1}+k_{2}})
\]
注意这里$\overline{\varPhi}_{\alpha}(p,a_{1},\cdots,a_{k_{1}})+\tilde{\varPhi}_{\alpha}(p,a_{k_{1}+1},\cdots,a_{k_{1}+k_{2}})$只是形式和. 显然$\varPhi_{\alpha}:U_{\alpha}\times \mr^{k_{1}+k_{2}} \longrightarrow \pi^{-1}(U_{\alpha})$为双射。

当$U_{\alpha}\cap U_{\beta}\neq\varnothing$时，则也有$\overline{U}_{\alpha}\cap\overline{U}_{\beta}\neq\varnothing,\tilde{U}_{\alpha}\cap\tilde{U}_{\beta}\neq\varnothing$. 则
\[
\begin{split}
&\varPhi^{ -1}_{\beta}\circ\varPhi_{\alpha}(p,a_{1},\cdots,a_{k_{1}},a_{k_{1}+1},\cdots,a_{k_{1}+k_{2}})=(p,b_{1},\cdots,b_{k_{1}},b_{k_{1}+1},\cdots,b_{k_{1}+k_{2}})\\
\Leftrightarrow&\varPhi_{\alpha}(p,a_{1},\cdots,a_{k_{1}},a_{k_{1}+1},\cdots,a_{k_{1}+k_{2}})=\varPhi_{\beta}(p,b_{1},\cdots,b_{k_{1}},b_{k_{1}+1},\cdots,b_{k_{1}+k_{2}})\\
\Leftrightarrow&\overline{\varPhi}_{\alpha}(p,a_{1},\cdots,a_{k_{1}})+\tilde{\varPhi}_{\alpha}(p,a_{k_{1}+1},\cdots,a_{k_{1}+k_{2}})=\overline{\varPhi}_{\beta}(p,b_{1},\cdots,b_{k_{1}})+\tilde{\varPhi}_{\beta}(p,b_{k_{1}+1},\cdots,b_{k_{1}+k_{2}})\\
\Leftrightarrow&\overline{\varPhi}_{\alpha}(p,a_{1},\cdots,a_{k_{1}})=\overline{\varPhi}_{\beta}(p,b_{1},\cdots,b_{k_{1}})\text{且}\tilde{\varPhi}_{\alpha}(p,a_{k_{1}+1},\cdots,a_{k_{1}+k_{2}})=\tilde{\varPhi}_{\beta}(p,b_{k_{1}+1},\cdots,b_{k_{1}+k_{2}})
%\Leftrightarrow&\overline{\varPhi}^{-1}_{\beta}\circ\overline{\varPhi}_{\alpha}(p,a_{1},\cdots,a_{k_{1}})=(p,b_{1},\cdots,b_{k_{1}})\text{且}\tilde{\varPhi}^{-1}_{\beta}\circ\tilde{\varPhi}_{\alpha}(p,a_{k_{1}+1},\cdots,a_{k_{1}+k_{2}})=(p,b_{k_{1}+1},\cdots,b_{k_{1}+k_{2}})
\end{split}
\]
因为
\[
\left( \begin{array}{ccc}{b_{1}} \\ {\vdots}  \\ {b_{k_{1}}} \end{array}\right)=\overline{t}_{\beta\alpha}(p) \left( \begin{array}{ccc}{a_{1}} \\ {\vdots}  \\ {a_{k_{1}}} \end{array}\right)\quad,\quad  \left( \begin{array}{ccc}{b_{k_{1}+1}} \\ {\vdots}  \\ {b_{k_{1}+k_{2}}} \end{array}\right)=\tilde{t}_{\beta\alpha}(p) \left( \begin{array}{ccc}{a_{k_{1}+1}} \\ {\vdots}  \\ {a_{k_{1}+k_{2}}} \end{array}\right)
\]
所以
\[
\left( \begin{array}{cccccc}{b_{1}} \\ {\vdots}  \\ {b_{k_{1}}}\\{b_{k_{1}+1}} \\ {\vdots}  \\ {b_{k_{1}+k_{2}}}  \end{array}\right)=\left( \begin{array}{cc}{\overline{t}_{\beta\alpha}(p)}&{0} \\ {0}&{\tilde{t}_{\beta\alpha}(p)}  \end{array}\right) \left( \begin{array}{cccccc}{a_{1}} \\ {\vdots}  \\ {a_{k_{1}}}\\{a_{k_{1}+1}} \\ {\vdots}  \\ {a_{k_{1}+k_{2}}}  \end{array}\right)
\]

因此，当$U_{\alpha}\cap U_{\beta}\neq\varnothing$时，我们定义$t_{\beta\alpha}:U_{\alpha}\cap U_{\beta}\longrightarrow GL_{k_{1}+k_{2}}(\mr)$为
\begin{equation}
	t_{\beta\alpha}\triangleq\left( \begin{array}{cc}{\overline{t}_{\beta\alpha}}&{0} \\ {0}&{\tilde{t}_{\beta\alpha}}  \end{array}\right)
\end{equation}
由于$\overline{t}_{\beta\alpha}:\overline{U}_{\alpha}\cap \overline{U}_{\beta}\longrightarrow GL_{k_{1}}(\mr),\tilde{t}_{\beta\alpha}:\tilde{U}_{\alpha}\cap \tilde{U}_{\beta}\longrightarrow GL_{k_{2}}(\mr)$都是光滑映射，所以如此定义的$t_{\beta\alpha}:U_{\alpha}\cap U_{\beta}\longrightarrow GL_{k_{1}+k_{2}}(\mr)$也是光滑映射。

可以验证以上给出的$\{\varPhi_{\alpha}:U_{\alpha}\times \mr^{k_{1}+k_{2}} \longrightarrow \pi^{-1}(U_{\alpha})\}$以及$\{ t_{\beta\alpha} \}$满足\rthm{thm:chapter2向量丛的等价定义}的所有要求，从而使得$(E,M,\pi)$成为$M$的一个新的秩为$k_{1}+k_{2}$的向量丛，式\requ{equ:}给出了它的过渡函数族。我们称这样得到的向量丛$(E,M,\pi)$为向量丛$(\overline{E},M,\overline{\pi})$与$(\tilde{E},M,\tilde{\pi})$的\textbf{直和}或者\textbf{Whitney和}，并记作$E\triangleq \overline{E}\oplus \tilde{E}$.

\subsection{直和的一个应用:外形式丛}

在\ref{s次外形式丛}节，我们介绍了$s$次外形式丛。在那里是给流形的每一点$p$粘上$\Lambda^{s}(T^{\ast}_{p}M)$作为该点处的向量空间。而我们知道\textbf{外代数}
\[
	\Lambda(T^{\ast}_{p}M)=\bigoplus_{s=0}^{m}\Lambda^{s}(T^{\ast}_{p}M),\;\text{\kaishu 其中}m=\operatorname{dim}(M)
\]
本身也构成一个向量空间，所以自然可以考虑给流形的每一点$p$赋予$\Lambda(T^{\ast}_{p}M)$作为该点处的向量空间。这一想法恰好契合了丛的直和这种由已知丛构造新丛的方法，于是我们令
\[
	\Lambda(T^{\ast}M)\triangleq \Lambda^{0}(T^{\ast}M)\oplus \Lambda^{1}(T^{\ast}M)\oplus \cdots\oplus\Lambda^{m}(T^{\ast}M)
\]

由$s$次外形式丛($s=0,1,\cdots,m$)像这样做直和得到的新的丛$\Lambda(T^{\ast}M)$，我们称之为\textbf{外形式丛}，注意名称前面不再加上“$s$次”的前缀，因为各次数的丛都已经包含在$\Lambda(T^{\ast}M)$里面了；或者你也可以叫它“外代数丛”。在外貌形态上，$\Lambda(T^{\ast}M)$就是流形$M$上各种阶数的反对称协变张量及其形式和的全体，即
\[
	\Lambda(T^{\ast}M)=\bigcup\limits_{p\in M}\Lambda(T^{\ast}_{p}M)
\]

这个丛可以看成是$s$次外形式丛的进一步推广，跟$s$次外形式丛一样，都是非常非常非常重要的丛！

\subsection{张量积}
假设已经给定两个秩分别为$k_{1},k_{2}$的向量丛$(\overline{E},M,\overline{\pi}),(\tilde{E},M,\tilde{\pi})$. 我们来构造一个新的向量丛$(E,M,\pi)$.

对任意$p\in M$，令$E_{p}\triangleq \overline{E}_{p}\otimes \tilde{E}_{p}$. 即给流形$M$的每一点赋予向量空间$\overline{E}_{p}$与$\tilde{E}_{p}$的\textbf{张量积空间}。再令$E\triangleq \bigcup\limits_{p\in M}E_{p}$. $\pi:E\longrightarrow M$定义为$\pi(X_{p})=p,\forall X_{p}\in E_{p}$. 这些都是很自然的构造。

由于$(\overline{E},M,\overline{\pi})$为向量丛，根据\rthm{thm:chapter2向量丛的性质}可知，存在$M$的开覆盖$\{\overline{U}_{\alpha} \}_{\alpha\in \Lambda}$以及双射族$\{\overline{\varPhi}_{\alpha}:\overline{U}_{\alpha}\times \mr^{k}\longrightarrow \overline{\pi}^{-1}(\overline{U}_{\alpha}) \}$，满足以下性质：
\begin{itemize}
	\item[(1)] $\overline{\varPhi}_{\alpha}(q,a)\in \overline{E}_{q},\forall q\in \overline{U}_{\alpha},\forall a\in \mr^{k_{1}}$.
	\item[(2)] $\overline{\varPhi}_{\alpha}(q,\;{\color{red}{\cdot}}\;):\mr^{k_{1}}\longrightarrow \overline{E}_{q}$为线性同构.
	\item[(3)] 当$\overline{U}_{\alpha}\cap \overline{U}_{\beta}\neq\varnothing$时，存在光滑映射$\overline{t}_{\beta\alpha}:\overline{U}_{\alpha}\cap \overline{U}_{\beta}\longrightarrow GL_{k_{1}}(\mr)$，使得对$\forall p\in \overline{U}_{\alpha}\cap \overline{U}_{\beta},\forall (a_{1},\cdots,a_{k_{1}})\in \mr^{k_{1}}$，若$\overline{\varPhi}_{\beta}^{-1}\circ\overline{\varPhi}_{\alpha}(a_{1},\cdots,a_{k_{1}})=(p,b_{1},\cdots,b_{k_{1}})$，则
	\[
	\left( \begin{array}{ccc}{b_{1}} \\ {\vdots}  \\ {b_{k_{1}}} \end{array}\right)=\overline{t}_{\beta\alpha}(p) \left( \begin{array}{ccc}{a_{1}} \\ {\vdots}  \\ {a_{k_{1}}} \end{array}\right)
	\]
\end{itemize}
同样的，由于$(\tilde{E},M,\tilde{\pi})$为向量丛，根据\rthm{thm:chapter2向量丛的性质}可知，存在$M$的开覆盖$\{\tilde{U}_{\alpha} \}_{\alpha\in \Lambda'}$以及双射族$\{\tilde{\varPhi}_{\alpha}:\tilde{U}_{\alpha}\times \mr^{k}\longrightarrow \tilde{\pi}^{-1}(\tilde{U}_{\alpha}) \}$，满足以下性质：
\begin{itemize}
	\item[(1)] $\tilde{\varPhi}_{\alpha}(q,a)\in \tilde{E}_{q},\forall q\in \tilde{U}_{\alpha},\forall a\in \mr^{k_{2}}$.
	\item[(2)] $\tilde{\varPhi}_{\alpha}(q,\;{\color{red}{\cdot}}\;):\mr^{k_{2}}\longrightarrow \tilde{E}_{q}$为线性同构.
	\item[(3)] 当$\tilde{U}_{\alpha}\cap \tilde{U}_{\beta}\neq\varnothing$时，存在光滑映射$\tilde{t}_{\beta\alpha}:\tilde{U}_{\alpha}\cap \tilde{U}_{\beta}\longrightarrow GL_{k_{2}}(\mr)$，使得对$\forall p\in \tilde{U}_{\alpha}\cap \tilde{U}_{\beta},\forall (a_{1},\cdots,a_{k_{2}})\in \mr^{k_{2}}$，若$\tilde{\varPhi}_{\beta}^{-1}\circ\tilde{\varPhi}_{\alpha}(a_{1},\cdots,a_{k_{2}})=(p,b_{1},\cdots,b_{k_{2}})$，则
	\[
	\left( \begin{array}{ccc}{b_{1}} \\ {\vdots}  \\ {b_{k_{2}}} \end{array}\right)=\tilde{t}_{\beta\alpha}(p) \left( \begin{array}{ccc}{a_{1}} \\ {\vdots}  \\ {a_{k_{2}}} \end{array}\right)
	\]
\end{itemize}

对$\forall p\in M$，存在包含点$p$的开集$\overline{U}_{\alpha}$和$\tilde{U}_{\alpha}$. 令$U_{\alpha}=\overline{U}_{\alpha}\cap \tilde{U}_{\alpha}$，并将映射$\overline{\varPhi}_{\alpha}$和$\tilde{\varPhi}_{\alpha}$限制在$U_{\alpha}$上来考虑。

假设$\{e_{1},e_{2},\cdots,e_{k_{1}} \}$为$\mr^{k_{1}}$中的标准基，$\{e'_{1},e'_{2},\cdots,e'_{k_{2}} \}$为$\mr^{k_{2}}$中的标准基。即$e_{i}$表示$\mr^{k_{1}}$中第$1\leqslant i\leqslant k_{1}$个分量为1，其余分量都为0的单位向量；$e'_{j}$表示$\mr^{k_{2}}$中第$1\leqslant j\leqslant k_{2}$个分量为1，其余分量都为0的单位向量。注意到$\overline{\varPhi}_{\alpha}(p,\;{\color{red}{\cdot}}\;):\mr^{k_{1}}\longrightarrow \overline{E}_{p}$为线性同构，所以$\{\overline{\varPhi}_{\alpha}(p,e_{\gamma_{ i }})|1\leqslant i\leqslant k_{1} \}$为$\overline{E}_{p}$的基；同理可知，$\{\tilde{\varPhi}_{\alpha}(p,e_{\gamma_{2}}) |1\leqslant j\leqslant k_{2}\}$为$\tilde{E}_{p}$的基。从而
\[
\{\overline{\varPhi}_{\alpha}(p,e_{i})\otimes \tilde{\varPhi}_{\alpha}(p,e'_{j})| 1\leqslant i\leqslant k_{1},1\leqslant j\leqslant k_{2}\} 
\]
就是$E_{p}=\overline{E}_{p}\otimes \tilde{E}_{p}$的基，并且默认按照字典顺序排列。

我们定义$\varPhi_{\alpha}:U_{\alpha}\times \mr^{k_{1}\times k_{2}} \longrightarrow \pi^{-1}(U_{\alpha})$为
\[
\varPhi_{\alpha}(p,a_{\eta},\cdots,a_{\gamma},\cdots,a_{\xi})\triangleq  \sum_{ \gamma\in \Gamma(k_{1},k_{2}) }a_{\gamma}\cdot\overline{\varPhi}_{\alpha}(p,e_{\gamma_{ 1 }})\otimes  \tilde{\varPhi}_{\alpha}(p,e'_{\gamma_{2}})
\]
这个映射的功能有点像前面张量丛那里的$\theta^{U}_{p}$，就是把一组数字$(a_{\eta},\cdots,a_{\gamma},\cdots,a_{\xi})\in \mr^{k_{1}\times k_{2}}$看作为某个张量在基下的坐标。
显然这样定义的映射$\varPhi_{\alpha}:U_{\alpha}\times \mr^{k_{1}+k_{2}} \longrightarrow \pi^{-1}(U_{\alpha})$为双射。

当$U_{\alpha}\cap U_{\beta}\neq\varnothing$时，则也有$\overline{U}_{\alpha}\cap\overline{U}_{\beta}\neq\varnothing,\tilde{U}_{\alpha}\cap\tilde{U}_{\beta}\neq\varnothing$. 则
\[
\begin{split}
&\varPhi^{ -1}_{\beta}\circ\varPhi_{\alpha}(p,a_{\eta},\cdots,a_{\gamma},\cdots,a_{\xi})=(p,b_{\eta},\cdots,b_{\gamma},\cdots,b_{\xi})\\
\Leftrightarrow&\varPhi_{\alpha}(p,a_{\eta},\cdots,a_{\gamma},\cdots,a_{\xi})=\varPhi_{\beta}(p,b_{\eta},\cdots,b_{\gamma},\cdots,b_{\xi})\\
\Leftrightarrow&\sum_{ \gamma\in \Gamma(k_{1},k_{2}) }a_{\gamma}\cdot\overline{\varPhi}_{\alpha}(p,e_{\gamma_{ 1 }})\otimes  \tilde{\varPhi}_{\alpha}(p,e'_{\gamma_{2}})=X_{P}=\sum_{ \eta\in \Gamma(k_{1},k_{2}) }a_{\eta}\cdot\overline{\varPhi}_{\beta}(p,e_{\eta_{ 1 }})\otimes  \tilde{\varPhi}_{\beta}(p,e'_{\eta_{2}})
\end{split}
\]
注意到$\{\overline{\varPhi}_{\alpha}(p,e_{i})\otimes \tilde{\varPhi}_{\alpha}(p,e'_{j})| 1\leqslant i\leqslant k_{1},1\leqslant j\leqslant k_{2}\} $与$\{\overline{\varPhi}_{\beta}(p,e_{i})\otimes \tilde{\varPhi}_{\beta}(p,e'_{j})| 1\leqslant i\leqslant k_{1},1\leqslant j\leqslant k_{2}\} $都是张量空间$E_{p}=\overline{E}_{p}\otimes \tilde{E}_{p}$的基，所以$(a_{\eta},\cdots,a_{\gamma},\cdots,a_{\xi})$与$(b_{\eta},\cdots,b_{\gamma},\cdots,b_{\xi})$之间的关系就是同一个张量$X_{p}$在不同基下的坐标之间的关系，为了求出这个坐标变换，我们需要写出基的过渡矩阵。设$\overline{E}_{p}$中基的过渡矩阵为$A$，即
\begin{equation}
	\left(\overline{\varPhi}_{\beta}(p,e_{1}),\cdots,\overline{\varPhi}_{\beta}(p,e_{k_{1}}) \right)=\left(\overline{\varPhi}_{\alpha}(p,e_{1}),\cdots,\overline{\varPhi}_{\alpha}(p,e_{k_{1}}) \right)\cdot A
\end{equation}
$\tilde{E}_{p}$中基的过渡矩阵为$B$，即
\[
\left(\tilde{\varPhi}_{\beta}(p,e'_{1}),\cdots,\tilde{\varPhi}_{\beta}(p,e'_{k_{2}}) \right)=\left(\tilde{\varPhi}_{\alpha}(p,e'_{1}),\cdots,\tilde{\varPhi}_{\alpha}(p,e'_{k_{2}}) \right)\cdot B
\]
则$E_{p}=\overline{E}_{p}\otimes \tilde{E}_{p}$中基的过渡矩阵为$A\otimes B$，即
\[
\{\overline{\varPhi}_{\beta}(p,e_{i})\otimes \tilde{\varPhi}_{\beta}(p,e'_{j})\}=\{\overline{\varPhi}_{\alpha}(p,e_{i})\otimes \tilde{\varPhi}_{\alpha}(p,e'_{j})\}\cdot A\otimes B
\]
因此
\[
\left( \begin{array}{ccccc}{\vdots} \\ {b_{\gamma}}  \\ {\vdots} \\ {\vdots}\\ {\vdots}\end{array}\right)=A^{-1}\otimes B^{-1}\left( \begin{array}{ccccc}{\vdots} \\ {\vdots}  \\ {\vdots} \\ {a_{\eta}}\\ {\vdots}\end{array}\right)
\]
下面再来看$A^{-1}$和$B^{-1}$究竟是什么？不妨设
\[
\overline{t}_{\beta\alpha}(p)=\left( \begin{array}{cccc}{c_{11}} & {c_{12}} & {\cdots} & {c_{1 n}} \\ {c_{21}} & {c_{22}} & {\cdots} & {c_{2 n}} \\ {\vdots} & {\vdots} & { } & {\vdots} \\ {c_{m 1}} & {c_{m 2}} & {\cdots} & {c_{m n}}\end{array}\right)
\]
我们来计算$\overline{\varPhi}_{\alpha}(p,e_{1})=\overline{\varPhi}_{\alpha}(p,1,0,\cdots,0)$. 设$\overline{\varPhi}_{\alpha}(p,e_{1})=\overline{\varPhi}_{\beta}(p,b_{1},\cdots,b_{k_{1}})$，则有
\[
\left( \begin{array}{ccc}{b_{1}} \\ {\vdots}  \\ {b_{k_{1}}} \end{array}\right)=\overline{t}_{\beta\alpha}(p) \left( \begin{array}{ccc}{a_{1}} \\ {\vdots}  \\ {a_{k_{1}}} \end{array}\right)=\left( \begin{array}{ccc}{c_{11}} \\ {\vdots}  \\ {c_{k_{1}1}} \end{array}\right)
\]
于是$\overline{\varPhi}_{\alpha}(p,e_{1})=\overline{\varPhi}_{\beta}(p,c_{11},\cdots,c_{k_{1}1})=c_{11}\overline{\varPhi}_{\beta}(p,e_{1})+\cdots+c_{k_{1}1}\overline{\varPhi}_{\beta}(p,e_{k_{1}})$.
依次计算$\overline{\varPhi}_{\alpha}(p,e_{2}),\cdots,\overline{\varPhi}_{\alpha}(p,e_{k_{1}})$，可以发现
\[
\left \{ \begin{array} { l } {\overline{\varPhi}_{\alpha}(p,e_{1})=\overline{\varPhi}_{\beta}(p,c_{11},\cdots,c_{k_{1}1})=c_{11}\overline{\varPhi}_{\beta}(p,e_{1})+\cdots+c_{k_{1}1}\overline{\varPhi}_{\beta}(p,e_{k_{1}}) } \\ { \overline{\varPhi}_{\alpha}(p,e_{2})=\overline{\varPhi}_{\beta}(p,c_{12},\cdots,c_{k_{1}2})=c_{12}\overline{\varPhi}_{\beta}(p,e_{1})+\cdots+c_{k_{1}2}\overline{\varPhi}_{\beta}(p,e_{k_{1}})  } \\{\vdots}\\ {  \overline{\varPhi}_{\alpha}(p,e_{k_{1}})=\overline{\varPhi}_{\beta}(p,c_{1k_{1}},\cdots,c_{k_{1}1})=c_{1k_{1}}\overline{\varPhi}_{\beta}(p,e_{1})+\cdots+c_{k_{1}1}\overline{\varPhi}_{\beta}(p,e_{k_{1}})  }\end{array} \right.
\]
也就是
\[
\left(\overline{\varPhi}_{\alpha}(p,e_{1}),\cdots,\overline{\varPhi}_{\alpha}(p,e_{k_{1}}) \right)=\left(\overline{\varPhi}_{\beta}(p,e_{1}),\cdots,\overline{\varPhi}_{\beta}(p,e_{k_{1}}) \right)\left( \begin{array}{cccc}{c_{11}} & {c_{12}} & {\cdots} & {c_{1 n}} \\ {c_{21}} & {c_{22}} & {\cdots} & {c_{2 n}} \\ {\vdots} & {\vdots} & { } & {\vdots} \\ {c_{m 1}} & {c_{m 2}} & {\cdots} & {c_{m n}}\end{array}\right)
\]
比照式\requ{equ:}便知，$A^{-1}=\overline{t}_{\beta\alpha}(p)$. 同理可知$B^{-1}=\tilde{t}_{\beta\alpha}(p)$. 所以我们就有
\[
\left( \begin{array}{ccccc}{\vdots} \\ {b_{\gamma}}  \\ {\vdots} \\ {\vdots}\\ {\vdots}\end{array}\right)=\overline{t}_{\beta\alpha}(p)\otimes \tilde{t}_{\beta\alpha}(p)\left( \begin{array}{ccccc}{\vdots} \\ {\vdots}  \\ {\vdots} \\ {a_{\eta}}\\ {\vdots}\end{array}\right)
\]

因此，当$U_{\alpha}\cap U_{\beta}\neq\varnothing$时，我们定义$t_{\beta\alpha}:U_{\alpha}\cap U_{\beta}\longrightarrow GL_{k_{1}+k_{2}}(\mr)$为
\begin{equation}
	t_{\beta\alpha}\triangleq\overline{t}_{\beta\alpha}\otimes \tilde{t}_{\beta\alpha}
\end{equation}
这里$\overline{t}_{\beta\alpha}\otimes \tilde{t}_{\beta\alpha}$表示矩阵的张量积。

由于$\overline{t}_{\beta\alpha}:\overline{U}_{\alpha}\cap \overline{U}_{\beta}\longrightarrow GL_{k_{1}}(\mr),\tilde{t}_{\beta\alpha}:\tilde{U}_{\alpha}\cap \tilde{U}_{\beta}\longrightarrow GL_{k_{2}}(\mr)$都是光滑映射，所以如此定义的$t_{\beta\alpha}:U_{\alpha}\cap U_{\beta}\longrightarrow GL_{k_{1}+k_{2}}(\mr)$也是光滑映射。

可以验证以上给出的$\{\varPhi_{\alpha}:U_{\alpha}\times \mr^{k_{1}\times k_{2}} \longrightarrow \pi^{-1}(U_{\alpha})\}$以及$\{ t_{\beta\alpha} \}$满足\rthm{thm:chapter2向量丛的等价定义}的所有要求，从而使得$(E,M,\pi)$成为$M$的一个新的秩为$k_{1}\times k_{2}$的向量丛，式\requ{equ:}给出了它的过渡函数族。我们称这样得到的向量丛$(E,M,\pi)$为向量丛$(\overline{E},M,\overline{\pi})$与$(\tilde{E},M,\tilde{\pi})$的\textbf{张量积}，并记作$E\triangleq \overline{E}\otimes \tilde{E}$.


以上介绍了三种常用的向量丛的运算。所有向量丛里面，最简单也是最重要的应该就是切丛了，比如切丛的对偶丛就是余切丛，再将切丛与余切丛做张量积就得到了\ref{张量丛}节中所讨论的的张量丛。

\subsection{拉回丛}

设$f:M^{m}\longrightarrow N^{n}$是一个光滑映射，并给定$N$上的一个秩为$k$的向量丛$(E,N,\pi)$. 利用映射$f$我们来给流形$M$诱导一个向量丛$(E^{\prime},M,\pi^{\prime})$.

对任意$p\in M$，令$E^{\prime}_{p}\triangleq \{p\}\times E_{f(p)}$. 即给流形$M$的每一点赋予向量空间$\{p\}\times E_{f(p)}$，注意这里赋予的向量空间是很有讲究的，必须是$\{p\}\times E_{f(p)}$而不能仅仅只是$ E_{f(p)}$.

再令$E^{\prime}\triangleq \bigcup\limits_{p\in M}E^{\prime}_{p}$，这其实等价于不交并
\[
E^{\prime}=\bigcup_{p\in M}E^{\prime}_{p}=\bigsqcup_{p\in M} E_{f(p)}=\{(p,X)|(p,X)\in M\times E\;\text{\kaishu 且}\;f(p)=\pi(X) \}
\]
显然$E^{\prime}\subset M\times E$. $\pi^{\prime}:E^{\prime}\longrightarrow M$定义为向第一个分量的投影，即$\pi^{\prime}\left(p,X  \right)=p,\forall X\in E_{f(p)}$.

由于$(E,N,\pi)$为向量丛，根据\rthm{thm:chapter2向量丛的性质}可知，存在$N$的开覆盖$\{V_{\alpha} \}_{\alpha\in \Lambda}$以及双射族$\{\varPsi_{\alpha}:V_{\alpha}\times \mr^{k}\longrightarrow \pi^{-1}(V_{\alpha}) \}$，满足以下性质：
\begin{itemize}
	\item[(1)] $\varPsi_{\alpha}(q,a)\in E_{q},\forall q\in V_{\alpha},\forall a\in \mr^{k}$.
	\item[(2)] $\varPsi_{\alpha}(q,\;{\color{red}{\cdot}}\;):\mr^{k}\longrightarrow E_{q}$为线性同构.
	\item[(3)] 当$V_{\alpha}\cap V_{\beta}\neq\varnothing$时，存在光滑映射$t_{\beta\alpha}:V_{\alpha}\cap V_{\beta}\longrightarrow GL_{k}(\mr)$，使得对$\forall p\in V_{\alpha}\cap V_{\beta},\forall (a_{1},\cdots,a_{k})\in \mr^{k}$，若$\varPsi_{\beta}^{-1}\circ\varPsi_{\alpha}(p,a_{1},\cdots,a_{k})=(p,b_{1},\cdots,b_{k})$，则
	\[
	\left( \begin{array}{ccc}{b_{1}} \\ {\vdots}  \\ {b_{k}} \end{array}\right)=t_{\beta\alpha}(p) \left( \begin{array}{ccc}{a_{1}} \\ {\vdots}  \\ {a_{k}} \end{array}\right)
	\]
\end{itemize}

根据以上条件，我们构造$U_{\alpha}\triangleq f^{-1} \left( V_{\alpha} \right)$. 注意到$f:M\longrightarrow N$是光滑映射，自然是连续映射，故$f^{-1} \left( V_{\alpha} \right)$是$M$上的开集。又因为$f(M)\subset N=\bigcup_{\alpha\in \Lambda}V_{\alpha}$，所以$\{U_{\alpha}\}_{\alpha\in \Lambda}$肯定也能覆盖$M$. 也就是说$\{U_{\alpha}\}_{\alpha\in \Lambda}$是$M$的一个开覆盖。

再定义映射$\varPhi_{\alpha} :U_{\alpha}\times \mr^{k}\longrightarrow \pi^{\prime -1}(U_{\alpha})$如下：$\forall p\in U_{\alpha},\forall (x_{1},\cdots,x_{k})\in \mr^{k}$
\[
\varPhi_{\alpha}(p,x_{1},\cdots,x_{k})\triangleq \left( p,\varPsi_{\alpha}(f(p),x_{1},\cdots,x_{k}) \right)
\]
显然$\varPhi_{\alpha} :U_{\alpha}\times \mr^{k}\longrightarrow \pi^{\prime -1}(U_{\alpha})$是双射。当$U_{\alpha}\cap U_{\beta}\neq\varnothing$时，设$\varPhi^{ -1}_{\beta}\circ\varPhi_{\alpha}(p,x_{1},\cdots,x_{k})=(p,y_{1},\cdots,y_{k})$. 则
\[
\begin{split}
&\varPhi^{ -1}_{\beta}\circ\varPhi_{\alpha}(p,x_{1},\cdots,x_{k})=(p,y_{1},\cdots,y_{k})\\
&\Leftrightarrow \varPhi_{\alpha}(p,x_{1},\cdots,x_{k})=\varPhi_{\beta}(p,y_{1},\cdots,y_{k})\\
&\Leftrightarrow  \left( p,\varPsi_{\alpha}(f(p),x_{1},\cdots,x_{k}) \right)= \left( p,\varPsi_{\beta}(f(p),y_{1},\cdots,y_{k}) \right)\\
&\Leftrightarrow \varPsi_{\alpha}(f(p),x_{1},\cdots,x_{k})=\varPsi_{\beta}(f(p),y_{1},\cdots,y_{k})\\
&\Leftrightarrow \varPsi^{-1}_{\beta}\circ\varPsi_{\alpha}(f(p),x_{1},\cdots,x_{k})=(f(p),y_{1},\cdots,y_{k})\\
\end{split}
\]
注意到，当$U_{\alpha}\cap U_{\beta}\neq\varnothing$时，那么也有$V_{\alpha}\cap V_{\beta}\neq\varnothing$. 所以上面推导过程中出现的$\varPsi^{-1}_{\beta}\circ\varPsi_{\alpha}$是有意义的。根据已知条件(3)，我们得到
\[
\begin{split}
	&\varPsi^{-1}_{\beta}\circ\varPsi_{\alpha}(f(p),x_{1},\cdots,x_{k})=(f(p),y_{1},\cdots,y_{k})\\
	&\Leftrightarrow \left( \begin{array}{ccc}{b_{1}} \\ {\vdots}  \\ {b_{k}} \end{array}\right)=t_{\beta\alpha}(f(p)) \left( \begin{array}{ccc}{a_{1}} \\ {\vdots}  \\ {a_{k}} \end{array}\right)
\end{split}
\]
改写一下，也就是
\[
\left( \begin{array}{ccc}{b_{1}} \\ {\vdots}  \\ {b_{k}} \end{array}\right)=(t_{\beta\alpha}\circ f)(p) \left( \begin{array}{ccc}{a_{1}} \\ {\vdots}  \\ {a_{k}} \end{array}\right)
\]

 因此，当$U_{\alpha}\cap U_{\beta}\neq\varnothing$时，我们定义$t^{\prime}_{\beta\alpha}:U_{\alpha}\cap U_{\beta}\longrightarrow GL_{k}(\mr)$为
 \begin{equation}
	 t^{\prime}_{\beta\alpha}\triangleq t_{\beta\alpha}\circ f
	 \label{拉回丛的过渡函数}
 \end{equation}
 
 由于$t_{\beta\alpha}:U_{\alpha}\cap U_{\beta}\longrightarrow GL_{k}(\mr)$和$f:M\longrightarrow N$都是光滑映射，所以如此定义的$t^{\prime}_{\beta\alpha}:U_{\alpha}\cap U_{\beta}\longrightarrow GL_{k}(\mr)$也是光滑映射。
 
 可以验证以上给出的$\{\varPhi^{\prime}_{\alpha} :U_{\alpha}\times \mr^{k}\longrightarrow \pi^{\prime -1}(U_{\alpha})\}$以及$\{ t^{\prime}_{\beta\alpha} \}$满足\rthm{thm:chapter2向量丛的等价定义}的所有要求，从而使得$(E^{\prime},M,\pi^{\prime})$成为$M$的一个秩为$k$的向量丛，式\requ{拉回丛的过渡函数}给出了它的过渡函数族。我们称$(E^{\prime},M,\pi^{\prime})$是由$(E,N,\pi)$通过光滑映射$f:M\longrightarrow N$诱导的向量丛，简称\textbf{诱导丛}(induced bundle). 由式\requ{拉回丛的过渡函数}确定的过渡函数族$t^{\prime}_{\beta\alpha}$也经常用拉回映射的记号来记：
\[
	 t^{\prime}_{\beta\alpha}= t_{\beta\alpha}\circ f=f^{\ast}(t_{\beta\alpha})
 \]
 正因如此，我们也把$(E^{\prime},M,\pi^{\prime})$称为\textbf{拉回丛}(pullback bundle)，并重新记
 \[
	f^{\ast}(E)\triangleq E^{\prime}
 \]

\section{光滑截面}

截面这个概念可以看成是函数的推广。

\begin{mydef}{向量丛的光滑截面}{chapter2向量丛的光滑截面}
设$M$是光滑流形，$(E,M,\pi)$是$M$上的一个秩为$k$的光滑实向量丛. 若光滑映射$s:M\longrightarrow E$满足$\pi\circ s= id_{M}$，其中$id_{M}:M\longrightarrow M$是恒等映射，则称$s$是向量丛$E$的一个光滑截面(smooth section). 向量丛$E$的全体光滑截面的集合记作$\Gamma(E)$.
\end{mydef}

定义中最核心的要求$\pi\circ s= id_{M}$，其{\color{red}{等价表述}}是：$\forall p\in M$，$s(p)\in E_{p}$.

截面之间也可以进行运算：对任意$ s_{1},s_{2}\in \Gamma(E)$，任意光滑函数$\varphi:M\longrightarrow \mr$，规定
\begin{equation}
\begin{split}
	\left( s_{1}+s_{2} \right)(p)&\triangleq s_{1}(p)+s_{2}(p)\\
	\left( \varphi s\right)(p)&\triangleq \varphi(p)\cdot s(p)
\end{split}
\label{截面之间的加法与数乘}
\end{equation}
其中$p\in M$. $\Gamma(E)$在上述两个运算下形成一个$C^{\infty}(M)$-模。你不必深究啥叫“模”，只要知道截面可以相加，也可以和光滑函数相乘即可。

对于一些常用的、特殊的向量丛，它们的截面也有特别的称呼。
\begin{itemize}
\item $(r,s)$型张量丛$T^{r}_{s}$的光滑截面，我们称之为流形$M$上的\textbf{光滑张量场}.
\item $s$次外形式丛$\Lambda^{s}(T^{\ast}M)$的光滑截面，我们称之为流形$M$上的\textbf{$s$次外微分形式}，并将$\Gamma(\Lambda^{s}(T^{\ast}M))$重新记为$A^{s}(M)$. 特别地，$A^{1}(M)=\Gamma(T^{\ast}M),A^{0}(M)=C^{\infty}(M)$.
\item 外形式丛$\Lambda(T^{\ast}M)$的光滑截面，我们称之为流形$M$上的\textbf{外微分形式}，并将$\Gamma(\Lambda(T^{\ast}M))$重新记为$A(M)$.
\end{itemize}
特别地，切丛的光滑截面称为流形$M$上的\textbf{光滑切向量场}。对于余切丛的光滑截面，你既可以把它叫作光滑余切向量场(因为$T^{\ast}M=T^{0}_{1}$)，也可以叫作$1$次外微分形式(因为$T^{\ast}M=\Lambda^{1}(T^{\ast}M)$). 不过还是把切丛的截面称为\textbf{$1$次外微分形式}或简称\textbf{$1$-形式}更普遍一些。

\begin{myprop}{}{}
对$\forall \omega\in A(M)$，则$\omega$可以表示为直和(形式和)：
\[
\omega=\omega^{0}+\omega^{1}+\cdots+\omega^{m}
\]
其中$\omega^{i}\in A^{i}(M),i=0,1,\cdots,m=\operatorname{dim}(M)$.
\end{myprop}
\begin{proof}
因为$\omega\in A(M)=\Gamma(\Lambda(T^{\ast}M))$，所以对任意点$p\in M$，$\omega(p)\in \Lambda(T^{\ast}_{p}M)$. 注意到
\[
	\Lambda(T^{\ast}_{p}M)=\bigoplus_{s=0}^{m}\Lambda^{s}(T^{\ast}_{p}M)
\]
因此$\Lambda(T^{\ast}_{p}M)$里面的元素都长成这个样子：$\forall \xi_{p}\in \Lambda(T^{\ast}_{p}M),\xi_{p}=\xi^{0}_{p}+\xi^{1}_{p}+\cdots +\xi^{m}_{p}$，其中$\xi^{i}_{p}\in \Lambda^{i}(T^{\ast}_{p}M),i=0,1,\cdots,m$. 而现在$\omega(p)\in \Lambda(T^{\ast}_{p}M)$，于是
\[
	\omega(p)=\omega^{0}(p)+\omega^{1}(p)+\cdots+\omega^{m}(p)
\]
其中$\omega^{i}(p)\in \Lambda^{i}(T^{\ast}_{p}M)$，再由点$p$的任意性便知，$\omega^{i}\in \Gamma(\Lambda^{i}(T^{\ast}M))=A^{i}(M),i=0,1,\cdots,m$. 从而
\[
	\omega=\omega^{0}+\omega^{1}+\cdots+\omega^{m}
\]
\end{proof}

上述性质告诉我们，截面空间$A(M)$实际上是截面空间$A^{s}(M)$的直和，即
\[
A(M)=\bigoplus_{s=0}^{m}A^{s}(M)
\]

拉回丛的截面是一个非常特殊而重要的概念，所以我们单独讲一下。设$f:M^{m}\longrightarrow N^{n}$是光滑映射，$N$上带有秩为$k$的实向量丛$(E,N,\pi)$，由它诱导了$M$上的拉回丛$f^{\ast}E$. 根据截面的\rdef{def:chapter2向量丛的光滑截面}，拉回丛$f^{\ast}E$的一个截面是指一个映射$\xi:M\longrightarrow f^{\ast}E$，满足$\xi(p)\in {\color{red}{f^{\ast}_{p}E=E_{f(p)}}}$. 也就是说对每一点$p\in M$，给它指定一个$E_{f(p)}$中的元素。

我们把拉回丛的截面又称为\textbf{$M$上沿着映射$f$定义的截面}。这里“沿着映射$f$”的意思是说截面值的分布范围只在像集$f(M)$上。 

特别地，如果已经给定向量丛$E$在$N$上的一个截面$s$，那么我们可以用$s$来诱导出拉回丛的一个截面。令
\[
	\xi(p)\triangleq s(f(p)),\;\forall p\in M
\]
显然$\xi$是$M$上的一个沿着映射$f$定义的截面。而且根据上式可以看出$\xi=f^{\ast}s$. 

\section{外微分算子}

本节我们主要在截面空间$A(M)$上考虑。首先可以通过逐点定义的方式在$A(M)$中引进外微分形式的外积运算。

\begin{mydef}{$A(M)$上的外积}{chapter2AM上的外积}
对$\forall \omega,\eta\in A(M)$，定义$\omega\wedge\eta:M\longrightarrow \Lambda(T^{\ast}M)$如下：
\[
\left( \omega\wedge\eta \right)(p)\triangleq \omega(p)\wedge\eta(p),\quad\forall p\in M
\]
则$\omega\wedge\eta\in \Gamma(\Lambda(T^{\ast}M))=A(M)$. 我们称$\omega\wedge\eta$为外微分形式$\omega$与$\eta$的外积.
\end{mydef}

注意到$\omega(p),\eta(p)\in \Lambda(T^{\ast}_{p}M)$，所以$\omega(p)\wedge\eta(p)$执行的是外代数$\Lambda(T^{\ast}_{p}M)$上的外积运算，从而是有意义的表达式，并且运算结果$\omega(p)\wedge\eta(p)$仍在外代数$\Lambda(T^{\ast}_{p}M)$中。

如果把\rdef{def:chapter2AM上的外积}规定的外积运算当作是$A(M)$中元素之间的乘法运算，再加上由\requ{截面之间的加法与数乘}规定的加法与数乘(乘的是函数)运算，则$A(M)$关于这三种运算构成一个代数。

在本科的微积分或者数学分析中，基本的研究对象是函数；与之相对的是流形上的微积分，其基本的研究对象是外微分形式。以下定理给出如何对微分形式“求微分”的方法，这是函数微分的推广。

\begin{mythm}{外微分算子的存在唯一性}{chapter2外微分算子的存在唯一性}
设$M^{m}$为光滑流形，则存在唯一的映射$\dd:A(M)\longrightarrow A(M)$使得以下条件成立：
\begin{itemize}
\item[(1)] 对$\forall \omega\in A^{s}(M)$，均有$\dd(\omega)\in A^{s+1}(M)$. 其中$s=0,1,\cdots,m$.
\item[(2)] 对$\forall \omega,\eta\in A(M)$，$\forall k_{1},k_{2}\in\mr$，均有$\dd (k_{1}\omega+k_{2}\eta)=k_{1}\dd (\omega)+k_{2}\dd(\eta)$.
\item[(3)] 对$\forall \omega\in A^{s}(M),\eta\in A(M)$，均有$\dd(\omega\wedge\eta)=\dd(\omega)\wedge\eta+(-1)^{s}\cdot \omega\wedge\dd(\eta)$.
\item[(4)] 当$f$是$M$上的光滑函数时，$\dd(f)=\dd f$. 这里等号右边$\dd f$表示函数$f$的微分.
\item[(5)] 当$f$是$M$上的光滑函数时，$\dd(\dd(f))=0$. 
\end{itemize}
\end{mythm}

为了证明这条定理，需要做一些准备工作。

\begin{mylem}{外微分算子的局部性}{chapter2外微分算子的局部性}
	假设满足\rthm{thm:chapter2外微分算子的存在唯一性}的映射$\dd:A(M)\longrightarrow A(M)$存在. 对任意$\omega\in A(M)$，任意开集$U\subset M$，只要$\left.\omega \right|_{U}=0$，便有$\left.\dd \omega \right|_{U}=0$.
\end{mylem}

\begin{proof}
	$\forall p\in U$，根据流形的局部紧致性(\rthm{thm:chapter1流形局部紧致})可知，点$p$处有紧致邻域$W$，使得$p\in W\subset U$. 由\rthm{thm:chapter1流形上的截断函数}又知，存在截断函数$f$满足
	\[
	\left. f \right|_{W}\equiv 1\quad\quad\quad \left.f \right|_{M\backslash U}\equiv 0
	\]

	考虑$f\omega$. 因为在$U$上$\omega=0$，故$\left. f\omega \right|_{U}=0$，而在$U$之外$f=0$，故$\left.f \omega\right|_{M\backslash U}= 0$. 因此在整个$M$上$f\omega\equiv 0$. 从而根据\rthm{thm:chapter2外微分算子的存在唯一性}的条件(3)得
	\[
	0=\dd (f\omega)=\dd f\wedge \omega +f\dd \omega
	\]
	特别地在$p$点有$\left( \dd f \right)_{p}\wedge \left( \omega \right)_{p} +f(p)\left( \dd \omega \right)_{p}=0$. 因为$\omega_{p}=0,f(p)=1$，所以$\left( \dd \omega \right)_{p}=0$. 由$p\in U$的任意性即知$\left.\dd \omega \right|_{U}=0$.
\end{proof}

因此根据上述引理，如果$\omega,\eta\in A(M)$限制在某开集$U$上相等，即$\left.\omega \right|_{U}=\left.\eta \right|_{U}$，则有$\left.\dd\omega \right|_{U}=\left.\dd\eta \right|_{U}$.


下面来证明\rthm{thm:chapter2外微分算子的存在唯一性}.
\begin{proof}
	
\end{proof}


\section{Stokes公式}

\section{联络}

\subsection{向量丛上的联络}

\begin{mydef}{}{chapter2联络}
	设$M$是一个光滑流形，$E$是$M$上的秩为$q$的实向量丛. 若映射$D:\Gamma(E)\longrightarrow \Gamma(T^{\ast}M\otimes E)$满足：
	\begin{itemize}
		\item[(1)] $D(s_{1}+s_{2})=Ds_{1}+Ds_{2},\forall s_{1},s_{2}\in \Gamma(E)$.
		\item[(2)] $D(f s)=\dd f\otimes s+f\cdot Ds,\forall s\in \Gamma(E),\forall f\in C^{\infty}(M)$.
	\end{itemize}
则称$D$为向量丛$E$上的一个联络(connection)，$Ds$称为截面$s$的绝对微分.
\end{mydef}

\begin{remark}
	有的书上在定义联络时会预先约定$\forall f\in C^{\infty}(M),{\color{red}{Df\triangleq \dd f}}$. 我们采用的这种定义实际上蕴含了这一点。因此上述定义的第(2)个条件其实就是Leibniz律。
\end{remark}

我们来分析一下$D$是怎么起作用的。首先，对任意一个截面$s\in \Gamma(E)$，映射$D$将$s$变成$Ds\in \Gamma(T^{\ast}M\otimes E)$. 也就是说$Ds$是张量积丛$T^{\ast}M\otimes E$的光滑截面。所以对任意$p\in M$，$(Ds)(p)\in T^{\ast}_{p}M\otimes E_{p}$，从而$(Ds)(p)$长成这个样子
\[
(Ds)(p)=\sum_{ \omega,z }a_{\omega z}(p)  \omega_{p}\otimes Z_{p}
\]
其中$a_{\omega z}(p)\in \mr$是系数，$\omega_{p}\in T^{\ast}_{p}M,Z_{p}\in E_{p}$. 由于点$p\in M$是任意取的，所以如果我们不局限在一个点来看的话，那么就有
\[
Ds=\sum_{ \omega,z }a_{\omega z}\cdot \omega\otimes Z
\]
这里系数$a_{\omega z}$是$M$上的光滑函数，即$a_{\omega z}\in C^{\infty}(M)$，而$\omega$则满足$\omega(p)\in T^{\ast}_{p}M,\forall p\in M$，这表明$\omega$是$T^{\ast}M$的光滑截面，即$\omega\in \Gamma(T^{\ast}M)$. 同理，$Z$满足$Z(p)\in E_{p}$，即$Z$是$E$的光滑截面，或者说$Z\in \Gamma(E)$. 这样我们就对$Ds$有了一点直观的感觉。

再来看$D$所满足的两个条件。在条件(2)中，取$f=k\in\mr$，则有$\dd k=0\Rightarrow D(k  s)=k  Ds$. 再结合条件(1)，这表明$D$是一个线性映射。特别地，$D(-s)=-Ds$，当$s=0$时便有$D(0)=-D0\Rightarrow D0=0$. 

$D$还具有局部性：若在开集$U$上$s_{1}=s_{2}$，则在$U$上有$Ds_{1}=Ds_{2}$.

\begin{mydef}{}{chapter2协变导数}
	设$X$是光滑流形$M$上的任一$C^{\infty}$向量场，即$X\in \Gamma(TM)$. 对$\forall s\in \Gamma(E)$，令
	\[
	D_{X}s\triangleq\langle Ds\;,X \rangle \in \Gamma(E)
	\]
	其中$\langle \bullet\;,\bullet\rangle $表示$T^{\ast}M$与$TM$之间的配合运算. 则称$D_{X}s$为截面$s$沿着向量场$X$的协变导数或共变导数.
\end{mydef}

记住，所有的配合运算都是指余切空间里的元素和切空间里的元素之间的配合运算。有了这一句话在心中，我们就可以切实抓住看似虚无缥缈的定义其中所蕴含的真意。对任意一点$p\in M$，设$(Ds)(p)=\sum_{ \omega,Z }a_{\omega Z}(p)  \omega_{p}\otimes Z_{p}$，其中$a_{\omega Z}(p)\in \mr,\omega_{p}\in T^{\ast}_{p}M,Z_{p}\in E_{p}$. $X(p)=X_{p}\in T_{p}M$. 则有
\[
(D_{X}s)(p)\triangleq\langle (Ds)(p)\;,X_{p} \rangle =\langle \sum_{ \omega,z }a_{\omega z}(p)  {\color{red}{\omega_{p}}}\otimes Z_{p},{\color{red}{X_{p}}} \rangle=\sum_{ \omega,z }a_{\omega z}(p)  {\color{red}{\omega_{p}(X_{p})}}  Z_{p}\in E_{p}
\]
因为$p$是任意的，所以
\[
(D_{X}s)=\langle Ds\;,X \rangle =\langle \sum_{ \omega,z }a_{\omega z}  {\color{red}{\omega}}\otimes Z\;,{\color{red}{X}} \rangle=\sum_{ \omega,z }a_{\omega z}  {\color{red}{\omega(X)}}  Z\in \Gamma(E)
\]
这里${\color{red}{\omega(X)}}(p)\triangleq{\color{red}{\omega_{p}(X_{p})}}$. 从而${\color{red}{\omega(X)}}$就是$M$上的光滑函数，即${\color{red}{\omega(X)}}\in C^{\infty}(M)$. 此即为定义中所说的“$\langle \bullet\;,\bullet\rangle $表示$T^{\ast}M$与$TM$之间的配合运算”的含义。

协变导数有下面这些重要的性质。
\begin{mythm}{}{chapter2协变导数的性质}
	\begin{itemize}
		\item[(1)] $D_{X+Y}(s)=D_{X}s+D_{Y}s,\forall X,Y\in \Gamma(TM)$.
		\item[(2)] $D_{fX}(s)=f(D_{X}s),\forall f\in C^{\infty}(M),\forall s\in \Gamma(E)$.
		\item[(3)] $D_{X}(s_{1}+s_{2})=D_{X}s_{1}+D_{X}s_{2},\forall s_{1},s_{2}\in\Gamma(E) $.
		\item[(4)] $D_{X}(fs)=(Xf)s+f(D_{X}s),\forall f\in C^{\infty}(M),\forall s\in \Gamma(E)$.
	\end{itemize}
\end{mythm}

注意到$Df=\dd f$，所以$D_{X}f=\langle Df,X \rangle=\langle\dd f,X \rangle=\dd f(X)=Xf$，所以第(4)条性质也可以改写成$D_{X}(fs)=(D_{X}f)s+f(D_{X}s)$，因此第(4)条性质本质上也表示Leibniz律。

下面我们来看联络的局部表示，这对于实际计算有最直接的帮助。

\begin{mydef}{}{chapter2局部标架场}
	设$E$是光滑流形$M$上秩为$q$的向量丛，$U\subset M$为一个开子集. 若$E$在$U$上的$q$个截面$\{s_{1},\cdots,s_{q} \}$满足：对任意$p\in U$，$\{s_{1}(p),\cdots,s_{q}(p) \}$为向量空间$E_{p}$的基，则称$\{s_{1},\cdots,s_{q} \}$为$E$在$U$上的一个局部标架场.
\end{mydef}

取定$T^{\ast}M$在$U$上的一个局部标架场$\{\theta^{1},\cdots,\theta^{m} \}$，取定$E$在$U$上的一个局部标架场$\{s_{1},\cdots,s_{q} \}$. 于是对任意$p\in U$，$\{\theta^{i}(p)\otimes s_{\alpha} (p)| 1\leqslant i\leqslant m,1\leqslant \alpha\leqslant q\}$为向量空间$T^{\ast}_{p}M\otimes E_{p}$的基，由于$p\in U$是任意的，所以$\{\theta^{i}\otimes s_{\alpha} | 1\leqslant i\leqslant m,1\leqslant \alpha\leqslant q\}$为丛$T^{\ast}M\otimes E$在$U$上的局部标架场。以后我们省略使用“对任意$p\in U$”这个步骤，而直接用函数、截面和标架场的语言来书写，这不会产生理解上的困难。

因为$Ds_{\alpha}\in \Gamma(T^{\ast}M\otimes E)$，所以$Ds_{\alpha}$就可以表示为
\[
Ds_{\alpha}=\sum_{\beta,i}\Gamma^{\beta}_{\alpha i}  \theta^{i}\otimes s_{\beta}
\]
其中$\Gamma^{\beta}_{\alpha i}\in C^{\infty}(U)$. 令${\color{red}{\omega^{\beta}_{\alpha}\triangleq \Gamma^{\beta}_{\alpha i}  \theta^{i}}}$，显然$\omega^{\beta}_{\alpha}\in \Gamma(T^{\ast}M)$. 我们把$\omega^{\beta}_{\alpha}$称为\textbf{联络1-形式}。于是$Ds_{\alpha}$又可以写成
\begin{equation}
	Ds_{\alpha}=\omega^{\beta}_{\alpha}\otimes s_{\beta}
\end{equation}
引进矩阵的记号，令
\[
S=\left( \begin{array}{ccc}{s_{1}} \\ {\vdots}  \\ {s_{q}} \end{array}\right)\quad DS=\left( \begin{array}{ccc}{Ds_{1}} \\ {\vdots}  \\ {Ds_{q}} \end{array}\right) \quad  \omega=\left( \begin{array}{cccc}{\omega^{1}_{1}} & {\omega^{2}_{1}}&{\cdots}&{\omega^{q}_{1}}\\ {\omega^{1}_{2}} & {\omega^{2}_{2}}&{\cdots}&{\omega^{q}_{2}}\\ {\vdots} & {\vdots}&{ }&{\vdots}\\{\omega^{1}_{q}} & {\omega^{2}_{q}}&{\cdots}&{\omega^{q}_{q}}\\ \end{array}\right)
\]
我们把矩阵$\omega$称为\textbf{联络矩阵}，它的元素是联络1-形式。利用联络矩阵，我们可以将上式简记为
\[
DS=\omega\otimes S
\]
特别要说明的是，$\omega\otimes S$不是前面介绍的矩阵张量积运算！我们这里的$\omega\otimes S$的含义是：你就按照普通的矩阵的乘法来算，只不过乘的时候用$\otimes$来乘，有点滥用记号了，不过读者只要明白了我们约定的含义，也就不会搞混了。
\begin{remark}
	如果我们把开集$U\subset M$直接取为流形$M$的参数化$\cx$，那么$T^{\ast}M$在$U$上自带一个局部标架场$\{\dd x^{1},\dd x^{2},\cdots,\dd x^{m} \}$. 此时你只需要给向量丛$E$指定局部标架场即可。
\end{remark}

联络矩阵$\omega$依赖于向量丛$E$在$U$上的局部标架场的选取。如果在开集$U$上另取$E$的局部标架场$\{s'_{1},\cdots,s'_{q} \}$，而$T^{\ast}M$在$U$上的局部标架场仍取作$\{\theta^{1},\cdots,\theta^{m} \}$. 此时相应的联络矩阵设为$\omega'$. 我们来看$\omega$和$\omega'$之间是什么关系？
\begin{mythm}{}{chapter2联络矩阵的变换公式}
	若$S'=A  S$，其中$A=(f_{ij})_{q\times q }$为可逆的矩阵函数。则$\omega'=(\dd A)  A^{-1}+A  \omega  A^{-1} $.
\end{mythm}
所谓“$A=(f_{ij})_{q\times q }$为可逆的矩阵函数”，也就是说，对$\forall p\in U$，$A(p)$给出了向量空间$E_{p}$中的基$\{s_{1}(p),\cdots,s_{q}(p) \}$到基$\{s'_{1}(p),\cdots,s'_{q}(p) \}$的过渡矩阵，即$S'(p)=A(p)S(p)$. 
\begin{proof}
	证明思路很简单，就是直接计算$DS'=D(A S)$. 
	
	不过$D(AS)$是矩阵形式的，我们还是要回到它的分量形式来看$D$具体是怎么作用的。$A  S$的第$k$个元素是$f_{k\beta}s_{\beta}$. 所以$D(f_{k\beta}s_{\beta})=\dd f_{k\beta}\otimes s_{\beta}+f_{k\beta}\cdot Ds_{\beta}$. 再恢复成矩阵形式，便得到
	\[D(A  S)=(\dd A)\otimes S+A\cdot (DS)\]
	其中$\dd A$表示对矩阵$A$的每个元素分别求外微分，即$\dd A=(\dd f_{ij})_{q\times q }$. 而$A\cdot (DS)=A\cdot (\omega\otimes S)$. 矩阵$A\cdot (\omega\otimes S)$中的第$k$元素是$f_{ki}\cdot (\omega^{j}_{i}\otimes S_{j})$，因为$f_{ki}$是函数，对任意一点$p\in U$，函数值$f_{ki}(p)$就是一个实数，所以$f_{ki}$放在哪都可以。因此$f_{ki}\cdot (\omega^{j}_{i}\otimes S_{j})=(f_{ki}\cdot \omega^{j}_{i})\otimes S_{j}$. 再恢复成矩阵形式
	\[
	\begin{split}
	D(AS)&=(\dd A)\otimes S+A\cdot (DS)\\
	&=(\dd A)\otimes S+A\cdot (\omega\otimes S)\\
	&=(\dd A)\otimes S+(A\cdot \omega)\otimes S\\
	&=(\dd A+A  \omega )\otimes S
	\end{split}
	\]
	而$S=A^{-1}  S'$，所以$D(AS)=(\dd A+A  \omega )\otimes (A^{-1}  S')$. 假设$A^{-1}=(g_{ij})_{q\times q}$. 同理我们来考察$(\dd A+A \omega )\otimes (A^{-1}  S')$的分量，它的第$k$个分量是$(\dd f_{kj}+f_{kl}\cdot \omega^{j}_{l})\otimes (g_{ji}\cdot s'_{i})$. 函数系数放哪都无所谓，所以
	\[
	\begin{split}
	(\dd f_{kj}+f_{kl}\cdot \omega^{j}_{l})\otimes (g_{ji}\cdot s'_{i})&=\dd f_{kj}\otimes  (g_{ji}\cdot s'_{i})+(f_{kl}\cdot \omega^{j}_{l})\otimes (g_{ji}\cdot s'_{i})\\
	&=(\dd f_{kj}\cdot g_{ji})  \otimes s'_{i}+(f_{kl}\cdot \omega^{j}_{l}\cdot g_{ji})\otimes s'_{i}\\
	&=(\dd f_{kj}\cdot g_{ji}+f_{kl}\cdot \omega^{j}_{l}\cdot g_{ji})\otimes s'_{i}
	\end{split}
	\]
	恢复成矩阵形式，也就是
	\[
	(\dd A+A  \omega )\otimes S=\left((\dd A)  A^{-1}+A   \omega  A^{-1} \right)\otimes S'
	\]
	
	把以上过程整理一下，我们得到	
	\[
	\begin{split}
	DS'&=D(AS)=(\dd A)\otimes S+A\cdot (DS)\\
	&=(\dd A)\otimes S+A\cdot (\omega\otimes S)\\
	&=(\dd A+A\omega )\otimes S\\
	&=(\dd A+A\omega )\otimes (A^{-1}  S')\\
	&=\underbrace{\left((\dd A)  A^{-1}+A   \omega  A^{-1} \right)}_{=\omega'}\otimes S'
	\end{split}
	\]
	从而$\omega'=(\dd A)  A^{-1}+A   \omega  A^{-1} $.
\end{proof}
 
在上述证明中，我们给出了如何在矩阵形式下计算联络的详细过程，核心就是去看分量是如何算的，再恢复成矩阵形式即可。以后我就不一一说明分量是如何算得了，而直接写矩阵表达式。

\begin{mydef}{}{chapter2曲率矩阵}
	设$\omega=(\omega_{\alpha}^{\beta})_{q\times q}$为联络矩阵，令
	\[
	\Omega=\dd \omega-\omega \wedge \omega
	\]
	这里$\omega \wedge \omega$表示按照普通矩阵的乘法相乘，只是乘的时候按$\wedge$来算. 我们称$\Omega$为联络$D$在$U$上的曲率矩阵.
\end{mydef}

因为$\omega$里面的元素是1-形式，所以如上定义的$\Omega$里面的元素都是2-形式。虽然名字叫曲率矩阵，但是目前我们只是从代数上定义的，还没有赋予它几何背景。与\rthm{thm:chapter2联络矩阵的变换公式}相应，我们可以给出曲率矩阵在向量丛的不同局部标架场下之间的关系。

\begin{mythm}{}{chapter2曲率矩阵的变换公式}
	若$S'=A  S$，其中$A=(f_{ij})_{q\times q }$为可逆的矩阵函数. 则$\Omega'=A \Omega  A^{-1}$.
\end{mythm}
\begin{proof}
	因为$\omega'=(\dd A)  A^{-1}+A   \omega  A^{-1} $，即$\omega'  A=\dd A+A   \omega$. 对该式两边求外微分
	\[
	\dd (\omega'  A)=\dd (\dd A+A   \omega)=\dd (\dd A)+\dd (A   \omega)=\dd (A   \omega)
	\]
	首先来计算$\dd (\omega' A)$. 先看分量$\dd (\omega'^{k}_{i}\cdot f_{kj})$. {\color{red}{因为$f_{kj}$是函数，也就是0-形式}}，所以我们有${\color{red}{\omega'^{k}_{i}\cdot f_{kj}=\omega'^{k}_{i}\wedge f_{kj}}}$. 于是
	\[
	\dd (\omega'^{k}_{i}\cdot f_{kj})=\dd(\omega'^{k}_{i}\wedge f_{kj})=\dd(\omega'^{k}_{i}){\color{red}{\wedge f_{kj}}}+(-1)^{1}\omega'^{k}_{i}\wedge \dd(f_{kj})=\dd(\omega'^{k}_{i}){\color{red}{\cdot f_{kj}}}-\omega'^{k}_{i}\wedge \dd(f_{kj})
	\]
	因此$\dd (\omega'  A)=(\dd\omega')  A-\omega'\wedge(\dd A)$. 
	
	其次来计算$\dd (A   \omega)$. 先看分量$\dd(f_{ik}\cdot \omega^{j}_{k} )$. 由于$f_{ik}$是0-形式，所以$f_{ik}\cdot \omega^{j}_{k} =(f_{ik}\wedge \omega^{j}_{k} )$. 于是
	\[
	\dd(f_{ik}\cdot \omega^{j}_{k} )=\dd(f_{ik}\wedge \omega^{j}_{k} )=\dd(f_{ik})\wedge \omega^{j}_{k}+(-1)^{0}{\color{red}{f_{ik}\wedge}} \dd (\omega^{j}_{k})=\dd(f_{ik})\wedge \omega^{j}_{k}+{\color{red}{f_{ik}\cdot}}  \dd (\omega^{j}_{k})
	\]
	因此$\dd (A   \omega)=(\dd A)\wedge \omega+A  (\dd\omega)$. 从而
	\[
	\begin{split}
	&(\dd\omega')  A-\omega'\wedge(\dd A)=(\dd A)\wedge \omega+A  (\dd\omega)\\
	\Leftrightarrow& (\dd\omega')  A-\omega'\wedge(\underbrace{\omega'  A-A  \omega}_{=\dd A})=(\underbrace{\omega'  A-A  \omega}_{\dd A})\wedge \omega+A  (\dd\omega)\\
	\Leftrightarrow& (\dd\omega')  A-\omega'\wedge(\omega'  A)+\omega'\wedge (A  \omega)=(\omega'  A)\wedge \omega-(A  \omega)\wedge\omega+A  (\dd\omega)\\
	\end{split}
	\]
	注意到$A$是0-形式组成的矩阵，所以$\omega'\wedge ({\color{red}{A\cdot \omega}})=\omega'\wedge ({\color{red}{A\wedge \omega}})=({\color{blue}{\omega'\wedge A}})\wedge \omega=({\color{blue}{\omega'\cdot A}})\wedge \omega$. 因此
	\[
	\begin{split}
	& (\dd\omega')  A-\omega'\wedge(\omega'  A)=-(A \omega)\wedge\omega+A  (\dd\omega)\\
	\Leftrightarrow& (\dd\omega'-\omega'\wedge\omega')   A=A  (\dd \omega-\omega\wedge\omega)\\
	\Leftrightarrow& \Omega'  A=A  \Omega\\
	\Leftrightarrow&\Omega'=A  \Omega  A^{-1}
	\end{split}
	\]
\end{proof}

\begin{mydef}{}{chapter2曲率算子}
	任意取定$X,Y\in \Gamma(TM)$. 设$\{s_{1},s_{2},\cdots ,s_{q} \}$为向量丛$E$在开集$U$上的局部标架场. \\
	定义映射$\rr(X,Y):\Gamma(E)\longrightarrow \Gamma(E)$为：对$\forall s=\lambda^{\alpha}s_{\alpha}\in \Gamma(E)$
	\[
	\rr(X,Y)s=\left(\lambda^{1},\lambda^{2},\cdots,\lambda^{q} \right)\left( \begin{array}{cccc}{\Omega^{1}_{1}(X,Y)} & {\Omega^{2}_{1}(X,Y)}&{\cdots}&{\Omega^{q}_{1}(X,Y)}\\ {\Omega^{1}_{2}(X,Y)} & {\Omega^{2}_{2}(X,Y)}&{\cdots}&{\Omega^{q}_{2}(X,Y)}\\ {\vdots} & {\vdots}&{ }&{\vdots}\\{\Omega^{1}_{q}(X,Y)} & {\Omega^{2}_{q}(X,Y)}&{\cdots}&{\Omega^{q}_{q}(X,Y)}\\ \end{array}\right)\left( \begin{array}{cccc}{s_{1}} \\{s_{2}}\\ {\vdots}  \\ {s_{q}} \end{array}\right)
	\]
	我们称映射$\rr(X,Y)$为曲率算子.
\end{mydef}
曲率算子$\rr(X,Y)$与向量丛$E$的局部标架场的选取无关，因而上述定义的良定的。事实上，如果$\{s'_{1},s'_{2},\cdots ,s'_{q} \}$为向量丛$E$在开集$U$上的另一个局部标架场，不妨设$S'=A  S$. 则有$\Omega'=A  \Omega  A^{-1}$以及$\left(\lambda'^{1},\cdots,\lambda'^{q} \right)=\left(\lambda^{1},\cdots,\lambda^{q} \right) A^{-1}$. 所以
\[
\begin{split}
&\left(\lambda'^{1},\lambda'^{2},\cdots,\lambda'^{q} \right)\left( \begin{array}{cccc}{\Omega'^{1}_{1}(X,Y)} & {\Omega'^{2}_{1}(X,Y)}&{\cdots}&{\Omega'^{q}_{1}(X,Y)}\\ {\Omega'^{1}_{2}(X,Y)} & {\Omega'^{2}_{2}(X,Y)}&{\cdots}&{\Omega'^{q}_{2}(X,Y)}\\ {\vdots} & {\vdots}&{ }&{\vdots}\\{\Omega'^{1}_{q}(X,Y)} & {\Omega'^{2}_{q}(X,Y)}&{\cdots}&{\Omega'^{q}_{q}(X,Y)}\\ \end{array}\right)\left( \begin{array}{cccc}{s'_{1}} \\{s'_{2}}\\ {\vdots}  \\ {s'_{q}} \end{array}\right)\\
&=\left(\lambda^{1},\cdots,\lambda^{q} \right)  A^{-1}  A 
\left( \begin{array}{cccc}{\Omega^{1}_{1}(X,Y)} & {\Omega^{2}_{1}(X,Y)}&{\cdots}&{\Omega^{q}_{1}(X,Y)}\\ {\Omega^{1}_{2}(X,Y)} & {\Omega^{2}_{2}(X,Y)}&{\cdots}&{\Omega^{q}_{2}(X,Y)}\\ {\vdots} & {\vdots}&{ }&{\vdots}\\{\Omega^{1}_{q}(X,Y)} & {\Omega^{2}_{q}(X,Y)}&{\cdots}&{\Omega^{q}_{q}(X,Y)}\\ \end{array}\right)
  A^{-1}  A\left( \begin{array}{cccc}{s_{1}} \\{s_{2}}\\ {\vdots}  \\ {s_{q}} \end{array}\right)\\
  &=\left(\lambda^{1},\lambda^{2},\cdots,\lambda^{q} \right)\left( \begin{array}{cccc}{\Omega^{1}_{1}(X,Y)} & {\Omega^{2}_{1}(X,Y)}&{\cdots}&{\Omega^{q}_{1}(X,Y)}\\ {\Omega^{1}_{2}(X,Y)} & {\Omega^{2}_{2}(X,Y)}&{\cdots}&{\Omega^{q}_{2}(X,Y)}\\ {\vdots} & {\vdots}&{ }&{\vdots}\\{\Omega^{1}_{q}(X,Y)} & {\Omega^{2}_{q}(X,Y)}&{\cdots}&{\Omega^{q}_{q}(X,Y)}\\ \end{array}\right)\left( \begin{array}{cccc}{s_{1}} \\{s_{2}}\\ {\vdots}  \\ {s_{q}} \end{array}\right)
\end{split}
\]
这就表明曲率算子$\rr(X,Y)$不依赖于向量丛$E$的局部标架场的选取。

我们把定义中的矩阵形式乘出来，就得到$\rr(X,Y)s=\lambda^{\alpha}\Omega^{\beta}_{\alpha}(X,Y)s_{\beta}$. 其中$\Omega^{\beta}_{\alpha}(X,Y)$是可以明确计算出来的。注意到$\Omega^{\beta}_{\alpha}=\dd \omega^{\beta}_{\alpha}-\omega_{\alpha}^{k}\wedge\omega_{k}^{\beta}$，所以
\[
\begin{split}
\Omega^{\beta}_{\alpha}(X,Y)&=\dd \omega^{\beta}_{\alpha}(X,Y)-\omega_{\alpha}^{k}\wedge\omega_{k}^{\beta}(X,Y)\\
&=X\left(\omega^{\beta}_{\alpha}(Y) \right)-Y\left(\omega^{\beta}_{\alpha}(X) \right)-\omega^{\beta}_{\alpha}\left([X,Y] \right)-\left| \begin{array}{cc}{\omega^{k}_{\alpha}(X)} & {\omega^{k}_{\alpha}(Y)}\\ {\omega^{\beta}_{k}(X)} & {\omega^{\beta}_{k}(Y)}\\ \end{array}\right|
\end{split}
\]

曲率算子有一些基本的性质。
\begin{mythm}{}{chapter2曲率算子的性质}
	\begin{itemize}
		\item[(1)] $\rr(X,Y):\Gamma(E)\longrightarrow \Gamma(E)$是线性映射.
		\item[(2)] $\rr(X,Y)=-\rr(Y,X)$.
		\item[(3)] $\rr(fX,Y)=f \rr(X,Y)=\rr(X,fY)$.
		\item[(4)] $\rr(X,Y)(fs)=f\rr(X,Y)(s)$.
	\end{itemize}
\end{mythm}
\begin{proof}
	(1)根据定义这是显然的。\\
	(2)因为
	\[
	\begin{split}
	\Omega^{\beta}_{\alpha}(fX,Y)&=fX\left(\omega^{\beta}_{\alpha}(Y) \right)-Y\left(\omega^{\beta}_{\alpha}(fX) \right)-\omega^{\beta}_{\alpha}\left([fX,Y] \right)-\left| \begin{array}{cc}{\omega^{k}_{\alpha}(fX)} & {\omega^{k}_{\alpha}(Y)}\\ {\omega^{\beta}_{k}(fX)} & {\omega^{\beta}_{k}(Y)}\\ \end{array}\right|\\
	&=Y\left(\omega^{\beta}_{\alpha}(X) \right)-X\left(\omega^{\beta}_{\alpha}(Y) \right)+\omega^{\beta}_{\alpha}\left([X,Y] \right)+\left| \begin{array}{cc}{\omega^{k}_{\alpha}(X)} & {\omega^{k}_{\alpha}(Y)}\\ {\omega^{\beta}_{k}(X)} & {\omega^{\beta}_{k}(Y)}\\ \end{array}\right|\\
	&=-\Omega^{\beta}_{\alpha}(X,Y)
	\end{split}
	\]
	所以对$\forall s=\lambda^{\alpha}s_{\alpha}\in \Gamma(E)$，$\rr(Y,X)s=\lambda^{\alpha}\Omega^{\beta}_{\alpha}(Y,X)s_{\beta}=-\lambda^{\alpha}\Omega^{\beta}_{\alpha}(X,Y)s_{\beta}$. 由$s$的任意性可知$\rr(X,Y)=-\rr(Y,X)$.\\
	(3)因为
	\[
	\begin{split}
	\Omega^{\beta}_{\alpha}(fX,Y)&=fX\left(\omega^{\beta}_{\alpha}(Y) \right)-Y\left(\omega^{\beta}_{\alpha}(fX) \right)-\omega^{\beta}_{\alpha}\left([fX,Y] \right)-\left| \begin{array}{cc}{\omega^{k}_{\alpha}(fX)} & {\omega^{k}_{\alpha}(Y)}\\ {\omega^{\beta}_{k}(fX)} & {\omega^{\beta}_{k}(Y)}\\ \end{array}\right|\\
	&=f\cdot X\left(\omega^{\beta}_{\alpha}(Y) \right)-f\cdot Y\left(\omega^{\beta}_{\alpha}(X) \right)-f\cdot \omega^{\beta}_{\alpha}\left([X,Y] \right)-f\cdot \left| \begin{array}{cc}{\omega^{k}_{\alpha}(X)} & {\omega^{k}_{\alpha}(Y)}\\ {\omega^{\beta}_{k}(X)} & {\omega^{\beta}_{k}(Y)}\\ \end{array}\right|\\
	&=f\cdot \Omega^{\beta}_{\alpha}(X,Y)
	\end{split}
	\]
	所以对$\forall s=\lambda^{\alpha}s_{\alpha}\in \Gamma(E)$，$\rr(fX,Y)s=\lambda^{\alpha}\Omega^{\beta}_{\alpha}(fX,Y)s_{\beta}=f\cdot \lambda^{\alpha}\Omega^{\beta}_{\alpha}(X,Y)s_{\beta}=f\cdot \rr(X,Y)s=f\rr(X,Y)s$. 由$s$的任意性可知$\rr(fX,Y)=f\cdot \rr(X,Y)$. 
	
	另一方面，$\rr(X,fY)=-\rr(fY,X)=-f\rr(Y,X)=f\rr(X,Y)$.
	
	(4)设$s=\lambda^{\alpha}\cdot s_{\alpha}$，则$fs=(f\lambda^{\alpha})\cdot s_{\alpha}$. 因此
	\[\rr(X,Y)(fs)=(f\lambda^{\alpha})\Omega_{\alpha}^{\beta}(X,Y)s_{\beta}=f(\lambda^{\alpha}\Omega_{\alpha}^{\beta}(X,Y)s_{\beta})=f\rr(X,Y)(s)\]
\end{proof}

下面这个定理是曲率算子的本质特征，很多书上也用它作为曲率算子的定义。
\begin{mythm}{曲率算子的等价定义}{chapter2曲率算子的等价定义}
	设$X,Y\in \Gamma(TM)$. 则曲率算子满足：
	\[
	\rr(X,Y)=D_{X}D_{Y}-D_{Y}D_{X}-D_{[X,Y]}
	\]
	即对任意$s\in \Gamma(E)$，$\rr(X,Y)s=D_{X}\left(D_{Y}s\right)-D_{Y}\left(D_{X}s\right)-D_{[X,Y]}s$.
\end{mythm}

\begin{proof}
	取定向量丛$E$在$U$上的一个局部标架场$\{s_{1},\cdots,s_{q} \}$. 任取$s\in \Gamma(E)$，设$s$在$U$上的局部表达式为$s=\lambda^{\alpha}s_{\alpha}$. 则直接计算可知
	\[
	\begin{split}
	D_{Y}s&=\langle Ds,Y \rangle=\langle D(\lambda^{\alpha}s_{\alpha}),Y \rangle=\langle \dd \lambda^{\alpha}\otimes s_{\alpha}+\lambda^{\alpha}Ds_{\alpha},Y \rangle\\
	&=\dd\lambda^{\alpha}(Y)  s_{\alpha}+\langle \lambda^{\alpha}\omega_{\alpha}^{\beta}\otimes s_{\beta},Y \rangle\\
	&=\dd\lambda^{\alpha}(Y)  s_{\alpha}+ \lambda^{\alpha}\omega_{\alpha}^{\beta}(Y) s_{\beta} \\
	&\xlongequal{\text{统一指标}}\dd\lambda^{\alpha}(Y)  s_{\alpha}+ \lambda^{\beta}\omega_{\beta}^{\alpha}(Y) s_{\alpha} \\
	&\xlongequal{\dd\lambda^{\alpha}(Y) =Y(\lambda^{\alpha})}\left(Y(\lambda^{\alpha}) +\lambda^{\beta}\omega_{\beta}^{\alpha}(Y) \right)s_{\alpha}\\
	&=\underbrace{\left(Y(\lambda^{\alpha}) +\lambda^{\beta}\omega_{\beta}^{\alpha}(Y) \right)}_{=\mu^{\alpha}}s_{\alpha}\\
	&\triangleq\mu^{\alpha}s_{\alpha}
	\end{split}
	\]
	则利用上面的计算过程我们又有
	\[
	\begin{split}
	D_{X}\left(D_{Y}s \right)&=D_{X}\left(\mu^{\alpha}s_{\alpha} \right)=\left(X(\mu^{\alpha}) +\mu^{\beta}\omega_{\beta}^{\alpha}(X) \right)s_{\alpha}\\
	&=\left[X(Y(\lambda^{\alpha})) +X(\lambda^{\alpha}\omega_{\beta}^{\alpha}(Y))+\left(Y(\lambda^{\beta}) +\lambda^{\gamma}\omega_{\gamma}^{\beta}(Y) \right)\omega_{\beta}^{\alpha}(X)\right] s_{\alpha}\\
	&=\left[{\color{green}{X(Y(\lambda^{\alpha}))}}+{\color{red}{\omega_{\beta}^{\alpha}(Y)X(\lambda^{\beta})}} +\lambda^{\beta}X\left(\omega_{\beta}^{\alpha}(Y) \right)+{\color{blue}{Y(\lambda^{\beta})\omega_{\beta}^{\alpha}(X)}}+\lambda^{\gamma}\omega_{\gamma}^{\beta}(Y)\omega_{\beta}^{\alpha}(X)\right]s_{\alpha}
	\end{split}
	\]
	同理可得
	\[
	D_{Y}\left(D_{X}s \right)=\left[{\color{seco}{Y(X(\lambda^{\alpha}))}}+{\color{blue}{\omega_{\beta}^{\alpha}(X)Y(\lambda^{\beta})}} +\lambda^{\beta}Y\left(\omega_{\beta}^{\alpha}(X) \right)+{\color{red}{X(\lambda^{\beta})\omega_{\beta}^{\alpha}(Y)}}+\lambda^{\gamma}\omega_{\gamma}^{\beta}(X)\omega_{\beta}^{\alpha}(Y)\right]s_{\alpha}
	\]
	以及
	\[
	D_{[X,Y]}s=\left([X,Y](\lambda^{\alpha}) +\lambda^{\beta}\omega_{\beta}^{\alpha}([X,Y]) \right)s_{\alpha}=\left({\color{green}{X(Y(\lambda^{\alpha}))}} -{\color{seco}{Y(X(\lambda^{\alpha}))}}+\lambda^{\beta}\omega_{\beta}^{\alpha}([X,Y]) \right)s_{\alpha}
	\]
	因此
    \[
    \begin{split}
    &D_{X}\left(D_{Y}s\right)-D_{Y}\left(D_{X}s\right)-D_{[X,Y]}s\\
    =&[\underbrace{\lambda^{\beta}X\left(\omega_{\beta}^{\alpha}(Y) \right)- \lambda^{\beta}Y\left(\omega_{\beta}^{\alpha}(X) \right) - \lambda^{\beta}\omega_{\beta}^{\alpha}([X,Y])}_{=\lambda^{\beta}\dd \omega_{\beta}^{\alpha}(X,Y)}+\underbrace{\lambda^{\gamma}\omega_{\gamma}^{\beta}(Y)\omega_{\beta}^{\alpha}(X)-\lambda^{\gamma}\omega_{\gamma}^{\beta}(X)\omega_{\beta}^{\alpha}(Y)}_{=-\lambda^{\gamma}\omega_{\gamma}^{\beta}\wedge\omega_{\beta}^{\alpha}(X,Y)} ]s_{\alpha}\\
    =&\left[\lambda^{\beta}\dd \omega_{\beta}^{\alpha}(X,Y)- \lambda^{\gamma}\omega_{\gamma}^{\beta}\wedge\omega_{\beta}^{\alpha}(X,Y) \right]s_{\alpha}\xlongequal{\text{统一指标}}\left[\lambda^{\beta}\dd \omega_{\beta}^{\alpha}(X,Y)- \lambda^{\beta}\omega_{\beta}^{\gamma}\wedge\omega_{\gamma}^{\alpha}(X,Y) \right]s_{\alpha}\\
    =&\lambda^{\beta}\left[\dd \omega_{\beta}^{\alpha}(X,Y)- \omega_{\beta}^{\gamma}\wedge\omega_{\gamma}^{\alpha}(X,Y) \right]s_{\alpha}=\lambda^{\beta}\Omega_{\beta}^{\alpha}(X,Y)s_{\alpha}=\rr(X,Y)s
    \end{split}
    \]
    从而$	\rr(X,Y)=D_{X}D_{Y}-D_{Y}D_{X}-D_{[X,Y]}$.
\end{proof}

不严谨地说，上述定理表明：交换两次“求导”的顺序，它们的差在某种程度上就泄露了该空间本身的曲率信息。利用这个等价定义，可以得到曲率算子的以下线性性质。

\begin{mythm}{}{chapter2曲率算子的双线性}
	\begin{itemize}
		\item[(1)] $\rr(X+Z,Y)=\rr(X,Y)+\rr(Z,Y)$.
		\item[(2)] $\rr(X,Y+Z)=\rr(X,Y)+\rr(X,Z)$.
	\end{itemize}
\end{mythm}
\begin{proof}
	利用\rthm{thm:chapter2曲率算子的等价定义}直接计算：对任意$s\in \Gamma(E)$
	\[
	\begin{split}
	\rr(X+Z,Y)s&=D_{X+Z}\left(D_{Y}s\right)-D_{Y}\left(D_{X+Z}s\right)-D_{[X+Z,Y]}s\\
	&=D_{X}\left(D_{Y}s\right)+D_{Z}\left(D_{Y}s\right)-D_{Y}\left( D_{X}s+D_{Z}s\right)-D_{[X,Y]+[Z,Y]}s\\
	&=D_{X}\left(D_{Y}s \right)+D_{Z}\left(D_{Y}s \right)-D_{Y}\left(D_{X}s \right)-D_{Y}\left(D_{Z}s \right)-D_{[X,Y]}s-D_{[Z,Y]}s\\
	&=\rr(X,Y)s+\rr(Z,Y)s
	\end{split}
	\]
	所以$\rr(X+Z,Y)=\rr(X,Y)+\rr(Z,Y)$.
	\[
	\begin{split}
	\rr(X,Y+Z)s&=D_{X}\left(D_{Y+Z}s\right)-D_{Y+Z}\left(D_{X}s\right)-D_{[X,Y+Z]}s\\
	&=D_{X}\left(D_{Y}s+D_{Z}s\right)-D_{Y}\left(D_{X}s\right)-D_{Z}\left( D_{X}s \right)-D_{[X,Y]+[X,Z]}s\\
	&=D_{X}\left(D_{Y}s \right)+D_{X}\left(D_{Z}s \right)-D_{Y}\left(D_{X}s \right)-D_{Z}\left(D_{X}s \right)-D_{[X,Y]}s-D_{[X,Z]}s\\
	&=\rr(X,Y)s+\rr(X,Z)s
	\end{split}
	\]
	所以$\rr(X,Y+Z)=\rr(X,Y)+\rr(X,Z)$.
\end{proof}


\begin{mythm}{Bianchi恒等式}{chapter2Bianchi恒等式}
	设$\omega$为联络$D$在开集$U$上的的联络矩阵，$\Omega$为相应的曲率矩阵. 则有
	\[
	\dd \Omega=\omega\wedge\Omega-\Omega\wedge\omega
	\]
\end{mythm}
\begin{proof}
	根据定义，$\Omega=\dd \omega-\omega\wedge\omega$. 对式子两边做外微分
	\[
	\dd \Omega=\dd(\dd \omega)-\dd (\omega\wedge\omega)=0-(\dd\omega\wedge\omega-\omega\wedge\dd\omega)=\omega\wedge\dd\omega-\dd\omega\wedge\omega
	\]
	代入$\dd \omega=\Omega+\omega\wedge\omega$继续计算
	\[
	\begin{split}
	\dd \Omega&=\omega\wedge(\Omega+\omega\wedge\omega)-(\Omega+\omega\wedge\omega)\wedge\omega\\
	&=\omega\wedge\Omega+\cancel{\omega\wedge(\omega\wedge\omega)}-\Omega\wedge\omega-\cancel{(\omega\wedge\omega)\wedge\omega}\\
	&=\omega\wedge\Omega-\Omega\wedge\omega
	\end{split}
	\]
\end{proof}

\begin{mydef}{平行截面}{chapter2平行截面}
	设$D$是向量丛$E$上的一个联络，$s\in \Gamma(E)$. 
	若$Ds=0$，则称$s$关于联络$D$是一个平行截面.
\end{mydef}

\subsection{仿射联络}

前面我们讨论的是一般向量丛上的联络，下面我们主要考虑一类特殊的联络：切丛上的联络。

\begin{mydef}{仿射联络}{chapter2仿射联络}
	设$M$是一个光滑流形. 我们把切丛$TM$上的联络$D$称为流形$M$上的仿射联络(affine connection).
\end{mydef}

在$M^{m}$的任意一个开子集$U$上，取定切丛$TM$的局部标架场$\{e_{1},e_{2},\cdots,e_{m} \}$，同时取对偶标架场$\{\theta^{1},\theta^{2},\cdots,\theta^{m} \}$作为$T^{\ast}M$在$U$上的局部标架，即$\theta^{i}(e_{j})=\delta^{i}_{j}$. 下面我们先准备一些必要的数据，再做进一步的讨论。

假设$De_{j}=\omega_{j}^{k}\otimes e_{k}$，其中$\omega_{j}^{k}=\Gamma^{k}_{ji}\theta^{i}$. 即联络矩阵为$\omega=(\omega_{j}^{k})_{m\times m}$. 在仿射联络下，我们还特别称$\Gamma^{k}_{ji}$为联络$D$在局部标架场$\{e_{1},e_{2},\cdots,e_{m} \}$下的\textbf{联络系数}。以上这些只需把前面一般向量丛上的有关数据照搬过来即可。除此以外，我们还需要一个不可或缺的关键数据：假设$[e_{i},e_{j}]=C_{ij}^{k}e_{k}$. 这里$[\;,\;]$表示括号积运算。

\begin{mythm}{}{chapter2仿射联络的系数变换规律}
	仿射联络系数$\{\Gamma^{k}_{ij}\}$不满足$(1,2)$型张量系数的变换规律.
\end{mythm}
\begin{proof}
	设在$U$上切丛$TM$的另一个局部标架场为$\{e^{\prime}_{1},e^{\prime}_{2},\cdots,e^{\prime}_{m} \}$，相应的对偶标架场为$\{\theta^{\prime1},\theta^{\prime2},\cdots,\theta^{\prime m} \}$，联络矩阵为$\omega^{\prime}$，联络系数为$\{\Gamma^{\prime k}_{ij}\}$. 记
	\[
		S=\left( \begin{array}{ccc}{e_{1}} \\ {\vdots}  \\ {e_{m}} \end{array}\right)\quad S^{\prime}=\left( \begin{array}{ccc}{e^{\prime}_{1}} \\ {\vdots}  \\ {e^{\prime}_{m}} \end{array}\right)
	\]
	再设$S^{\prime}=AS$. 其中
	\[
	A=\left( \begin{array}{cccc}{f_{1}^{1}} & {f_{1}^{2}}&{\cdots}&{f_{1}^{m}}\\ {f_{2}^{1}} & {f_{2}^{2}}&{\cdots}&{f_{2}^{m}}\\ {\vdots} & {\vdots}&{ }&{\vdots}\\{f_{m}^{1}} & {f_{m}^{2}}&{\cdots}&{f_{m}^{m}}\\ \end{array}\right)\quad A^{-1}=\left( \begin{array}{cccc}{g_{1}^{1}} & {g_{1}^{2}}&{\cdots}&{g_{1}^{m}}\\ {g_{2}^{1}} & {g_{2}^{2}}&{\cdots}&{g_{2}^{m}}\\ {\vdots} & {\vdots}&{ }&{\vdots}\\{g_{m}^{1}} & {g_{m}^{2}}&{\cdots}&{g_{m}^{m}}\\ \end{array}\right)
	\]
	则根据\rthm{thm:chapter2联络矩阵的变换公式}可知，$\omega^{\prime}=(\dd A)A^{-1}+A\omega A^{-1}$. 即$\omega^{\prime k}_{i}=\dd (f_{i}^{p})g_{p}^{k}+f_{i}^{p}\omega_{p}^{q}g_{q}^{k}$. 由于$f_{i}^{p}$是0-形式，所以$\dd f_{i}^{p}$是1-形式，即$\dd f_{i}^{p}\in \Gamma(T^{\ast}M)$. 而$\{\theta^{\prime1},\theta^{\prime2},\cdots,\theta^{\prime m} \}$是$T^{\ast}M$的局部标架场，所以{\color{blue}{可设$\dd f_{i}^{p}=h_{ij}^{p}\theta^{\prime j}$}}. 因此
	\[
	\omega^{\prime k}_{i}=\dd (f_{i}^{p})g_{p}^{k}+f_{i}^{p}\omega_{p}^{q}g_{q}^{k}
	=(h_{ij}^{p}\theta^{\prime j})g_{p}^{k}+f_{i}^{p}(\Gamma_{pr}^{q}\theta^{r})g_{q}^{k}
	\]
	注意到$S^{\prime}=AS$，从而易知$\left(\theta^{ 1},\theta^{ 2},\cdots,\theta^{  m}\right)=\left(\theta^{\prime1},\theta^{\prime2},\cdots,\theta^{\prime m}\right)A$. 于是有${\color{blue}{\theta^{r}=f^{r}_{j}\theta^{\prime j}}}$. 因此 
	\[
	\omega^{\prime k}_{i}=h_{ij}^{p}g_{p}^{k}\theta^{\prime j}+f_{i}^{p}g_{q}^{k}\Gamma_{pr}^{q}f^{r}_{j}\theta^{\prime j}
	=\underbrace{\left(h_{ij}^{p}g_{p}^{k}+f_{i}^{p}g_{q}^{k}\Gamma_{pr}^{q}f^{r}_{j} \right)}_{=\Gamma^{\prime k}_{ij}}\theta^{\prime j}
	\]
所以我们得到
\begin{equation}
	{\color{red}{\Gamma^{\prime k}_{ij}}}=h_{ij}^{p}g_{p}^{k}+f_{i}^{p}g_{q}^{k}{\color{red}{\Gamma_{pr}^{q}}}f^{r}_{j}
\end{equation}
这不符合$(1,2)$型张量系数的变换规律。
\end{proof}

一般向量丛上联络的曲率矩阵和曲率算子，都可以原封不动地照搬到仿射联络的情境下。但是在仿射联络下，事情会变得更加有趣。

回忆一下曲率方阵的定义，$\Omega=\dd \omega-\omega\wedge \omega$，我们可以算出其中第$i$行第$j$列的元素$\Omega_{i}^{j}=\dd \omega_{i}^{j}-\omega_{i}^{k}\wedge\omega^{j}_{k}$. 这是一个2-形式，我们姑且把它称为\textbf{曲率形式}。

再把曲率算子也搬过来。因为这里我们考虑的是切丛，局部标架场是$\{ e_{1},e_{2},\cdots,e_{m}\}$，所以曲率算子的定义就是：任意取定$X,Y\in \Gamma(TM)$，对$\forall Z=Z^{i}e_{i}\in \Gamma(TM)$
\[
\rr(X,Y)Z=Z^{i}\Omega_{i}^{j}(X,Y)e_{j}
\]

第一个有趣的事情发生于暴力计算$\Omega_{i}^{j}(X,Y)$的时候。在$U$上，假设$X=X^{k}e_{k}$，$Y=Y^{l}e_{l}$，则有
\[
\begin{split}
[X,Y]=[X^{k}e_{k},Y^{l}e_{l}]&=X^{k}\left(e_{k}(Y^{l}) \right)e_{l}-Y^{l}\left(e_{l}(X^{k}) \right)e_{k}+X^{k}Y^{l}[e_{k},e_{l}]\\
&=X^{k}\left(e_{k}(Y^{l}) \right)e_{l}-Y^{l}\left(e_{l}(X^{k}) \right)e_{k}+X^{k}Y^{l}C_{kl}^{t}e_{t}\\
&\xlongequal{\text{统一指标}}{\color{red}{X^{k}}}\left({\color{red}{e_{k}}}(Y^{l}) \right)e_{l}-{\color{red}{Y^{k}}}\left({\color{red}{e_{k}}}(X^{l}) \right)e_{l}+X^{k}Y^{t}C_{kt}^{l}e_{l}\\
&={\color{red}{X}}(Y^{l})e_{l}-{\color{red}{Y}}(X^{l})e_{l}+X^{k}Y^{t}C_{kt}^{l}e_{l}
\end{split}
\]

下面我们就来计算$\Omega_{i}^{j}(X,Y)$. 注意到$\Omega_{i}^{j}(X,Y)=\dd \omega_{i}^{j}(X,Y)-\omega_{i}^{k}\wedge\omega^{j}_{k}(X,Y)$，所以
\[
\begin{split}
&\Omega_{i}^{j}(X,Y)=\dd \omega_{i}^{j}(X,Y)-\omega_{i}^{k}\wedge\omega^{j}_{k}(X,Y)\\
&=X\left(\omega_{i}^{j}(Y) \right)-Y\left(\omega_{i}^{j}(X)\right)-\omega_{i}^{j}\left([X,Y]\right)-\left| \begin{array}{cc}{\omega^{k}_{i}(X)} & {\omega^{k}_{i}(Y)}\\ {\omega^{j}_{k}(X)} & {\omega^{j}_{k}(Y)}\\ \end{array}\right|\\
&=X\left(\Gamma^{j}_{il}\theta^{l}(Y) \right)-Y\left(\Gamma^{j}_{il}\theta^{l}(X)\right)-\Gamma^{j}_{il}\theta^{l}\left([X,Y]\right)-\left[\Gamma^{k}_{ir}\Gamma_{ks}^{j}\theta^{r}(X)\theta^{s}(Y)-\Gamma^{j}_{kr}\Gamma^{k}_{is}\theta^{r}(X)\theta^{s}(Y) \right]\\
&=X\left(\Gamma^{j}_{il}Y^{l} \right)-Y\left(\Gamma^{j}_{il}X^{l}\right)-\Gamma^{j}_{il}\left[X(Y^{l})-Y(X^{l})+X^{k}Y^{t}C_{kt}^{l}\right]  -\left[\Gamma^{k}_{ir}\Gamma_{ks}^{j}X^{r}Y^{s}-\Gamma^{j}_{kr}\Gamma^{k}_{is}X^{r}Y^{s} \right]\\
&= Y^{l}X(\Gamma^{j}_{il})+\cancel{\Gamma^{j}_{il}X(Y^{l})}-X^{l}Y(\Gamma^{j}_{il})-\cancel{\Gamma^{j}_{il}Y(X^{l})}-\cancel{\Gamma^{j}_{il}X(Y^{l})}+\cancel{\Gamma^{j}_{il}Y(X^{l})}-\Gamma^{j}_{il}X^{k}Y^{t}C^{l}_{kt}\\
&\;\;\;\;\;\quad\quad\quad\quad\quad\quad\quad\quad\quad\quad\quad\quad\quad\quad\quad\quad\quad\quad\quad\quad\quad-\Gamma^{k}_{ir}\Gamma_{ks}^{j}X^{r}Y^{s}+\Gamma^{j}_{kr}\Gamma^{k}_{is}X^{r}Y^{s}\\
&=Y^{l}X(\Gamma^{j}_{il}) -X^{l}Y(\Gamma^{j}_{il})   -\Gamma^{j}_{il}X^{k}Y^{t}C^{l}_{kt} -\Gamma^{k}_{ir}\Gamma_{ks}^{j}X^{r}Y^{s}+\Gamma^{j}_{kr}\Gamma^{k}_{is}X^{r}Y^{s}\\
&=Y^{l}X^{r}e_{r}(\Gamma^{j}_{il}) -X^{l}Y^{s}e_{s}(\Gamma^{j}_{il})   -\Gamma^{j}_{il}X^{k}Y^{t}C^{l}_{kt} -\Gamma^{k}_{ir}\Gamma_{ks}^{j}X^{r}Y^{s}+\Gamma^{j}_{kr}\Gamma^{k}_{is}X^{r}Y^{s}\\
 &\xlongequal{\text{统一指标}}Y^{s}X^{r}e_{r}(\Gamma^{j}_{is}) -X^{r}Y^{s}e_{s}(\Gamma^{j}_{ir})   -\Gamma^{j}_{il}X^{r}Y^{s}C^{l}_{rs} -\Gamma^{k}_{ir}\Gamma_{ks}^{j}X^{r}Y^{s}+\Gamma^{j}_{kr}\Gamma^{k}_{is}X^{r}Y^{s}\\
&=X^{r}\left[ e_{r}(\Gamma^{j}_{is})-e_{s}(\Gamma^{j}_{ir})-\Gamma^{j}_{il} C^{l}_{rs}  - \Gamma^{k}_{ir}\Gamma_{ks}^{j}+\Gamma^{j}_{kr}\Gamma^{k}_{is}\right]Y^{s}
\end{split}
\]

这样我们就得到了仿射联络下的一个极其重要的量，令
\begin{equation}
R^{j}_{irs}\triangleq e_{r}(\Gamma^{j}_{is})-e_{s}(\Gamma^{j}_{ir})-\Gamma^{j}_{il} C^{l}_{rs}  - \Gamma^{k}_{ir}\Gamma_{ks}^{j}+\Gamma^{j}_{kr}\Gamma^{k}_{is}
\end{equation}
其中$C_{rs}^{l}=\theta^{l}\left([e_{r},e_{s}]\right) $. 并且$\Omega_{i}^{j}(X,Y)=X^{r}R^{j}_{irs}Y^{s}$.
\begin{mythm}{}{chapter2Rijkl的性质}
	\begin{itemize}
		\item[(1)] $R^{j}_{irs}=-R^{j}_{isr}$.
		\item[(2)] $\rr(e_{r},e_{s})e_{i}=R^{j}_{irs}e_{j}$. 这里$\rr(e_{r},e_{s})e_{i}$是指曲率算子的作用.
		\item[(3)] $\{R^{j}_{irs} \}$满足$(1,3)$型张量系数的变换规律.
	\end{itemize}
\end{mythm}
\begin{proof}
	(1)\;因为$R^{j}_{isr}= e_{s}(\Gamma^{j}_{ir})-e_{r}(\Gamma^{j}_{is})-\Gamma^{j}_{il} C^{l}_{sr}  - \Gamma^{k}_{is}\Gamma_{kr}^{j}+\Gamma^{j}_{ks}\Gamma^{k}_{ir}$，再注意到括号积的性质
	\[
	C_{sr}^{l}=\theta^{l}\left([e_{s},e_{r}]\right) =\theta^{l}\left(-[e_{r},e_{s}]\right)=-C_{rs}^{l}
	\]
	因此就有$R^{j}_{irs}=-R^{j}_{isr}$.
	
	(2)\; 根据曲率算子的定义结合之前的计算可知，$\rr(e_{r},e_{s})e_{i}=\Omega_{i}^{j}(e_{r},e_{s})e_{j}=R^{j}_{irs}e_{j}$. 
	
	(3)\;设在$U$上切丛$TM$的另一个局部标架场为$\{e^{\prime}_{1},e^{\prime}_{2},\cdots,e^{\prime}_{m} \}$，相应的对偶标架场为$\{\theta^{\prime1},\theta^{\prime2},\cdots,\theta^{\prime m} \}$，联络系数为$\{\Gamma^{\prime k}_{ij}\}$. 根据(2)性质可知，$\rr(e^{\prime}_{r},e^{\prime}_{s})e^{\prime}_{i}=R^{\prime j}_{irs}e^{\prime}_{j}$. 记
	\[
	S=\left( \begin{array}{ccc}{e_{1}} \\ {\vdots}  \\ {e_{m}} \end{array}\right)\quad S^{\prime}=\left( \begin{array}{ccc}{e^{\prime}_{1}} \\ {\vdots}  \\ {e^{\prime}_{m}} \end{array}\right)
	\]
	再设$S^{\prime}=AS$. 其中
	\[
	A=\left( \begin{array}{cccc}{f_{1}^{1}} & {f_{1}^{2}}&{\cdots}&{f_{1}^{m}}\\ {f_{2}^{1}} & {f_{2}^{2}}&{\cdots}&{f_{2}^{m}}\\ {\vdots} & {\vdots}&{ }&{\vdots}\\{f_{m}^{1}} & {f_{m}^{2}}&{\cdots}&{f_{m}^{m}}\\ \end{array}\right)\quad A^{-1}=\left( \begin{array}{cccc}{g_{1}^{1}} & {g_{1}^{2}}&{\cdots}&{g_{1}^{m}}\\ {g_{2}^{1}} & {g_{2}^{2}}&{\cdots}&{g_{2}^{m}}\\ {\vdots} & {\vdots}&{ }&{\vdots}\\{g_{m}^{1}} & {g_{m}^{2}}&{\cdots}&{g_{m}^{m}}\\ \end{array}\right)
	\]
	则$e^{\prime}_{r}=f^{k}_{r}e_{k}\;,\;e^{\prime}_{s}=f^{l}_{s}e_{l}\;,\;e^{\prime}_{i}=f^{h}_{i}e_{h}\;,\;e^{\prime}_{j}=f^{p}_{j}e_{p}$. 从而
	\[
	\begin{split}
	 R^{\prime j}_{irs}e^{\prime}_{j}&=\rr(e^{\prime}_{r},e^{\prime}_{s})e^{\prime}_{i}=\rr(f^{k}_{r}e_{k},f^{l}_{s}e_{l})(f^{h}_{i}e_{h})\\
	 &=f^{h}_{i}\left(\rr(f^{k}_{r}e_{k},f^{l}_{s}e_{l}) e_{h} \right)\\
	 &=f^{h}_{i}f^{k}_{r}f^{l}_{s}\left(\rr(e_{k},e_{l}) e_{h} \right)\\
	 &=f^{h}_{i}f^{k}_{r}f^{l}_{s}R_{hkl}^{p}e_{p}
	\end{split}
	\]
	因此，$R^{\prime j}_{irs}e^{\prime}_{j}=f^{h}_{i}f^{k}_{r}f^{l}_{s}R_{hkl}^{p}e_{p}$. 再带入$e^{\prime}_{j}=f^{p}_{j}e_{p}$可得
	$R^{\prime j}_{irs}f^{p}_{j}e_{p}=f^{h}_{i}f^{k}_{r}f^{l}_{s}R_{hkl}^{p}e_{p}$. 两边再同乘$g_{p}^{j}$便得
	\[
	R^{\prime j}_{irs} e_{p}=g_{p}^{j}f^{h}_{i}f^{k}_{r}f^{l}_{s}R_{hkl}^{p}e_{p}
	\]
	于是$R^{\prime j}_{irs} =g_{p}^{j}f^{h}_{i}f^{k}_{r}f^{l}_{s}R_{hkl}^{p}$. 而这正是$(1,3)$型张量分量的变换规律！
\end{proof}

因为$\{R^{j}_{irs} \}$满足$(1,3)$型张量系数的变换规律，所以我们可以以其作为系数构造一个$(1,3)$型张量。不过——在这里我们要对以前所介绍的张量的符号做一些澄清和进一步的说明。

在第2.1.4节我们介绍$(r,s)$型混合张量时，对于张量系数的书写，由于采用了多重指标的记号，所以写得比较简洁。例如对于张量$X\in T^{r}_{s}$，它在基下的第$\gamma$个坐标分量，我们写作为$X_{\gamma}$. 其中$\gamma=(\overline{\gamma_{1}},\cdots,\overline{\gamma_{r }},\underline{\gamma_{1}},\cdots,\underline{\gamma_{s }})\in \Gamma^{r }_{s }$表示多重指标。当然，我们所用的多重指标的记号实际上跟大多数教材里用的记号没有本质区别，因为你可以把$X_{\gamma}$改写成那种常见的形式：
\[
X_{\gamma}=X^{\overline{\gamma_{1}} \;\overline{\gamma_{2}}\;\cdots\; \overline{\gamma_{r }}}_{\underline{\gamma_{1}}\;\underline{\gamma_{2}} \;\cdots \;\underline{\gamma_{s }}}
\]

由于以后我们涉及的张量的指标并不很多很复杂，所以也没必要再用多重指标的记号来记。就比如上面所说的$R^{j}_{irs}$，我们没必要把它重新写成“$R_{\gamma}$，其中$\gamma=(j,i,r,s)$”这种臃肿的样子。要记住，使用或者发明自己的记号的原则是：一切都是为了方便起见。所以在有些情形下，如果使用多重指标的记号来得更方便，那你就用多重指标；如果使用$X^{i_{1}\cdots i_{r}}_{j_{1}\cdots j_{s}}$这种记号更方便，那就用这种记号。

但是形如“$X^{i_{1}\cdots i_{r}}_{j_{1}\cdots j_{s}}$”这种记号也有缺陷，某些情况下甚至可以说很糟糕(terrible notation). 主要原因在于，这种记号不利于指标的升降运算。升降指标的技术在微分几何里大量使用，后面我们会具体讲。所以有必要对“$X^{i_{1}\cdots i_{r}}_{j_{1}\cdots j_{s}}$”这种写法做进一步的改进。改进的方法是：引入空格，强调指标的顺序，从而严格区分定义域顺序不同的张量。我们用一个具体的例子来说明。
\begin{example}[张量指标的新约定]
	\label{张量指标的新约定}
	假设$\tensor{T}{^{ab}_{cde}}\;,\tensor{T}{_{cd}^{a}_{e}^{b}}$与$\tensor{T}{^a_c^{b}_{de}}$都符合$(2,3)$型张量系数的变换规律，但我们认为它们所表示的$(2,3)$型张量不是同一个。
	\soln
	
	设向量空间$V$的基是$\{v_{1},\cdots,v_{n} \}$，$V^{\ast}$中的对偶基是$\{v^{1},\cdots,v^{n} \}$. 我们约定$\tensor{T}{^{ab}_{cde}}$所表示的张量$X\in V\otimes V\otimes V^{\ast} \otimes V^{\ast}\otimes V^{\ast}$，即
	\[X=\tensor{T}{^{ab}_{cde}} v_{a}\otimes v_{b}\otimes v^{c}\otimes v^{d}\otimes v^{d}\]
	约定$\tensor{T}{_{cd}^{a}_{e}^{b}}$所表示的张量$Y\in V^{\ast}\otimes V^{\ast}\otimes V\otimes V^{\ast}\otimes V$，即
	\[
	Y=\tensor{T}{_{cd}^{a}_{e}^{b}} v^{c}\otimes v^{d}\otimes v_{a}\otimes v^{e}\otimes v_{b}
	\]
	而约定$\tensor{T}{^a_c^{b}_{de}}$所表示的张量$Z\in V\otimes V^{\ast}\otimes V\otimes V^{\ast}\otimes V^{\ast}$，即
	\[
	Z=\tensor{T}{^a_c^{b}_{de}}v_{a}\otimes v^{c}\otimes v_{b}\otimes v^{d}\otimes v^{e}
	\]
\end{example}

通过上述例子，读者应该能看出我们所做约定的含义。也就是说，不再笼统地认为$(r,s)$型张量都来自于
$\underbrace{V\otimes \cdots\otimes V}_{r\text{个}}\otimes \underbrace{V^{\ast}\otimes \cdots \otimes V^{\ast}}_{s\text{个}}$这个张量空间，我们约定现在的张量空间是带有顺序的！例如$V\otimes V\otimes V^{\ast}\neq V\otimes V^{\ast}\otimes V$，虽然它们都是$(2,1)$型张量空间。

在新约定下，每一个张量分量的指标正上方或者正下方都有一个空格位置，这就是我们为以后做升降运算而预留的。

现在回到本节内容。我们已经证明了$\{R^{j}_{irs} \}$满足$(1,3)$型张量系数的变换规律，因此它是流形$M$上某个$(1,3)$型张量(场)的坐标分量。根据上面所做的约定，我们知道$(1,3)$型张量空间在带有顺序的意义下是不同的。所以我们要构造以$R^{j}_{irs}$为系数的张量，就要考虑所构造的这个张量所在张量空间的构成顺序。这里我们遵循陈省身《微分几何讲义》一书的习惯，
以$\tensor{R}{_i^j_{rs}}$这个顺序来构造我们想要的$(1,3)$型张量(场). 即重新记
\begin{equation}
\tensor{R}{_i^j_{rs}}\triangleq e_{r}(\Gamma^{j}_{is})-e_{s}(\Gamma^{j}_{ir})-\Gamma^{j}_{il} C^{l}_{rs}  - \Gamma^{k}_{ir}\Gamma_{ks}^{j}+\Gamma^{j}_{kr}\Gamma^{k}_{is}
\end{equation}
如此一来，\rthm{thm:chapter2Rijkl的性质}也就可以重新表述为
\begin{mythm}{}{chapter2Rijkl的性质的另一种表述}
	\begin{itemize}
		\item[(1)] $\tensor{R}{_i^j_{rs}}=-\tensor{R}{_i^j_{sr}}$.
		\item[(2)] $\rr(e_{r},e_{s})e_{i}=\tensor{R}{_i^j_{rs}}e_{j}$. 这里$\rr(e_{r},e_{s})e_{i}$是指曲率算子的作用.
		\item[(3)] $\{\tensor{R}{_i^j_{rs}} \}$满足$(1,3)$型张量系数的变换规律.
	\end{itemize}
\end{mythm}

\begin{remark}
	陈省身的书上在后面做指标升降运算时，默认约定都是关于第二个位置升降。既然我们遵循他的习惯，那么反映在这里，便是以$\tensor{R}{_i^j_{rs}}$作为$(1,3)$型张量的系数。当然，也有很多书里采用关于第一个位置升降指标的习惯，这种约定实际上对应于采用$\tensor{R}{^j_{irs}}$作为张量的系数。这两种约定本质上没有什么不同，无非是构造出的张量的定义域顺序不同而已。
\end{remark}

\begin{mydef}{曲率张量(场)}{chapter2曲率张量}
	在$M^{m}$的任意一个开子集$U$上，取定切丛$TM$的局部标架场$\{e_{1},e_{2},\cdots,e_{m} \}$，同时取对偶标架场$\{\theta^{1},\theta^{2},\cdots,\theta^{m} \}$作为$T^{\ast}M$在$U$上的局部标架. 并设联络系数为$\Gamma^{k}_{ij}$. 令
	\[
	R\triangleq\tensor{R}{_i^j_{rs}}\theta^{i}\otimes e_{j}\otimes \theta^{r}\otimes \theta^{s}
	\]
	则$R\in \Gamma(T^{\ast}M\otimes TM\otimes T^{\ast}M\otimes T^{\ast}M)$，它是流形$M$上的一个$(1,3)$型张量场. 我们称其为仿射联络$D$的曲率张量场，常常就简称为曲率张量.
\end{mydef}

曲率张量作为四重线性函数，与曲率算子有如下关系。
\begin{mythm}{}{chapter2曲率张量与曲率算子的关系}
	对任意$X,Y,Z\in \Gamma(TM)$，均有
	\[\rr(X,Y)Z=R(Z,\quad,X,Y)\]
	其中曲率张量的作用定义为：$R(Z,\quad,X,Y)\triangleq \tensor{R}{_i^j_{rs}}\theta^{i}(Z)\theta^{r}(X)\theta^{s}(Y)e_{j}$.
\end{mythm}

\begin{proof}
	直接在局部标架场下计算。
	一方面，$\rr(X,Y)Z=Z^{i}\Omega_{i}^{j}(X,Y)e_{j}=Z^{i} X^{r}\tensor{R}{_i^j_{rs}}Y^{s} e_{j}$. 另一方面，$R(Z,\quad,X,Y)=\tensor{R}{_i^j_{rs}}\theta^{i}(Z)    \theta^{r}(X)  \theta^{s}(Y)e_{j}=\tensor{R}{_i^j_{rs}} Z^{i}X^{r}Y^{s}e_{j}$. 
	
	因此$\rr(X,Y)Z=R(Z,\quad,X,Y)$.
\end{proof}

\begin{mythm}{}{chapter2曲率作用在局部标架上}
	$\rr(e_{r},e_{s})e_{i}=\tensor{R}{_i^j_{rs}}e_{j}=R(e_{i},\quad,e_{r},e_{s})$.
\end{mythm}
 \begin{proof}
根据以上定理，显然成立。
 \end{proof}

以上讨论了仿射联络框架下的曲率算子和曲率张量的性质。在仿射联络下，我们还有一个一般向量丛上未必有的概念，接下来我们就给出相关概念和性质。

在$M^{m}$的任意一个开子集$U$上，取定切丛$TM$的局部标架场$\{e_{1},e_{2},\cdots,e_{m} \}$，同时取对偶标架场$\{\theta^{1},\theta^{2},\cdots,\theta^{m} \}$作为$T^{\ast}M$在$U$上的局部标架，即$\theta^{i}(e_{j})=\delta^{i}_{j}$. 设联络矩阵为$\omega=(\omega_{j}^{k})_{m\times m}$，联络系数为$\Gamma^{k}_{ji}$. 并且假设$[e_{i},e_{j}]=C_{ij}^{k}e_{k}$.  

\begin{mydef}{挠率形式}{chapter2挠率形式}
	在$M^{m}$的任意一个开子集$U$上，取定切丛$TM$的局部标架场$\{e_{1},e_{2},\cdots,e_{m} \}$，同时取对偶标架场$\{\theta^{1},\theta^{2},\cdots,\theta^{m} \}$作为$T^{\ast}M$在$U$上的局部标架. 设联络矩阵为$\omega=(\omega_{j}^{k})_{m\times m}$. 令
	\[
	\Omega^{\alpha}\triangleq\dd \theta^{\alpha}-\theta^{k}\wedge \omega_{k}^{\alpha}
	\]
	我们称$\Omega^{\alpha}$为仿射联络$D$在局部标架场下的挠率形式.
\end{mydef}

显然挠率形式是一个2-形式。而且从它的定义可以看出，这的确是一个无法定义在一般向量丛上的概念，因为在式子$\dd \theta^{\alpha}-\theta^{\color{red}{k}}\wedge \omega_{\color{red}{k}}^{\alpha}$中，要对$k$求和，这就要求所考虑的向量丛的秩与余切丛$T^{\ast}M$的秩相同，而仿射联络所考虑的向量丛为切丛，显然与余切丛有相同的秩。

类似于有了曲率形式(即曲率矩阵中的元素)就可以定义曲率算子，我们也可以用挠率形式来定义“挠率算子”(这个名词并不是标准术语，是我胡诌的...好像也没在别的书上看到过有这个说法。不过问题不大，在这本笔记范围中，这个概念还是可以做到自洽的)。
\begin{mydef}{挠率算子(非正式的)}{chapter2挠率算子}
	在$M^{m}$的任意一个开子集$U$上，取定切丛$TM$的局部标架场$\{e_{1},e_{2},\cdots,e_{m} \}$，同时取对偶标架场$\{\theta^{1},\theta^{2},\cdots,\theta^{m} \}$作为$T^{\ast}M$在$U$上的局部标架. 定义映射$\T:\Gamma(TM)\times \Gamma(TM)\longrightarrow \Gamma(TM)$为：对$\forall X,Y\in \Gamma(TM)$
	\[
	\T(X,Y)=\Omega^{i}(X,Y)e_{i}
	\]
\end{mydef}

挠率算子有一些基本的性质。
\begin{mythm}{}{chapter2挠率算子的性质}
	\begin{itemize}
		\item[(1)] $\T(X,Y)=-\T(Y,X)$.
		\item[(2)] $\T(fX,Y)=f \T(X,Y)=\T(X,fY)$.
	\end{itemize}
\end{mythm}
\begin{proof}
	(1)\; 根据定义可知，$\T(Y,X)=\Omega^{i}(Y,X)e_{i}= \left(\dd \theta^{i}(Y,X)-\theta^{k}\wedge\omega_{k}^{i}(Y,X)\right)e_{i}$. 注意到
	\[
	\begin{split}
	 \Omega^{i}(Y,X)&= \dd \theta^{i}(Y,X)-\theta^{k}\wedge \omega_{k}^{i}(Y,X) \\
	 &=Y\left(\theta^{i}(X) \right)-X\left(\theta^{i}(Y) \right)-\theta^{i}\left([Y,X] \right)-\left| \begin{array}{cc}{\theta^{k}(Y)} & {\theta^{k}(X)}\\ {\omega_{k}^{i}(Y)} & {\omega_{k}^{i}(X)}\\ \end{array}\right|\\
	 &= -X\left(\theta^{i}(Y) \right)+Y\left(\theta^{i}(X) \right)+\theta^{i}\left([X,Y] \right)+\left| \begin{array}{cc}{\theta^{k}(Y)} & {\theta^{k}(X)}\\ {\omega_{k}^{i}(Y)} & {\omega_{k}^{i}(X)}\\ \end{array}\right|\\
	 &=-\dd \theta^{i}(X,Y)+\theta^{k}\wedge \omega_{k}^{i}(X,Y) \\
	 &=-\Omega^{i}(X,Y)
	\end{split}
	\]
	从而，$\T(Y,X) = \left(\dd \theta^{i}(Y,X)-\theta^{k}\wedge\omega_{k}^{i}\right)e_{i}=-\left(\dd \theta^{i}(X,Y)-\theta^{k}\wedge \omega_{k}^{i}(X,Y) \right)e_{i} =-\T(X,Y)$.
	
	(2)\;因为$\T(fX,Y)=\Omega^{i}(fX,Y)e_{i}= \left(\dd \theta^{i}(fX,Y)-\theta^{k}\wedge\omega_{k}^{i}(fX,Y)\right)e_{i}$，而
	\[
	\begin{split}
	 \Omega^{i}(fX,Y)&=\dd \theta^{i}(fX,Y)-\theta^{k}\wedge\omega_{k}^{i}(fX,Y)\\
	 &=fX\left(\theta^{i}(Y) \right)-Y\left(\theta^{i}(fX) \right)-\theta^{i}\left([fX,Y] \right)-\left| \begin{array}{cc}{\theta^{k}(fX)} & {\theta^{k}(Y)}\\ {\omega_{k}^{i}(fX)} & {\omega_{k}^{i}(Y)}\\ \end{array}\right|\\
	 &= fX\left(\theta^{i}(Y) \right)-fY\left(\theta^{i}(X) \right)-f\theta^{i}\left([X,Y] \right)-f\left| \begin{array}{cc}{\theta^{k}(X)} & {\theta^{k}(Y)}\\ {\omega_{k}^{i}(X)} & {\omega_{k}^{i}(Y)}\\ \end{array}\right|\\
	 &= f \Omega^{i}(X,Y) 
	\end{split}
	\]
	所以$\T(fX,Y) =f\T(X,Y)$. 又根据(1)可知，$\T(X,fY)=-\T(fY,X)=-f\T(Y,X)=f\T(X,Y)$.
\end{proof}

\begin{mythm}{挠率算子的等价定义}{chapter2挠率算子的等价定义}
		设$X,Y\in \Gamma(TM)$. 则挠率算子满足：$\T(X,Y)=D_{X}Y-D_{Y}X-[X,Y]$.
\end{mythm}

\begin{proof}
	 设$X=X^{i}e_{i},Y=Y^{j}e_{j}$. 则$DX=D(X^{i}e_{i})=\dd X^{i}\otimes e_{i}+X^{i}D(e_{i})=\dd X^{i}\otimes e_{i}+X^{i}\omega_{i}^{j}\otimes e_{j}$.
	 于是
	 \[
	 \begin{split}
	 D_{Y}X&=\langle DX,Y \rangle=\langle \dd X^{i}\otimes e_{i}+X^{i}\omega_{i}^{j}\otimes e_{j},Y\rangle\\
	 &=\dd X^{i}(Y)e_{i}+X^{i}\omega_{i}^{j}(Y)e_{j}\\
	 &=Y(X^{i})e_{i}+X^{i}\omega_{i}^{j}(Y)e_{j}\\
	 &\xlongequal{\text{统一指标}}{\color{green}{Y(X^{i})e_{i}}}+X^{j}\omega_{j}^{i}(Y)e_{i}
	 \end{split}
	 \]
	 同理可知，$D_{X}Y={\color{red}{X(Y^{i})e_{i}}}+Y^{j}\omega_{j}^{i}(X)e_{i}$. 另外
	 \[
	 \begin{split}
	 [X,Y]=[X^{i}e_{i},Y^{j}e_{j}]&=X^{i}\left(e_{i}(Y^{j}) \right)e_{j}-Y^{j}\left(e_{j}(X^{i}) \right)e_{i}+X^{i}Y^{j}[e_{i},e_{j}]\\
	 &=X^{i}\left(e_{i}(Y^{j}) \right)e_{j}-Y^{j}\left(e_{j}(X^{i}) \right)e_{i}+X^{i}Y^{j}C_{ij}^{t}e_{t}\\
	 &\xlongequal{\text{统一指标}} X^{j}\left(\ e_{j}(Y^{i}) \right)e_{i}-Y^{j}\left(e_{j}(X^{i}) \right)e_{i}+X^{t}Y^{j}C_{tj}^{i}e_{i}\\
	 &={\color{red}{X(Y^{i})e_{i}}}{\color{green}{-Y(X^{i})e_{i}}}+X^{t}Y^{j}C_{tj}^{i}e_{i}
	 \end{split}
	 \]
	 从而，$D_{X}Y-D_{Y}X-[X,Y]=\left(Y^{j}\omega_{j}^{i}(X)-X^{j}\omega_{j}^{i}(Y)-X^{t}Y^{j}C_{tj}^{i}\right)e_{i}$. 另一方面根据\rdef{def:chapter2挠率算子}，$\T(X,Y)=\Omega^{i}(X,Y)e_{i}$. 下面来计算$\Omega^{i}(X,Y)$.
	 \begin{equation}
	 \begin{split}
	 &\Omega ^{i}(X,Y)=\dd \theta^{i}(X,Y)-\theta^{k}\wedge\omega^{i}_{k}(X,Y)\\
	 &=X\left(\theta^{i}(Y) \right)-Y\left(\theta^{i}(X)\right)-\theta^{i}\left([X,Y]\right)-\left| \begin{array}{cc}{\theta^{k}(X)} & {\theta^{k}(Y)}\\ {\omega^{i}_{k}(X)} & {\omega^{i}_{k}(Y)}\\ \end{array}\right|\\
	 &=\cancel{X\left(Y^{i} \right)}-\cancel{Y\left(X^{i}\right)}-\cancel{X(Y^{i})}+\cancel{Y(X^{i})}-X^{t}Y^{j}C_{tj}^{i}-\left[X^{k}\omega^{i}_{k}(Y)-\omega^{i}_{k}(X)Y^{k} \right]\\
	 &=\omega^{i}_{k}(X)Y^{k}-X^{k}\omega^{i}_{k}(Y)-X^{t}Y^{j}C_{tj}^{i}\\
	 &\xlongequal{\text{统一指标}}\omega^{i}_{j}(X)Y^{j}-X^{j}\omega^{i}_{j}(Y)-X^{t}Y^{j}C_{tj}^{i}
	 \end{split}
	 \end{equation}
	 因此$\T(X,Y)=\left(\omega^{i}_{j}(X)Y^{j}-X^{j}\omega^{i}_{j}(Y)-X^{t}Y^{j}C_{tj}^{i} \right)e_{i}$. 这表明$\T(X,Y)=D_{X}Y-D_{Y}X-[X,Y]$.
\end{proof}



\begin{mythm}{}{chapter2挠率算子的双线性}
	\begin{itemize}
		\item[(1)] $\T(X+Z,Y) =\T(X,Y) +\T(Z,Y)$. 
		\item[(2)] $\T(X,Y+Z)=\T(X,Y)+\T(X,Z)$.
	\end{itemize}
\end{mythm}

\begin{proof}
	利用\rthm{thm:chapter2挠率算子的等价定义}直接计算：对任意$X,Y,Z\in \Gamma(TM)$
	\[
	\begin{split}
	\T(X+Z,Y)&=D_{X+Z}Y-D_{Y}(X+Z)-[X+Z,Y]\\
	&=D_{X}Y+D_{Z}Y-D_{Y}X-D_{Y}Z-[X,Y]-[Z,Y]\\
	&=D_{X}Y-D_{Y}X-[X,Y]+D_{Z}Y-D_{Y}Z-[Z,Y]\\
	&=\T(X,Y)+\T(Z,Y)
	\end{split}
	\]
	所以$\T(X+Z,Y) =\T(X,Y) +\T(Z,Y)$.
	\[
	\begin{split}
	\T(X,Y+Z)&=D_{X}(Y+Z)-D_{Y+Z}X-[X,Y+Z]\\
	&=D_{X}Y+D_{X}Z-D_{Y}X-D_{Z}X-[X,Y]-[X,Z]\\
	&=D_{X}Y-D_{Y}X-[X,Y]+D_{X}Z-D_{Z}X-[X,Z]\\
	&=\T(X,Y)+\T(X,Z)
	\end{split}
	\]
	所以$\T(X,Y+Z)=\T(X,Y)+\T(X,Z)$.
\end{proof}

在式\requ{equ:}中我们得到这样一个计算结果：$\Omega ^{i}(X,Y)=\omega^{i}_{j}(X)Y^{j}-X^{j}\omega^{i}_{j}(Y)-X^{t}Y^{j}C_{tj}^{i}$. 将这个式子进一步计算并调整相应的指标，我们会得到
\begin{equation}
	\begin{split}
		\Omega ^{i}(X,Y)&=\omega^{i}_{j}(X)Y^{j}-X^{j}\omega^{i}_{j}(Y)-X^{t}Y^{j}C_{tj}^{i}\\
		&=\Gamma_{jk}^{i}X^{k}Y^{j}-X^{j}\Gamma_{jk}^{i}Y^{k}-X^{t}Y^{j}C_{tj}^{i}\\
		&\xlongequal{\text{统一指标}}\Gamma_{jk}^{i}X^{k}Y^{j}-X^{k}\Gamma_{kj}^{i}Y^{j}-X^{k}Y^{j}C_{kj}^{i}\\
		&=X^{k}\left(\Gamma_{jk}^{i} -\Gamma_{kj}^{i}-C_{kj}^{i}\right)Y^{j}
	\end{split}
\end{equation}
类似于曲率张量，在这里我们获得了仿射联络下的另一个极其重要的量，令
\begin{equation}
	T_{kj}^{i}\triangleq \Gamma_{jk}^{i} -\Gamma_{kj}^{i}-C_{kj}^{i}
\end{equation}
其中$C_{kj}^{i}=\theta^{i}\left([e_{k},e_{j}] \right)$. 并且$\Omega ^{i}(X,Y)=X^{k}T_{kj}^{i}Y^{j}$.

\begin{mythm}{}{chapter2Tijk的性质}
	\begin{itemize}
		\item[(1)] $T^{i}_{kj}=-T^{i}_{jk}$.
		\item[(2)] $\T(e_{k},e_{j})=T^{i}_{kj}e_{i}$. 这里$\T(e_{k},e_{j})$是指挠率算子的作用.
		\item[(3)] $\{T^{i}_{kj} \}$满足$(1,2)$型张量系数的变换规律.
	\end{itemize}
\end{mythm}
\begin{proof}
	(1)\;因为$T^{i}_{jk}=\Gamma_{kj}^{i} -\Gamma_{jk}^{i}-C_{jk}^{i}$，再注意到括号积的性质
	\[
	C_{jk}^{i}=\theta^{i}\left([e_{j},e_{k}] \right)=\theta^{i}\left(-[e_{k},e_{j}] \right)=-C_{kj}^{i}
	\]
	因此$T^{i}_{kj}=-T^{i}_{jk}$.
	
	(2)\;根据挠率算子的定义和之前的计算可知，$\T(e_{k},e_{j})=\Omega ^{i}(e_{k},e_{j})e_{i}=T^{i}_{kj}e_{i}$.
	
	(3)\;设在$U$上切丛$TM$的另一个局部标架场为$\{e^{\prime}_{1},e^{\prime}_{2},\cdots,e^{\prime}_{m} \}$，相应的对偶标架场为$\{\theta^{\prime1},\theta^{\prime2},\cdots,\theta^{\prime m} \}$，联络系数为$\{\Gamma^{\prime k}_{ij}\}$. 根据(2)性质可知，$\T(e^{\prime}_{k},e^{\prime}_{j})=T^{\prime i}_{kj}e^{\prime}_{i}$. 记
	\[
	S=\left( \begin{array}{ccc}{e_{1}} \\ {\vdots}  \\ {e_{m}} \end{array}\right)\quad S^{\prime}=\left( \begin{array}{ccc}{e^{\prime}_{1}} \\ {\vdots}  \\ {e^{\prime}_{m}} \end{array}\right)
	\]
	再设$S^{\prime}=AS$. 其中
	\[
	A=\left( \begin{array}{cccc}{f_{1}^{1}} & {f_{1}^{2}}&{\cdots}&{f_{1}^{m}}\\ {f_{2}^{1}} & {f_{2}^{2}}&{\cdots}&{f_{2}^{m}}\\ {\vdots} & {\vdots}&{ }&{\vdots}\\{f_{m}^{1}} & {f_{m}^{2}}&{\cdots}&{f_{m}^{m}}\\ \end{array}\right)\quad A^{-1}=\left( \begin{array}{cccc}{g_{1}^{1}} & {g_{1}^{2}}&{\cdots}&{g_{1}^{m}}\\ {g_{2}^{1}} & {g_{2}^{2}}&{\cdots}&{g_{2}^{m}}\\ {\vdots} & {\vdots}&{ }&{\vdots}\\{g_{m}^{1}} & {g_{m}^{2}}&{\cdots}&{g_{m}^{m}}\\ \end{array}\right)
	\]
	则$\;e^{\prime}_{k}=f^{l}_{k}e_{l}\;,\;e^{\prime}_{i}=f^{h}_{i}e_{h}\;,\;e^{\prime}_{j}=f^{p}_{j}e_{p}$. 从而
	\[
	T^{\prime i}_{kj}e^{\prime}_{i}=\T(e^{\prime}_{k},e^{\prime}_{j})
	=\T(f^{l}_{k}e_{l},f^{p}_{j}e_{p}) =f^{l}_{k}f^{p}_{j}\T(e_{l},e_{p})=f^{l}_{k}f^{p}_{j}T^{h}_{lp}e_{h}
	\]
	因此，$T^{\prime i}_{kj}{\color{red}{e^{\prime}_{i}}}=f^{l}_{k}f^{p}_{j}T^{h}_{lp}e_{h}$. 再带入${\color{red}{e^{\prime}_{i}=f^{h}_{i}e_{h}}}$可得
	$T^{\prime i}_{kj}{\color{red}{f^{h}_{i}e_{h}}}=f^{l}_{k}f^{p}_{j}T^{h}_{lp}e_{h}$. 两边再同乘$g_{h}^{i}$便得
	\[
	T^{\prime i}_{kj}e_{h}=g_{h}^{i}f^{l}_{k}f^{p}_{j}T^{h}_{lp}e_{h}
	\]
	于是$T^{\prime i}_{kj} =g_{h}^{i}f^{l}_{k}f^{p}_{j}T^{h}_{lp}$. 而这正是$(1,2)$型张量分量的变换规律！
\end{proof}

既然$\{T^{i}_{kj} \}$满足$(1,2)$型张量系数的变换规律，那么就可以以其作为系数构造一个$(1,2)$型张量。但是根据之前所讲的，\textbf{张量空间的构成顺序是有讲究的}。我们这里约定所构造的$(1,2)$型张量的指标顺序如下：
\begin{equation}
\tensor{T}{^i_{kj}}\triangleq \Gamma_{jk}^{i} -\Gamma_{kj}^{i}-C_{kj}^{i}
\end{equation}
如此一来，\rthm{thm:chapter2Tijk的性质}也就可以重新表述为
\begin{mythm}{}{chapter2Tijk的性质的另一种表述}
	\begin{itemize}
		\item[(1)] $\tensor{T}{^i_{kj}}=-\tensor{T}{^i_{jk}}$.
		\item[(2)] $\T(e_{k},e_{j})=\tensor{T}{^i_{kj}}e_{i}$. 这里$\T(e_{k},e_{j})$是指挠率算子的作用.
		\item[(3)] $\{\tensor{T}{^i_{kj}} \}$满足$(1,2)$型张量系数的变换规律.
	\end{itemize}
\end{mythm}

\begin{mydef}{挠率张量(场)}{chapter2挠率张量}
	在$M^{m}$的任意一个开子集$U$上，取定切丛$TM$的局部标架场$\{e_{1},e_{2},\cdots,e_{m} \}$，同时取对偶标架场$\{\theta^{1},\theta^{2},\cdots,\theta^{m} \}$作为$T^{\ast}M$在$U$上的局部标架. 并设联络系数为$\Gamma^{k}_{ij}$. 令
	\[
	T\triangleq\tensor{T}{^i_{kj}}e_{i}\otimes \theta^{k}\otimes \theta^{j}
	\]
	则$T\in \Gamma(TM\otimes T^{\ast}M\otimes T^{\ast}M)$，它是流形$M$上的一个$(1,2)$型张量场. 我们称其为仿射联络$D$的挠率张量场，常常就简称为挠率张量.
\end{mydef}

挠率张量作为三重线性函数，与挠率算子有如下关系。
\begin{mythm}{}{chapter2挠率算子与挠率张量的关系}
	对任意$X,Y\in \Gamma(TM)$，均有
	\[
	\T(X,Y)=T(\quad,X,Y)
	\]
	其中挠率张量的作用定义为：$T(\quad,X,Y)\triangleq \tensor{T}{^i_{kj}} \theta^{k}(X) \theta^{j}(Y)e_{i}$.
\end{mythm}
\begin{proof}
	一方面，$\T(X,Y)=\Omega ^{i}(X,Y)e_{i}=X^{k}\tensor{T}{^i_{kj}}Y^{j}e_{i}$. 另一方面，
	\[
		T(\quad,X,Y)= \tensor{T}{^i_{kj}} \theta^{k}(X) \theta^{j}(Y)e_{i}=\tensor{T}{^i_{kj}} X^{k} Y^{j}e_{i}
	\]
	因此$\T(X,Y)=T(\quad,X,Y)$.
\end{proof}

\begin{mythm}{}{chapter2挠率在局部标架上的作用}
	$\T(e_{k},e_{j})=\tensor{T}{^i_{kj}}e_{i}=T(\quad,X,Y)$.
\end{mythm}
\begin{proof}
	根据以上定理，显然成立。
\end{proof}

\begin{mythm}{Cartan结构方程}{chapter2Cartan结构方程}
	设$D$是光滑流形$M^{m}$上的仿射联络. 在$M$的任意一个开子集$U$上，取定切丛$TM$的局部标架场$\{e_{1},e_{2},\cdots,e_{m} \}$，同时取对偶标架场$\{\theta^{1},\theta^{2},\cdots,\theta^{m} \}$作为$T^{\ast}M$在$U$上的局部标架. 则有：(1)\;$\Omega_{i}^{j}=\frac{1}{2}\tensor{R}{_i^j_{rs}}\theta^{r}\wedge\theta^{s}$；(2)\;$\Omega^{i}=\frac{1}{2}\tensor{T}{^i_{kj}}\theta^{k}\wedge\theta^{j}$.
\end{mythm}
\begin{proof}
	只要证明对任意$X,Y\in \Gamma(TM)$，方程左右两边作用在$(X,Y)$上的值相等。而这又只需去说明对$\forall p,q\in \{1,2,\cdots,m \}$，方程左右两边在$(e_{p},e_{q})$上的值相等。
	
	(1)\;前面已经计算过$\Omega_{i}^{j}(X,Y)=X^{r}\tensor{R}{_i^j_{rs}}Y^{s}$. 这里直接拿来用便有$\Omega_{i}^{j}(e_{p},e_{q})=\tensor{R}{_i^j_{pq}}$. 另一方面
	\[
	\begin{split}
	\frac{1}{2}\tensor{R}{_i^j_{rs}}\theta^{r}\wedge\theta^{s}(e_{p},e_{q})&=\frac{1}{2}\tensor{R}{_i^j_{rs}}\left| \begin{array}{cc}{\theta^{r}(e_{p})} & {\theta^{r}(e_{q})}\\ {\theta^{s}(e_{p})} & {\theta^{s}(e_{q})}\\ \end{array}\right|=\frac{1}{2}\tensor{R}{_i^j_{rs}}\left| \begin{array}{cc}{\delta_{p}^{r}} & {\delta_{q}^{r}}\\ {\delta_{p}^{s}} & {\delta_{q}^{s}}\\ \end{array}\right|\\
	&=\frac{1}{2}\tensor{R}{_i^j_{rs}}(\delta_{p}^{r}\delta_{q}^{s}-\delta_{p}^{s}\delta_{q}^{r})=\frac{1}{2}\tensor{R}{_i^j_{pq}}-\frac{1}{2}\tensor{R}{_i^j_{qp}}\\
	&=\frac{1}{2}\tensor{R}{_i^j_{pq}}+\frac{1}{2}\tensor{R}{_i^j_{pq}}=\tensor{R}{_i^j_{pq}}
	\end{split}
	\]
	因此，$\Omega_{i}^{j}=\frac{1}{2}\tensor{R}{_i^j_{rs}}\theta^{r}\wedge\theta^{s}$.
	
	(2)\;前面也已经算过$\Omega ^{i}(X,Y)=X^{k}T_{kj}^{i}Y^{j}$. 于是$\Omega ^{i}(e_{p},e_{q})=T_{pq}^{i}$. 另一方面
	\[
	\begin{split}
	\frac{1}{2}\tensor{T}{^i_{kj}}\theta^{k}\wedge\theta^{j}(e_{p},e_{q})&=\frac{1}{2}\tensor{T}{^i_{kj}}\left| \begin{array}{cc}{\theta^{k}(e_{p})} & {\theta^{k}(e_{q})}\\ {\theta^{j}(e_{p})} & {\theta^{j}(e_{q})}\\ \end{array}\right|=\frac{1}{2}\tensor{T}{^i_{kj}}\left| \begin{array}{cc}{\delta_{p}^{k}} & {\delta_{q}^{k}}\\ {\delta_{p}^{j}} & {\delta_{q}^{j}}\\ \end{array}\right|\\
	&=\frac{1}{2}\tensor{T}{^i_{kj}}(\delta_{p}^{k}\delta_{q}^{j}-\delta_{p}^{j}\delta_{q}^{k})=\frac{1}{2}\tensor{T}{^i_{pq}}-\frac{1}{2}\tensor{T}{^i_{qp}}\\
	&=\frac{1}{2}\tensor{T}{^i_{pq}}+\frac{1}{2}\tensor{T}{^i_{pq}}=\tensor{T}{^i_{pq}}
	\end{split}
	\]
	因此，$\Omega^{i}=\frac{1}{2}\tensor{T}{^i_{kj}}\theta^{k}\wedge\theta^{j}$.
\end{proof}

\begin{mydef}{无挠联络}{chapter2无挠联络}
	若仿射联络$D$的挠率张量$T=0$，则称$D$是无挠联络或对称联络(torsion free).
\end{mydef}
无挠联络总是存在的。

\begin{mythm}{第一Bianchi恒等式}{chapter2第一Bianchi恒等式}
	若$D$是无挠仿射联络，则有第一Bianchi恒等式：$\forall i,j,k,l\in\{1,2,\cdots,m \}$
	\[
	\tensor{R}{_i^j_{kl}}+\tensor{R}{_k^j_{li}}+\tensor{R}{_l^j_{ik}}=0
	\]
\end{mythm}
\begin{proof}
	取定切丛$TM$的局部标架场$\{e_{1},e_{2},\cdots,e_{m} \}$，同时取对偶标架场$\{\theta^{1},\theta^{2},\cdots,\theta^{m} \}$作为$T^{\ast}M$的局部标架. 根据\rthm{thm:chapter2Rijkl的性质的另一种表述}(2)我们有
	\[
	\begin{split}
	\rr(e_{k},e_{l})e_{i}&=\tensor{R}{_i^j_{kl}}e_{j}\\
	\rr(e_{l},e_{i})e_{k}&=\tensor{R}{_k^j_{li}}e_{j}\\
	\rr(e_{i},e_{k})e_{l}&=\tensor{R}{_l^j_{ik}}e_{j}
	\end{split}
	\]
	而根据曲率算子的等价定义，又有
	\[
	\begin{split}
	\rr(e_{k},e_{l})e_{i}&=D_{e_{k}}\left(D_{e_{l}} e_{i}\right)-D_{e_{l}}\left(D_{e_{k}} e_{i}\right)-D_{[e_{k},e_{l}]}e_{i}\\
	\rr(e_{l},e_{i})e_{k}&=D_{e_{l}}\left(D_{e_{i}} e_{k}\right)-D_{e_{i}}\left(D_{e_{l}} e_{k}\right)-D_{[e_{l},e_{i}]}e_{k}\\
	\rr(e_{i},e_{k})e_{l}&=D_{e_{i}}\left(D_{e_{k}} e_{l}\right)-D_{e_{k}}\left(D_{e_{i}} e_{l}\right)-D_{[e_{i},e_{k}]}e_{l}
	\end{split}
	\]
	所以
	\begin{equation}
	\begin{split}
	\rr(e_{k},e_{l})e_{i}+\rr(e_{l},e_{i})e_{k}+\rr(e_{i},e_{k})e_{l}
	&=D_{e_{k}}\left({\color{red}{D_{e_{l}} e_{i}-D_{e_{i}} e_{l}}}\right)-D_{[e_{l},e_{i}]}e_{k}\\
	&+D_{e_{l}}\left({\color{red}{D_{e_{i}} e_{k}-D_{e_{k}} e_{i}}}\right)-D_{[e_{i},e_{k}]}e_{l}\\
	&+D_{e_{i}}\left({\color{red}{D_{e_{k}} e_{l}-D_{e_{l}} e_{k}}}\right)-D_{[e_{k},e_{l}]}e_{i}
	\end{split}
	\end{equation}
	
	由于$D$是无挠联络，所以挠率张量为0，即有
	\[
	\begin{split}
	{\color{red}{D_{e_{k}}e_{l}-D_{e_{l}}e_{k}}}&=[e_{k},e_{l}]\\
	{\color{red}{D_{e_{i}}e_{k}-D_{e_{k}}e_{i}}}&=[e_{i},e_{k}]\\
	{\color{red}{D_{e_{l}}e_{i}-D_{e_{i}}e_{l}}}&=[e_{l},e_{i}]\\
	\end{split}
	\]
	带入式\requ{equ:}，得到
	\[
	\begin{split}
	\rr(e_{k},e_{l})e_{i}+\rr(e_{l},e_{i})e_{k}+\rr(e_{i},e_{k})e_{l}
	&=D_{e_{k}}\left([e_{l},e_{i}]\right)-D_{[e_{l},e_{i}]}e_{k}\\
	&+D_{e_{l}}\left([e_{i},e_{k}]\right)-D_{[e_{i},e_{k}]}e_{l}\\
	&+D_{e_{i}}\left([e_{k},e_{l}]\right)-D_{[e_{k},e_{l}]}e_{i}
	\end{split}
	\]
	然后再用一次无挠的性质
	\[
	\rr(e_{k},e_{l})e_{i}+\rr(e_{l},e_{i})e_{k}+\rr(e_{i},e_{k})e_{l}=[e_{k},[e_{l},e_{i}]]+[e_{l},[e_{i},e_{k}]]+[e_{i},[e_{k},e_{l}]]=0	
	\]
	最后一个等号是因为\textbf{括号积满足Jacobi恒等式}。因此$\left(\tensor{R}{_i^j_{kl}}+\tensor{R}{_k^j_{li}}+\tensor{R}{_l^j_{ik}}\right)e_{j}=0$
	由于$\{e_{1},e_{2},\cdots,e_{m}\}$是局部标架场，从而在局部处处线性无关，因此便有
	\[
	\tensor{R}{_i^j_{kl}}+\tensor{R}{_k^j_{li}}+\tensor{R}{_l^j_{ik}}=0
	\]
\end{proof}

最后要说明的是，以上所有这些都是在一般的局部标架场情形下考虑和计算的。而流形天生拥有局部坐标，所以实际上我们可以在选取局部坐标卡$\cx$后，把切丛$TM$的局部标架场直接取作自然标架场$\left\{  \frac { \partial } { \partial x _ { 1 } } ,\frac { \partial } { \partial x _ { 2 } }, \cdots , \frac { \partial } { \partial x _ { m } }  \right\}$，在余切丛$T^{\ast}M$上相应的取$\left\{\dd x_{1},\dd x_{2},\cdots,\dd x_{m} \right\}$作为局部标架场。

我们之前的所有讨论基于的是一般标架场，当然也适用于这一特殊的标架场。而且很多推导和计算在自然标架场下面进行其实更简单，就比如说$[\frac { \partial } { \partial x _ { i } } ,\frac { \partial } { \partial x _ { j } }]=0,\forall i,j$. 读者可以自己试着用自然标架场来把联络这一节的内容全部复现一遍，只需替换关键数据：
\begin{itemize}
	\item $\{e_{1},e_{2},\cdots,e_{m} \}=\left\{  \frac { \partial } { \partial x _ { 1 } } ,\frac { \partial } { \partial x _ { 2 } }, \cdots , \frac { \partial } { \partial x _ { m } }  \right\}$
	\item $\{\theta^{1},\theta^{2},\cdots,\theta^{m} \}=\left\{\dd x_{1},\dd x_{2},\cdots,\dd x_{m} \right\}$
	\item $[\frac { \partial } { \partial x _ { i } } ,\frac { \partial } { \partial x _ { j } }]=0,\forall i,j\in \{1,2,\cdots,m\}$
\end{itemize}

\subsection{仿射联络的一些计算}

这一节我们主要给出计算一般$(r,s)$型张量场的绝对微分的方法，以及由交换两次求导顺序产生的Ricci恒等式。

不过在此之前，我们先来考虑如何利用已知的联络定义对偶丛、直和丛与张量积丛上的联络？这是在计算$(r,s)$型张量场的绝对微分时不可避免会遇到的问题。注意，以下\textbf{我们将使用同一个符号$D$来表示不同的联络}，读者\textbf{只需看每个$D$作用在什么元素上，就可以看出它具体表示的是哪个向量丛上的联络}，所以这并不会产生歧义。

\begin{mydef}{对偶丛上的诱导联络}{chapter2对偶对偶丛上的诱导联络}
	设$M$是一个光滑流形，$(E,D)$是$M$上的一个带有联络$D$的实向量丛. $E^{\ast}$为$E$的对偶丛. 若映射$D:\Gamma(E^{\ast})\longrightarrow \Gamma(T^{\ast}M\otimes E^{\ast})$满足：
	\begin{itemize}
		\item[(1)] $D(s^{\ast}_{1}+s^{\ast}_{2})=Ds^{\ast}_{1}+Ds^{\ast}_{2},\forall s^{\ast}_{1},s^{\ast}_{2}\in \Gamma(E^{\ast})$.
		\item[(2)] $D(f s^{\ast})=\dd f\otimes s^{\ast}+f\cdot Ds^{\ast},\forall s^{\ast}\in \Gamma(E^{\ast}),\forall f\in C^{\infty}(M)$.
		\item[(3)] 对$\forall s\in \Gamma(E),\forall s^{\ast}\in \Gamma(E^{\ast})$，有${\color{red}{\langle D s^{\ast}\;,s \rangle=\dd \langle s^{\ast}\;,s \rangle-\langle s^{\ast} \;,Ds \rangle}}$. 
	\end{itemize}
	则称$D:\Gamma(E^{\ast})\longrightarrow \Gamma(T^{\ast}M\otimes E^{\ast})$为向量丛$E^{\ast}$上的由$(E,D)$所诱导的联络.
\end{mydef}

设$\{s_{1},\cdots,s_{q} \}$是向量丛$E$的一个局部标架场，若记$\{s^{\ast 1},\cdots,s^{\ast q} \}$为相应的对偶标架场，即满足$s^{\ast \beta }(s_{\alpha})=\delta_{\alpha}^{\beta}$. 我们来看对偶丛上的诱导联络在对偶标架场下的联络矩阵长什么样子？
设$\omega=(\omega_{\alpha}^{\beta})$为$E$上的联络$D$在$\{s_{1},\cdots,s_{q} \}$下的联络矩阵，$\omega^{\ast}=(\omega^{\ast}_{ij})$为$E^{\ast}$上的诱导联络$D$在$\{s^{\ast 1},\cdots,s^{\ast q} \}$下的联络矩阵。要说明的是，这里我们把矩阵$\omega^{\ast}$里第$i$行第$j$列的元素暂时先记作$\omega^{\ast}_{ij}$这个样子，等稍后确定了$\omega^{\ast}$和$\omega$的关系，再像$\omega$中的元素那样，把$\omega^{\ast}_{ij}$用上下指标来记。所以符号上暂时会显得有些怪异。

则我们有$Ds_{\alpha}=\omega_{\alpha}^{i }\otimes s_{i}$以及$Ds^{\ast \beta }=\omega^{\ast}_{\beta j}\otimes s^{\ast j}$. 根据诱导联络定义的第(3)个条件，我们又有
\[
\dd \langle s^{\ast\beta}\;,s_{\alpha} \rangle=\langle Ds^{\ast\beta}\;,s_{\alpha} \rangle+\langle s^{\ast\beta} \;,Ds_{\alpha} \rangle
\]
注意到$\langle s^{\ast\beta}\;,s_{\alpha} \rangle=s^{\ast \beta }(s_{\alpha})=\delta_{\alpha}^{\beta}$，所以$\dd \langle s^{\ast\beta}\;,s_{\alpha} \rangle=0$. 又因为
\[
\begin{split}
&\langle Ds^{\ast\beta}\;,s_{\alpha} \rangle=\langle \omega^{\ast}_{\beta j}\otimes s^{\ast j}\;,s_{\alpha} \rangle=\omega^{\ast}_{\beta j} s^{\ast j}(s_{\alpha})=\omega^{\ast}_{\beta j}\delta^{j}_{\alpha}=\omega^{\ast}_{\beta \alpha}\\
&\langle s^{\ast\beta} \;,Ds_{\alpha} \rangle=\langle s^{\ast\beta} \;,\omega_{\alpha}^{i }\otimes s_{i} \rangle=\omega_{\alpha}^{i } s^{\ast\beta}(s_{i})=\omega_{\alpha}^{i }\delta^{\beta}_{i}=\omega_{\alpha}^{\beta }
\end{split}
\]
所以
\[
0=\omega^{\ast}_{\beta \alpha}+\omega_{\alpha}^{\beta }
\]
上式表明：矩阵$\omega^{\ast}$的第$\beta$行第$\alpha$列的元素与矩阵$\omega$的第$\alpha$行第$\beta$列的元素之和等于0. 这也就是说$0=\omega^{\ast}+\omega^{T}$，从而
\[
\omega^{\ast}=-\omega^{T}
\]
因此，对偶丛上的诱导联络在对偶标架场下的联络矩阵$\omega^{\ast}$恰好是$E$上的联络在标架场下的联络矩阵$\omega$的负转置。正是因为存在这样的关系，我们把$\omega^{\ast}$中第$i$行第$j$列的元素$\omega^{\ast}_{ij}$重新记作$\omega^{\ast i}_{j}$，即用上指标表示行，下指标表示列；这与$\omega$中的元素用上指标表示列，下指标表示行的约定恰好相反，以此体现$\omega^{\ast}$与$\omega$之间的这种转置关系。
\[
\omega^{\ast}=\left( \begin{array}{cccc}{-\omega^{1}_{1}} & {-\omega^{1}_{2}}&{\cdots}&{-\omega^{1}_{q}}\\ {-\omega^{2}_{1}} & {-\omega^{2}_{2}}&{\cdots}&{-\omega^{2}_{q}}\\ {\vdots} & {\vdots}&{ }&{\vdots}\\{-\omega^{q}_{1}} & {-\omega^{q}_{2}}&{\cdots}&{-\omega^{q}_{q}}\\ \end{array}\right)
\]

\begin{mydef}{直和丛上的诱导联络}{chapter2直和丛上的诱导联络}
	设$M$是一个光滑流形，$(E_{1},D)$与$(E_{2},D)$是$M$上的两个分别带有联络$D$(记成同一个符号)的实向量丛. 若映射$D:\Gamma(E_{1}\oplus E_{2})\longrightarrow \Gamma(T^{\ast}M\otimes (E_{1}\oplus E_{2}))$满足：
	\begin{itemize}
		\item[(1)] $D(s +s^{\prime})=Ds  +Ds^{\prime},\forall s ,s^{\prime}\in \Gamma(E_{1}\oplus E_{2})$.
		\item[(2)] $D(f s )=\dd f\otimes s +f\cdot Ds ,\forall s \in \Gamma(E_{1}\oplus E_{2}),\forall f\in C^{\infty}(M)$.
		\item[(3)] 对$\forall s_{1}\in \Gamma(E_{1}),\forall s_{2}\in \Gamma(E_{2})$，则$s_{1}\oplus s_{2}\in \Gamma(E_{1}\oplus E_{2})$，有$	{\color{red}{D(s_{1}\oplus s_{2})= Ds_{1}\oplus Ds_{2}}}$.
	\end{itemize}
	则称$D:\Gamma(E_{1}\oplus E_{2})\longrightarrow \Gamma(T^{\ast}M\otimes (E_{1}\oplus E_{2}))$为向量丛$E_{1}\oplus E_{2}$上的由$(E_{1},D)$和$(E_{2},D)$所诱导的联络.
\end{mydef}

\begin{mydef}{张量积丛上的诱导联络}{chapter2张量积丛上的诱导联络}
	设$M$是一个光滑流形，$(E_{1},D)$与$(E_{2},D)$是$M$上的两个分别带有联络$D$(记成同一个符号)的实向量丛. 若映射$D:\Gamma(E_{1}\otimes E_{2})\longrightarrow \Gamma(T^{\ast}M\otimes (E_{1}\otimes E_{2}))$满足：
	\begin{itemize}
		\item[(1)] $D(s +s^{\prime})=Ds  +Ds^{\prime},\forall s ,s^{\prime}\in \Gamma(E_{1}\otimes E_{2})$.
		\item[(2)] $D(f s )=\dd f\otimes s +f\cdot Ds ,\forall s \in \Gamma(E_{1}\otimes E_{2}),\forall f\in C^{\infty}(M)$.
		\item[(3)] 对$\forall s_{1}\in \Gamma(E_{1}),\forall s_{2}\in \Gamma(E_{2})$，则$s_{1}\otimes s_{2}\in \Gamma(E_{1}\otimes E_{2})$，有
				\[
				{\color{red}{D(s_{1}\otimes s_{2})= Ds_{1}\otimes s_{2}+ s_{1}\otimes Ds_{2}}}
				\] 
	\end{itemize}
	则称$D:\Gamma(E_{1}\otimes E_{2})\longrightarrow \Gamma(T^{\ast}M\otimes (E_{1}\otimes E_{2}))$为向量丛$E_{1}\otimes E_{2}$上的由$(E_{1},D)$和$(E_{2},D)$所诱导的联络.
\end{mydef}

\begin{remark}
	关于张量积丛上的诱导联络，我们要额外做一个约定：在使用
	\[
	D(s_{1}\otimes s_{2})\triangleq Ds_{1}\otimes s_{2}+ s_{1}\otimes Ds_{2}
	\]
	这一性质的时候，要把其中的$s_{1}\otimes Ds_{2}$的1-形式部分提出来放到最左边，从而使得计算结果在形式上与定义$D:\Gamma(E_{1}\otimes E_{2})\longrightarrow \Gamma(T^{\ast}M\otimes (E_{1}\otimes E_{2}))$保持一致。
\end{remark}

举个例子，假设$\{s^{1}_{1},\cdots,s^{1}_{q} \}$是向量丛$E_{1}$的在开集$U$上的局部标架场，$\{s^{2}_{1},\cdots,s^{2}_{r} \}$是向量丛$E_{2}$在$U$上的局部标架场，那么容易看出$\{s^{1}_{i}\otimes s^{2}_{j}|1\leqslant i\leqslant q,1\leqslant j\leqslant r \}$是$E_{1}\otimes E_{2}$在$U$上的局部标架场。再分别设联络矩阵为$\overline{\omega}$和$\tilde{  \omega } $. 我们来计算$D(s^{1}_{i}\otimes s^{2}_{j})$.
\[
\begin{split}
D(s^{1}_{i}\otimes s^{2}_{j})&=Ds^{1}_{i}\otimes s^{2}_{j}+s^{1}_{i}\otimes Ds^{2}_{j}\\
&=\left(\overline{\omega}^{k}_{i}\otimes s^{1}_{k} \right)\otimes s^{2}_{j}+s^{1}_{i}\otimes \left({\color{red}{\tilde{  \omega } ^{l}_{j}}}\otimes s^{2}_{l} \right)\\
&\xlongequal{\text{约定}}\left(\overline{\omega}^{k}_{i}\otimes s^{1}_{k} \right)\otimes s^{2}_{j}+{\color{red}{\tilde{  \omega } ^{l}_{j}}}\otimes s^{1}_{i}\otimes s^{2}_{l} \\
&=\overline{\omega}^{k}_{i}\otimes \left(s^{1}_{k}  \otimes s^{2}_{j}\right)+\tilde{  \omega } ^{l}_{j}\otimes \left(s^{1}_{i}\otimes s^{2}_{l}\right) \\
\end{split}
\]
红色部分就是计算$s^{1}_{i}\otimes Ds^{2}_{j}$时所产生的$1$-形式部分。按照约定，我们把它放到最左边。这样计算的结果就跟定义完全吻合了，即$D(s^{1}_{i}\otimes s^{2}_{j})\in \Gamma(T^{\ast}M\otimes(E_{1}\otimes E_{2} ))$. 这个约定在下面计算绝对微分时要用到。

\begin{mydef}{拉回丛上的诱导联络}{chapter2拉回丛上的诱导联络}
	设$M,N$为两个光滑流形，$f:M\longrightarrow N$为光滑映射. $(E,\nabla)$是$N$上一个带有联络$\nabla$的实向量丛. $f^{\ast}E$为$E$的拉回丛. 若映射$D:\Gamma(f^{\ast}E)\longrightarrow \Gamma(T^{\ast}M\otimes f^{\ast}E)$满足：
	\begin{itemize}
		\item[(1)] $D(\xi+\eta)=D\xi+D\eta,\forall \xi,\eta\in \Gamma(f^{\ast}E)$.
		\item[(2)] $D(g\xi)=\dd g\otimes \xi+g\cdot D\xi,\forall \xi\in \Gamma(f^{\ast}E),\forall g\in C^{\infty}(M)$.
		\item[(3)] 对$\forall s\in \Gamma(E),\forall X\in \Gamma(TM)$，则$f^{\ast}s\in \Gamma(f^{\ast}E)$，有
			\[
				{\color{red}{\langle D(f^{\ast}s)\;,X\rangle=\langle f^{\ast}\left( \nabla s \right) \;,f_{\ast}(X) \rangle}}
			\] 
	\end{itemize}
	则称$D:\Gamma(f^{\ast}E)\longrightarrow \Gamma(T^{\ast}M\otimes f^{\ast}E)$为拉回丛$f^{\ast}E$上的由$(E,\nabla)$所诱导的联络.
\end{mydef}

有了诱导联络的概念，我们就可以把仿射联络从切丛推广到一般的张量丛上：首先利用\rdef{def:chapter2对偶对偶丛上的诱导联络}，把切丛$TM$上的联络(即仿射联络)诱导至余切丛$T^{\ast}M$上；其次注意到$(r,s)$型张量丛$T^{r}_{s}$是$TM$与$T^{\ast}M$做张量积得到的；最后利用\rdef{def:chapter2张量积丛上的诱导联络}，得到$T^{r}_{s}$上的诱导联络，这个诱导联络就称为\textbf{张量丛$T^{r}_{s}$上的仿射联络}。

下面我们来计算张量场在仿射联络下的绝对微分。先看一些简单的情形。
\begin{example}[函数的绝对微分]
	\label{函数的绝对微分}
     显然对$\forall f\in C^{\infty}(M)$，$Df=\dd f$. 因此光滑函数的绝对微分就是它的普通微分(这是函数外微分$\dd f$的定义)。而且如果把函数看成$(0,0)$型张量场，则$Df=\dd f\in \Gamma(T^{\ast}M)$就是$(0,1)$型张量场.
\end{example}

\begin{example}[$(1,0)$型]
	\label{切向量场的绝对微分}
	设$D:\Gamma(TM)\rightarrow \Gamma(T^{\ast}M\otimes TM)$为仿射联络. 对任意$Y\in \Gamma(TM)$，设$Y$在局部坐标系$\cx$下的局部表示为$Y=\tensor{Y}{^i}\frac { \partial } { \partial x _ { i } }$. 计算绝对微分$DY$.
	\soln
	
	\[
	\begin{split}
	DY&=D\left( Y^{i}\frac { \partial } { \partial x _ { i } }\right)=\dd Y^{i}\otimes \frac { \partial } { \partial x _ { i } }+Y^{i}D\left(\frac { \partial } { \partial x _ { i } }\right) \\
	&=\dd Y^{i}\otimes \frac { \partial } { \partial x _ { i } }+Y^{i}\omega_{i}^{k}\otimes\frac { \partial } { \partial x _ { k } }\\
	&\xlongequal{\text{统一指标}}\dd Y^{i}\otimes \frac { \partial } { \partial x _ { i } }+Y^{k}\omega_{k}^{i}\otimes\frac { \partial } { \partial x _ { i } }\\
	&=\left(\dd Y^{i}+Y^{k}\omega_{k}^{i}\right)\otimes\frac { \partial } { \partial x _ { i } }\\
	&=\left(\frac { \partial Y^{i}} { \partial x _ { j } }\dd x_{j}+Y^{k}\Gamma_{kj}^{i}\dd x_{j}\right)\otimes\frac { \partial } { \partial x _ { i } }\\
	&=\underbrace{\left(\frac { \partial Y^{i}} { \partial x _ { j } }+Y^{k}\Gamma_{kj}^{i}\right)}_{\left(\tensor{Y}{^i}\right)^{\prime}_{ j}}\dd x_{j}\otimes\frac { \partial } { \partial x _ { i } }\\
	&=\left(\tensor{Y}{^i}\right)^{\prime}_{ j}{\color{red}{\dd x_{j}}}\otimes\frac { \partial } { \partial x _ { i } }
	\end{split}
	\]
	
可见$(1,0)$型张量场(即切向量场)的绝对微分是$(1,1)$型张量场，并且$DY\in \Gamma({\color{red}{T^{\ast}M}}\otimes TM)$. 
\end{example}

\begin{example}[$(0,1$型]
	\label{余切向量场的绝对微分}
	设$D:\Gamma(T^{\ast}M)\rightarrow \Gamma(T^{\ast}M\otimes T^{\ast}M)$为仿射联络. 对任意$\alpha\in \Gamma(T^{\ast}M)$，设$\alpha$在局部坐标系$\cx$下的局部表示为$\alpha=\tensor{\alpha}{_i}\dd x^{i}$. 计算绝对微分$D\alpha$.
	\soln
	
	\[
	\begin{split}
	D\alpha&=D\left(\tensor{\alpha}{_i}\dd x^{i} \right)=\dd \tensor{\alpha}{_i}\otimes \dd x^{i}+\tensor{\alpha}{_i}D\left(\dd x^{i}\right) \\
	&\xlongequal{\text{诱导联络}}\dd \tensor{\alpha}{_i}\otimes \dd x^{i}+\tensor{\alpha}{_i}({\color{blue}{-\omega_{k}^{i}}})\otimes\dd x^{k}\\
	&\xlongequal{\text{统一指标}}\dd \tensor{\alpha}{_i}\otimes \dd x^{i}{\color{blue}{-}}\tensor{\alpha}{_k}{\color{blue}{\omega_{i}^{k}}}\otimes\dd x^{i}\\
	&=\left(\dd \tensor{\alpha}{_i} {\color{blue}{-}}\tensor{\alpha}{_k}{\color{blue}{\omega_{i}^{k}}}\right)\otimes\dd x^{i}\\
	&=\left(\frac { \partial \tensor{\alpha}{_i}} { \partial x ^{ j } }\dd x^{j} {\color{blue}{-}}\tensor{\alpha}{_k}{\color{blue}{\Gamma_{ij}^{k}\dd x^{j}}}\right)\otimes\dd x^{i}\\
	&=\underbrace{\left(\frac { \partial \tensor{\alpha}{_i}} { \partial x ^{ j } } -\tensor{\alpha}{_k}\Gamma_{ij}^{k}\right)}_{\left(\tensor{\alpha}{_{i}}\right)^{\prime}_{ j}}\dd x^{j}\otimes\dd x^{i}\\
	&=\left(\tensor{\alpha}{_{i}}\right)^{\prime}_{ j}{\color{red}{\dd x^{j}}}\otimes\dd x^{i}
	\end{split}
	\]
	
	可见$(0,1)$型张量场的绝对微分是$(0,2)$型张量场，并且$D\alpha\in \Gamma({\color{red}{T^{\ast}M}}\otimes T^{\ast}M)$.
\end{example}

下面我们把问题一般化：设$D:\Gamma(T^{r}_{s})\longrightarrow \Gamma(T^{\ast}M\otimes T^{r}_{s})$为张量丛$T^{r}_{s}$上的(诱导)仿射联络. 任取$t\in \Gamma(T^{r}_{s})$，设$t$在局部坐标系$\cx$下的局部表示为
		\[
		t=\tensor{t}{^{i_{1}i_{2}\cdots i_{r}}_{j_{1}j_{2}\cdots j_{s}}}\frac { \partial } { \partial x ^ { i_{1} } }\otimes\frac { \partial } { \partial x ^ { i_{2} } }\otimes\cdots\otimes\frac { \partial } { \partial x ^ { i_{r} } }\otimes\dd x^{j_{1}}\otimes\dd x^{j_{2}}\otimes\cdots\otimes\dd x^{j_{s}}
		\]
我们来计算绝对微分$Dt$. 

首先根据联络的定义得到
\[
\begin{split}
Dt&=D\left(\tensor{t}{^{i_{1}i_{2}\cdots i_{r}}_{j_{1}j_{2}\cdots j_{s}}}\frac { \partial } { \partial x ^ { i_{1} } }\otimes\frac { \partial } { \partial x ^ { i_{2} } }\otimes\cdots\otimes\frac { \partial } { \partial x ^ { i_{r} } }\otimes\dd x^{j_{1}}\otimes\dd x^{j_{2}}\otimes\cdots\otimes\dd x^{j_{s}} \right)\\
&=\dd \left(\tensor{t}{^{i_{1}i_{2}\cdots i_{r}}_{j_{1}j_{2}\cdots j_{s}}}\right) \otimes \frac { \partial } { \partial x ^ { i_{1} } }\otimes\frac { \partial } { \partial x ^ { i_{2} } }\otimes\cdots\otimes\frac { \partial } { \partial x ^ { i_{r} } }\otimes\dd x^{j_{1}}\otimes\dd x^{j_{2}}\otimes\cdots\otimes\dd x^{j_{s}} \\
&\quad+\tensor{t}{^{i_{1}i_{2}\cdots i_{r}}_{j_{1}j_{2}\cdots j_{s}}}D\left(\frac { \partial } { \partial x ^ { i_{1} } }\otimes\frac { \partial } { \partial x ^ { i_{2} } }\otimes\cdots\otimes\frac { \partial } { \partial x ^ { i_{r} } }\otimes \dd x^{j_{1}}\otimes\dd x^{j_{2}}\otimes\cdots\otimes \dd x^{j_{s}}\right)
\end{split}
\]

其次根据张量积丛上诱导联络的定义得到
\[
\begin{split}
Dt&={\color{red}{\dd \left(\tensor{t}{^{i_{1}i_{2}\cdots i_{r}}_{j_{1}j_{2}\cdots j_{s}}}\right)}} \otimes \frac { \partial } { \partial x ^ { i_{1} } }\otimes\frac { \partial } { \partial x ^ { i_{2} } }\otimes\cdots\otimes\frac { \partial } { \partial x ^ { i_{r} } }\otimes\dd x^{j_{1}}\otimes\dd x^{j_{2}}\otimes\cdots\otimes\dd x^{j_{s}} \quad\quad\;\;\;\circled{$\star$}\\
&\quad+\tensor{t}{^{i_{1}i_{2}\cdots i_{r}}_{j_{1}j_{2}\cdots j_{s}}}{\color{red}{D\left(\frac { \partial } { \partial x ^ { i_{1} } }\right)}}\otimes\frac { \partial } { \partial x ^ { i_{2} } }\otimes\cdots\otimes\frac { \partial } { \partial x ^ { i_{r} } }\otimes\dd x^{j_{1}}\otimes\dd x^{j_{2}}\otimes\cdots\otimes\dd x^{j_{s}} \quad\quad\;\;\;\circled{$i_{1}$}\\
&\quad+\tensor{t}{^{i_{1}i_{2}\cdots i_{r}}_{j_{1}j_{2}\cdots j_{s}}}\frac { \partial } { \partial x ^ { i_{1} } }\otimes {\color{red}{D\left(\frac { \partial } { \partial x ^ { i_{2} } }\right)}}\otimes\cdots\otimes\frac { \partial } { \partial x ^ { i_{r} } }\otimes\dd x^{j_{1}}\otimes\dd x^{j_{2}}\otimes\cdots\otimes\dd x^{j_{s}} \quad\quad\;\;\;\circled{$i_{2}$}\\
&\quad\vdots\\
&\quad+\tensor{t}{^{i_{1}i_{2}\cdots i_{r}}_{j_{1}j_{2}\cdots j_{s}}}\frac { \partial } { \partial x ^ { i_{1} } }\otimes\frac { \partial } { \partial x ^ { i_{2} } }\otimes\cdots\otimes {\color{red}{D\left(\frac { \partial } { \partial x ^ { i_{r} } }\right)}}\otimes \dd x^{j_{1}}\otimes\dd x^{j_{2}}\otimes\cdots\otimes \dd x^{j_{s}} \quad\quad\;\;\;\circled{$i_{r}$}\\
&\quad+\tensor{t}{^{i_{1}i_{2}\cdots i_{r}}_{j_{1}j_{2}\cdots j_{s}}}\frac { \partial } { \partial x ^ { i_{1} } }\otimes\frac { \partial } { \partial x ^ { i_{2} } }\otimes\cdots\otimes\frac { \partial } { \partial x ^ { i_{r} } }\otimes {\color{red}{D\left(\dd x^{j_{1}}\right)}}\otimes\dd x^{j_{2}}\otimes\cdots\otimes\dd x^{j_{s}} \quad\quad\quad\circled{$j_{1}$}\\
&\quad+\tensor{t}{^{i_{1}i_{2}\cdots i_{r}}_{j_{1}j_{2}\cdots j_{s}}}\frac { \partial } { \partial x ^ { i_{1} } }\otimes\frac { \partial } { \partial x ^ { i_{2} } }\otimes\cdots\otimes\frac { \partial } { \partial x ^ { i_{r} } }\otimes \dd x^{j_{1}}\otimes {\color{red}{D\left(\dd x^{j_{2}}\right)}}\otimes\cdots\otimes\dd x^{j_{s}} \quad\quad\quad\circled{$j_{2}$}\\
&\quad\vdots\\
&\quad+\tensor{t}{^{i_{1}i_{2}\cdots i_{r}}_{j_{1}j_{2}\cdots j_{s}}}\frac { \partial } { \partial x ^ { i_{1} } }\otimes\frac { \partial } { \partial x ^ { i_{2} } }\otimes\cdots\otimes\frac { \partial } { \partial x ^ { i_{r} } }\otimes \dd x^{j_{1}}\otimes\dd x^{j_{2}}\otimes\cdots\otimes {\color{red}{D\left(\dd x^{j_{s}} \right)}}\quad\quad\quad\circled{$j_{s}$}\\
\end{split}
\]
由于上面这个式子太长了，所以我们拆开来一项一项看，我已经给每一项都标了序号。第一项，也就是标了\circled{$\star$}的这一项，是之后各项计算及调整指标的参考标准。
\[
\begin{split}
\circled{$i_{1}$}&=\tensor{t}{^{i_{1}i_{2}\cdots i_{r}}_{j_{1}j_{2}\cdots j_{s}}}\left[{\color{red}{\omega_{i_{1}}^{l}\otimes\frac { \partial } { \partial x ^ { l} }}}\right]\otimes\frac { \partial } { \partial x ^ { i_{2} } }\otimes\cdots\otimes\frac { \partial } { \partial x ^ { i_{r} } }\otimes\dd x^{j_{1}}\otimes\dd x^{j_{2}}\otimes\cdots\otimes\dd x^{j_{s}} \\
&\xlongequal{l\leftrightarrow i_{1}}\tensor{t}{^{li_{2}\cdots i_{r}}_{j_{1}j_{2}\cdots j_{s}}}\omega_{l}^{i_{1}}\otimes\frac { \partial } { \partial x ^ { i_{1}} }\otimes\frac { \partial } { \partial x ^ { i_{2} } }\otimes\cdots\otimes\frac { \partial } { \partial x ^ { i_{r} } }\otimes\dd x^{j_{1}}\otimes\dd x^{j_{2}}\otimes\cdots\otimes\dd x^{j_{s}} 
\end{split}
\]
解释一下：$l\leftrightarrow i_{1}$表示将指标$l$与指标$i_{1}$对调。之所以这样调整指标，是因为我们\textbf{以\circled{$\star$}这一项的指标顺序作为参考标准}。后面各项的指标调整也遵循这一原则。
\[
\begin{split}
\circled{$i_{2}$}&=\tensor{t}{^{i_{1}i_{2}\cdots i_{r}}_{j_{1}j_{2}\cdots j_{s}}}\frac { \partial } { \partial x ^ { i_{1} } }\otimes \left[{\color{red}{\omega_{i_{2}}^{l}\otimes\frac { \partial } { \partial x ^ { l} }}}\right]\otimes\cdots\otimes\frac { \partial } { \partial x ^ { i_{r} } }\otimes\dd x^{j_{1}}\otimes\dd x^{j_{2}}\otimes\cdots\otimes\dd x^{j_{s}}\\
&=\tensor{t}{^{i_{1}i_{2}\cdots i_{r}}_{j_{1}j_{2}\cdots j_{s}}}{\color{red}{\omega_{i_{2}}^{l}}}\otimes\frac { \partial } { \partial x ^ { i_{1} } }\otimes {\color{red}{\frac { \partial } { \partial x ^ { l} }}}\otimes\cdots\otimes\frac { \partial } { \partial x ^ { i_{r} } }\otimes\dd x^{j_{1}}\otimes\dd x^{j_{2}}\otimes\cdots\otimes\dd x^{j_{s}}\\
&\xlongequal{l\leftrightarrow i_{2}}\tensor{t}{^{i_{1}l\cdots i_{r}}_{j_{1}j_{2}\cdots j_{s}}}\omega_{l}^{i_{2}}\otimes\frac { \partial } { \partial x ^ { i_{1} } }\otimes \frac { \partial } { \partial x ^ { i_{2}} }\otimes\cdots\otimes\frac { \partial } { \partial x ^ { i_{r} } }\otimes\dd x^{j_{1}}\otimes\dd x^{j_{2}}\otimes\cdots\otimes\dd x^{j_{s}}
\end{split}
\]
注意上面第二个等号是约定：“\textbf{1-形式部分放到最左边}”。\circled{$i_{3}$}$,\cdots,$\circled{$i_{r}$}可类似计算出来。
\[
\begin{split}
\circled{$i_{r}$}&=\tensor{t}{^{i_{1}i_{2}\cdots i_{r}}_{j_{1}j_{2}\cdots j_{s}}}\frac { \partial } { \partial x ^ { i_{1} } }\otimes\frac { \partial } { \partial x ^ { i_{2} } }\otimes\cdots\otimes \left[{\color{red}{\omega_{i_{r}}^{l}\otimes\frac { \partial } { \partial x ^ { l} }}}\right]\otimes \dd x^{j_{1}}\otimes\dd x^{j_{2}}\otimes\cdots\otimes \dd x^{j_{s}}\\
&=\tensor{t}{^{i_{1}i_{2}\cdots i_{r}}_{j_{1}j_{2}\cdots j_{s}}}{\color{red}{\omega_{i_{r}}^{l}}}\otimes\frac { \partial } { \partial x ^ { i_{1} } }\otimes\frac { \partial } { \partial x ^ { i_{2} } }\otimes\cdots\otimes {\color{red}{\frac { \partial } { \partial x ^ { l} }}}\otimes \dd x^{j_{1}}\otimes\dd x^{j_{2}}\otimes\cdots\otimes \dd x^{j_{s}}\\
&\xlongequal{l\leftrightarrow i_{r}}\tensor{t}{^{i_{1}i_{2}\cdots l}_{j_{1}j_{2}\cdots j_{s}}}\omega_{l}^{i_{r}}\otimes\frac { \partial } { \partial x ^ { i_{1} } }\otimes\frac { \partial } { \partial x ^ { i_{2} } }\otimes\cdots\otimes \frac { \partial } { \partial x ^ { i_{r}} }\otimes \dd x^{j_{1}}\otimes\dd x^{j_{2}}\otimes\cdots\otimes \dd x^{j_{s}}
\end{split}
\]

利用对偶联络的性质，我们同理可计算\circled{$j_{1}$}$,\cdots,$\circled{$j_{s}$}.
\[
\begin{split}
	\circled{$j_{1}$}&=\tensor{t}{^{i_{1}i_{2}\cdots i_{r}}_{j_{1}j_{2}\cdots j_{s}}}\frac { \partial } { \partial x ^ { i_{1} } }\otimes\frac { \partial } { \partial x ^ { i_{2} } }\otimes\cdots\otimes\frac { \partial } { \partial x ^ { i_{r} } }\otimes \left[  {\color{red}{-\omega_{l}^{j_{1}}\otimes \dd x^{l}}}\right]\otimes\dd x^{j_{2}}\otimes\cdots\otimes\dd x^{j_{s}}  \\
	&={\color{red}{-}}\tensor{t}{^{i_{1}i_{2}\cdots i_{r}}_{j_{1}j_{2}\cdots j_{s}}}{\color{red}{\omega_{l}^{j_{1}}}}\otimes\frac { \partial } { \partial x ^ { i_{1} } }\otimes\frac { \partial } { \partial x ^ { i_{2} } }\otimes\cdots\otimes\frac { \partial } { \partial x ^ { i_{r} } }\otimes  {\color{red}{ \dd x^{l}}}\otimes\dd x^{j_{2}}\otimes\cdots\otimes\dd x^{j_{s}}  \\
	&\xlongequal{l\leftrightarrow j_{1}}-\tensor{t}{^{i_{1}i_{2}\cdots i_{r}}_{l j_{2}\cdots j_{s}}}\omega_{j_{1}}^{l}\otimes\frac { \partial } { \partial x ^ { i_{1} } }\otimes\frac { \partial } { \partial x ^ { i_{2} } }\otimes\cdots\otimes\frac { \partial } { \partial x ^ { i_{r} } }\otimes   \dd x^{j_{1}}\otimes\dd x^{j_{2}}\otimes\cdots\otimes\dd x^{j_{s}}  \\
\end{split}
\]
\[
\begin{split}
\circled{$j_{2}$}&=\tensor{t}{^{i_{1}i_{2}\cdots i_{r}}_{j_{1}j_{2}\cdots j_{s}}}\frac { \partial } { \partial x ^ { i_{1} } }\otimes\frac { \partial } { \partial x ^ { i_{2} } }\otimes\cdots\otimes\frac { \partial } { \partial x ^ { i_{r} } }\otimes \dd x^{j_{1}}\otimes \left[  {\color{red}{-\omega_{l}^{j_{2}}\otimes \dd x^{l}}}\right]\otimes\cdots\otimes\dd x^{j_{s}}  \\
&={\color{red}{-}}\tensor{t}{^{i_{1}i_{2}\cdots i_{r}}_{j_{1}j_{2}\cdots j_{s}}}{\color{red}{\omega_{l}^{j_{2}}}}\otimes\frac { \partial } { \partial x ^ { i_{1} } }\otimes\frac { \partial } { \partial x ^ { i_{2} } }\otimes\cdots\otimes\frac { \partial } { \partial x ^ { i_{r} } }\otimes  \dd x^{j_{1}}\otimes{\color{red}{ \dd x^{l}}}\otimes\cdots\otimes\dd x^{j_{s}}  \\
&\xlongequal{l\leftrightarrow j_{2}}-\tensor{t}{^{i_{1}i_{2}\cdots i_{r}}_{j_{1} l\cdots j_{s}}}\omega_{j_{2}}^{l}\otimes\frac { \partial } { \partial x ^ { i_{1} } }\otimes\frac { \partial } { \partial x ^ { i_{2} } }\otimes\cdots\otimes\frac { \partial } { \partial x ^ { i_{r} } }\otimes   \dd x^{j_{1}}\otimes\dd x^{j_{2}}\otimes\cdots\otimes\dd x^{j_{s}}  \\
\end{split}
\]
\[
\begin{split}
\circled{$j_{s}$}&=\tensor{t}{^{i_{1}i_{2}\cdots i_{r}}_{j_{1}j_{2}\cdots j_{s}}}\frac { \partial } { \partial x ^ { i_{1} } }\otimes\frac { \partial } { \partial x ^ { i_{2} } }\otimes\cdots\otimes\frac { \partial } { \partial x ^ { i_{r} } }\otimes \dd x^{j_{1}}\otimes\dd x^{j_{2}}\otimes\cdots\otimes \left[  {\color{red}{-\omega_{l}^{j_{s}}\otimes \dd x^{l}}}\right]\\
&={\color{red}{-}}\tensor{t}{^{i_{1}i_{2}\cdots i_{r}}_{j_{1}j_{2}\cdots j_{s}}}{\color{red}{\omega_{l}^{j_{s}}}}\otimes\frac { \partial } { \partial x ^ { i_{1} } }\otimes\frac { \partial } { \partial x ^ { i_{2} } }\otimes\cdots\otimes\frac { \partial } { \partial x ^ { i_{r} } }\otimes  \dd x^{j_{1}}\otimes\dd x^{j_{2}}\otimes\cdots\otimes{\color{red}{ \dd x^{l}}}  \\
&\xlongequal{l\leftrightarrow j_{s}}-\tensor{t}{^{i_{1}i_{2}\cdots i_{r}}_{j_{1} j_{2}\cdots l}}\omega^{l}_{j_{s}}\otimes\frac { \partial } { \partial x ^ { i_{1} } }\otimes\frac { \partial } { \partial x ^ { i_{2} } }\otimes\cdots\otimes\frac { \partial } { \partial x ^ { i_{r} } }\otimes   \dd x^{j_{1}}\otimes\dd x^{j_{2}}\otimes\cdots\otimes\dd x^{j_{s}}  \\
\end{split}
\]
因此
\[
\begin{split}
Dt&=\circled{$\star$}+\circled{$i_{1}$}+\circled{$i_{2}$}+\cdots+\circled{$i_{r}$}+\circled{$j_{1}$}+\circled{$j_{2}$}+\cdots+\circled{$j_{s}$}\\
&=\left[\dd \left(\tensor{t}{^{i_{1}i_{2}\cdots i_{r}}_{j_{1}j_{2}\cdots j_{s}}}\right)+\sum_{a=1}^{r}\tensor{t}{^{i_{1}\cdots l\cdots i_{r}}_{j_{1}j_{2}\cdots j_{s}}}\omega_{l}^{i_{a}}-\sum_{b=1}^{s}\tensor{t}{^{i_{1}i_{2}\cdots i_{r}}_{j_{1}\cdots l\cdots j_{s}}}\omega_{j_{b}}^{l}\right]\\
&\quad\quad\quad\quad\quad\quad\quad\;\quad\otimes\left[\frac { \partial } { \partial x ^ { i_{1} } }\otimes\frac { \partial } { \partial x ^ { i_{2} } }\otimes\cdots\otimes\frac { \partial } { \partial x ^ { i_{r} } }\otimes   \dd x^{j_{1}}\otimes\dd x^{j_{2}}\otimes\cdots\otimes\dd x^{j_{s}} \right]\\
&=\left[\frac{{\partial}(\tensor{t}{^{i_{1}i_{2}\cdots i_{r}}_{j_{1}j_{2}\cdots j_{s}}})}{\partial x^{k}}+\sum_{a=1}^{r}\tensor{t}{^{i_{1}\cdots l\cdots i_{r}}_{j_{1}j_{2}\cdots j_{s}}}\Gamma_{lk}^{i_{a}}\dd x^{k}-\sum_{b=1}^{s}\tensor{t}{^{i_{1}i_{2}\cdots i_{r}}_{j_{1}\cdots l\cdots j_{s}}}\Gamma_{j_{b}k}^{l}\dd x^{k}\right]\\
&\quad\quad\quad\quad\quad\quad\quad\;\quad\otimes\left[\frac { \partial } { \partial x ^ { i_{1} } }\otimes\frac { \partial } { \partial x ^ { i_{2} } }\otimes\cdots\otimes\frac { \partial } { \partial x ^ { i_{r} } }\otimes   \dd x^{j_{1}}\otimes\dd x^{j_{2}}\otimes\cdots\otimes\dd x^{j_{s}} \right]\\
&=\left[\frac{{\partial}(\tensor{t}{^{i_{1}i_{2}\cdots i_{r}}_{j_{1}j_{2}\cdots j_{s}}})}{\partial x^{k}}+\sum_{a=1}^{r}\tensor{t}{^{i_{1}\cdots l\cdots i_{r}}_{j_{1}j_{2}\cdots j_{s}}}\Gamma_{lk}^{i_{a}}-\sum_{b=1}^{s}\tensor{t}{^{i_{1}i_{2}\cdots i_{r}}_{j_{1}\cdots l\cdots j_{s}}}\Gamma_{j_{b}k}^{l}\right]\\
&\quad\quad\quad\quad\quad\quad\;\quad\dd x^{k}\otimes\left[\frac { \partial } { \partial x ^ { i_{1} } }\otimes\frac { \partial } { \partial x ^ { i_{2} } }\otimes\cdots\otimes\frac { \partial } { \partial x ^ { i_{r} } }\otimes   \dd x^{j_{1}}\otimes\dd x^{j_{2}}\otimes\cdots\otimes\dd x^{j_{s}} \right]\\
&=\left(\tensor{t}{^{i_{1}i_{2}\cdots i_{r}}_{j_{1}j_{2}\cdots j_{s}}}\right)^{\prime}_{ k}{\color{red}{\dd x^{k}}}\otimes\left[\frac { \partial } { \partial x ^ { i_{1} } }\otimes\frac { \partial } { \partial x ^ { i_{2} } }\otimes\cdots\otimes\frac { \partial } { \partial x ^ { i_{r} } }\otimes   \dd x^{j_{1}}\otimes\dd x^{j_{2}}\otimes\cdots\otimes\dd x^{j_{s}} \right]
\end{split}
\]
这里我们约定
\begin{equation}
\left(\tensor{t}{^{i_{1}i_{2}\cdots i_{r}}_{j_{1}j_{2}\cdots j_{s}}}\right)^{\prime}_{ k}\triangleq \frac{{\partial}(\tensor{t}{^{i_{1}i_{2}\cdots i_{r}}_{j_{1}j_{2}\cdots j_{s}}})}{\partial x^{k}}+\sum_{a=1}^{r}\tensor{t}{^{i_{1}\cdots l\cdots i_{r}}_{j_{1}j_{2}\cdots j_{s}}}\Gamma_{lk}^{i_{a}}-\sum_{b=1}^{s}\tensor{t}{^{i_{1}i_{2}\cdots i_{r}}_{j_{1}\cdots l\cdots j_{s}}}\Gamma_{j_{b}k}^{l}
\end{equation}

以上计算其实就是根据定义一步一步进行，唯一需要注意的地方只有那个约定：1-形式部分放在最左边。这个约定在最终的计算结果里面就体现为把计算过程中新产生的${\color{red}{\dd x^{k}}}$写在原有的基底张量的最前面，从而达到与定义$Dt\in \Gamma({\color{red}{T^{\ast}M}}\otimes T^{r}_{s})$保持一致的目的。

从绝对微分$Dt$的系数表达式\requ{equ:}来看，$\left(\tensor{t}{^{i_{1}i_{2}\cdots i_{r}}_{j_{1}j_{2}\cdots j_{s}}}\right)^{\prime}_{ k}$与普通欧氏空间中的偏导数相比，只是多了后面含有$\Gamma_{jk}^{i}$的这一部分。正是多出来的这一部分揭露了弯曲空间与平直空间的本质不同。

上面的计算结果也表明，$Dt$是一个$(r,s+1)$型张量，但是就物理学的要求来说，我们不太希望看到张量类型的改变，最好能消掉多出来的${\color{red}{\dd x^{k}}}$. 我们可以通过再对$Dt$做一次配合运算$\langle X,Dt\rangle ,X\in \Gamma(TM)$来消除额外的${\color{red}{\dd x^{k}}}$. 这样得到的新的张量也就是$D_{X}t$，称为张量场$t$沿着切向量场$X$的协变导数。
设$X=X^{l}\frac { \partial } { \partial x ^ { l } }$，让$X$与${\color{red}{\dd x^{k}}}$做配合，便有
\[
	\begin{split}
	&D_{X}t=\langle X,Dt\rangle\\
	&={\color{red}{\dd x^{k}(X)}}\left(\tensor{t}{^{i_{1}i_{2}\cdots i_{r}}_{j_{1}j_{2}\cdots j_{s}}}\right)^{\prime}_{ k}\frac { \partial } { \partial x ^ { i_{1} } }\otimes\frac { \partial } { \partial x ^ { i_{2} } }\otimes\cdots\otimes\frac { \partial } { \partial x ^ { i_{r} } }\otimes   \dd x^{j_{1}}\otimes\dd x^{j_{2}}\otimes\cdots\otimes\dd x^{j_{s}} \\
	&={\color{red}{X^{k}}}\left(\tensor{t}{^{i_{1}i_{2}\cdots i_{r}}_{j_{1}j_{2}\cdots j_{s}}}\right)^{\prime}_{ k}\frac { \partial } { \partial x ^ { i_{1} } }\otimes\frac { \partial } { \partial x ^ { i_{2} } }\otimes\cdots\otimes\frac { \partial } { \partial x ^ { i_{r} } }\otimes   \dd x^{j_{1}}\otimes\dd x^{j_{2}}\otimes\cdots\otimes\dd x^{j_{s}} \\
	\end{split}
\]

\begin{mythm}{第二Bianchi恒等式}{chapter2第二Bianchi恒等式}
	若$D$是无挠仿射联络，则有第二Bianchi恒等式：$\forall i,j,k,l,h\in\{1,2,\cdots,m \}$
	\[
	\left(\tensor{R}{_i^j_{kl}}\right)^{\prime}_{h}+\left(\tensor{R}{_i^j_{lh}}\right)^{\prime}_{k}+\left(\tensor{R}{_i^j_{hk}}\right)^{\prime}_{l}=0
	\]
\end{mythm}
\begin{proof}
	\[
	\begin{split}
	\left(\tensor{R}{_i^j_{kl}}\right)^{\prime}_{h}&=\frac{\partial \tensor{R}{_i^j_{kl}}}{\partial x^{h}}+\tensor{R}{_i^p_{kl}}\Gamma_{ph}^{j} -\tensor{R}{_p^j_{kl}}\Gamma_{ih}^{p}{\color{yellow}{-\tensor{R}{_i^j_{pl}}\Gamma_{kh}^{p}}}-{\color{blue}{\tensor{R}{_i^j_{kp}}\Gamma_{lh}^{p}}}\\
	\left(\tensor{R}{_i^j_{lh}}\right)^{\prime}_{k}&=\frac{\partial \tensor{R}{_i^j_{lh}}}{\partial x^{k}}+\tensor{R}{_i^p_{lh}}\Gamma_{pk}^{j} -\tensor{R}{_p^j_{lh}}\Gamma_{ik}^{p}{\color{green}{-\tensor{R}{_i^j_{ph}}\Gamma_{lk}^{p}}}-{\color{yellow}{\tensor{R}{_i^j_{lp}}\Gamma_{hk}^{p}}}\\
	\left(\tensor{R}{_i^j_{hk}}\right)^{\prime}_{l}&=\frac{\partial \tensor{R}{_i^j_{hk}}}{\partial x^{l}}+\tensor{R}{_i^p_{hk}}\Gamma_{pl}^{j} -\tensor{R}{_p^j_{hk}}\Gamma_{il}^{p}{\color{blue}{-\tensor{R}{_i^j_{pk}}\Gamma_{hl}^{p}}}-{\color{green}{\tensor{R}{_i^j_{hp}}\Gamma_{kl}^{p}}}
	\end{split}
	\]
	注意在自然标架场下，$[\frac { \partial } { \partial x _ { i } } ,\frac { \partial } { \partial x _ { j } }]=0$. 所以由联络的无挠性可得$\Gamma_{ik}^{j}=\Gamma_{ki}^{j}$. 又因为$\tensor{R}{_i^j_{kl}}=-\tensor{R}{_i^j_{lk}}$. 所以
	\[
	\begin{split}
	\left(\tensor{R}{_i^j_{kl}}\right)^{\prime}_{h}+\left(\tensor{R}{_i^j_{lh}}\right)^{\prime}_{k}+\left(\tensor{R}{_i^j_{hk}}\right)^{\prime}_{l}&=\frac{\partial \tensor{R}{_i^j_{kl}}}{\partial x^{h}}+\tensor{R}{_i^p_{kl}}\Gamma_{ph}^{j} -\tensor{R}{_p^j_{kl}}\Gamma_{ih}^{p}\\
	&+\frac{\partial \tensor{R}{_i^j_{lh}}}{\partial x^{k}}+\tensor{R}{_i^p_{lh}}\Gamma_{pk}^{j} -\tensor{R}{_p^j_{lh}}\Gamma_{ik}^{p}\\
	&+\frac{\partial \tensor{R}{_i^j_{hk}}}{\partial x^{l}}+\tensor{R}{_i^p_{hk}}\Gamma_{pl}^{j} -\tensor{R}{_p^j_{hk}}\Gamma_{il}^{p}
	\end{split}
	\]
	根据Bianchi恒等式(\rthm{thm:chapter2Bianchi恒等式})，我们有$\dd \Omega_{i}^{j}=\omega_{i}^{p}\wedge\Omega_{p}^{j}-\Omega_{i}^{p}\wedge\omega_{p}^{j}$. 代入Cartan结构方程，我们可作如下计算
	\[
	\begin{split}
	&\dd \left(\frac{1}{2} \tensor{R}{_i^j_{kl}}\dd x^{k}\wedge\dd x^{l}\right)=\Gamma_{ih}^{p}\dd x^{h}\wedge\left(\frac{1}{2} \tensor{R}{_p^j_{kl}}\dd x^{k}\wedge\dd x^{l} \right)-\left(\frac{1}{2} \tensor{R}{_i^p_{kl}}\dd x^{k}\wedge\dd x^{l} \right)\wedge\Gamma_{ph}^{j}\dd x^{h}\\
	&\Longrightarrow \frac{\partial \tensor{R}{_i^j_{kl}}}{\partial x^{h}}\dd x^{h}\wedge\dd x^{k}\wedge\dd x^{l}=\left(\tensor{R}{_p^j_{kl}}\Gamma_{ih}^{p}-\tensor{R}{_i^p_{kl}}\Gamma_{ph}^{j}\right)\dd x^{h}\wedge\dd x^{k}\wedge\dd x^{l}\\
	&\Longrightarrow \left(\frac{\partial \tensor{R}{_i^j_{kl}}}{\partial x^{h}}+\tensor{R}{_i^p_{kl}}\Gamma_{ph}^{j} -\tensor{R}{_p^j_{kl}}\Gamma_{ih}^{p}\right)\dd x^{h}\wedge\dd x^{k}\wedge\dd x^{l}=0
	\end{split}
	\]
	同理可得
	\[
	\begin{split}
	\left(\frac{\partial \tensor{R}{_i^j_{lh}}}{\partial x^{k}}+\tensor{R}{_i^p_{lh}}\Gamma_{pk}^{j} -\tensor{R}{_p^j_{lh}}\Gamma_{ik}^{p}\right)\dd x^{k}\wedge\dd x^{l}\wedge\dd x^{h}&=0\\
	\left(\frac{\partial \tensor{R}{_i^j_{hk}}}{\partial x^{l}}+\tensor{R}{_i^p_{hk}}\Gamma_{pl}^{j} -\tensor{R}{_p^j_{hk}}\Gamma_{il}^{p}\right)\dd x^{l}\wedge\dd x^{h}\wedge\dd x^{k}&=0
	\end{split}
	\]
	把外微分式都调整成$\dd x^{h}\wedge\dd x^{k}\wedge\dd x^{l}$的顺序，则有
	\[
	\begin{cases}
%	\begin{split}
	\left(\frac{\partial \tensor{R}{_i^j_{kl}}}{\partial x^{h}}+\tensor{R}{_i^p_{kl}}\Gamma_{ph}^{j} -\tensor{R}{_p^j_{kl}}\Gamma_{ih}^{p}\right)\dd x^{h}\wedge\dd x^{k}\wedge\dd x^{l}=0\\
	\left(\frac{\partial \tensor{R}{_i^j_{lh}}}{\partial x^{k}}+\tensor{R}{_i^p_{lh}}\Gamma_{pk}^{j} -\tensor{R}{_p^j_{lh}}\Gamma_{ik}^{p}\right)\dd x^{h}\wedge\dd x^{k}\wedge\dd x^{l}=0\\
	\left(\frac{\partial \tensor{R}{_i^j_{hk}}}{\partial x^{l}}+\tensor{R}{_i^p_{hk}}\Gamma_{pl}^{j} -\tensor{R}{_p^j_{hk}}\Gamma_{il}^{p}\right)\dd x^{h}\wedge\dd x^{k}\wedge\dd x^{l}=0
%	\end{split}
	\end{cases}
	\]
	三式相加便得到
	\[
	\left[\left(\tensor{R}{_i^j_{kl}}\right)^{\prime}_{h}+\left(\tensor{R}{_i^j_{lh}}\right)^{\prime}_{k}+\left(\tensor{R}{_i^j_{hk}}\right)^{\prime}_{l}\right]\dd x^{h}\wedge\dd x^{k}\wedge\dd x^{l}=0
	\]
	注意到$\left(\tensor{R}{_i^j_{kl}}\right)^{\prime}_{h}+\left(\tensor{R}{_i^j_{lh}}\right)^{\prime}_{k}+\left(\tensor{R}{_i^j_{hk}}\right)^{\prime}_{l}$关于指标$k,l,h$有反对称性，而$\dd x^{h}\wedge\dd x^{k}\wedge\dd x^{l}$也关于指标$k,l,h$反对称，所以我们必有
	\[
	\left(\tensor{R}{_i^j_{kl}}\right)^{\prime}_{h}+\left(\tensor{R}{_i^j_{lh}}\right)^{\prime}_{k}+\left(\tensor{R}{_i^j_{hk}}\right)^{\prime}_{l}=0
	\]
\end{proof}

下面我们转而来考虑交换两次求导顺序得到的结果是否相同的问题。从曲率算子的等价定义(\rthm{thm:chapter2曲率算子的等价定义})我们可以看出一般情况下答案是否定的，并且按照不同顺序连续求两次“导数”所得结果的差可以用曲率来描述，即
\[
D_{X}D_{Y}-D_{Y}D_{X}=\rr(X,Y)+D_{[X,Y]}
\]
现在我们着重在局部坐标下来计算，从而给出具体的表达式，即重要的Ricci恒等式。

首先约定一个记号：
\[
	\left(\tensor{t}{^{i_{1}i_{2}\cdots i_{r}}_{j_{1}j_{2}\cdots j_{s}}}\right)^{\prime\prime}_{ kl}\triangleq\left(\left(\tensor{t}{^{i_{1}i_{2}\cdots i_{r}}_{j_{1}j_{2}\cdots j_{s}}}\right)^{\prime}_{ k}\right)^{\prime}_{l}
\]
我们可以依此类推约定更高阶“导数”的记号。所谓Ricci恒等式其实就是问
\[
	\left(\tensor{t}{^{i_{1}i_{2}\cdots i_{r}}_{j_{1}j_{2}\cdots j_{s}}}\right)^{\prime\prime}_{ kl}-\left(\tensor{t}{^{i_{1}i_{2}\cdots i_{r}}_{j_{1}j_{2}\cdots j_{s}}}\right)^{\prime\prime}_{ lk}=\;\mbf{?} 
\]
\begin{mythm}{Ricci恒等式}{chapter2Ricci恒等式}
	设$D$是光滑流形$M$上的无挠仿射联络，$(r,s)$型张量场$t$在坐标卡$\cx$下的局部表示为
	\[
	t=\tensor{t}{^{i_{1}i_{2}\cdots i_{r}}_{j_{1}j_{2}\cdots j_{s}}}\frac { \partial } { \partial x ^ { i_{1} } }\otimes\frac { \partial } { \partial x ^ { i_{2} } }\otimes\cdots\otimes\frac { \partial } { \partial x ^ { i_{r} } }\otimes\dd x^{j_{1}}\otimes\dd x^{j_{2}}\otimes\cdots\otimes\dd x^{j_{s}}
	\]
	则有
	\[
	\begin{split}
		&\left(\tensor{t}{^{i_{1}i_{2}\cdots i_{r}}_{j_{1}j_{2}\cdots j_{s}}}\right)^{\prime\prime}_{ kl}-\left(\tensor{t}{^{i_{1}i_{2}\cdots i_{r}}_{j_{1}j_{2}\cdots j_{s}}}\right)^{\prime\prime}_{ lk}\\
		&=-\sum_{a=1}^{r} \tensor{t}{^{i_{1}\cdots q\cdots i_{r}}_{j_{1}j_{2}\cdots j_{s}}}\tensor{R}{_q^{i_{a}}_{kl}} -\sum_{b=1}^{s} \tensor{t}{^{i_{1}i_{2}\cdots i_{r}}_{j_{1}\cdots q\cdots j_{s}}}\tensor{R}{_{j_{b}}^{q}_{lk}} 
	\end{split}
	\]
\end{mythm}
\begin{proof}
\[
\begin{split}
\left( \tensor{t}{^{i_{1}i_{2}\cdots i_{r}}_{j_{1}j_{2}\cdots j_{s}}} \right)^{\prime\prime}_{kl}=\left( \left( \tensor{t}{^{i_{1}i_{2}\cdots i_{r}}_{j_{1}j_{2}\cdots j_{s}}} \right)^{\prime}_{k} \right)^{\prime}_{l}=&\frac{{\partial}\left( \tensor{t}{^{i_{1}i_{2}\cdots i_{r}}_{j_{1}j_{2}\cdots j_{s}}} \right)^{\prime}_{k} }{\partial x^{l}}+\sum_{c=1}^{r}\left( \tensor{t}{^{i_{1}\cdots p\cdots i_{r}}_{j_{1}j_{2}\cdots j_{s}}} \right)^{\prime}_{k} \Gamma_{pl}^{i_{c}}\\
-&\left( \tensor{t}{^{i_{1}i_{2}\cdots i_{r}}_{j_{1}j_{2}\cdots j_{s}}} \right)^{\prime}_{p}\Gamma_{kl}^{p}-\sum_{d=1}^{s}\left( \tensor{t}{^{i_{1}i_{2}\cdots i_{r}}_{j_{1}\cdots p\cdots  j_{s}}} \right)^{\prime}_{k} \Gamma_{j_{d}l}^{p}
\end{split}
\]
下面来分别计算各项：
\[
\begin{split}
&\frac{\partial \left( \tensor{t}{^{i_{1}i_{2}\cdots i_{r}}_{j_{1}j_{2}\cdots j_{s}}} \right)^{\prime}_{k} }{\partial x^{l}}\\
&=\frac{{\partial}^{2}(\tensor{t}{^{i_{1}i_{2}\cdots i_{r}}_{j_{1}j_{2}\cdots j_{s}}})}{\partial x^{l}\partial x^{k}}+\sum_{c=1}^{r}\frac{\partial  }{\partial x^{l}}\left( \tensor{t}{^{i_{1}\cdots q\cdots i_{r}}_{j_{1}j_{2}\cdots j_{s}}}\Gamma_{qk}^{i_{c}} \right)-\sum_{d=1}^{s}\frac{\partial  }{\partial x^{l}}\left( \tensor{t}{^{i_{1}i_{2}\cdots i_{r}}_{j_{1}\cdots q\cdots j_{s}}}\Gamma_{j_{d}k}^{q} \right)\\
&=\frac{{\partial}^{2}(\tensor{t}{^{i_{1}i_{2}\cdots i_{r}}_{j_{1}j_{2}\cdots j_{s}}})}{\partial x^{l}\partial x^{k}}+\sum_{c=1}^{r}\tensor{t}{^{i_{1}\cdots q\cdots i_{r}}_{j_{1}j_{2}\cdots j_{s}}}\frac{\partial  }{\partial x^{l}} \left( \Gamma_{qk}^{i_{c}} \right) +\sum_{c=1}^{r}\Gamma_{qk}^{i_{c}}\frac{\partial  }{\partial x^{l}}\left( \tensor{t}{^{i_{1}\cdots q\cdots i_{r}}_{j_{1}j_{2}\cdots j_{s}}} \right)\\
& \quad\quad\quad\quad\quad\quad\quad\quad-\sum_{d=1}^{s} \tensor{t}{^{i_{1}i_{2}\cdots i_{r}}_{j_{1}\cdots q\cdots j_{s}}}\frac{\partial  }{\partial x^{l}}\left(\Gamma_{j_{d}k}^{q} \right)-\sum_{d=1}^{s}\Gamma_{j_{d}k}^{q}  \frac{\partial  }{\partial x^{l}}\left(\tensor{t}{^{i_{1}i_{2}\cdots i_{r}}_{j_{1}\cdots q\cdots j_{s}}}\right)\\
\end{split}
\]
因为
\[
\begin{split}
	&\left( \tensor{t}{^{i_{1}\cdots p\cdots i_{r}}_{j_{1}j_{2}\cdots j_{s}}} \right)^{\prime}_{k}\\
	&=\frac{{\partial}(\tensor{t}{^{i_{1}\cdots p\cdots i_{r}}_{j_{1}j_{2}\cdots j_{s}}} )}{\partial x^{k}}+\tensor{t}{^{i_{1}\cdots q\cdots i_{r}}_{j_{1}j_{2}\cdots j_{s}}} \Gamma_{qk}^{p}+\sum_{\substack{a=1,\\ a\neq c}}^{r}\tensor{t}{^{i_{1}\cdots q\cdots p\cdots i_{r}}_{j_{1}j_{2}\cdots j_{s}}}\Gamma_{qk}^{i_{a}}-\sum_{b=1}^{s}\tensor{t}{^{i_{1}\cdots p\cdots i_{r}}_{j_{1}\cdots q\cdots j_{s}}}\Gamma_{j_{b}k}^{q}
\end{split}
\]
所以
\[
\begin{split}
	&\sum_{c=1}^{r}\left( \tensor{t}{^{i_{1}\cdots p\cdots i_{r}}_{j_{1}j_{2}\cdots j_{s}}} \right)^{\prime}_{k} \Gamma_{pl}^{i_{c}}\\
	&=\sum_{c=1}^{r}\frac{{\partial}(\tensor{t}{^{i_{1}\cdots p\cdots i_{r}}_{j_{1}j_{2}\cdots j_{s}}} )}{\partial x^{k}}\Gamma_{pl}^{i_{c}}+\sum_{c=1}^{r}\tensor{t}{^{i_{1}\cdots q\cdots i_{r}}_{j_{1}j_{2}\cdots j_{s}}} \Gamma_{qk}^{p}\Gamma_{pl}^{i_{c}}+\sum_{c=1}^{r}\sum_{\substack{a=1,\\ a\neq c}}^{r}\tensor{t}{^{i_{1}\cdots q\cdots p\cdots i_{r}}_{j_{1}j_{2}\cdots j_{s}}}\Gamma_{qk}^{i_{a}}\Gamma_{pl}^{i_{c}}\\
	&\;\quad\quad\quad\quad\quad\quad\quad\quad\quad\quad\quad\quad\quad\quad\quad\quad\quad\quad\quad\quad\quad\quad-\sum_{c=1}^{r}\sum_{b=1}^{s}\tensor{t}{^{i_{1}\cdots p\cdots i_{r}}_{j_{1}\cdots q\cdots j_{s}}}\Gamma_{j_{b}k}^{q}  \Gamma_{pl}^{i_{c}}
\end{split}
\]
因为
\[
\begin{split}
	&\left( \tensor{t}{^{i_{1}i_{2}\cdots  i_{r}}_{j_{1}\cdots p\cdots j_{s}}} \right)^{\prime}_{k}\\
	&=\frac{{\partial}( \tensor{t}{^{i_{1}i_{2}\cdots  i_{r}}_{j_{1}\cdots p\cdots j_{s}}})}{\partial x^{k}}+\sum_{a=1}^{r} \tensor{t}{^{i_{1}\cdots q \cdots i_{r}}_{j_{1}\cdots p\cdots j_{s}}}\Gamma_{qk}^{i_{a}}-\tensor{t}{^{i_{1}i_{2}\cdots  i_{r}}_{j_{1}\cdots q\cdots j_{s}}} \Gamma_{pk}^{q}-\sum_{\substack{b=1,\\ b\neq d}}^{s}\tensor{t}{^{i_{1}i_{2}\cdots  i_{r}}_{j_{1}\cdots q\cdots p\cdots j_{s}}}\Gamma_{j_{b}k}^{q}
\end{split}
\]
所以
\[
\begin{split}
	&\sum_{d=1}^{s}\left( \tensor{t}{^{i_{1}i_{2}\cdots i_{r}}_{j_{1}\cdots p\cdots  j_{s}}} \right)^{\prime}_{k} \Gamma_{j_{d}l}^{p}\\
	&=\sum_{d=1}^{s}\frac{{\partial}( \tensor{t}{^{i_{1}i_{2}\cdots  i_{r}}_{j_{1}\cdots p\cdots j_{s}}})}{\partial x^{k}}\Gamma_{j_{d}l}^{p}-\sum_{d=1}^{s}\tensor{t}{^{i_{1}i_{2}\cdots  i_{r}}_{j_{1}\cdots q\cdots j_{s}}} \Gamma_{pk}^{q}\Gamma_{j_{d}l}^{p}-\sum_{d=1}^{s}\sum_{\substack{b=1,\\ b\neq d}}^{s}\tensor{t}{^{i_{1}i_{2}\cdots  i_{r}}_{j_{1}\cdots q\cdots p\cdots j_{s}}}\Gamma_{j_{b}k}^{q}\Gamma_{j_{d}l}^{p}\\
	&\;\quad\quad\quad\quad\quad\quad\quad\quad\quad\quad\quad\quad\quad\quad\quad\quad\quad\quad\quad\quad\quad\quad+\sum_{d=1}^{s}\sum_{a=1}^{r} \tensor{t}{^{i_{1}\cdots q \cdots i_{r}}_{j_{1}\cdots p\cdots j_{s}}}\Gamma_{qk}^{i_{a}}  \Gamma_{j_{d}l}^{p}
\end{split}
\]
所以 
\[
\begin{split}
&\left( \tensor{t}{^{i_{1}i_{2}\cdots i_{r}}_{j_{1}j_{2}\cdots j_{s}}} \right)^{\prime\prime}_{kl}\\
&=\frac{{\partial}^{2}(\tensor{t}{^{i_{1}i_{2}\cdots i_{r}}_{j_{1}j_{2}\cdots j_{s}}})}{\partial x^{l}\partial x^{k}}+\sum_{c=1}^{r}\tensor{t}{^{i_{1}\cdots q\cdots i_{r}}_{j_{1}j_{2}\cdots j_{s}}}\frac{\partial  }{\partial x^{l}} \left( \Gamma_{qk}^{i_{c}} \right) +\sum_{c=1}^{r}\Gamma_{qk}^{i_{c}}\frac{\partial  }{\partial x^{l}}\left( \tensor{t}{^{i_{1}\cdots q\cdots i_{r}}_{j_{1}j_{2}\cdots j_{s}}} \right)\\
& \quad\quad\quad\quad\quad\quad\quad\quad-\sum_{d=1}^{s} \tensor{t}{^{i_{1}i_{2}\cdots i_{r}}_{j_{1}\cdots q\cdots j_{s}}}\frac{\partial  }{\partial x^{l}}\left(\Gamma_{j_{d}k}^{q} \right)-\sum_{d=1}^{s}\Gamma_{j_{d}k}^{q}  \frac{\partial  }{\partial x^{l}}\left(\tensor{t}{^{i_{1}i_{2}\cdots i_{r}}_{j_{1}\cdots q\cdots j_{s}}}\right)\\
&+\sum_{c=1}^{r}\frac{{\partial}(\tensor{t}{^{i_{1}\cdots p\cdots i_{r}}_{j_{1}j_{2}\cdots j_{s}}} )}{\partial x^{k}}\Gamma_{pl}^{i_{c}}+\sum_{c=1}^{r}\tensor{t}{^{i_{1}\cdots q\cdots i_{r}}_{j_{1}j_{2}\cdots j_{s}}} \Gamma_{qk}^{p}\Gamma_{pl}^{i_{c}}+\sum_{c=1}^{r}\sum_{\substack{a=1,\\ a\neq c}}^{r}\tensor{t}{^{i_{1}\cdots q\cdots p\cdots i_{r}}_{j_{1}j_{2}\cdots j_{s}}}\Gamma_{qk}^{i_{a}}\Gamma_{pl}^{i_{c}}\\
&\;\quad\quad\quad\quad\quad\quad\quad\quad\quad\quad\quad\quad\quad\quad\quad\quad\quad\quad\quad\quad\quad\quad-\sum_{c=1}^{r}\sum_{b=1}^{s}\tensor{t}{^{i_{1}\cdots p\cdots i_{r}}_{j_{1}\cdots q\cdots j_{s}}}\Gamma_{j_{b}k}^{q}  \Gamma_{pl}^{i_{c}}\\
&-\left( \tensor{t}{^{i_{1}i_{2}\cdots i_{r}}_{j_{1}j_{2}\cdots j_{s}}} \right)^{\prime}_{p}\Gamma_{kl}^{p}\\
&-\sum_{d=1}^{s}\frac{{\partial}( \tensor{t}{^{i_{1}i_{2}\cdots  i_{r}}_{j_{1}\cdots p\cdots j_{s}}})}{\partial x^{k}}\Gamma_{j_{d}l}^{p}+\sum_{d=1}^{s}\tensor{t}{^{i_{1}i_{2}\cdots  i_{r}}_{j_{1}\cdots q\cdots j_{s}}} \Gamma_{pk}^{q}\Gamma_{j_{d}l}^{p}+\sum_{d=1}^{s}\sum_{\substack{b=1,\\ b\neq d}}^{s}\tensor{t}{^{i_{1}i_{2}\cdots  i_{r}}_{j_{1}\cdots q\cdots p\cdots j_{s}}}\Gamma_{j_{b}k}^{q}\Gamma_{j_{d}l}^{p}\\
&\;\quad\quad\quad\quad\quad\quad\quad\quad\quad\quad\quad\quad\quad\quad\quad\quad\quad\quad\quad\quad\quad\quad-\sum_{d=1}^{s}\sum_{a=1}^{r} \tensor{t}{^{i_{1}\cdots q \cdots i_{r}}_{j_{1}\cdots p\cdots j_{s}}}\Gamma_{qk}^{i_{a}}  \Gamma_{j_{d}l}^{p}
\end{split}
\]
将 $ k $与 $ l $互换便得
\[
\begin{split}
&\left( \tensor{t}{^{i_{1}i_{2}\cdots i_{r}}_{j_{1}j_{2}\cdots j_{s}}} \right)^{\prime\prime}_{lk}\\
&=\frac{{\partial}^{2}(\tensor{t}{^{i_{1}i_{2}\cdots i_{r}}_{j_{1}j_{2}\cdots j_{s}}})}{\partial x^{k}\partial x^{l}}+\sum_{c=1}^{r}\tensor{t}{^{i_{1}\cdots q\cdots i_{r}}_{j_{1}j_{2}\cdots j_{s}}}\frac{\partial  }{\partial x^{k}} \left( \Gamma_{ql}^{i_{c}} \right) +\sum_{c=1}^{r}\Gamma_{ql}^{i_{c}}\frac{\partial  }{\partial x^{k}}\left( \tensor{t}{^{i_{1}\cdots q\cdots i_{r}}_{j_{1}j_{2}\cdots j_{s}}} \right)\\
& \quad\quad\quad\quad\quad\quad\quad\quad-\sum_{d=1}^{s} \tensor{t}{^{i_{1}i_{2}\cdots i_{r}}_{j_{1}\cdots q\cdots j_{s}}}\frac{\partial  }{\partial x^{k}}\left(\Gamma_{j_{d}l}^{q} \right)-\sum_{d=1}^{s}\Gamma_{j_{d}l}^{q}  \frac{\partial  }{\partial x^{k}}\left(\tensor{t}{^{i_{1}i_{2}\cdots i_{r}}_{j_{1}\cdots q\cdots j_{s}}}\right)\\
&+\sum_{c=1}^{r}\frac{{\partial}(\tensor{t}{^{i_{1}\cdots p\cdots i_{r}}_{j_{1}j_{2}\cdots j_{s}}} )}{\partial x^{l}}\Gamma_{pk}^{i_{c}}+\sum_{c=1}^{r}\tensor{t}{^{i_{1}\cdots q\cdots i_{r}}_{j_{1}j_{2}\cdots j_{s}}} \Gamma_{ql}^{p}\Gamma_{pk}^{i_{c}}+\sum_{c=1}^{r}\sum_{\substack{a=1,\\ a\neq c}}^{r}\tensor{t}{^{i_{1}\cdots q\cdots p\cdots i_{r}}_{j_{1}j_{2}\cdots j_{s}}}\Gamma_{ql}^{i_{a}}\Gamma_{pk}^{i_{c}}\\
&\;\quad\quad\quad\quad\quad\quad\quad\quad\quad\quad\quad\quad\quad\quad\quad\quad\quad\quad\quad\quad\quad\quad-\sum_{c=1}^{r}\sum_{b=1}^{s}\tensor{t}{^{i_{1}\cdots p\cdots i_{r}}_{j_{1}\cdots q\cdots j_{s}}}\Gamma_{j_{b}l}^{q}  \Gamma_{pk}^{i_{c}}\\
&-\left( \tensor{t}{^{i_{1}i_{2}\cdots i_{r}}_{j_{1}j_{2}\cdots j_{s}}} \right)^{\prime}_{p}\Gamma_{lk}^{p}\\
&-\sum_{d=1}^{s}\frac{{\partial}( \tensor{t}{^{i_{1}i_{2}\cdots  i_{r}}_{j_{1}\cdots p\cdots j_{s}}})}{\partial x^{l}}\Gamma_{j_{d}k}^{p}+\sum_{d=1}^{s}\tensor{t}{^{i_{1}i_{2}\cdots  i_{r}}_{j_{1}\cdots q\cdots j_{s}}} \Gamma_{pl}^{q}\Gamma_{j_{d}k}^{p}+\sum_{d=1}^{s}\sum_{\substack{b=1,\\ b\neq d}}^{s}\tensor{t}{^{i_{1}i_{2}\cdots  i_{r}}_{j_{1}\cdots q\cdots p\cdots j_{s}}}\Gamma_{j_{b}l}^{q}\Gamma_{j_{d}k}^{p}\\
&\;\quad\quad\quad\quad\quad\quad\quad\quad\quad\quad\quad\quad\quad\quad\quad\quad\quad\quad\quad\quad\quad\quad-\sum_{d=1}^{s}\sum_{a=1}^{r} \tensor{t}{^{i_{1}\cdots q \cdots i_{r}}_{j_{1}\cdots p\cdots j_{s}}}\Gamma_{ql}^{i_{a}}  \Gamma_{j_{d}k}^{p}
\end{split}
\]

将上面所得到的两个式子相减，由于式子太长，我只把相减过程中可以配对消掉的写出来(有些需要统一求和的哑指标，这一点也容易看出来)：
\[
	-\left( \tensor{t}{^{i_{1}i_{2}\cdots i_{r}}_{j_{1}j_{2}\cdots j_{s}}} \right)^{\prime}_{p}\Gamma_{kl}^{p}+\left( \tensor{t}{^{i_{1}i_{2}\cdots i_{r}}_{j_{1}j_{2}\cdots j_{s}}} \right)^{\prime}_{p}\Gamma_{lk}^{p}=0
\]
\[
	\frac{{\partial}^{2}(\tensor{t}{^{i_{1}i_{2}\cdots i_{r}}_{j_{1}j_{2}\cdots j_{s}}})}{\partial x^{l}\partial x^{k}}-\frac{{\partial}^{2}(\tensor{t}{^{i_{1}i_{2}\cdots i_{r}}_{j_{1}j_{2}\cdots j_{s}}})}{\partial x^{k}\partial x^{l}}=0
\]
\[
	\sum_{c=1}^{r}\Gamma_{qk}^{i_{c}}\frac{\partial  }{\partial x^{l}}\left( \tensor{t}{^{i_{1}\cdots q\cdots i_{r}}_{j_{1}j_{2}\cdots j_{s}}} \right)-\sum_{c=1}^{r}\frac{{\partial}(\tensor{t}{^{i_{1}\cdots p\cdots i_{r}}_{j_{1}j_{2}\cdots j_{s}}} )}{\partial x^{l}}\Gamma_{pk}^{i_{c}}=0
\]
\[
	-\sum_{d=1}^{s}\Gamma_{j_{d}k}^{q}  \frac{\partial  }{\partial x^{l}}\left(\tensor{t}{^{i_{1}i_{2}\cdots i_{r}}_{j_{1}\cdots q\cdots j_{s}}}\right)+\sum_{d=1}^{s}\frac{{\partial}( \tensor{t}{^{i_{1}i_{2}\cdots  i_{r}}_{j_{1}\cdots p\cdots j_{s}}})}{\partial x^{l}}\Gamma_{j_{d}k}^{p}=0
\]
\[
	\sum_{c=1}^{r}\frac{{\partial}(\tensor{t}{^{i_{1}\cdots p\cdots i_{r}}_{j_{1}j_{2}\cdots j_{s}}} )}{\partial x^{k}}\Gamma_{pl}^{i_{c}}-\sum_{c=1}^{r}\Gamma_{ql}^{i_{c}}\frac{\partial  }{\partial x^{k}}\left( \tensor{t}{^{i_{1}\cdots q\cdots i_{r}}_{j_{1}j_{2}\cdots j_{s}}} \right)=0
\]
\[
	-\sum_{c=1}^{r}\sum_{b=1}^{s}\tensor{t}{^{i_{1}\cdots p\cdots i_{r}}_{j_{1}\cdots q\cdots j_{s}}}\Gamma_{j_{b}k}^{q}  \Gamma_{pl}^{i_{c}}+\sum_{d=1}^{s}\sum_{a=1}^{r} \tensor{t}{^{i_{1}\cdots q \cdots i_{r}}_{j_{1}\cdots p\cdots j_{s}}}\Gamma_{ql}^{i_{a}}  \Gamma_{j_{d}k}^{p}=0
\]
\[
	-\sum_{d=1}^{s}\frac{{\partial}( \tensor{t}{^{i_{1}i_{2}\cdots  i_{r}}_{j_{1}\cdots p\cdots j_{s}}})}{\partial x^{k}}\Gamma_{j_{d}l}^{p}+\sum_{d=1}^{s}\Gamma_{j_{d}l}^{q}  \frac{\partial  }{\partial x^{k}}\left(\tensor{t}{^{i_{1}i_{2}\cdots i_{r}}_{j_{1}\cdots q\cdots j_{s}}}\right)=0
\]
\[
	-\sum_{d=1}^{s}\sum_{a=1}^{r} \tensor{t}{^{i_{1}\cdots q \cdots i_{r}}_{j_{1}\cdots p\cdots j_{s}}}\Gamma_{qk}^{i_{a}}  \Gamma_{j_{d}l}^{p}+\sum_{c=1}^{r}\sum_{b=1}^{s}\tensor{t}{^{i_{1}\cdots p\cdots i_{r}}_{j_{1}\cdots q\cdots j_{s}}}\Gamma_{j_{b}l}^{q}  \Gamma_{pk}^{i_{c}}=0
\]
\[
	\sum_{c=1}^{r}\sum_{\substack{a=1,\\ a\neq c}}^{r}\tensor{t}{^{i_{1}\cdots q\cdots p\cdots i_{r}}_{j_{1}j_{2}\cdots j_{s}}}\Gamma_{qk}^{i_{a}}\Gamma_{pl}^{i_{c}}-\sum_{c=1}^{r}\sum_{\substack{a=1,\\ a\neq c}}^{r}\tensor{t}{^{i_{1}\cdots q\cdots p\cdots i_{r}}_{j_{1}j_{2}\cdots j_{s}}}\Gamma_{ql}^{i_{a}}\Gamma_{pk}^{i_{c}}=0
\]
\[
\sum_{d=1}^{s}\sum_{\substack{b=1,\\ b\neq d}}^{s}\tensor{t}{^{i_{1}i_{2}\cdots  i_{r}}_{j_{1}\cdots q\cdots p\cdots j_{s}}}\Gamma_{j_{b}k}^{q}\Gamma_{j_{d}l}^{p}-\sum_{d=1}^{s}\sum_{\substack{b=1,\\ b\neq d}}^{s}\tensor{t}{^{i_{1}i_{2}\cdots  i_{r}}_{j_{1}\cdots q\cdots p\cdots j_{s}}}\Gamma_{j_{b}l}^{q}\Gamma_{j_{d}k}^{p}=0
\]
于是我们得到
\[
\begin{split}
	&\left(\tensor{t}{^{i_{1}i_{2}\cdots i_{r}}_{j_{1}j_{2}\cdots j_{s}}}\right)^{\prime\prime}_{ kl}-\left(\tensor{t}{^{i_{1}i_{2}\cdots i_{r}}_{j_{1}j_{2}\cdots j_{s}}}\right)^{\prime\prime}_{ lk}\\
	&=\sum_{c=1}^{r}\tensor{t}{^{i_{1}\cdots q\cdots i_{r}}_{j_{1}j_{2}\cdots j_{s}}}\frac{\partial  }{\partial x^{l}} \left( \Gamma_{qk}^{i_{c}} \right)-\sum_{d=1}^{s} \tensor{t}{^{i_{1}i_{2}\cdots i_{r}}_{j_{1}\cdots q\cdots j_{s}}}\frac{\partial  }{\partial x^{l}}\left(\Gamma_{j_{d}k}^{q} \right)\\
	&+\sum_{c=1}^{r}\tensor{t}{^{i_{1}\cdots q\cdots i_{r}}_{j_{1}j_{2}\cdots j_{s}}} \Gamma_{qk}^{p}\Gamma_{pl}^{i_{c}}+\sum_{d=1}^{s}\tensor{t}{^{i_{1}i_{2}\cdots  i_{r}}_{j_{1}\cdots q\cdots j_{s}}} \Gamma_{pk}^{q}\Gamma_{j_{d}l}^{p}\\
	&-\sum_{c=1}^{r}\tensor{t}{^{i_{1}\cdots q\cdots i_{r}}_{j_{1}j_{2}\cdots j_{s}}}\frac{\partial  }{\partial x^{k}} \left( \Gamma_{ql}^{i_{c}} \right) +\sum_{d=1}^{s} \tensor{t}{^{i_{1}i_{2}\cdots i_{r}}_{j_{1}\cdots q\cdots j_{s}}}\frac{\partial  }{\partial x^{k}}\left(\Gamma_{j_{d}l}^{q} \right)\\
	&-\sum_{c=1}^{r}\tensor{t}{^{i_{1}\cdots q\cdots i_{r}}_{j_{1}j_{2}\cdots j_{s}}} \Gamma_{ql}^{p}\Gamma_{pk}^{i_{c}}-\sum_{d=1}^{s}\tensor{t}{^{i_{1}i_{2}\cdots  i_{r}}_{j_{1}\cdots q\cdots j_{s}}} \Gamma_{pl}^{q}\Gamma_{j_{d}k}^{p}\\
	&=\sum_{c=1}^{r}\tensor{t}{^{i_{1}\cdots q\cdots i_{r}}_{j_{1}j_{2}\cdots j_{s}}}\left(\frac{\partial  }{\partial x^{l}} \left( \Gamma_{qk}^{i_{c}} \right)-\frac{\partial  }{\partial x^{k}} \left( \Gamma_{ql}^{i_{c}} \right)- \Gamma_{ql}^{p}\Gamma_{pk}^{i_{c}}+ \Gamma_{qk}^{p}\Gamma_{pl}^{i_{c}} \right)\\
	&-\sum_{d=1}^{s} \tensor{t}{^{i_{1}i_{2}\cdots i_{r}}_{j_{1}\cdots q\cdots j_{s}}}\left( \frac{\partial  }{\partial x^{l}}\left(\Gamma_{j_{d}k}^{q} \right)- \frac{\partial  }{\partial x^{k}}\left(\Gamma_{j_{d}l}^{q} \right)-\Gamma_{pk}^{q}\Gamma_{j_{d}l}^{p}+\Gamma_{pl}^{q}\Gamma_{j_{d}k}^{p}\right)
\end{split}
\]
再根据式\requ{equ:}，便有
\begin{equation}
\begin{split}
	&\left(\tensor{t}{^{i_{1}i_{2}\cdots i_{r}}_{j_{1}j_{2}\cdots j_{s}}}\right)^{\prime\prime}_{ kl}-\left(\tensor{t}{^{i_{1}i_{2}\cdots i_{r}}_{j_{1}j_{2}\cdots j_{s}}}\right)^{\prime\prime}_{ lk}\\
	&=\sum_{c=1}^{r} \tensor{t}{^{i_{1}\cdots q\cdots i_{r}}_{j_{1}j_{2}\cdots j_{s}}}\tensor{R}{_q^{i_{c}}_{lk}} -\sum_{d=1}^{s} \tensor{t}{^{i_{1}i_{2}\cdots i_{r}}_{j_{1}\cdots q\cdots j_{s}}}\tensor{R}{_{j_{d}}^{q}_{lk}} \\
	&=-\sum_{a=1}^{r} \tensor{t}{^{i_{1}\cdots q\cdots i_{r}}_{j_{1}j_{2}\cdots j_{s}}}\tensor{R}{_q^{i_{a}}_{kl}} -\sum_{b=1}^{s} \tensor{t}{^{i_{1}i_{2}\cdots i_{r}}_{j_{1}\cdots q\cdots j_{s}}}\tensor{R}{_{j_{b}}^{q}_{lk}} 
\end{split}
\end{equation}
\end{proof}

\begin{remark}
要强调一点：我们这里所给的Ricci恒等式是在{\color{red}{自然标架场}}(由局部坐标卡给出)下计算得到的，利用了$[\frac { \partial } { \partial x _ { i } } ,\frac { \partial } { \partial x _ { j } }]=C_{ij}^{k}\frac { \partial } { \partial x _ { k } }=0,\forall i,j$以及联络的无挠性。如果在一般局部标架场下的进行计算，则所得结果中还会含有 $ C_{ij}^{k} $.
\end{remark}

